
\section{Variable importance}
\label{sec:rh_fi}
\ifdefined\rhfipNOrigins\else% AUTO-GENERATED by writeup/make_rehaul_fi_pooled.py -- do not edit; re-run the script.
\providecommand{\ser}[2]{#2}\providecommand{\serx}[3]{#3}
% \rhfip... macros for writeup/rehaul/08_feature_importance.tex (results/tree_shap_w3000_2014start_fomcfix_2026-10-26, the
% FOMC-patched design, its linear reference inputs/fiw_gamma_fomcfix.npz; the window study's ...W... macros:
% results/fi_by_window_2014start_2026-10-22, pre-patch; the final pooled runs, W = 3,000 trading days).
\providecommand{\rhfipNOrigins}{}\renewcommand{\rhfipNOrigins}{126}  % monthly origins (re-estimations explained)
\providecommand{\rhfipNRows}{}\renewcommand{\rhfipNRows}{33{,}970}  % bars explained, every bar of the evaluation span
\providecommand{\rhfipOriginEvery}{}\renewcommand{\rhfipOriginEvery}{21}  % trading days between origins
\providecommand{\rhfipLastBlock}{}\renewcommand{\rhfipLastBlock}{18}  % trading days of the last origin's block
\providecommand{\rhfipFirstOrigin}{}\renewcommand{\rhfipFirstOrigin}{2~January~2014}  % first origin (10:00 bar)
\providecommand{\rhfipLastOrigin}{}\renewcommand{\rhfipLastOrigin}{5~April~2024}  % last origin
\providecommand{\rhfipLastDay}{}\renewcommand{\rhfipLastDay}{30~April~2024}  % day of the last bar explained
\providecommand{\rhfipWindowDays}{}\renewcommand{\rhfipWindowDays}{3{,}000}  % rolling window, trading days
\providecommand{\rhfipWindowBars}{}\renewcommand{\rhfipWindowBars}{39{,}000}  % rolling window, bars
\providecommand{\rhfipNCols}{}\renewcommand{\rhfipNCols}{564}  % predictors of the full set
\providecommand{\rhfipNGroups}{}\renewcommand{\rhfipNGroups}{48}  % predictor groups: 46 exogenous series + HAR components + calendar dummies
\providecommand{\rhfipNExog}{}\renewcommand{\rhfipNExog}{46}  % exogenous series
\providecommand{\rhfipNPerSeries}{}\renewcommand{\rhfipNPerSeries}{12}  % predictors per exogenous series: value and availability averages at six horizons
\providecommand{\rhfipNStd}{}\renewcommand{\rhfipNStd}{12}  % standard groups
\providecommand{\rhfipNNonStd}{}\renewcommand{\rhfipNNonStd}{36}  % non-standard series
\providecommand{\rhfipNStdExog}{}\renewcommand{\rhfipNStdExog}{10}  % standard exogenous series
\providecommand{\rhfipTuneLast}{}\renewcommand{\rhfipTuneLast}{31~December~2013}  % last day of the last tuning fold (fold 3)
\providecommand{\rhfipRidgePen}{}\renewcommand{\rhfipRidgePen}{\ensuremath{10^{3}}}
\providecommand{\rhfipLassoPen}{}\renewcommand{\rhfipLassoPen}{\ensuremath{10^{-4}}}
\providecommand{\rhfipEnetPen}{}\renewcommand{\rhfipEnetPen}{\ensuremath{10^{-3}}}
\providecommand{\rhfipHarPf}{}\renewcommand{\rhfipHarPf}{0.01}  % penalty factor, HAR block
\providecommand{\rhfipBootBlock}{}\renewcommand{\rhfipBootBlock}{12}  % bootstrap block, origins
\providecommand{\rhfipBootB}{}\renewcommand{\rhfipBootB}{2{,}000}
\providecommand{\rhfipCiLevel}{}\renewcommand{\rhfipCiLevel}{95}  % percentile intervals of both rules (RULE_shap.txt, RULE_fiw.txt)
\providecommand{\rhfipShapeMin}{}\renewcommand{\rhfipShapeMin}{0.5}  % M3 selection: share of at least this many percent
\providecommand{\rhfipTopN}{}\renewcommand{\rhfipTopN}{10}  % length of the top lists compared (m2_rank.csv top10_overlap)
\providecommand{\rhfipSlopeUnit}{}\renewcommand{\rhfipSlopeUnit}{1{,}000}  % trading days per unit of the window slope
\providecommand{\rhfipNMom}{}\renewcommand{\rhfipNMom}{seven}  % the futures' realized moments (standard)
\providecommand{\rhfipLgbmGateEq}{}\renewcommand{\rhfipLgbmGateEq}{126}  % origins bit-equal to the merged every-bar walk
\providecommand{\rhfipXgbGateEq}{}\renewcommand{\rhfipXgbGateEq}{126}  % origins bit-equal to the merged every-bar walk
\providecommand{\rhfipRfGateEq}{}\renewcommand{\rhfipRfGateEq}{126}  % origins bit-equal to the merged every-bar walk
\providecommand{\rhfipRfGateAtShap}{}\renewcommand{\rhfipRfGateAtShap}{126}  % forest origins the walk had written when the SHAP gate ran
\providecommand{\rhfipXgbAddMax}{}\renewcommand{\rhfipXgbAddMax}{\ensuremath{3.7\times 10^{-6}}}  % XGBoost additivity, float32
\providecommand{\rhfipXgbAddTol}{}\renewcommand{\rhfipXgbAddTol}{\ensuremath{10^{-5}}}
\providecommand{\rhfipAddMaxDouble}{}\renewcommand{\rhfipAddMaxDouble}{\ensuremath{5.1\times 10^{-14}}}
\providecommand{\rhfipLinGateMax}{}\renewcommand{\rhfipLinGateMax}{\ensuremath{4.4\times 10^{-16}}}  % linear replay vs walk forecasts, W = 3,000
\providecommand{\rhfipInterRows}{}\renewcommand{\rhfipInterRows}{1{,}008}  % M5 rows
\providecommand{\rhfipInterPerOrigin}{}\renewcommand{\rhfipInterPerOrigin}{8}  % M5 rows per origin
\providecommand{\rhfipFomcRows}{}\renewcommand{\rhfipFomcRows}{1{,}014}
\providecommand{\rhfipFomcDaysShap}{}\renewcommand{\rhfipFomcDaysShap}{78}  % FOMC statement days in the SHAP subset (M4; releases.parquet only)
\providecommand{\rhfipFomcDays}{}\renewcommand{\rhfipFomcDays}{81}  % scheduled FOMC statement days in the span (loader calendar)
\providecommand{\rhfipFomcDaysMissing}{}\renewcommand{\rhfipFomcDaysMissing}{3}  % statement days the SHAP subset leaves out (counted with the other days; the patched rerun's reducer keeps releases.parquet's days)
\providecommand{\rhfipNHolSessions}{}\renewcommand{\rhfipNHolSessions}{64}  % NYSE-holiday futures sessions among the explained bars (not scored in Table 1)
\providecommand{\rhfipNHolBars}{}\renewcommand{\rhfipNHolBars}{443}  % their regular-hours bars
\providecommand{\rhfipVixLo}{}\renewcommand{\rhfipVixLo}{13.96}
\providecommand{\rhfipVixHi}{}\renewcommand{\rhfipVixHi}{19.02}
\providecommand{\rhfipLgbmHar}{}\renewcommand{\rhfipLgbmHar}{51.7}
\providecommand{\rhfipLgbmHarLo}{}\renewcommand{\rhfipLgbmHarLo}{49.5}
\providecommand{\rhfipLgbmHarHi}{}\renewcommand{\rhfipLgbmHarHi}{53.7}
\providecommand{\rhfipLgbmMom}{}\renewcommand{\rhfipLgbmMom}{22.3}
\providecommand{\rhfipLgbmNs}{}\renewcommand{\rhfipLgbmNs}{21.8}
\providecommand{\rhfipLgbmNsLo}{}\renewcommand{\rhfipLgbmNsLo}{20.4}
\providecommand{\rhfipLgbmNsHi}{}\renewcommand{\rhfipLgbmNsHi}{23.4}
\providecommand{\rhfipLgbmNsEx}{}\renewcommand{\rhfipLgbmNsEx}{45.9}
\providecommand{\rhfipLgbmNsExLo}{}\renewcommand{\rhfipLgbmNsExLo}{44.2}
\providecommand{\rhfipLgbmNsExHi}{}\renewcommand{\rhfipLgbmNsExHi}{47.7}
\providecommand{\rhfipXgbHar}{}\renewcommand{\rhfipXgbHar}{52.0}
\providecommand{\rhfipXgbHarLo}{}\renewcommand{\rhfipXgbHarLo}{50.0}
\providecommand{\rhfipXgbHarHi}{}\renewcommand{\rhfipXgbHarHi}{53.8}
\providecommand{\rhfipXgbMom}{}\renewcommand{\rhfipXgbMom}{23.6}
\providecommand{\rhfipXgbNs}{}\renewcommand{\rhfipXgbNs}{20.6}
\providecommand{\rhfipXgbNsLo}{}\renewcommand{\rhfipXgbNsLo}{19.2}
\providecommand{\rhfipXgbNsHi}{}\renewcommand{\rhfipXgbNsHi}{22.2}
\providecommand{\rhfipXgbNsEx}{}\renewcommand{\rhfipXgbNsEx}{43.6}
\providecommand{\rhfipXgbNsExLo}{}\renewcommand{\rhfipXgbNsExLo}{41.8}
\providecommand{\rhfipXgbNsExHi}{}\renewcommand{\rhfipXgbNsExHi}{45.7}
\providecommand{\rhfipRfHar}{}\renewcommand{\rhfipRfHar}{57.6}
\providecommand{\rhfipRfHarLo}{}\renewcommand{\rhfipRfHarLo}{55.4}
\providecommand{\rhfipRfHarHi}{}\renewcommand{\rhfipRfHarHi}{59.6}
\providecommand{\rhfipRfMom}{}\renewcommand{\rhfipRfMom}{20.7}
\providecommand{\rhfipRfNs}{}\renewcommand{\rhfipRfNs}{18.6}
\providecommand{\rhfipRfNsLo}{}\renewcommand{\rhfipRfNsLo}{17.2}
\providecommand{\rhfipRfNsHi}{}\renewcommand{\rhfipRfNsHi}{20.3}
\providecommand{\rhfipRfNsEx}{}\renewcommand{\rhfipRfNsEx}{44.2}
\providecommand{\rhfipRfNsExLo}{}\renewcommand{\rhfipRfNsExLo}{42.1}
\providecommand{\rhfipRfNsExHi}{}\renewcommand{\rhfipRfNsExHi}{46.5}
\providecommand{\rhfipRidgeNsEx}{}\renewcommand{\rhfipRidgeNsEx}{68.8}
\providecommand{\rhfipRidgeNs}{}\renewcommand{\rhfipRidgeNs}{55.6}
\providecommand{\rhfipRidgeHar}{}\renewcommand{\rhfipRidgeHar}{17.5}
\providecommand{\rhfipLassoNsEx}{}\renewcommand{\rhfipLassoNsEx}{82.5}
\providecommand{\rhfipLassoNs}{}\renewcommand{\rhfipLassoNs}{76.5}
\providecommand{\rhfipLassoHar}{}\renewcommand{\rhfipLassoHar}{6.7}
\providecommand{\rhfipEnetNsEx}{}\renewcommand{\rhfipEnetNsEx}{60.0}
\providecommand{\rhfipEnetNs}{}\renewcommand{\rhfipEnetNs}{42.9}
\providecommand{\rhfipEnetHar}{}\renewcommand{\rhfipEnetHar}{26.4}
\providecommand{\rhfipHarMin}{}\renewcommand{\rhfipHarMin}{51.7}
\providecommand{\rhfipHarMax}{}\renewcommand{\rhfipHarMax}{57.6}
\providecommand{\rhfipMomMin}{}\renewcommand{\rhfipMomMin}{20.7}
\providecommand{\rhfipMomMax}{}\renewcommand{\rhfipMomMax}{23.6}
\providecommand{\rhfipNsExMin}{}\renewcommand{\rhfipNsExMin}{43.6}
\providecommand{\rhfipNsExMax}{}\renewcommand{\rhfipNsExMax}{45.9}
\providecommand{\rhfipSentMin}{}\renewcommand{\rhfipSentMin}{0.6}
\providecommand{\rhfipSentMax}{}\renewcommand{\rhfipSentMax}{0.8}
\providecommand{\rhfipVdMin}{}\renewcommand{\rhfipVdMin}{0.2}
\providecommand{\rhfipVdMax}{}\renewcommand{\rhfipVdMax}{0.3}
\providecommand{\rhfipFomcBucketMin}{}\renewcommand{\rhfipFomcBucketMin}{1.4}
\providecommand{\rhfipFomcBucketMax}{}\renewcommand{\rhfipFomcBucketMax}{2.3}
\providecommand{\rhfipVvixMin}{}\renewcommand{\rhfipVvixMin}{0.2}
\providecommand{\rhfipVvixMax}{}\renewcommand{\rhfipVvixMax}{0.3}
\providecommand{\rhfipVixThreeMMin}{}\renewcommand{\rhfipVixThreeMMin}{0.3}
\providecommand{\rhfipVixThreeMMax}{}\renewcommand{\rhfipVixThreeMMax}{0.7}
\providecommand{\rhfipFomcDayMin}{}\renewcommand{\rhfipFomcDayMin}{0.1}
\providecommand{\rhfipFomcDayMax}{}\renewcommand{\rhfipFomcDayMax}{0.1}
\providecommand{\rhfipNsOne}{}\renewcommand{\rhfipNsOne}{\ensuremath{\ser{sumabsret_vwstock}{\mathrm{RAV}^{\mathrm{VW}}}}}
\providecommand{\rhfipNsOneMin}{}\renewcommand{\rhfipNsOneMin}{2.4}
\providecommand{\rhfipNsOneMax}{}\renewcommand{\rhfipNsOneMax}{3.9}
\providecommand{\rhfipNsOneLgbm}{}\renewcommand{\rhfipNsOneLgbm}{3.5}
\providecommand{\rhfipNsOneXgb}{}\renewcommand{\rhfipNsOneXgb}{3.9}
\providecommand{\rhfipNsOneRf}{}\renewcommand{\rhfipNsOneRf}{2.4}
\providecommand{\rhfipNsTwo}{}\renewcommand{\rhfipNsTwo}{\ensuremath{\ser{sumabsret_ewstock}{\mathrm{RAV}^{\mathrm{EW}}}}}
\providecommand{\rhfipNsTwoMin}{}\renewcommand{\rhfipNsTwoMin}{1.9}
\providecommand{\rhfipNsTwoMax}{}\renewcommand{\rhfipNsTwoMax}{3.9}
\providecommand{\rhfipNsTwoLgbm}{}\renewcommand{\rhfipNsTwoLgbm}{3.9}
\providecommand{\rhfipNsTwoXgb}{}\renewcommand{\rhfipNsTwoXgb}{3.8}
\providecommand{\rhfipNsTwoRf}{}\renewcommand{\rhfipNsTwoRf}{1.9}
\providecommand{\rhfipNsThree}{}\renewcommand{\rhfipNsThree}{\ensuremath{\ser{fomc_until_inv}{\mathrm{FOMC}^{\mathrm{next}}}}}
\providecommand{\rhfipNsThreeMin}{}\renewcommand{\rhfipNsThreeMin}{1.1}
\providecommand{\rhfipNsThreeMax}{}\renewcommand{\rhfipNsThreeMax}{2.0}
\providecommand{\rhfipNsThreeLgbm}{}\renewcommand{\rhfipNsThreeLgbm}{2.0}
\providecommand{\rhfipNsThreeXgb}{}\renewcommand{\rhfipNsThreeXgb}{2.0}
\providecommand{\rhfipNsThreeRf}{}\renewcommand{\rhfipNsThreeRf}{1.1}
\providecommand{\rhfipNsFour}{}\renewcommand{\rhfipNsFour}{\ensuremath{\ser{ofi_vwstock}{\mathrm{OFI}^{\mathrm{VW}}}}}
\providecommand{\rhfipNsFourMin}{}\renewcommand{\rhfipNsFourMin}{1.4}
\providecommand{\rhfipNsFourMax}{}\renewcommand{\rhfipNsFourMax}{1.8}
\providecommand{\rhfipNsFourLgbm}{}\renewcommand{\rhfipNsFourLgbm}{1.8}
\providecommand{\rhfipNsFourXgb}{}\renewcommand{\rhfipNsFourXgb}{1.7}
\providecommand{\rhfipNsFourRf}{}\renewcommand{\rhfipNsFourRf}{1.4}
\providecommand{\rhfipNsFive}{}\renewcommand{\rhfipNsFive}{\ensuremath{\ser{effspread_vwstock}{\mathrm{EffSpread}^{\mathrm{VW}}}}}
\providecommand{\rhfipNsFiveMin}{}\renewcommand{\rhfipNsFiveMin}{1.2}
\providecommand{\rhfipNsFiveMax}{}\renewcommand{\rhfipNsFiveMax}{1.3}
\providecommand{\rhfipNsFiveLgbm}{}\renewcommand{\rhfipNsFiveLgbm}{1.2}
\providecommand{\rhfipNsFiveXgb}{}\renewcommand{\rhfipNsFiveXgb}{1.2}
\providecommand{\rhfipNsFiveRf}{}\renewcommand{\rhfipNsFiveRf}{1.3}
\providecommand{\rhfipTreeTopOne}{}\renewcommand{\rhfipTreeTopOne}{\ensuremath{\ser{sumabsret}{\mathrm{RAV}}}}
\providecommand{\rhfipTreeTopTwo}{}\renewcommand{\rhfipTreeTopTwo}{\ensuremath{\ser{sumret}{r}}}
\providecommand{\rhfipLgbmTopOneEx}{}\renewcommand{\rhfipLgbmTopOneEx}{20.3}
\providecommand{\rhfipXgbTopOneEx}{}\renewcommand{\rhfipXgbTopOneEx}{24.7}
\providecommand{\rhfipRfTopOneEx}{}\renewcommand{\rhfipRfTopOneEx}{29.2}
\providecommand{\rhfipRidgeTopOne}{}\renewcommand{\rhfipRidgeTopOne}{\ensuremath{\ser{fomc_until_inv}{\mathrm{FOMC}^{\mathrm{next}}}}}
\providecommand{\rhfipRidgeTopOneEx}{}\renewcommand{\rhfipRidgeTopOneEx}{9.0}
\providecommand{\rhfipLgbmNUsedAll}{}\renewcommand{\rhfipLgbmNUsedAll}{42}  % exogenous series split on at every origin
\providecommand{\rhfipLgbmUsedMin}{}\renewcommand{\rhfipLgbmUsedMin}{58}
\providecommand{\rhfipXgbNUsedAll}{}\renewcommand{\rhfipXgbNUsedAll}{43}  % exogenous series split on at every origin
\providecommand{\rhfipXgbUsedMin}{}\renewcommand{\rhfipXgbUsedMin}{72}
\providecommand{\rhfipRfNUsedAll}{}\renewcommand{\rhfipRfNUsedAll}{46}  % exogenous series split on at every origin
\providecommand{\rhfipRfUsedMin}{}\renewcommand{\rhfipRfUsedMin}{126}
\providecommand{\rhfipLgbmDRidge}{}\renewcommand{\rhfipLgbmDRidge}{\ensuremath{-}22.9}
\providecommand{\rhfipLgbmDRidgeLo}{}\renewcommand{\rhfipLgbmDRidgeLo}{\ensuremath{-}24.8}
\providecommand{\rhfipLgbmDRidgeHi}{}\renewcommand{\rhfipLgbmDRidgeHi}{\ensuremath{-}21.1}
\providecommand{\rhfipLgbmDLasso}{}\renewcommand{\rhfipLgbmDLasso}{\ensuremath{-}36.6}
\providecommand{\rhfipLgbmDLassoLo}{}\renewcommand{\rhfipLgbmDLassoLo}{\ensuremath{-}39.1}
\providecommand{\rhfipLgbmDLassoHi}{}\renewcommand{\rhfipLgbmDLassoHi}{\ensuremath{-}34.4}
\providecommand{\rhfipLgbmDEnet}{}\renewcommand{\rhfipLgbmDEnet}{\ensuremath{-}14.1}
\providecommand{\rhfipLgbmDEnetLo}{}\renewcommand{\rhfipLgbmDEnetLo}{\ensuremath{-}15.7}
\providecommand{\rhfipLgbmDEnetHi}{}\renewcommand{\rhfipLgbmDEnetHi}{\ensuremath{-}12.7}
\providecommand{\rhfipXgbDRidge}{}\renewcommand{\rhfipXgbDRidge}{\ensuremath{-}25.1}
\providecommand{\rhfipXgbDRidgeLo}{}\renewcommand{\rhfipXgbDRidgeLo}{\ensuremath{-}27.0}
\providecommand{\rhfipXgbDRidgeHi}{}\renewcommand{\rhfipXgbDRidgeHi}{\ensuremath{-}23.3}
\providecommand{\rhfipXgbDLasso}{}\renewcommand{\rhfipXgbDLasso}{\ensuremath{-}38.9}
\providecommand{\rhfipXgbDLassoLo}{}\renewcommand{\rhfipXgbDLassoLo}{\ensuremath{-}41.4}
\providecommand{\rhfipXgbDLassoHi}{}\renewcommand{\rhfipXgbDLassoHi}{\ensuremath{-}36.6}
\providecommand{\rhfipXgbDEnet}{}\renewcommand{\rhfipXgbDEnet}{\ensuremath{-}16.3}
\providecommand{\rhfipXgbDEnetLo}{}\renewcommand{\rhfipXgbDEnetLo}{\ensuremath{-}17.8}
\providecommand{\rhfipXgbDEnetHi}{}\renewcommand{\rhfipXgbDEnetHi}{\ensuremath{-}14.9}
\providecommand{\rhfipRfDRidge}{}\renewcommand{\rhfipRfDRidge}{\ensuremath{-}24.5}
\providecommand{\rhfipRfDRidgeLo}{}\renewcommand{\rhfipRfDRidgeLo}{\ensuremath{-}26.3}
\providecommand{\rhfipRfDRidgeHi}{}\renewcommand{\rhfipRfDRidgeHi}{\ensuremath{-}22.9}
\providecommand{\rhfipRfDLasso}{}\renewcommand{\rhfipRfDLasso}{\ensuremath{-}38.3}
\providecommand{\rhfipRfDLassoLo}{}\renewcommand{\rhfipRfDLassoLo}{\ensuremath{-}40.5}
\providecommand{\rhfipRfDLassoHi}{}\renewcommand{\rhfipRfDLassoHi}{\ensuremath{-}36.3}
\providecommand{\rhfipRfDEnet}{}\renewcommand{\rhfipRfDEnet}{\ensuremath{-}15.7}
\providecommand{\rhfipRfDEnetLo}{}\renewcommand{\rhfipRfDEnetLo}{\ensuremath{-}17.0}
\providecommand{\rhfipRfDEnetHi}{}\renewcommand{\rhfipRfDEnetHi}{\ensuremath{-}14.4}
\providecommand{\rhfipNDLess}{}\renewcommand{\rhfipNDLess}{27}  % tree-minus-linear D with the 95% interval below 0
\providecommand{\rhfipNDAll}{}\renewcommand{\rhfipNDAll}{27}  % D comparisons: 3 trees x 3 linear models x 3 denominators
\providecommand{\rhfipLgbmSpRidge}{}\renewcommand{\rhfipLgbmSpRidge}{0.57}
\providecommand{\rhfipLgbmTopRidge}{}\renewcommand{\rhfipLgbmTopRidge}{6}
\providecommand{\rhfipLgbmSpLasso}{}\renewcommand{\rhfipLgbmSpLasso}{0.49}
\providecommand{\rhfipLgbmTopLasso}{}\renewcommand{\rhfipLgbmTopLasso}{5}
\providecommand{\rhfipLgbmSpEnet}{}\renewcommand{\rhfipLgbmSpEnet}{0.66}
\providecommand{\rhfipLgbmTopEnet}{}\renewcommand{\rhfipLgbmTopEnet}{6}
\providecommand{\rhfipXgbSpRidge}{}\renewcommand{\rhfipXgbSpRidge}{0.55}
\providecommand{\rhfipXgbTopRidge}{}\renewcommand{\rhfipXgbTopRidge}{6}
\providecommand{\rhfipXgbSpLasso}{}\renewcommand{\rhfipXgbSpLasso}{0.47}
\providecommand{\rhfipXgbTopLasso}{}\renewcommand{\rhfipXgbTopLasso}{5}
\providecommand{\rhfipXgbSpEnet}{}\renewcommand{\rhfipXgbSpEnet}{0.63}
\providecommand{\rhfipXgbTopEnet}{}\renewcommand{\rhfipXgbTopEnet}{6}
\providecommand{\rhfipRfSpRidge}{}\renewcommand{\rhfipRfSpRidge}{0.53}
\providecommand{\rhfipRfTopRidge}{}\renewcommand{\rhfipRfTopRidge}{5}
\providecommand{\rhfipRfSpLasso}{}\renewcommand{\rhfipRfSpLasso}{0.43}
\providecommand{\rhfipRfTopLasso}{}\renewcommand{\rhfipRfTopLasso}{4}
\providecommand{\rhfipRfSpEnet}{}\renewcommand{\rhfipRfSpEnet}{0.60}
\providecommand{\rhfipRfTopEnet}{}\renewcommand{\rhfipRfTopEnet}{5}
\providecommand{\rhfipSpLgbmXgb}{}\renewcommand{\rhfipSpLgbmXgb}{0.99}
\providecommand{\rhfipSpLgbmRf}{}\renewcommand{\rhfipSpLgbmRf}{0.90}
\providecommand{\rhfipSpXgbRf}{}\renewcommand{\rhfipSpXgbRf}{0.89}
\providecommand{\rhfipRidgeRankFomcUntil}{}\renewcommand{\rhfipRidgeRankFomcUntil}{1}
\providecommand{\rhfipTreeRankFomcUntilMin}{}\renewcommand{\rhfipTreeRankFomcUntilMin}{6}
\providecommand{\rhfipTreeRankFomcUntilMax}{}\renewcommand{\rhfipTreeRankFomcUntilMax}{10}
\providecommand{\rhfipRidgeRankFomcDay}{}\renewcommand{\rhfipRidgeRankFomcDay}{9}
\providecommand{\rhfipTreeRankFomcDayMin}{}\renewcommand{\rhfipTreeRankFomcDayMin}{39}
\providecommand{\rhfipTreeRankFomcDayMax}{}\renewcommand{\rhfipTreeRankFomcDayMax}{43}
\providecommand{\rhfipFomcUntil}{}\renewcommand{\rhfipFomcUntil}{\ensuremath{\ser{fomc_until_inv}{\mathrm{FOMC}^{\mathrm{next}}}}}
\providecommand{\rhfipFomcDay}{}\renewcommand{\rhfipFomcDay}{\ensuremath{\ser{fomc_day}{\mathrm{FOMC}^{\mathrm{day}}}}}
\providecommand{\rhfipAbsVw}{}\renewcommand{\rhfipAbsVw}{\ensuremath{\ser{sumabsret_vwstock}{\mathrm{RAV}^{\mathrm{VW}}}}}
\providecommand{\rhfipOfi}{}\renewcommand{\rhfipOfi}{\ensuremath{\ser{ofi_vwstock}{\mathrm{OFI}^{\mathrm{VW}}}}}
\providecommand{\rhfipLgbmNeff}{}\renewcommand{\rhfipLgbmNeff}{11.0}
\providecommand{\rhfipLgbmNHalf}{}\renewcommand{\rhfipLgbmNHalf}{4}
\providecommand{\rhfipLgbmNEighty}{}\renewcommand{\rhfipLgbmNEighty}{13}
\providecommand{\rhfipXgbNeff}{}\renewcommand{\rhfipXgbNeff}{9.1}
\providecommand{\rhfipXgbNHalf}{}\renewcommand{\rhfipXgbNHalf}{4}
\providecommand{\rhfipXgbNEighty}{}\renewcommand{\rhfipXgbNEighty}{11}
\providecommand{\rhfipRfNeff}{}\renewcommand{\rhfipRfNeff}{8.9}
\providecommand{\rhfipRfNHalf}{}\renewcommand{\rhfipRfNHalf}{4}
\providecommand{\rhfipRfNEighty}{}\renewcommand{\rhfipRfNEighty}{16}
\providecommand{\rhfipRidgeNeff}{}\renewcommand{\rhfipRidgeNeff}{28.5}
\providecommand{\rhfipRidgeNHalf}{}\renewcommand{\rhfipRidgeNHalf}{11}
\providecommand{\rhfipRidgeNEighty}{}\renewcommand{\rhfipRidgeNEighty}{26}
\providecommand{\rhfipLassoNeff}{}\renewcommand{\rhfipLassoNeff}{5.1}
\providecommand{\rhfipLassoNHalf}{}\renewcommand{\rhfipLassoNHalf}{2}
\providecommand{\rhfipLassoNEighty}{}\renewcommand{\rhfipLassoNEighty}{11}
\providecommand{\rhfipEnetNeff}{}\renewcommand{\rhfipEnetNeff}{16.2}
\providecommand{\rhfipEnetNHalf}{}\renewcommand{\rhfipEnetNHalf}{6}
\providecommand{\rhfipEnetNEighty}{}\renewcommand{\rhfipEnetNEighty}{15}
\providecommand{\rhfipYearFirst}{}\renewcommand{\rhfipYearFirst}{2014}
\providecommand{\rhfipYearLast}{}\renewcommand{\rhfipYearLast}{2024}
\providecommand{\rhfipLgbmNsYearMin}{}\renewcommand{\rhfipLgbmNsYearMin}{18.1}
\providecommand{\rhfipLgbmNsYearMinYr}{}\renewcommand{\rhfipLgbmNsYearMinYr}{2021}
\providecommand{\rhfipLgbmNsYearMax}{}\renewcommand{\rhfipLgbmNsYearMax}{25.6}
\providecommand{\rhfipLgbmNsYearMaxYr}{}\renewcommand{\rhfipLgbmNsYearMaxYr}{2015}
\providecommand{\rhfipLgbmHarYearMax}{}\renewcommand{\rhfipLgbmHarYearMax}{57.4}
\providecommand{\rhfipLgbmHarYearMaxYr}{}\renewcommand{\rhfipLgbmHarYearMaxYr}{2021}
\providecommand{\rhfipLgbmHarYearMin}{}\renewcommand{\rhfipLgbmHarYearMin}{44.6}
\providecommand{\rhfipLgbmHarYearMinYr}{}\renewcommand{\rhfipLgbmHarYearMinYr}{2017}
\providecommand{\rhfipXgbNsYearMin}{}\renewcommand{\rhfipXgbNsYearMin}{17.2}
\providecommand{\rhfipXgbNsYearMinYr}{}\renewcommand{\rhfipXgbNsYearMinYr}{2021}
\providecommand{\rhfipXgbNsYearMax}{}\renewcommand{\rhfipXgbNsYearMax}{24.3}
\providecommand{\rhfipXgbNsYearMaxYr}{}\renewcommand{\rhfipXgbNsYearMaxYr}{2015}
\providecommand{\rhfipXgbHarYearMax}{}\renewcommand{\rhfipXgbHarYearMax}{57.0}
\providecommand{\rhfipXgbHarYearMaxYr}{}\renewcommand{\rhfipXgbHarYearMaxYr}{2021}
\providecommand{\rhfipXgbHarYearMin}{}\renewcommand{\rhfipXgbHarYearMin}{45.8}
\providecommand{\rhfipXgbHarYearMinYr}{}\renewcommand{\rhfipXgbHarYearMinYr}{2017}
\providecommand{\rhfipRfNsYearMin}{}\renewcommand{\rhfipRfNsYearMin}{15.2}
\providecommand{\rhfipRfNsYearMinYr}{}\renewcommand{\rhfipRfNsYearMinYr}{2021}
\providecommand{\rhfipRfNsYearMax}{}\renewcommand{\rhfipRfNsYearMax}{22.9}
\providecommand{\rhfipRfNsYearMaxYr}{}\renewcommand{\rhfipRfNsYearMaxYr}{2015}
\providecommand{\rhfipRfHarYearMax}{}\renewcommand{\rhfipRfHarYearMax}{62.3}
\providecommand{\rhfipRfHarYearMaxYr}{}\renewcommand{\rhfipRfHarYearMaxYr}{2019}
\providecommand{\rhfipRfHarYearMin}{}\renewcommand{\rhfipRfHarYearMin}{50.4}
\providecommand{\rhfipRfHarYearMinYr}{}\renewcommand{\rhfipRfHarYearMinYr}{2017}
\providecommand{\rhfipNsYearMinYr}{}\renewcommand{\rhfipNsYearMinYr}{2021}  % the year of every tree's lowest non-standard share
\providecommand{\rhfipDipFirst}{}\renewcommand{\rhfipDipFirst}{2018}  % first of the years in which every tree's non-standard share is below its pooled share
\providecommand{\rhfipDipLast}{}\renewcommand{\rhfipDipLast}{2022}
\providecommand{\rhfipRidgeNsExYearMin}{}\renewcommand{\rhfipRidgeNsExYearMin}{65.1}
\providecommand{\rhfipRidgeNsExYearMinYr}{}\renewcommand{\rhfipRidgeNsExYearMinYr}{2021}
\providecommand{\rhfipRidgeNsExYearMax}{}\renewcommand{\rhfipRidgeNsExYearMax}{72.4}
\providecommand{\rhfipRidgeNsExYearMaxYr}{}\renewcommand{\rhfipRidgeNsExYearMaxYr}{2017}
\providecommand{\rhfipLassoNsExYearMin}{}\renewcommand{\rhfipLassoNsExYearMin}{77.1}
\providecommand{\rhfipLassoNsExYearMinYr}{}\renewcommand{\rhfipLassoNsExYearMinYr}{2021}
\providecommand{\rhfipLassoNsExYearMax}{}\renewcommand{\rhfipLassoNsExYearMax}{87.1}
\providecommand{\rhfipLassoNsExYearMaxYr}{}\renewcommand{\rhfipLassoNsExYearMaxYr}{2017}
\providecommand{\rhfipEnetNsExYearMin}{}\renewcommand{\rhfipEnetNsExYearMin}{56.1}
\providecommand{\rhfipEnetNsExYearMinYr}{}\renewcommand{\rhfipEnetNsExYearMinYr}{2019}
\providecommand{\rhfipEnetNsExYearMax}{}\renewcommand{\rhfipEnetNsExYearMax}{65.2}
\providecommand{\rhfipEnetNsExYearMaxYr}{}\renewcommand{\rhfipEnetNsExYearMaxYr}{2014}
\providecommand{\rhfipLgbmNsExYearMin}{}\renewcommand{\rhfipLgbmNsExYearMin}{42.3}
\providecommand{\rhfipLgbmNsExYearMax}{}\renewcommand{\rhfipLgbmNsExYearMax}{50.1}
\providecommand{\rhfipXgbNsExYearMin}{}\renewcommand{\rhfipXgbNsExYearMin}{39.4}
\providecommand{\rhfipXgbNsExYearMax}{}\renewcommand{\rhfipXgbNsExYearMax}{48.9}
\providecommand{\rhfipRfNsExYearMin}{}\renewcommand{\rhfipRfNsExYearMin}{39.5}
\providecommand{\rhfipRfNsExYearMax}{}\renewcommand{\rhfipRfNsExYearMax}{50.9}
\providecommand{\rhfipLgbmNShape}{}\renewcommand{\rhfipLgbmNShape}{11}  % non-standard series at >= 0.5% (M3)
\providecommand{\rhfipLgbmNMono}{}\renewcommand{\rhfipLgbmNMono}{7}
\providecommand{\rhfipLgbmNIrr}{}\renewcommand{\rhfipLgbmNIrr}{3}
\providecommand{\rhfipLgbmNStep}{}\renewcommand{\rhfipLgbmNStep}{1}
\providecommand{\rhfipLgbmNFlat}{}\renewcommand{\rhfipLgbmNFlat}{0}
\providecommand{\rhfipLgbmNThresh}{}\renewcommand{\rhfipLgbmNThresh}{0}
\providecommand{\rhfipLgbmNU}{}\renewcommand{\rhfipLgbmNU}{0}
\providecommand{\rhfipLgbmNAsym}{}\renewcommand{\rhfipLgbmNAsym}{9}
\providecommand{\rhfipXgbNShape}{}\renewcommand{\rhfipXgbNShape}{10}  % non-standard series at >= 0.5% (M3)
\providecommand{\rhfipXgbNMono}{}\renewcommand{\rhfipXgbNMono}{5}
\providecommand{\rhfipXgbNIrr}{}\renewcommand{\rhfipXgbNIrr}{3}
\providecommand{\rhfipXgbNStep}{}\renewcommand{\rhfipXgbNStep}{0}
\providecommand{\rhfipXgbNFlat}{}\renewcommand{\rhfipXgbNFlat}{1}
\providecommand{\rhfipXgbNThresh}{}\renewcommand{\rhfipXgbNThresh}{1}
\providecommand{\rhfipXgbNU}{}\renewcommand{\rhfipXgbNU}{0}
\providecommand{\rhfipXgbNAsym}{}\renewcommand{\rhfipXgbNAsym}{8}
\providecommand{\rhfipRfNShape}{}\renewcommand{\rhfipRfNShape}{13}  % non-standard series at >= 0.5% (M3)
\providecommand{\rhfipRfNMono}{}\renewcommand{\rhfipRfNMono}{10}
\providecommand{\rhfipRfNIrr}{}\renewcommand{\rhfipRfNIrr}{3}
\providecommand{\rhfipRfNStep}{}\renewcommand{\rhfipRfNStep}{0}
\providecommand{\rhfipRfNFlat}{}\renewcommand{\rhfipRfNFlat}{0}
\providecommand{\rhfipRfNThresh}{}\renewcommand{\rhfipRfNThresh}{0}
\providecommand{\rhfipRfNU}{}\renewcommand{\rhfipRfNU}{0}
\providecommand{\rhfipRfNAsym}{}\renewcommand{\rhfipRfNAsym}{13}
\providecommand{\rhfipNUThreshAll}{}\renewcommand{\rhfipNUThreshAll}{1}  % U-shaped, inverted-U or threshold profiles in any tree
\providecommand{\rhfipAbsVwLgbmLo}{}\renewcommand{\rhfipAbsVwLgbmLo}{\ensuremath{-}1.4}
\providecommand{\rhfipAbsVwLgbmHi}{}\renewcommand{\rhfipAbsVwLgbmHi}{3.9}
\providecommand{\rhfipAbsVwXgbLo}{}\renewcommand{\rhfipAbsVwXgbLo}{\ensuremath{-}1.5}
\providecommand{\rhfipAbsVwXgbHi}{}\renewcommand{\rhfipAbsVwXgbHi}{4.3}
\providecommand{\rhfipAbsVwRfLo}{}\renewcommand{\rhfipAbsVwRfLo}{\ensuremath{-}0.9}
\providecommand{\rhfipAbsVwRfHi}{}\renewcommand{\rhfipAbsVwRfHi}{2.3}
\providecommand{\rhfipOfiLgbmLo}{}\renewcommand{\rhfipOfiLgbmLo}{0.8}
\providecommand{\rhfipOfiLgbmHi}{}\renewcommand{\rhfipOfiLgbmHi}{\ensuremath{-}1.1}
\providecommand{\rhfipOfiXgbLo}{}\renewcommand{\rhfipOfiXgbLo}{0.8}
\providecommand{\rhfipOfiXgbHi}{}\renewcommand{\rhfipOfiXgbHi}{\ensuremath{-}0.9}
\providecommand{\rhfipOfiRfLo}{}\renewcommand{\rhfipOfiRfLo}{0.8}
\providecommand{\rhfipOfiRfHi}{}\renewcommand{\rhfipOfiRfHi}{\ensuremath{-}0.7}
\providecommand{\rhfipUntilLgbmLo}{}\renewcommand{\rhfipUntilLgbmLo}{\ensuremath{-}0.3}
\providecommand{\rhfipUntilLgbmHi}{}\renewcommand{\rhfipUntilLgbmHi}{2.9}
\providecommand{\rhfipUntilXgbLo}{}\renewcommand{\rhfipUntilXgbLo}{0.2}
\providecommand{\rhfipUntilXgbHi}{}\renewcommand{\rhfipUntilXgbHi}{2.4}
\providecommand{\rhfipUntilRfLo}{}\renewcommand{\rhfipUntilRfLo}{\ensuremath{-}0.0}
\providecommand{\rhfipUntilRfHi}{}\renewcommand{\rhfipUntilRfHi}{1.5}
\providecommand{\rhfipUntilHorizon}{}\renewcommand{\rhfipUntilHorizon}{25}  % horizon of fomc_until_inv's dominant column
\providecommand{\rhfipUntilHorizonXgb}{}\renewcommand{\rhfipUntilHorizonXgb}{1}  % the same in XGBoost (differs from the other trees')
\providecommand{\rhfipBarFirst}{}\renewcommand{\rhfipBarFirst}{10:00}
\providecommand{\rhfipBarLast}{}\renewcommand{\rhfipBarLast}{16:00}
\providecommand{\rhfipBarPeakA}{}\renewcommand{\rhfipBarPeakA}{14:30}
\providecommand{\rhfipBarPeakB}{}\renewcommand{\rhfipBarPeakB}{15:00}
\providecommand{\rhfipNBars}{}\renewcommand{\rhfipNBars}{13}
\providecommand{\rhfipLgbmNsFomc}{}\renewcommand{\rhfipLgbmNsFomc}{42.6}
\providecommand{\rhfipLgbmNsOther}{}\renewcommand{\rhfipLgbmNsOther}{20.9}
\providecommand{\rhfipLgbmUntilFomc}{}\renewcommand{\rhfipLgbmUntilFomc}{25.4}
\providecommand{\rhfipLgbmUntilOther}{}\renewcommand{\rhfipLgbmUntilOther}{1.06}
\providecommand{\rhfipLgbmFomcBucketFomc}{}\renewcommand{\rhfipLgbmFomcBucketFomc}{27.3}
\providecommand{\rhfipLgbmFomcBucketOther}{}\renewcommand{\rhfipLgbmFomcBucketOther}{1.3}
\providecommand{\rhfipLgbmHarFomc}{}\renewcommand{\rhfipLgbmHarFomc}{39.2}
\providecommand{\rhfipLgbmHarOther}{}\renewcommand{\rhfipLgbmHarOther}{52.2}
\providecommand{\rhfipLgbmNsFirstBar}{}\renewcommand{\rhfipLgbmNsFirstBar}{23.9}
\providecommand{\rhfipLgbmNsLastBar}{}\renewcommand{\rhfipLgbmNsLastBar}{20.5}
\providecommand{\rhfipLgbmNsAftMin}{}\renewcommand{\rhfipLgbmNsAftMin}{24.3}
\providecommand{\rhfipLgbmNsAftMax}{}\renewcommand{\rhfipLgbmNsAftMax}{25.2}
\providecommand{\rhfipLgbmNsMidMin}{}\renewcommand{\rhfipLgbmNsMidMin}{20.1}
\providecommand{\rhfipLgbmNsMidMax}{}\renewcommand{\rhfipLgbmNsMidMax}{21.9}
\providecommand{\rhfipLgbmFomcBucketAft}{}\renewcommand{\rhfipLgbmFomcBucketAft}{5.5}
\providecommand{\rhfipLgbmFomcBucketOtherBars}{}\renewcommand{\rhfipLgbmFomcBucketOtherBars}{1.4}  % median FOMC-announcement share over the other 11 bars
\providecommand{\rhfipLgbmFomcDiff}{}\renewcommand{\rhfipLgbmFomcDiff}{\ensuremath{+}21.7}
\providecommand{\rhfipLgbmFomcDiffLo}{}\renewcommand{\rhfipLgbmFomcDiffLo}{\ensuremath{+}19.3}
\providecommand{\rhfipLgbmFomcDiffHi}{}\renewcommand{\rhfipLgbmFomcDiffHi}{\ensuremath{+}24.1}
\providecommand{\rhfipLgbmLateDiff}{}\renewcommand{\rhfipLgbmLateDiff}{\ensuremath{-}3.4}
\providecommand{\rhfipLgbmLateDiffLo}{}\renewcommand{\rhfipLgbmLateDiffLo}{\ensuremath{-}4.3}
\providecommand{\rhfipLgbmLateDiffHi}{}\renewcommand{\rhfipLgbmLateDiffHi}{\ensuremath{-}2.5}
\providecommand{\rhfipLgbmVixDiff}{}\renewcommand{\rhfipLgbmVixDiff}{\ensuremath{-}3.2}
\providecommand{\rhfipLgbmVixDiffLo}{}\renewcommand{\rhfipLgbmVixDiffLo}{\ensuremath{-}5.3}
\providecommand{\rhfipLgbmVixDiffHi}{}\renewcommand{\rhfipLgbmVixDiffHi}{\ensuremath{-}1.1}
\providecommand{\rhfipLgbmVixHar}{}\renewcommand{\rhfipLgbmVixHar}{\ensuremath{+}5.1}
\providecommand{\rhfipLgbmVixHarLo}{}\renewcommand{\rhfipLgbmVixHarLo}{\ensuremath{+}1.8}
\providecommand{\rhfipLgbmVixHarHi}{}\renewcommand{\rhfipLgbmVixHarHi}{\ensuremath{+}8.4}
\providecommand{\rhfipLgbmVixDiffEx}{}\renewcommand{\rhfipLgbmVixDiffEx}{\ensuremath{-}2.3}
\providecommand{\rhfipLgbmVixDiffExLo}{}\renewcommand{\rhfipLgbmVixDiffExLo}{\ensuremath{-}4.5}
\providecommand{\rhfipLgbmVixDiffExHi}{}\renewcommand{\rhfipLgbmVixDiffExHi}{\ensuremath{-}0.2}
\providecommand{\rhfipXgbNsFomc}{}\renewcommand{\rhfipXgbNsFomc}{41.1}
\providecommand{\rhfipXgbNsOther}{}\renewcommand{\rhfipXgbNsOther}{19.7}
\providecommand{\rhfipXgbUntilFomc}{}\renewcommand{\rhfipXgbUntilFomc}{23.9}
\providecommand{\rhfipXgbUntilOther}{}\renewcommand{\rhfipXgbUntilOther}{1.10}
\providecommand{\rhfipXgbFomcBucketFomc}{}\renewcommand{\rhfipXgbFomcBucketFomc}{25.6}
\providecommand{\rhfipXgbFomcBucketOther}{}\renewcommand{\rhfipXgbFomcBucketOther}{1.3}
\providecommand{\rhfipXgbHarFomc}{}\renewcommand{\rhfipXgbHarFomc}{39.1}
\providecommand{\rhfipXgbHarOther}{}\renewcommand{\rhfipXgbHarOther}{52.5}
\providecommand{\rhfipXgbNsFirstBar}{}\renewcommand{\rhfipXgbNsFirstBar}{22.5}
\providecommand{\rhfipXgbNsLastBar}{}\renewcommand{\rhfipXgbNsLastBar}{19.5}
\providecommand{\rhfipXgbNsAftMin}{}\renewcommand{\rhfipXgbNsAftMin}{23.5}
\providecommand{\rhfipXgbNsAftMax}{}\renewcommand{\rhfipXgbNsAftMax}{23.7}
\providecommand{\rhfipXgbNsMidMin}{}\renewcommand{\rhfipXgbNsMidMin}{19.1}
\providecommand{\rhfipXgbNsMidMax}{}\renewcommand{\rhfipXgbNsMidMax}{20.4}
\providecommand{\rhfipXgbFomcBucketAft}{}\renewcommand{\rhfipXgbFomcBucketAft}{5.8}
\providecommand{\rhfipXgbFomcBucketOtherBars}{}\renewcommand{\rhfipXgbFomcBucketOtherBars}{1.3}  % median FOMC-announcement share over the other 11 bars
\providecommand{\rhfipXgbFomcDiff}{}\renewcommand{\rhfipXgbFomcDiff}{\ensuremath{+}21.4}
\providecommand{\rhfipXgbFomcDiffLo}{}\renewcommand{\rhfipXgbFomcDiffLo}{\ensuremath{+}19.3}
\providecommand{\rhfipXgbFomcDiffHi}{}\renewcommand{\rhfipXgbFomcDiffHi}{\ensuremath{+}23.5}
\providecommand{\rhfipXgbLateDiff}{}\renewcommand{\rhfipXgbLateDiff}{\ensuremath{-}3.1}
\providecommand{\rhfipXgbLateDiffLo}{}\renewcommand{\rhfipXgbLateDiffLo}{\ensuremath{-}4.0}
\providecommand{\rhfipXgbLateDiffHi}{}\renewcommand{\rhfipXgbLateDiffHi}{\ensuremath{-}2.3}
\providecommand{\rhfipXgbVixDiff}{}\renewcommand{\rhfipXgbVixDiff}{\ensuremath{-}3.0}
\providecommand{\rhfipXgbVixDiffLo}{}\renewcommand{\rhfipXgbVixDiffLo}{\ensuremath{-}5.1}
\providecommand{\rhfipXgbVixDiffHi}{}\renewcommand{\rhfipXgbVixDiffHi}{\ensuremath{-}0.9}
\providecommand{\rhfipXgbVixHar}{}\renewcommand{\rhfipXgbVixHar}{\ensuremath{+}3.9}
\providecommand{\rhfipXgbVixHarLo}{}\renewcommand{\rhfipXgbVixHarLo}{\ensuremath{+}0.4}
\providecommand{\rhfipXgbVixHarHi}{}\renewcommand{\rhfipXgbVixHarHi}{\ensuremath{+}7.2}
\providecommand{\rhfipXgbVixDiffEx}{}\renewcommand{\rhfipXgbVixDiffEx}{\ensuremath{-}3.2}
\providecommand{\rhfipXgbVixDiffExLo}{}\renewcommand{\rhfipXgbVixDiffExLo}{\ensuremath{-}5.7}
\providecommand{\rhfipXgbVixDiffExHi}{}\renewcommand{\rhfipXgbVixDiffExHi}{\ensuremath{-}0.8}
\providecommand{\rhfipRfNsFomc}{}\renewcommand{\rhfipRfNsFomc}{28.9}
\providecommand{\rhfipRfNsOther}{}\renewcommand{\rhfipRfNsOther}{18.3}
\providecommand{\rhfipRfUntilFomc}{}\renewcommand{\rhfipRfUntilFomc}{12.4}
\providecommand{\rhfipRfFomcBucketFomc}{}\renewcommand{\rhfipRfFomcBucketFomc}{13.0}
\providecommand{\rhfipRfFomcBucketOther}{}\renewcommand{\rhfipRfFomcBucketOther}{1.0}
\providecommand{\rhfipRfHarFomc}{}\renewcommand{\rhfipRfHarFomc}{50.8}
\providecommand{\rhfipRfHarOther}{}\renewcommand{\rhfipRfHarOther}{57.9}
\providecommand{\rhfipRfNsFirstBar}{}\renewcommand{\rhfipRfNsFirstBar}{20.1}
\providecommand{\rhfipRfNsLastBar}{}\renewcommand{\rhfipRfNsLastBar}{17.8}
\providecommand{\rhfipRfNsAftMin}{}\renewcommand{\rhfipRfNsAftMin}{19.5}
\providecommand{\rhfipRfNsAftMax}{}\renewcommand{\rhfipRfNsAftMax}{19.9}
\providecommand{\rhfipRfNsMidMin}{}\renewcommand{\rhfipRfNsMidMin}{17.8}
\providecommand{\rhfipRfNsMidMax}{}\renewcommand{\rhfipRfNsMidMax}{18.6}
\providecommand{\rhfipRfFomcBucketAft}{}\renewcommand{\rhfipRfFomcBucketAft}{3.1}
\providecommand{\rhfipRfFomcBucketOtherBars}{}\renewcommand{\rhfipRfFomcBucketOtherBars}{1.0}  % median FOMC-announcement share over the other 11 bars
\providecommand{\rhfipRfFomcDiff}{}\renewcommand{\rhfipRfFomcDiff}{\ensuremath{+}10.6}
\providecommand{\rhfipRfFomcDiffLo}{}\renewcommand{\rhfipRfFomcDiffLo}{\ensuremath{+}8.3}
\providecommand{\rhfipRfFomcDiffHi}{}\renewcommand{\rhfipRfFomcDiffHi}{\ensuremath{+}13.1}
\providecommand{\rhfipRfLateDiff}{}\renewcommand{\rhfipRfLateDiff}{\ensuremath{-}2.3}
\providecommand{\rhfipRfLateDiffLo}{}\renewcommand{\rhfipRfLateDiffLo}{\ensuremath{-}3.1}
\providecommand{\rhfipRfLateDiffHi}{}\renewcommand{\rhfipRfLateDiffHi}{\ensuremath{-}1.6}
\providecommand{\rhfipRfVixDiff}{}\renewcommand{\rhfipRfVixDiff}{\ensuremath{-}3.3}
\providecommand{\rhfipRfVixDiffLo}{}\renewcommand{\rhfipRfVixDiffLo}{\ensuremath{-}5.1}
\providecommand{\rhfipRfVixDiffHi}{}\renewcommand{\rhfipRfVixDiffHi}{\ensuremath{-}1.4}
\providecommand{\rhfipRfVixHar}{}\renewcommand{\rhfipRfVixHar}{\ensuremath{+}4.2}
\providecommand{\rhfipRfVixHarLo}{}\renewcommand{\rhfipRfVixHarLo}{\ensuremath{+}1.0}
\providecommand{\rhfipRfVixHarHi}{}\renewcommand{\rhfipRfVixHarHi}{\ensuremath{+}7.3}
\providecommand{\rhfipRfVixDiffEx}{}\renewcommand{\rhfipRfVixDiffEx}{\ensuremath{-}3.7}
\providecommand{\rhfipRfVixDiffExLo}{}\renewcommand{\rhfipRfVixDiffExLo}{\ensuremath{-}6.3}
\providecommand{\rhfipRfVixDiffExHi}{}\renewcommand{\rhfipRfVixDiffExHi}{\ensuremath{-}1.2}
\providecommand{\rhfipLgbmInterVar}{}\renewcommand{\rhfipLgbmInterVar}{33.1}
\providecommand{\rhfipLgbmInterCol}{}\renewcommand{\rhfipLgbmInterCol}{45.3}
\providecommand{\rhfipLgbmInterNs}{}\renewcommand{\rhfipLgbmInterNs}{49.6}
\providecommand{\rhfipLgbmInterStd}{}\renewcommand{\rhfipLgbmInterStd}{26.8}
\providecommand{\rhfipLgbmInterStdEx}{}\renewcommand{\rhfipLgbmInterStdEx}{40.2}  % descriptive: standard exogenous series only (no HAR, no calendar)
\providecommand{\rhfipLgbmInterHar}{}\renewcommand{\rhfipLgbmInterHar}{17.1}
\providecommand{\rhfipLgbmInterWeakMin}{}\renewcommand{\rhfipLgbmInterWeakMin}{54}
\providecommand{\rhfipLgbmInterWeakMax}{}\renewcommand{\rhfipLgbmInterWeakMax}{66}
\providecommand{\rhfipLgbmPairOne}{}\renewcommand{\rhfipLgbmPairOne}{1.25}
\providecommand{\rhfipLgbmPairTwo}{}\renewcommand{\rhfipLgbmPairTwo}{0.86}
\providecommand{\rhfipLgbmCellHarHar}{}\renewcommand{\rhfipLgbmCellHarHar}{48.0}
\providecommand{\rhfipXgbInterVar}{}\renewcommand{\rhfipXgbInterVar}{30.9}
\providecommand{\rhfipXgbInterCol}{}\renewcommand{\rhfipXgbInterCol}{42.5}
\providecommand{\rhfipXgbInterNs}{}\renewcommand{\rhfipXgbInterNs}{45.7}
\providecommand{\rhfipXgbInterStd}{}\renewcommand{\rhfipXgbInterStd}{26.0}
\providecommand{\rhfipXgbInterStdEx}{}\renewcommand{\rhfipXgbInterStdEx}{37.8}  % descriptive: standard exogenous series only (no HAR, no calendar)
\providecommand{\rhfipXgbInterHar}{}\renewcommand{\rhfipXgbInterHar}{17.6}
\providecommand{\rhfipXgbInterWeakMin}{}\renewcommand{\rhfipXgbInterWeakMin}{54}
\providecommand{\rhfipXgbInterWeakMax}{}\renewcommand{\rhfipXgbInterWeakMax}{63}
\providecommand{\rhfipXgbPairOne}{}\renewcommand{\rhfipXgbPairOne}{1.23}
\providecommand{\rhfipXgbPairTwo}{}\renewcommand{\rhfipXgbPairTwo}{0.83}
\providecommand{\rhfipXgbCellHarHar}{}\renewcommand{\rhfipXgbCellHarHar}{47.4}
\providecommand{\rhfipPairOneA}{}\renewcommand{\rhfipPairOneA}{HAR components}
\providecommand{\rhfipPairOneB}{}\renewcommand{\rhfipPairOneB}{\ensuremath{\ser{sumabsret_vwstock}{\mathrm{RAV}^{\mathrm{VW}}}}}
\providecommand{\rhfipPairTwoA}{}\renewcommand{\rhfipPairTwoA}{HAR components}
\providecommand{\rhfipPairTwoB}{}\renewcommand{\rhfipPairTwoB}{\ensuremath{\ser{sumabsret_ewstock}{\mathrm{RAV}^{\mathrm{EW}}}}}
\providecommand{\rhfipLgbmPairsHar}{}\renewcommand{\rhfipLgbmPairsHar}{7}
\providecommand{\rhfipXgbPairsHar}{}\renewcommand{\rhfipXgbPairsHar}{7}
\providecommand{\rhfipWMin}{}\renewcommand{\rhfipWMin}{1{,}000}
\providecommand{\rhfipWMax}{}\renewcommand{\rhfipWMax}{3{,}500}
\providecommand{\rhfipNWindows}{}\renewcommand{\rhfipNWindows}{6}
\providecommand{\rhfipSOneMeasures}{}\renewcommand{\rhfipSOneMeasures}{105}  % measure x common-penalty fits of the window study (P1 and P2)
\providecommand{\rhfipSOneRises}{}\renewcommand{\rhfipSOneRises}{0}  % ... of which rise with W
\providecommand{\rhfipSOneFalls}{}\renewcommand{\rhfipSOneFalls}{102}
\providecommand{\rhfipSOneFlat}{}\renewcommand{\rhfipSOneFlat}{3}
\providecommand{\rhfipSOneNotSupported}{}\renewcommand{\rhfipSOneNotSupported}{14}  % common-penalty verdicts NOT SUPPORTED
\providecommand{\rhfipSOneVerdicts}{}\renewcommand{\rhfipSOneVerdicts}{14}
\providecommand{\rhfipRidgeNeffWMin}{}\renewcommand{\rhfipRidgeNeffWMin}{29.2}
\providecommand{\rhfipRidgeNeffWMain}{}\renewcommand{\rhfipRidgeNeffWMain}{29.2}
\providecommand{\rhfipRidgeNeffWMax}{}\renewcommand{\rhfipRidgeNeffWMax}{28.2}
\providecommand{\rhfipRidgeNeffSlope}{}\renewcommand{\rhfipRidgeNeffSlope}{\ensuremath{-}0.10}
\providecommand{\rhfipRidgeNeffSlopeLo}{}\renewcommand{\rhfipRidgeNeffSlopeLo}{\ensuremath{-}0.25}
\providecommand{\rhfipRidgeNeffSlopeHi}{}\renewcommand{\rhfipRidgeNeffSlopeHi}{0.12}
\providecommand{\rhfipRidgeSOneLoNeffWMin}{}\renewcommand{\rhfipRidgeSOneLoNeffWMin}{33.9}
\providecommand{\rhfipRidgeSOneLoNeffWMax}{}\renewcommand{\rhfipRidgeSOneLoNeffWMax}{28.2}
\providecommand{\rhfipRidgeSOneHiNeffWMin}{}\renewcommand{\rhfipRidgeSOneHiNeffWMin}{29.2}
\providecommand{\rhfipRidgeSOneHiNeffWMax}{}\renewcommand{\rhfipRidgeSOneHiNeffWMax}{23.5}
\providecommand{\rhfipLassoSOneVWMin}{}\renewcommand{\rhfipLassoSOneVWMin}{38.7}
\providecommand{\rhfipLassoSOneVWMax}{}\renewcommand{\rhfipLassoSOneVWMax}{31.9}
\providecommand{\rhfipLassoSOneCWMin}{}\renewcommand{\rhfipLassoSOneCWMin}{213.3}
\providecommand{\rhfipLassoSOneCWMax}{}\renewcommand{\rhfipLassoSOneCWMax}{177.5}
\providecommand{\rhfipLassoWalkVWMin}{}\renewcommand{\rhfipLassoWalkVWMin}{38.8}
\providecommand{\rhfipLassoWalkVWMax}{}\renewcommand{\rhfipLassoWalkVWMax}{41.9}
\providecommand{\rhfipLassoSOneLo}{}\renewcommand{\rhfipLassoSOneLo}{\ensuremath{10^{-3}}}
\providecommand{\rhfipRidgeSOneLo}{}\renewcommand{\rhfipRidgeSOneLo}{\ensuremath{10^{3}}}
\providecommand{\rhfipRidgeSOneHi}{}\renewcommand{\rhfipRidgeSOneHi}{\ensuremath{10^{4}}}
\providecommand{\rhfipEnetPenWMin}{}\renewcommand{\rhfipEnetPenWMin}{\ensuremath{10^{-2}}}
\providecommand{\rhfipEnetPenWNext}{}\renewcommand{\rhfipEnetPenWNext}{\ensuremath{10^{-3}}}
\providecommand{\rhfipWNext}{}\renewcommand{\rhfipWNext}{1{,}500}
\providecommand{\rhfipRidgePOneWalk}{}\renewcommand{\rhfipRidgePOneWalk}{not supported}
\providecommand{\rhfipRidgePTwoWalk}{}\renewcommand{\rhfipRidgePTwoWalk}{not supported}
\providecommand{\rhfipLassoPOneWalk}{}\renewcommand{\rhfipLassoPOneWalk}{partly}
\providecommand{\rhfipLassoPTwoWalk}{}\renewcommand{\rhfipLassoPTwoWalk}{not supported}
\providecommand{\rhfipEnetPOneWalk}{}\renewcommand{\rhfipEnetPOneWalk}{supported}
\providecommand{\rhfipEnetPTwoWalk}{}\renewcommand{\rhfipEnetPTwoWalk}{supported}
\fi
\ifdefined\rhabLgbmFullAllPct\else% AUTO-GENERATED by writeup/make_rehaul_ablations.py -- do not edit; re-run the script.
% results/ablations_2014start_w3000 (RULE.txt): column-subset arms of the final 13-bar design,
% frozen hyperparameters, every bar scored as Table 1; DM on daily (average) QLIKE.
\newcommand{\rhabEnetAddEwAllQlike}{0.1871}  % QLIKE enet add_ew all
\newcommand{\rhabEnetAddEwAllMse}{1.47}  % MSE x 1e10 enet add_ew all
\newcommand{\rhabEnetAddEwAllPct}{$-$1.1}  % QLIKE change % vs har
\newcommand{\rhabEnetAddEwAllDm}{$-$2.81}  % DM vs the reference arm
\newcommand{\rhabEnetAddEwAllP}{0.005}  % p-value
\newcommand{\rhabEnetAddEwCloseQlike}{0.1249}  % QLIKE enet add_ew c16
\newcommand{\rhabEnetAddEwCloseMse}{2.10}  % MSE x 1e10 enet add_ew c16
\newcommand{\rhabEnetAddEwClosePct}{$-$1.5}  % QLIKE change % vs har
\newcommand{\rhabEnetAddEwCloseDm}{$-$1.37}  % DM vs the reference arm
\newcommand{\rhabEnetAddEwCloseP}{0.170}  % p-value
\newcommand{\rhabEnetAddEwCloseLateQlike}{0.1120}  % QLIKE enet add_ew c16_2020
\newcommand{\rhabEnetAddEwCloseLateMse}{4.56}  % MSE x 1e10 enet add_ew c16_2020
\newcommand{\rhabEnetAddEwCloseLatePct}{$-$1.5}  % QLIKE change % vs har
\newcommand{\rhabEnetAddEwCloseLateDm}{$-$0.90}  % DM vs the reference arm
\newcommand{\rhabEnetAddEwCloseLateP}{0.369}  % p-value
\newcommand{\rhabEnetAddFomcAllQlike}{0.1847}  % QLIKE enet add_fomc all
\newcommand{\rhabEnetAddFomcAllMse}{1.13}  % MSE x 1e10 enet add_fomc all
\newcommand{\rhabEnetAddFomcAllPct}{$-$2.4}  % QLIKE change % vs har
\newcommand{\rhabEnetAddFomcAllDm}{$-$5.85}  % DM vs the reference arm
\newcommand{\rhabEnetAddFomcAllP}{\ensuremath{<}0.001}  % p-value
\newcommand{\rhabEnetAddFomcCloseQlike}{0.1274}  % QLIKE enet add_fomc c16
\newcommand{\rhabEnetAddFomcCloseMse}{1.85}  % MSE x 1e10 enet add_fomc c16
\newcommand{\rhabEnetAddFomcClosePct}{0.4}  % QLIKE change % vs har
\newcommand{\rhabEnetAddFomcCloseDm}{1.66}  % DM vs the reference arm
\newcommand{\rhabEnetAddFomcCloseP}{0.097}  % p-value
\newcommand{\rhabEnetAddFomcCloseLateQlike}{0.1136}  % QLIKE enet add_fomc c16_2020
\newcommand{\rhabEnetAddFomcCloseLateMse}{3.96}  % MSE x 1e10 enet add_fomc c16_2020
\newcommand{\rhabEnetAddFomcCloseLatePct}{0.0}  % QLIKE change % vs har
\newcommand{\rhabEnetAddFomcCloseLateDm}{$-$0.09}  % DM vs the reference arm
\newcommand{\rhabEnetAddFomcCloseLateP}{0.927}  % p-value
\newcommand{\rhabEnetAddIvAllQlike}{0.1894}  % QLIKE enet add_iv all
\newcommand{\rhabEnetAddIvAllMse}{1.32}  % MSE x 1e10 enet add_iv all
\newcommand{\rhabEnetAddIvAllPct}{0.1}  % QLIKE change % vs har
\newcommand{\rhabEnetAddIvAllDm}{0.23}  % DM vs the reference arm
\newcommand{\rhabEnetAddIvAllP}{0.817}  % p-value
\newcommand{\rhabEnetAddIvCloseQlike}{0.1272}  % QLIKE enet add_iv c16
\newcommand{\rhabEnetAddIvCloseMse}{1.76}  % MSE x 1e10 enet add_iv c16
\newcommand{\rhabEnetAddIvClosePct}{0.3}  % QLIKE change % vs har
\newcommand{\rhabEnetAddIvCloseDm}{0.28}  % DM vs the reference arm
\newcommand{\rhabEnetAddIvCloseP}{0.780}  % p-value
\newcommand{\rhabEnetAddIvCloseLateQlike}{0.1152}  % QLIKE enet add_iv c16_2020
\newcommand{\rhabEnetAddIvCloseLateMse}{3.75}  % MSE x 1e10 enet add_iv c16_2020
\newcommand{\rhabEnetAddIvCloseLatePct}{1.4}  % QLIKE change % vs har
\newcommand{\rhabEnetAddIvCloseLateDm}{0.70}  % DM vs the reference arm
\newcommand{\rhabEnetAddIvCloseLateP}{0.485}  % p-value
\newcommand{\rhabEnetAddLiqAllQlike}{0.1881}  % QLIKE enet add_liq all
\newcommand{\rhabEnetAddLiqAllMse}{1.16}  % MSE x 1e10 enet add_liq all
\newcommand{\rhabEnetAddLiqAllPct}{$-$0.6}  % QLIKE change % vs har
\newcommand{\rhabEnetAddLiqAllDm}{$-$1.54}  % DM vs the reference arm
\newcommand{\rhabEnetAddLiqAllP}{0.124}  % p-value
\newcommand{\rhabEnetAddLiqCloseQlike}{0.1299}  % QLIKE enet add_liq c16
\newcommand{\rhabEnetAddLiqCloseMse}{1.99}  % MSE x 1e10 enet add_liq c16
\newcommand{\rhabEnetAddLiqClosePct}{2.4}  % QLIKE change % vs har
\newcommand{\rhabEnetAddLiqCloseDm}{1.77}  % DM vs the reference arm
\newcommand{\rhabEnetAddLiqCloseP}{0.077}  % p-value
\newcommand{\rhabEnetAddLiqCloseLateQlike}{0.1192}  % QLIKE enet add_liq c16_2020
\newcommand{\rhabEnetAddLiqCloseLateMse}{4.29}  % MSE x 1e10 enet add_liq c16_2020
\newcommand{\rhabEnetAddLiqCloseLatePct}{4.9}  % QLIKE change % vs har
\newcommand{\rhabEnetAddLiqCloseLateDm}{1.90}  % DM vs the reference arm
\newcommand{\rhabEnetAddLiqCloseLateP}{0.058}  % p-value
\newcommand{\rhabEnetAddRmAllQlike}{0.1842}  % QLIKE enet add_rm all
\newcommand{\rhabEnetAddRmAllMse}{1.72}  % MSE x 1e10 enet add_rm all
\newcommand{\rhabEnetAddRmAllPct}{$-$2.7}  % QLIKE change % vs har
\newcommand{\rhabEnetAddRmAllDm}{$-$4.25}  % DM vs the reference arm
\newcommand{\rhabEnetAddRmAllP}{\ensuremath{<}0.001}  % p-value
\newcommand{\rhabEnetAddRmCloseQlike}{0.1223}  % QLIKE enet add_rm c16
\newcommand{\rhabEnetAddRmCloseMse}{2.92}  % MSE x 1e10 enet add_rm c16
\newcommand{\rhabEnetAddRmClosePct}{$-$3.6}  % QLIKE change % vs har
\newcommand{\rhabEnetAddRmCloseDm}{$-$2.25}  % DM vs the reference arm
\newcommand{\rhabEnetAddRmCloseP}{0.024}  % p-value
\newcommand{\rhabEnetAddRmCloseLateQlike}{0.1147}  % QLIKE enet add_rm c16_2020
\newcommand{\rhabEnetAddRmCloseLateMse}{6.56}  % MSE x 1e10 enet add_rm c16_2020
\newcommand{\rhabEnetAddRmCloseLatePct}{1.0}  % QLIKE change % vs har
\newcommand{\rhabEnetAddRmCloseLateDm}{0.33}  % DM vs the reference arm
\newcommand{\rhabEnetAddRmCloseLateP}{0.741}  % p-value
\newcommand{\rhabEnetAddSentAllQlike}{0.1893}  % QLIKE enet add_sent all
\newcommand{\rhabEnetAddSentAllMse}{1.11}  % MSE x 1e10 enet add_sent all
\newcommand{\rhabEnetAddSentAllPct}{0.0}  % QLIKE change % vs har
\newcommand{\rhabEnetAddSentAllDm}{0.04}  % DM vs the reference arm
\newcommand{\rhabEnetAddSentAllP}{0.971}  % p-value
\newcommand{\rhabEnetAddSentCloseQlike}{0.1300}  % QLIKE enet add_sent c16
\newcommand{\rhabEnetAddSentCloseMse}{1.98}  % MSE x 1e10 enet add_sent c16
\newcommand{\rhabEnetAddSentClosePct}{2.5}  % QLIKE change % vs har
\newcommand{\rhabEnetAddSentCloseDm}{3.62}  % DM vs the reference arm
\newcommand{\rhabEnetAddSentCloseP}{\ensuremath{<}0.001}  % p-value
\newcommand{\rhabEnetAddSentCloseLateQlike}{0.1194}  % QLIKE enet add_sent c16_2020
\newcommand{\rhabEnetAddSentCloseLateMse}{4.27}  % MSE x 1e10 enet add_sent c16_2020
\newcommand{\rhabEnetAddSentCloseLatePct}{5.0}  % QLIKE change % vs har
\newcommand{\rhabEnetAddSentCloseLateDm}{2.88}  % DM vs the reference arm
\newcommand{\rhabEnetAddSentCloseLateP}{0.004}  % p-value
\newcommand{\rhabEnetAddVdAllQlike}{0.1900}  % QLIKE enet add_vd all
\newcommand{\rhabEnetAddVdAllMse}{1.11}  % MSE x 1e10 enet add_vd all
\newcommand{\rhabEnetAddVdAllPct}{0.4}  % QLIKE change % vs har
\newcommand{\rhabEnetAddVdAllDm}{2.92}  % DM vs the reference arm
\newcommand{\rhabEnetAddVdAllP}{0.003}  % p-value
\newcommand{\rhabEnetAddVdCloseQlike}{0.1277}  % QLIKE enet add_vd c16
\newcommand{\rhabEnetAddVdCloseMse}{1.96}  % MSE x 1e10 enet add_vd c16
\newcommand{\rhabEnetAddVdClosePct}{0.6}  % QLIKE change % vs har
\newcommand{\rhabEnetAddVdCloseDm}{1.85}  % DM vs the reference arm
\newcommand{\rhabEnetAddVdCloseP}{0.065}  % p-value
\newcommand{\rhabEnetAddVdCloseLateQlike}{0.1160}  % QLIKE enet add_vd c16_2020
\newcommand{\rhabEnetAddVdCloseLateMse}{4.21}  % MSE x 1e10 enet add_vd c16_2020
\newcommand{\rhabEnetAddVdCloseLatePct}{2.1}  % QLIKE change % vs har
\newcommand{\rhabEnetAddVdCloseLateDm}{2.57}  % DM vs the reference arm
\newcommand{\rhabEnetAddVdCloseLateP}{0.010}  % p-value
\newcommand{\rhabEnetAddVwAllQlike}{0.1872}  % QLIKE enet add_vw all
\newcommand{\rhabEnetAddVwAllMse}{1.99}  % MSE x 1e10 enet add_vw all
\newcommand{\rhabEnetAddVwAllPct}{$-$1.1}  % QLIKE change % vs har
\newcommand{\rhabEnetAddVwAllDm}{$-$2.03}  % DM vs the reference arm
\newcommand{\rhabEnetAddVwAllP}{0.042}  % p-value
\newcommand{\rhabEnetAddVwCloseQlike}{0.1256}  % QLIKE enet add_vw c16
\newcommand{\rhabEnetAddVwCloseMse}{2.37}  % MSE x 1e10 enet add_vw c16
\newcommand{\rhabEnetAddVwClosePct}{$-$1.0}  % QLIKE change % vs har
\newcommand{\rhabEnetAddVwCloseDm}{$-$0.75}  % DM vs the reference arm
\newcommand{\rhabEnetAddVwCloseP}{0.454}  % p-value
\newcommand{\rhabEnetAddVwCloseLateQlike}{0.1134}  % QLIKE enet add_vw c16_2020
\newcommand{\rhabEnetAddVwCloseLateMse}{5.22}  % MSE x 1e10 enet add_vw c16_2020
\newcommand{\rhabEnetAddVwCloseLatePct}{$-$0.2}  % QLIKE change % vs har
\newcommand{\rhabEnetAddVwCloseLateDm}{$-$0.13}  % DM vs the reference arm
\newcommand{\rhabEnetAddVwCloseLateP}{0.897}  % p-value
\newcommand{\rhabEnetFullAllQlike}{0.1778}  % QLIKE enet full all
\newcommand{\rhabEnetFullAllMse}{2.43}  % MSE x 1e10 enet full all
\newcommand{\rhabEnetFullAllPct}{$-$6.0}  % QLIKE change % vs har
\newcommand{\rhabEnetFullAllDm}{$-$6.98}  % DM vs the reference arm
\newcommand{\rhabEnetFullAllP}{\ensuremath{<}0.001}  % p-value
\newcommand{\rhabEnetFullCloseQlike}{0.1262}  % QLIKE enet full c16
\newcommand{\rhabEnetFullCloseMse}{3.94}  % MSE x 1e10 enet full c16
\newcommand{\rhabEnetFullClosePct}{$-$0.5}  % QLIKE change % vs har
\newcommand{\rhabEnetFullCloseDm}{$-$0.21}  % DM vs the reference arm
\newcommand{\rhabEnetFullCloseP}{0.832}  % p-value
\newcommand{\rhabEnetFullCloseLateQlike}{0.1234}  % QLIKE enet full c16_2020
\newcommand{\rhabEnetFullCloseLateMse}{9.01}  % MSE x 1e10 enet full c16_2020
\newcommand{\rhabEnetFullCloseLatePct}{8.5}  % QLIKE change % vs har
\newcommand{\rhabEnetFullCloseLateDm}{1.91}  % DM vs the reference arm
\newcommand{\rhabEnetFullCloseLateP}{0.056}  % p-value
\newcommand{\rhabEnetHarAllQlike}{0.1893}  % QLIKE enet har all
\newcommand{\rhabEnetHarAllMse}{1.11}  % MSE x 1e10 enet har all
\newcommand{\rhabEnetHarCloseQlike}{0.1269}  % QLIKE enet har c16
\newcommand{\rhabEnetHarCloseMse}{1.97}  % MSE x 1e10 enet har c16
\newcommand{\rhabEnetHarCloseLateQlike}{0.1137}  % QLIKE enet har c16_2020
\newcommand{\rhabEnetHarCloseLateMse}{4.24}  % MSE x 1e10 enet har c16_2020
\newcommand{\rhabEnetHarqAllQlike}{0.1893}  % QLIKE enet harq all
\newcommand{\rhabEnetHarqAllMse}{1.16}  % MSE x 1e10 enet harq all
\newcommand{\rhabEnetHarqAllPct}{0.0}  % QLIKE change % vs har
\newcommand{\rhabEnetHarqAllDm}{0.63}  % DM vs the reference arm
\newcommand{\rhabEnetHarqAllP}{0.531}  % p-value
\newcommand{\rhabEnetHarqCloseQlike}{0.1267}  % QLIKE enet harq c16
\newcommand{\rhabEnetHarqCloseMse}{1.80}  % MSE x 1e10 enet harq c16
\newcommand{\rhabEnetHarqClosePct}{$-$0.1}  % QLIKE change % vs har
\newcommand{\rhabEnetHarqCloseDm}{$-$1.31}  % DM vs the reference arm
\newcommand{\rhabEnetHarqCloseP}{0.191}  % p-value
\newcommand{\rhabEnetHarqCloseLateQlike}{0.1134}  % QLIKE enet harq c16_2020
\newcommand{\rhabEnetHarqCloseLateMse}{3.84}  % MSE x 1e10 enet harq c16_2020
\newcommand{\rhabEnetHarqCloseLatePct}{$-$0.3}  % QLIKE change % vs har
\newcommand{\rhabEnetHarqCloseLateDm}{$-$1.15}  % DM vs the reference arm
\newcommand{\rhabEnetHarqCloseLateP}{0.251}  % p-value
\newcommand{\rhabEnetLharAllQlike}{0.1841}  % QLIKE enet lhar all
\newcommand{\rhabEnetLharAllMse}{1.71}  % MSE x 1e10 enet lhar all
\newcommand{\rhabEnetLharAllPct}{$-$2.7}  % QLIKE change % vs har
\newcommand{\rhabEnetLharAllDm}{$-$6.66}  % DM vs the reference arm
\newcommand{\rhabEnetLharAllP}{\ensuremath{<}0.001}  % p-value
\newcommand{\rhabEnetLharCloseQlike}{0.1217}  % QLIKE enet lhar c16
\newcommand{\rhabEnetLharCloseMse}{2.76}  % MSE x 1e10 enet lhar c16
\newcommand{\rhabEnetLharClosePct}{$-$4.1}  % QLIKE change % vs har
\newcommand{\rhabEnetLharCloseDm}{$-$3.16}  % DM vs the reference arm
\newcommand{\rhabEnetLharCloseP}{0.002}  % p-value
\newcommand{\rhabEnetLharCloseLateQlike}{0.1124}  % QLIKE enet lhar c16_2020
\newcommand{\rhabEnetLharCloseLateMse}{6.19}  % MSE x 1e10 enet lhar c16_2020
\newcommand{\rhabEnetLharCloseLatePct}{$-$1.1}  % QLIKE change % vs har
\newcommand{\rhabEnetLharCloseLateDm}{$-$0.45}  % DM vs the reference arm
\newcommand{\rhabEnetLharCloseLateP}{0.656}  % p-value
\newcommand{\rhabEnetLoEwAllQlike}{0.1776}  % QLIKE enet lo_ew all
\newcommand{\rhabEnetLoEwAllMse}{2.54}  % MSE x 1e10 enet lo_ew all
\newcommand{\rhabEnetLoEwAllPct}{$-$0.1}  % QLIKE change % vs full
\newcommand{\rhabEnetLoEwAllDm}{$-$1.41}  % DM vs the reference arm
\newcommand{\rhabEnetLoEwAllP}{0.157}  % p-value
\newcommand{\rhabEnetLoEwCloseQlike}{0.1267}  % QLIKE enet lo_ew c16
\newcommand{\rhabEnetLoEwCloseMse}{4.02}  % MSE x 1e10 enet lo_ew c16
\newcommand{\rhabEnetLoEwClosePct}{0.3}  % QLIKE change % vs full
\newcommand{\rhabEnetLoEwCloseDm}{1.18}  % DM vs the reference arm
\newcommand{\rhabEnetLoEwCloseP}{0.238}  % p-value
\newcommand{\rhabEnetLoEwCloseLateQlike}{0.1235}  % QLIKE enet lo_ew c16_2020
\newcommand{\rhabEnetLoEwCloseLateMse}{9.21}  % MSE x 1e10 enet lo_ew c16_2020
\newcommand{\rhabEnetLoEwCloseLatePct}{0.1}  % QLIKE change % vs full
\newcommand{\rhabEnetLoEwCloseLateDm}{0.23}  % DM vs the reference arm
\newcommand{\rhabEnetLoEwCloseLateP}{0.822}  % p-value
\newcommand{\rhabEnetLoFomcAllQlike}{0.1820}  % QLIKE enet lo_fomc all
\newcommand{\rhabEnetLoFomcAllMse}{2.42}  % MSE x 1e10 enet lo_fomc all
\newcommand{\rhabEnetLoFomcAllPct}{2.3}  % QLIKE change % vs full
\newcommand{\rhabEnetLoFomcAllDm}{5.74}  % DM vs the reference arm
\newcommand{\rhabEnetLoFomcAllP}{\ensuremath{<}0.001}  % p-value
\newcommand{\rhabEnetLoFomcCloseQlike}{0.1261}  % QLIKE enet lo_fomc c16
\newcommand{\rhabEnetLoFomcCloseMse}{3.99}  % MSE x 1e10 enet lo_fomc c16
\newcommand{\rhabEnetLoFomcClosePct}{$-$0.1}  % QLIKE change % vs full
\newcommand{\rhabEnetLoFomcCloseDm}{$-$0.51}  % DM vs the reference arm
\newcommand{\rhabEnetLoFomcCloseP}{0.607}  % p-value
\newcommand{\rhabEnetLoFomcCloseLateQlike}{0.1234}  % QLIKE enet lo_fomc c16_2020
\newcommand{\rhabEnetLoFomcCloseLateMse}{9.13}  % MSE x 1e10 enet lo_fomc c16_2020
\newcommand{\rhabEnetLoFomcCloseLatePct}{0.1}  % QLIKE change % vs full
\newcommand{\rhabEnetLoFomcCloseLateDm}{0.13}  % DM vs the reference arm
\newcommand{\rhabEnetLoFomcCloseLateP}{0.894}  % p-value
\newcommand{\rhabEnetLoIvAllQlike}{0.1778}  % QLIKE enet lo_iv all
\newcommand{\rhabEnetLoIvAllMse}{2.49}  % MSE x 1e10 enet lo_iv all
\newcommand{\rhabEnetLoIvAllPct}{0.0}  % QLIKE change % vs full
\newcommand{\rhabEnetLoIvAllDm}{$-$0.13}  % DM vs the reference arm
\newcommand{\rhabEnetLoIvAllP}{0.898}  % p-value
\newcommand{\rhabEnetLoIvCloseQlike}{0.1243}  % QLIKE enet lo_iv c16
\newcommand{\rhabEnetLoIvCloseMse}{4.03}  % MSE x 1e10 enet lo_iv c16
\newcommand{\rhabEnetLoIvClosePct}{$-$1.5}  % QLIKE change % vs full
\newcommand{\rhabEnetLoIvCloseDm}{$-$1.51}  % DM vs the reference arm
\newcommand{\rhabEnetLoIvCloseP}{0.132}  % p-value
\newcommand{\rhabEnetLoIvCloseLateQlike}{0.1193}  % QLIKE enet lo_iv c16_2020
\newcommand{\rhabEnetLoIvCloseLateMse}{9.21}  % MSE x 1e10 enet lo_iv c16_2020
\newcommand{\rhabEnetLoIvCloseLatePct}{$-$3.3}  % QLIKE change % vs full
\newcommand{\rhabEnetLoIvCloseLateDm}{$-$1.72}  % DM vs the reference arm
\newcommand{\rhabEnetLoIvCloseLateP}{0.086}  % p-value
\newcommand{\rhabEnetLoLiqAllQlike}{0.1779}  % QLIKE enet lo_liq all
\newcommand{\rhabEnetLoLiqAllMse}{2.43}  % MSE x 1e10 enet lo_liq all
\newcommand{\rhabEnetLoLiqAllPct}{0.1}  % QLIKE change % vs full
\newcommand{\rhabEnetLoLiqAllDm}{0.35}  % DM vs the reference arm
\newcommand{\rhabEnetLoLiqAllP}{0.724}  % p-value
\newcommand{\rhabEnetLoLiqCloseQlike}{0.1251}  % QLIKE enet lo_liq c16
\newcommand{\rhabEnetLoLiqCloseMse}{4.29}  % MSE x 1e10 enet lo_liq c16
\newcommand{\rhabEnetLoLiqClosePct}{$-$0.9}  % QLIKE change % vs full
\newcommand{\rhabEnetLoLiqCloseDm}{$-$1.54}  % DM vs the reference arm
\newcommand{\rhabEnetLoLiqCloseP}{0.123}  % p-value
\newcommand{\rhabEnetLoLiqCloseLateQlike}{0.1223}  % QLIKE enet lo_liq c16_2020
\newcommand{\rhabEnetLoLiqCloseLateMse}{9.84}  % MSE x 1e10 enet lo_liq c16_2020
\newcommand{\rhabEnetLoLiqCloseLatePct}{$-$0.9}  % QLIKE change % vs full
\newcommand{\rhabEnetLoLiqCloseLateDm}{$-$0.92}  % DM vs the reference arm
\newcommand{\rhabEnetLoLiqCloseLateP}{0.358}  % p-value
\newcommand{\rhabEnetLoRmAllQlike}{0.1813}  % QLIKE enet lo_rm all
\newcommand{\rhabEnetLoRmAllMse}{2.00}  % MSE x 1e10 enet lo_rm all
\newcommand{\rhabEnetLoRmAllPct}{1.9}  % QLIKE change % vs full
\newcommand{\rhabEnetLoRmAllDm}{4.38}  % DM vs the reference arm
\newcommand{\rhabEnetLoRmAllP}{\ensuremath{<}0.001}  % p-value
\newcommand{\rhabEnetLoRmCloseQlike}{0.1285}  % QLIKE enet lo_rm c16
\newcommand{\rhabEnetLoRmCloseMse}{2.35}  % MSE x 1e10 enet lo_rm c16
\newcommand{\rhabEnetLoRmClosePct}{1.8}  % QLIKE change % vs full
\newcommand{\rhabEnetLoRmCloseDm}{1.31}  % DM vs the reference arm
\newcommand{\rhabEnetLoRmCloseP}{0.190}  % p-value
\newcommand{\rhabEnetLoRmCloseLateQlike}{0.1217}  % QLIKE enet lo_rm c16_2020
\newcommand{\rhabEnetLoRmCloseLateMse}{5.18}  % MSE x 1e10 enet lo_rm c16_2020
\newcommand{\rhabEnetLoRmCloseLatePct}{$-$1.3}  % QLIKE change % vs full
\newcommand{\rhabEnetLoRmCloseLateDm}{$-$0.62}  % DM vs the reference arm
\newcommand{\rhabEnetLoRmCloseLateP}{0.538}  % p-value
\newcommand{\rhabEnetLoSentAllQlike}{0.1782}  % QLIKE enet lo_sent all
\newcommand{\rhabEnetLoSentAllMse}{2.43}  % MSE x 1e10 enet lo_sent all
\newcommand{\rhabEnetLoSentAllPct}{0.2}  % QLIKE change % vs full
\newcommand{\rhabEnetLoSentAllDm}{1.49}  % DM vs the reference arm
\newcommand{\rhabEnetLoSentAllP}{0.136}  % p-value
\newcommand{\rhabEnetLoSentCloseQlike}{0.1253}  % QLIKE enet lo_sent c16
\newcommand{\rhabEnetLoSentCloseMse}{3.94}  % MSE x 1e10 enet lo_sent c16
\newcommand{\rhabEnetLoSentClosePct}{$-$0.7}  % QLIKE change % vs full
\newcommand{\rhabEnetLoSentCloseDm}{$-$2.12}  % DM vs the reference arm
\newcommand{\rhabEnetLoSentCloseP}{0.034}  % p-value
\newcommand{\rhabEnetLoSentCloseLateQlike}{0.1212}  % QLIKE enet lo_sent c16_2020
\newcommand{\rhabEnetLoSentCloseLateMse}{9.01}  % MSE x 1e10 enet lo_sent c16_2020
\newcommand{\rhabEnetLoSentCloseLatePct}{$-$1.7}  % QLIKE change % vs full
\newcommand{\rhabEnetLoSentCloseLateDm}{$-$2.21}  % DM vs the reference arm
\newcommand{\rhabEnetLoSentCloseLateP}{0.027}  % p-value
\newcommand{\rhabEnetLoVdAllQlike}{0.1777}  % QLIKE enet lo_vd all
\newcommand{\rhabEnetLoVdAllMse}{2.43}  % MSE x 1e10 enet lo_vd all
\newcommand{\rhabEnetLoVdAllPct}{$-$0.1}  % QLIKE change % vs full
\newcommand{\rhabEnetLoVdAllDm}{$-$1.13}  % DM vs the reference arm
\newcommand{\rhabEnetLoVdAllP}{0.260}  % p-value
\newcommand{\rhabEnetLoVdCloseQlike}{0.1252}  % QLIKE enet lo_vd c16
\newcommand{\rhabEnetLoVdCloseMse}{3.94}  % MSE x 1e10 enet lo_vd c16
\newcommand{\rhabEnetLoVdClosePct}{$-$0.8}  % QLIKE change % vs full
\newcommand{\rhabEnetLoVdCloseDm}{$-$1.94}  % DM vs the reference arm
\newcommand{\rhabEnetLoVdCloseP}{0.053}  % p-value
\newcommand{\rhabEnetLoVdCloseLateQlike}{0.1207}  % QLIKE enet lo_vd c16_2020
\newcommand{\rhabEnetLoVdCloseLateMse}{9.01}  % MSE x 1e10 enet lo_vd c16_2020
\newcommand{\rhabEnetLoVdCloseLatePct}{$-$2.1}  % QLIKE change % vs full
\newcommand{\rhabEnetLoVdCloseLateDm}{$-$2.18}  % DM vs the reference arm
\newcommand{\rhabEnetLoVdCloseLateP}{0.029}  % p-value
\newcommand{\rhabEnetLoVwAllQlike}{0.1777}  % QLIKE enet lo_vw all
\newcommand{\rhabEnetLoVwAllMse}{2.08}  % MSE x 1e10 enet lo_vw all
\newcommand{\rhabEnetLoVwAllPct}{$-$0.1}  % QLIKE change % vs full
\newcommand{\rhabEnetLoVwAllDm}{$-$0.65}  % DM vs the reference arm
\newcommand{\rhabEnetLoVwAllP}{0.516}  % p-value
\newcommand{\rhabEnetLoVwCloseQlike}{0.1254}  % QLIKE enet lo_vw c16
\newcommand{\rhabEnetLoVwCloseMse}{3.77}  % MSE x 1e10 enet lo_vw c16
\newcommand{\rhabEnetLoVwClosePct}{$-$0.7}  % QLIKE change % vs full
\newcommand{\rhabEnetLoVwCloseDm}{$-$1.66}  % DM vs the reference arm
\newcommand{\rhabEnetLoVwCloseP}{0.097}  % p-value
\newcommand{\rhabEnetLoVwCloseLateQlike}{0.1218}  % QLIKE enet lo_vw c16_2020
\newcommand{\rhabEnetLoVwCloseLateMse}{8.59}  % MSE x 1e10 enet lo_vw c16_2020
\newcommand{\rhabEnetLoVwCloseLatePct}{$-$1.3}  % QLIKE change % vs full
\newcommand{\rhabEnetLoVwCloseLateDm}{$-$2.65}  % DM vs the reference arm
\newcommand{\rhabEnetLoVwCloseLateP}{0.008}  % p-value
\newcommand{\rhabEnetSharAllQlike}{0.1877}  % QLIKE enet shar all
\newcommand{\rhabEnetSharAllMse}{1.38}  % MSE x 1e10 enet shar all
\newcommand{\rhabEnetSharAllPct}{$-$0.9}  % QLIKE change % vs har
\newcommand{\rhabEnetSharAllDm}{$-$3.68}  % DM vs the reference arm
\newcommand{\rhabEnetSharAllP}{\ensuremath{<}0.001}  % p-value
\newcommand{\rhabEnetSharCloseQlike}{0.1265}  % QLIKE enet shar c16
\newcommand{\rhabEnetSharCloseMse}{1.83}  % MSE x 1e10 enet shar c16
\newcommand{\rhabEnetSharClosePct}{$-$0.3}  % QLIKE change % vs har
\newcommand{\rhabEnetSharCloseDm}{$-$0.54}  % DM vs the reference arm
\newcommand{\rhabEnetSharCloseP}{0.586}  % p-value
\newcommand{\rhabEnetSharCloseLateQlike}{0.1141}  % QLIKE enet shar c16_2020
\newcommand{\rhabEnetSharCloseLateMse}{3.92}  % MSE x 1e10 enet shar c16_2020
\newcommand{\rhabEnetSharCloseLatePct}{0.3}  % QLIKE change % vs har
\newcommand{\rhabEnetSharCloseLateDm}{0.34}  % DM vs the reference arm
\newcommand{\rhabEnetSharCloseLateP}{0.731}  % p-value
\newcommand{\rhabEnetSharqlAllQlike}{0.1840}  % QLIKE enet sharql all
\newcommand{\rhabEnetSharqlAllMse}{1.69}  % MSE x 1e10 enet sharql all
\newcommand{\rhabEnetSharqlAllPct}{$-$2.8}  % QLIKE change % vs har
\newcommand{\rhabEnetSharqlAllDm}{$-$6.95}  % DM vs the reference arm
\newcommand{\rhabEnetSharqlAllP}{\ensuremath{<}0.001}  % p-value
\newcommand{\rhabEnetSharqlCloseQlike}{0.1216}  % QLIKE enet sharql c16
\newcommand{\rhabEnetSharqlCloseMse}{2.64}  % MSE x 1e10 enet sharql c16
\newcommand{\rhabEnetSharqlClosePct}{$-$4.1}  % QLIKE change % vs har
\newcommand{\rhabEnetSharqlCloseDm}{$-$3.46}  % DM vs the reference arm
\newcommand{\rhabEnetSharqlCloseP}{\ensuremath{<}0.001}  % p-value
\newcommand{\rhabEnetSharqlCloseLateQlike}{0.1121}  % QLIKE enet sharql c16_2020
\newcommand{\rhabEnetSharqlCloseLateMse}{5.90}  % MSE x 1e10 enet sharql c16_2020
\newcommand{\rhabEnetSharqlCloseLatePct}{$-$1.4}  % QLIKE change % vs har
\newcommand{\rhabEnetSharqlCloseLateDm}{$-$0.64}  % DM vs the reference arm
\newcommand{\rhabEnetSharqlCloseLateP}{0.524}  % p-value
\newcommand{\rhabLassoAddEwAllQlike}{0.1877}  % QLIKE lasso add_ew all
\newcommand{\rhabLassoAddEwAllMse}{1.47}  % MSE x 1e10 lasso add_ew all
\newcommand{\rhabLassoAddEwAllPct}{$-$0.8}  % QLIKE change % vs har
\newcommand{\rhabLassoAddEwAllDm}{$-$2.01}  % DM vs the reference arm
\newcommand{\rhabLassoAddEwAllP}{0.045}  % p-value
\newcommand{\rhabLassoAddEwCloseQlike}{0.1257}  % QLIKE lasso add_ew c16
\newcommand{\rhabLassoAddEwCloseMse}{2.13}  % MSE x 1e10 lasso add_ew c16
\newcommand{\rhabLassoAddEwClosePct}{$-$1.0}  % QLIKE change % vs har
\newcommand{\rhabLassoAddEwCloseDm}{$-$0.81}  % DM vs the reference arm
\newcommand{\rhabLassoAddEwCloseP}{0.417}  % p-value
\newcommand{\rhabLassoAddEwCloseLateQlike}{0.1123}  % QLIKE lasso add_ew c16_2020
\newcommand{\rhabLassoAddEwCloseLateMse}{4.65}  % MSE x 1e10 lasso add_ew c16_2020
\newcommand{\rhabLassoAddEwCloseLatePct}{$-$1.2}  % QLIKE change % vs har
\newcommand{\rhabLassoAddEwCloseLateDm}{$-$0.70}  % DM vs the reference arm
\newcommand{\rhabLassoAddEwCloseLateP}{0.485}  % p-value
\newcommand{\rhabLassoAddFomcAllQlike}{0.1820}  % QLIKE lasso add_fomc all
\newcommand{\rhabLassoAddFomcAllMse}{1.13}  % MSE x 1e10 lasso add_fomc all
\newcommand{\rhabLassoAddFomcAllPct}{$-$3.8}  % QLIKE change % vs har
\newcommand{\rhabLassoAddFomcAllDm}{$-$5.76}  % DM vs the reference arm
\newcommand{\rhabLassoAddFomcAllP}{\ensuremath{<}0.001}  % p-value
\newcommand{\rhabLassoAddFomcCloseQlike}{0.1296}  % QLIKE lasso add_fomc c16
\newcommand{\rhabLassoAddFomcCloseMse}{1.83}  % MSE x 1e10 lasso add_fomc c16
\newcommand{\rhabLassoAddFomcClosePct}{2.2}  % QLIKE change % vs har
\newcommand{\rhabLassoAddFomcCloseDm}{4.33}  % DM vs the reference arm
\newcommand{\rhabLassoAddFomcCloseP}{\ensuremath{<}0.001}  % p-value
\newcommand{\rhabLassoAddFomcCloseLateQlike}{0.1135}  % QLIKE lasso add_fomc c16_2020
\newcommand{\rhabLassoAddFomcCloseLateMse}{3.92}  % MSE x 1e10 lasso add_fomc c16_2020
\newcommand{\rhabLassoAddFomcCloseLatePct}{$-$0.1}  % QLIKE change % vs har
\newcommand{\rhabLassoAddFomcCloseLateDm}{$-$0.16}  % DM vs the reference arm
\newcommand{\rhabLassoAddFomcCloseLateP}{0.875}  % p-value
\newcommand{\rhabLassoAddIvAllQlike}{0.1887}  % QLIKE lasso add_iv all
\newcommand{\rhabLassoAddIvAllMse}{1.91}  % MSE x 1e10 lasso add_iv all
\newcommand{\rhabLassoAddIvAllPct}{$-$0.3}  % QLIKE change % vs har
\newcommand{\rhabLassoAddIvAllDm}{$-$0.79}  % DM vs the reference arm
\newcommand{\rhabLassoAddIvAllP}{0.430}  % p-value
\newcommand{\rhabLassoAddIvCloseQlike}{0.1263}  % QLIKE lasso add_iv c16
\newcommand{\rhabLassoAddIvCloseMse}{1.82}  % MSE x 1e10 lasso add_iv c16
\newcommand{\rhabLassoAddIvClosePct}{$-$0.4}  % QLIKE change % vs har
\newcommand{\rhabLassoAddIvCloseDm}{$-$0.38}  % DM vs the reference arm
\newcommand{\rhabLassoAddIvCloseP}{0.705}  % p-value
\newcommand{\rhabLassoAddIvCloseLateQlike}{0.1144}  % QLIKE lasso add_iv c16_2020
\newcommand{\rhabLassoAddIvCloseLateMse}{3.92}  % MSE x 1e10 lasso add_iv c16_2020
\newcommand{\rhabLassoAddIvCloseLatePct}{0.6}  % QLIKE change % vs har
\newcommand{\rhabLassoAddIvCloseLateDm}{0.29}  % DM vs the reference arm
\newcommand{\rhabLassoAddIvCloseLateP}{0.768}  % p-value
\newcommand{\rhabLassoAddLiqAllQlike}{0.1888}  % QLIKE lasso add_liq all
\newcommand{\rhabLassoAddLiqAllMse}{1.17}  % MSE x 1e10 lasso add_liq all
\newcommand{\rhabLassoAddLiqAllPct}{$-$0.3}  % QLIKE change % vs har
\newcommand{\rhabLassoAddLiqAllDm}{$-$0.55}  % DM vs the reference arm
\newcommand{\rhabLassoAddLiqAllP}{0.583}  % p-value
\newcommand{\rhabLassoAddLiqCloseQlike}{0.1305}  % QLIKE lasso add_liq c16
\newcommand{\rhabLassoAddLiqCloseMse}{2.02}  % MSE x 1e10 lasso add_liq c16
\newcommand{\rhabLassoAddLiqClosePct}{2.9}  % QLIKE change % vs har
\newcommand{\rhabLassoAddLiqCloseDm}{1.86}  % DM vs the reference arm
\newcommand{\rhabLassoAddLiqCloseP}{0.063}  % p-value
\newcommand{\rhabLassoAddLiqCloseLateQlike}{0.1206}  % QLIKE lasso add_liq c16_2020
\newcommand{\rhabLassoAddLiqCloseLateMse}{4.37}  % MSE x 1e10 lasso add_liq c16_2020
\newcommand{\rhabLassoAddLiqCloseLatePct}{6.1}  % QLIKE change % vs har
\newcommand{\rhabLassoAddLiqCloseLateDm}{2.00}  % DM vs the reference arm
\newcommand{\rhabLassoAddLiqCloseLateP}{0.045}  % p-value
\newcommand{\rhabLassoAddRmAllQlike}{0.1852}  % QLIKE lasso add_rm all
\newcommand{\rhabLassoAddRmAllMse}{1.81}  % MSE x 1e10 lasso add_rm all
\newcommand{\rhabLassoAddRmAllPct}{$-$2.1}  % QLIKE change % vs har
\newcommand{\rhabLassoAddRmAllDm}{$-$2.64}  % DM vs the reference arm
\newcommand{\rhabLassoAddRmAllP}{0.008}  % p-value
\newcommand{\rhabLassoAddRmCloseQlike}{0.1234}  % QLIKE lasso add_rm c16
\newcommand{\rhabLassoAddRmCloseMse}{3.11}  % MSE x 1e10 lasso add_rm c16
\newcommand{\rhabLassoAddRmClosePct}{$-$2.8}  % QLIKE change % vs har
\newcommand{\rhabLassoAddRmCloseDm}{$-$1.41}  % DM vs the reference arm
\newcommand{\rhabLassoAddRmCloseP}{0.160}  % p-value
\newcommand{\rhabLassoAddRmCloseLateQlike}{0.1176}  % QLIKE lasso add_rm c16_2020
\newcommand{\rhabLassoAddRmCloseLateMse}{7.01}  % MSE x 1e10 lasso add_rm c16_2020
\newcommand{\rhabLassoAddRmCloseLatePct}{3.5}  % QLIKE change % vs har
\newcommand{\rhabLassoAddRmCloseLateDm}{0.90}  % DM vs the reference arm
\newcommand{\rhabLassoAddRmCloseLateP}{0.368}  % p-value
\newcommand{\rhabLassoAddSentAllQlike}{0.1895}  % QLIKE lasso add_sent all
\newcommand{\rhabLassoAddSentAllMse}{1.11}  % MSE x 1e10 lasso add_sent all
\newcommand{\rhabLassoAddSentAllPct}{0.1}  % QLIKE change % vs har
\newcommand{\rhabLassoAddSentAllDm}{0.60}  % DM vs the reference arm
\newcommand{\rhabLassoAddSentAllP}{0.550}  % p-value
\newcommand{\rhabLassoAddSentCloseQlike}{0.1310}  % QLIKE lasso add_sent c16
\newcommand{\rhabLassoAddSentCloseMse}{1.98}  % MSE x 1e10 lasso add_sent c16
\newcommand{\rhabLassoAddSentClosePct}{3.2}  % QLIKE change % vs har
\newcommand{\rhabLassoAddSentCloseDm}{4.34}  % DM vs the reference arm
\newcommand{\rhabLassoAddSentCloseP}{\ensuremath{<}0.001}  % p-value
\newcommand{\rhabLassoAddSentCloseLateQlike}{0.1214}  % QLIKE lasso add_sent c16_2020
\newcommand{\rhabLassoAddSentCloseLateMse}{4.27}  % MSE x 1e10 lasso add_sent c16_2020
\newcommand{\rhabLassoAddSentCloseLatePct}{6.8}  % QLIKE change % vs har
\newcommand{\rhabLassoAddSentCloseLateDm}{3.64}  % DM vs the reference arm
\newcommand{\rhabLassoAddSentCloseLateP}{\ensuremath{<}0.001}  % p-value
\newcommand{\rhabLassoAddVdAllQlike}{0.1903}  % QLIKE lasso add_vd all
\newcommand{\rhabLassoAddVdAllMse}{1.12}  % MSE x 1e10 lasso add_vd all
\newcommand{\rhabLassoAddVdAllPct}{0.6}  % QLIKE change % vs har
\newcommand{\rhabLassoAddVdAllDm}{3.05}  % DM vs the reference arm
\newcommand{\rhabLassoAddVdAllP}{0.002}  % p-value
\newcommand{\rhabLassoAddVdCloseQlike}{0.1272}  % QLIKE lasso add_vd c16
\newcommand{\rhabLassoAddVdCloseMse}{1.95}  % MSE x 1e10 lasso add_vd c16
\newcommand{\rhabLassoAddVdClosePct}{0.3}  % QLIKE change % vs har
\newcommand{\rhabLassoAddVdCloseDm}{0.65}  % DM vs the reference arm
\newcommand{\rhabLassoAddVdCloseP}{0.516}  % p-value
\newcommand{\rhabLassoAddVdCloseLateQlike}{0.1169}  % QLIKE lasso add_vd c16_2020
\newcommand{\rhabLassoAddVdCloseLateMse}{4.20}  % MSE x 1e10 lasso add_vd c16_2020
\newcommand{\rhabLassoAddVdCloseLatePct}{2.9}  % QLIKE change % vs har
\newcommand{\rhabLassoAddVdCloseLateDm}{3.03}  % DM vs the reference arm
\newcommand{\rhabLassoAddVdCloseLateP}{0.002}  % p-value
\newcommand{\rhabLassoAddVwAllQlike}{0.1881}  % QLIKE lasso add_vw all
\newcommand{\rhabLassoAddVwAllMse}{2.04}  % MSE x 1e10 lasso add_vw all
\newcommand{\rhabLassoAddVwAllPct}{$-$0.6}  % QLIKE change % vs har
\newcommand{\rhabLassoAddVwAllDm}{$-$1.02}  % DM vs the reference arm
\newcommand{\rhabLassoAddVwAllP}{0.309}  % p-value
\newcommand{\rhabLassoAddVwCloseQlike}{0.1262}  % QLIKE lasso add_vw c16
\newcommand{\rhabLassoAddVwCloseMse}{2.38}  % MSE x 1e10 lasso add_vw c16
\newcommand{\rhabLassoAddVwClosePct}{$-$0.5}  % QLIKE change % vs har
\newcommand{\rhabLassoAddVwCloseDm}{$-$0.35}  % DM vs the reference arm
\newcommand{\rhabLassoAddVwCloseP}{0.726}  % p-value
\newcommand{\rhabLassoAddVwCloseLateQlike}{0.1136}  % QLIKE lasso add_vw c16_2020
\newcommand{\rhabLassoAddVwCloseLateMse}{5.24}  % MSE x 1e10 lasso add_vw c16_2020
\newcommand{\rhabLassoAddVwCloseLatePct}{$-$0.1}  % QLIKE change % vs har
\newcommand{\rhabLassoAddVwCloseLateDm}{$-$0.04}  % DM vs the reference arm
\newcommand{\rhabLassoAddVwCloseLateP}{0.968}  % p-value
\newcommand{\rhabLassoFullAllQlike}{0.1775}  % QLIKE lasso full all
\newcommand{\rhabLassoFullAllMse}{2.37}  % MSE x 1e10 lasso full all
\newcommand{\rhabLassoFullAllPct}{$-$6.2}  % QLIKE change % vs har
\newcommand{\rhabLassoFullAllDm}{$-$4.72}  % DM vs the reference arm
\newcommand{\rhabLassoFullAllP}{\ensuremath{<}0.001}  % p-value
\newcommand{\rhabLassoFullCloseQlike}{0.1310}  % QLIKE lasso full c16
\newcommand{\rhabLassoFullCloseMse}{3.57}  % MSE x 1e10 lasso full c16
\newcommand{\rhabLassoFullClosePct}{3.3}  % QLIKE change % vs har
\newcommand{\rhabLassoFullCloseDm}{1.04}  % DM vs the reference arm
\newcommand{\rhabLassoFullCloseP}{0.300}  % p-value
\newcommand{\rhabLassoFullCloseLateQlike}{0.1301}  % QLIKE lasso full c16_2020
\newcommand{\rhabLassoFullCloseLateMse}{8.13}  % MSE x 1e10 lasso full c16_2020
\newcommand{\rhabLassoFullCloseLatePct}{14.4}  % QLIKE change % vs har
\newcommand{\rhabLassoFullCloseLateDm}{2.51}  % DM vs the reference arm
\newcommand{\rhabLassoFullCloseLateP}{0.012}  % p-value
\newcommand{\rhabLassoHarAllQlike}{0.1893}  % QLIKE lasso har all
\newcommand{\rhabLassoHarAllMse}{1.11}  % MSE x 1e10 lasso har all
\newcommand{\rhabLassoHarCloseQlike}{0.1269}  % QLIKE lasso har c16
\newcommand{\rhabLassoHarCloseMse}{1.97}  % MSE x 1e10 lasso har c16
\newcommand{\rhabLassoHarCloseLateQlike}{0.1137}  % QLIKE lasso har c16_2020
\newcommand{\rhabLassoHarCloseLateMse}{4.24}  % MSE x 1e10 lasso har c16_2020
\newcommand{\rhabLassoHarqAllQlike}{0.1893}  % QLIKE lasso harq all
\newcommand{\rhabLassoHarqAllMse}{1.17}  % MSE x 1e10 lasso harq all
\newcommand{\rhabLassoHarqAllPct}{0.0}  % QLIKE change % vs har
\newcommand{\rhabLassoHarqAllDm}{0.61}  % DM vs the reference arm
\newcommand{\rhabLassoHarqAllP}{0.544}  % p-value
\newcommand{\rhabLassoHarqCloseQlike}{0.1267}  % QLIKE lasso harq c16
\newcommand{\rhabLassoHarqCloseMse}{1.78}  % MSE x 1e10 lasso harq c16
\newcommand{\rhabLassoHarqClosePct}{$-$0.2}  % QLIKE change % vs har
\newcommand{\rhabLassoHarqCloseDm}{$-$1.37}  % DM vs the reference arm
\newcommand{\rhabLassoHarqCloseP}{0.171}  % p-value
\newcommand{\rhabLassoHarqCloseLateQlike}{0.1133}  % QLIKE lasso harq c16_2020
\newcommand{\rhabLassoHarqCloseLateMse}{3.80}  % MSE x 1e10 lasso harq c16_2020
\newcommand{\rhabLassoHarqCloseLatePct}{$-$0.3}  % QLIKE change % vs har
\newcommand{\rhabLassoHarqCloseLateDm}{$-$1.18}  % DM vs the reference arm
\newcommand{\rhabLassoHarqCloseLateP}{0.238}  % p-value
\newcommand{\rhabLassoLharAllQlike}{0.1842}  % QLIKE lasso lhar all
\newcommand{\rhabLassoLharAllMse}{1.71}  % MSE x 1e10 lasso lhar all
\newcommand{\rhabLassoLharAllPct}{$-$2.7}  % QLIKE change % vs har
\newcommand{\rhabLassoLharAllDm}{$-$6.30}  % DM vs the reference arm
\newcommand{\rhabLassoLharAllP}{\ensuremath{<}0.001}  % p-value
\newcommand{\rhabLassoLharCloseQlike}{0.1219}  % QLIKE lasso lhar c16
\newcommand{\rhabLassoLharCloseMse}{2.80}  % MSE x 1e10 lasso lhar c16
\newcommand{\rhabLassoLharClosePct}{$-$3.9}  % QLIKE change % vs har
\newcommand{\rhabLassoLharCloseDm}{$-$2.77}  % DM vs the reference arm
\newcommand{\rhabLassoLharCloseP}{0.006}  % p-value
\newcommand{\rhabLassoLharCloseLateQlike}{0.1130}  % QLIKE lasso lhar c16_2020
\newcommand{\rhabLassoLharCloseLateMse}{6.28}  % MSE x 1e10 lasso lhar c16_2020
\newcommand{\rhabLassoLharCloseLatePct}{$-$0.6}  % QLIKE change % vs har
\newcommand{\rhabLassoLharCloseLateDm}{$-$0.21}  % DM vs the reference arm
\newcommand{\rhabLassoLharCloseLateP}{0.836}  % p-value
\newcommand{\rhabLassoLoEwAllQlike}{0.1770}  % QLIKE lasso lo_ew all
\newcommand{\rhabLassoLoEwAllMse}{2.54}  % MSE x 1e10 lasso lo_ew all
\newcommand{\rhabLassoLoEwAllPct}{$-$0.3}  % QLIKE change % vs full
\newcommand{\rhabLassoLoEwAllDm}{$-$1.46}  % DM vs the reference arm
\newcommand{\rhabLassoLoEwAllP}{0.144}  % p-value
\newcommand{\rhabLassoLoEwCloseQlike}{0.1310}  % QLIKE lasso lo_ew c16
\newcommand{\rhabLassoLoEwCloseMse}{3.82}  % MSE x 1e10 lasso lo_ew c16
\newcommand{\rhabLassoLoEwClosePct}{0.0}  % QLIKE change % vs full
\newcommand{\rhabLassoLoEwCloseDm}{0.00}  % DM vs the reference arm
\newcommand{\rhabLassoLoEwCloseP}{0.999}  % p-value
\newcommand{\rhabLassoLoEwCloseLateQlike}{0.1289}  % QLIKE lasso lo_ew c16_2020
\newcommand{\rhabLassoLoEwCloseLateMse}{8.72}  % MSE x 1e10 lasso lo_ew c16_2020
\newcommand{\rhabLassoLoEwCloseLatePct}{$-$0.9}  % QLIKE change % vs full
\newcommand{\rhabLassoLoEwCloseLateDm}{$-$0.76}  % DM vs the reference arm
\newcommand{\rhabLassoLoEwCloseLateP}{0.446}  % p-value
\newcommand{\rhabLassoLoFomcAllQlike}{0.1839}  % QLIKE lasso lo_fomc all
\newcommand{\rhabLassoLoFomcAllMse}{2.39}  % MSE x 1e10 lasso lo_fomc all
\newcommand{\rhabLassoLoFomcAllPct}{3.6}  % QLIKE change % vs full
\newcommand{\rhabLassoLoFomcAllDm}{5.10}  % DM vs the reference arm
\newcommand{\rhabLassoLoFomcAllP}{\ensuremath{<}0.001}  % p-value
\newcommand{\rhabLassoLoFomcCloseQlike}{0.1293}  % QLIKE lasso lo_fomc c16
\newcommand{\rhabLassoLoFomcCloseMse}{3.62}  % MSE x 1e10 lasso lo_fomc c16
\newcommand{\rhabLassoLoFomcClosePct}{$-$1.4}  % QLIKE change % vs full
\newcommand{\rhabLassoLoFomcCloseDm}{$-$3.21}  % DM vs the reference arm
\newcommand{\rhabLassoLoFomcCloseP}{0.001}  % p-value
\newcommand{\rhabLassoLoFomcCloseLateQlike}{0.1293}  % QLIKE lasso lo_fomc c16_2020
\newcommand{\rhabLassoLoFomcCloseLateMse}{8.25}  % MSE x 1e10 lasso lo_fomc c16_2020
\newcommand{\rhabLassoLoFomcCloseLatePct}{$-$0.6}  % QLIKE change % vs full
\newcommand{\rhabLassoLoFomcCloseLateDm}{$-$0.92}  % DM vs the reference arm
\newcommand{\rhabLassoLoFomcCloseLateP}{0.358}  % p-value
\newcommand{\rhabLassoLoIvAllQlike}{0.1778}  % QLIKE lasso lo_iv all
\newcommand{\rhabLassoLoIvAllMse}{2.43}  % MSE x 1e10 lasso lo_iv all
\newcommand{\rhabLassoLoIvAllPct}{0.2}  % QLIKE change % vs full
\newcommand{\rhabLassoLoIvAllDm}{0.34}  % DM vs the reference arm
\newcommand{\rhabLassoLoIvAllP}{0.731}  % p-value
\newcommand{\rhabLassoLoIvCloseQlike}{0.1284}  % QLIKE lasso lo_iv c16
\newcommand{\rhabLassoLoIvCloseMse}{3.90}  % MSE x 1e10 lasso lo_iv c16
\newcommand{\rhabLassoLoIvClosePct}{$-$2.0}  % QLIKE change % vs full
\newcommand{\rhabLassoLoIvCloseDm}{$-$1.66}  % DM vs the reference arm
\newcommand{\rhabLassoLoIvCloseP}{0.096}  % p-value
\newcommand{\rhabLassoLoIvCloseLateQlike}{0.1251}  % QLIKE lasso lo_iv c16_2020
\newcommand{\rhabLassoLoIvCloseLateMse}{8.92}  % MSE x 1e10 lasso lo_iv c16_2020
\newcommand{\rhabLassoLoIvCloseLatePct}{$-$3.8}  % QLIKE change % vs full
\newcommand{\rhabLassoLoIvCloseLateDm}{$-$1.77}  % DM vs the reference arm
\newcommand{\rhabLassoLoIvCloseLateP}{0.076}  % p-value
\newcommand{\rhabLassoLoLiqAllQlike}{0.1774}  % QLIKE lasso lo_liq all
\newcommand{\rhabLassoLoLiqAllMse}{2.39}  % MSE x 1e10 lasso lo_liq all
\newcommand{\rhabLassoLoLiqAllPct}{$-$0.1}  % QLIKE change % vs full
\newcommand{\rhabLassoLoLiqAllDm}{$-$0.21}  % DM vs the reference arm
\newcommand{\rhabLassoLoLiqAllP}{0.831}  % p-value
\newcommand{\rhabLassoLoLiqCloseQlike}{0.1285}  % QLIKE lasso lo_liq c16
\newcommand{\rhabLassoLoLiqCloseMse}{4.03}  % MSE x 1e10 lasso lo_liq c16
\newcommand{\rhabLassoLoLiqClosePct}{$-$1.9}  % QLIKE change % vs full
\newcommand{\rhabLassoLoLiqCloseDm}{$-$2.17}  % DM vs the reference arm
\newcommand{\rhabLassoLoLiqCloseP}{0.030}  % p-value
\newcommand{\rhabLassoLoLiqCloseLateQlike}{0.1272}  % QLIKE lasso lo_liq c16_2020
\newcommand{\rhabLassoLoLiqCloseLateMse}{9.23}  % MSE x 1e10 lasso lo_liq c16_2020
\newcommand{\rhabLassoLoLiqCloseLatePct}{$-$2.2}  % QLIKE change % vs full
\newcommand{\rhabLassoLoLiqCloseLateDm}{$-$1.99}  % DM vs the reference arm
\newcommand{\rhabLassoLoLiqCloseLateP}{0.047}  % p-value
\newcommand{\rhabLassoLoRmAllQlike}{0.1798}  % QLIKE lasso lo_rm all
\newcommand{\rhabLassoLoRmAllMse}{2.07}  % MSE x 1e10 lasso lo_rm all
\newcommand{\rhabLassoLoRmAllPct}{1.3}  % QLIKE change % vs full
\newcommand{\rhabLassoLoRmAllDm}{2.07}  % DM vs the reference arm
\newcommand{\rhabLassoLoRmAllP}{0.039}  % p-value
\newcommand{\rhabLassoLoRmCloseQlike}{0.1311}  % QLIKE lasso lo_rm c16
\newcommand{\rhabLassoLoRmCloseMse}{2.27}  % MSE x 1e10 lasso lo_rm c16
\newcommand{\rhabLassoLoRmClosePct}{0.0}  % QLIKE change % vs full
\newcommand{\rhabLassoLoRmCloseDm}{0.01}  % DM vs the reference arm
\newcommand{\rhabLassoLoRmCloseP}{0.990}  % p-value
\newcommand{\rhabLassoLoRmCloseLateQlike}{0.1255}  % QLIKE lasso lo_rm c16_2020
\newcommand{\rhabLassoLoRmCloseLateMse}{4.99}  % MSE x 1e10 lasso lo_rm c16_2020
\newcommand{\rhabLassoLoRmCloseLatePct}{$-$3.5}  % QLIKE change % vs full
\newcommand{\rhabLassoLoRmCloseLateDm}{$-$1.28}  % DM vs the reference arm
\newcommand{\rhabLassoLoRmCloseLateP}{0.200}  % p-value
\newcommand{\rhabLassoLoSentAllQlike}{0.1776}  % QLIKE lasso lo_sent all
\newcommand{\rhabLassoLoSentAllMse}{2.37}  % MSE x 1e10 lasso lo_sent all
\newcommand{\rhabLassoLoSentAllPct}{0.1}  % QLIKE change % vs full
\newcommand{\rhabLassoLoSentAllDm}{0.41}  % DM vs the reference arm
\newcommand{\rhabLassoLoSentAllP}{0.679}  % p-value
\newcommand{\rhabLassoLoSentCloseQlike}{0.1294}  % QLIKE lasso lo_sent c16
\newcommand{\rhabLassoLoSentCloseMse}{3.60}  % MSE x 1e10 lasso lo_sent c16
\newcommand{\rhabLassoLoSentClosePct}{$-$1.3}  % QLIKE change % vs full
\newcommand{\rhabLassoLoSentCloseDm}{$-$2.54}  % DM vs the reference arm
\newcommand{\rhabLassoLoSentCloseP}{0.011}  % p-value
\newcommand{\rhabLassoLoSentCloseLateQlike}{0.1262}  % QLIKE lasso lo_sent c16_2020
\newcommand{\rhabLassoLoSentCloseLateMse}{8.21}  % MSE x 1e10 lasso lo_sent c16_2020
\newcommand{\rhabLassoLoSentCloseLatePct}{$-$3.0}  % QLIKE change % vs full
\newcommand{\rhabLassoLoSentCloseLateDm}{$-$2.62}  % DM vs the reference arm
\newcommand{\rhabLassoLoSentCloseLateP}{0.009}  % p-value
\newcommand{\rhabLassoLoVdAllQlike}{0.1773}  % QLIKE lasso lo_vd all
\newcommand{\rhabLassoLoVdAllMse}{2.35}  % MSE x 1e10 lasso lo_vd all
\newcommand{\rhabLassoLoVdAllPct}{$-$0.1}  % QLIKE change % vs full
\newcommand{\rhabLassoLoVdAllDm}{$-$0.63}  % DM vs the reference arm
\newcommand{\rhabLassoLoVdAllP}{0.531}  % p-value
\newcommand{\rhabLassoLoVdCloseQlike}{0.1299}  % QLIKE lasso lo_vd c16
\newcommand{\rhabLassoLoVdCloseMse}{3.56}  % MSE x 1e10 lasso lo_vd c16
\newcommand{\rhabLassoLoVdClosePct}{$-$0.9}  % QLIKE change % vs full
\newcommand{\rhabLassoLoVdCloseDm}{$-$2.10}  % DM vs the reference arm
\newcommand{\rhabLassoLoVdCloseP}{0.036}  % p-value
\newcommand{\rhabLassoLoVdCloseLateQlike}{0.1275}  % QLIKE lasso lo_vd c16_2020
\newcommand{\rhabLassoLoVdCloseLateMse}{8.12}  % MSE x 1e10 lasso lo_vd c16_2020
\newcommand{\rhabLassoLoVdCloseLatePct}{$-$2.0}  % QLIKE change % vs full
\newcommand{\rhabLassoLoVdCloseLateDm}{$-$2.12}  % DM vs the reference arm
\newcommand{\rhabLassoLoVdCloseLateP}{0.034}  % p-value
\newcommand{\rhabLassoLoVwAllQlike}{0.1769}  % QLIKE lasso lo_vw all
\newcommand{\rhabLassoLoVwAllMse}{2.07}  % MSE x 1e10 lasso lo_vw all
\newcommand{\rhabLassoLoVwAllPct}{$-$0.3}  % QLIKE change % vs full
\newcommand{\rhabLassoLoVwAllDm}{$-$1.09}  % DM vs the reference arm
\newcommand{\rhabLassoLoVwAllP}{0.274}  % p-value
\newcommand{\rhabLassoLoVwCloseQlike}{0.1302}  % QLIKE lasso lo_vw c16
\newcommand{\rhabLassoLoVwCloseMse}{3.61}  % MSE x 1e10 lasso lo_vw c16
\newcommand{\rhabLassoLoVwClosePct}{$-$0.6}  % QLIKE change % vs full
\newcommand{\rhabLassoLoVwCloseDm}{$-$1.06}  % DM vs the reference arm
\newcommand{\rhabLassoLoVwCloseP}{0.289}  % p-value
\newcommand{\rhabLassoLoVwCloseLateQlike}{0.1277}  % QLIKE lasso lo_vw c16_2020
\newcommand{\rhabLassoLoVwCloseLateMse}{8.22}  % MSE x 1e10 lasso lo_vw c16_2020
\newcommand{\rhabLassoLoVwCloseLatePct}{$-$1.8}  % QLIKE change % vs full
\newcommand{\rhabLassoLoVwCloseLateDm}{$-$1.96}  % DM vs the reference arm
\newcommand{\rhabLassoLoVwCloseLateP}{0.050}  % p-value
\newcommand{\rhabLassoSharAllQlike}{0.1877}  % QLIKE lasso shar all
\newcommand{\rhabLassoSharAllMse}{1.39}  % MSE x 1e10 lasso shar all
\newcommand{\rhabLassoSharAllPct}{$-$0.8}  % QLIKE change % vs har
\newcommand{\rhabLassoSharAllDm}{$-$3.52}  % DM vs the reference arm
\newcommand{\rhabLassoSharAllP}{\ensuremath{<}0.001}  % p-value
\newcommand{\rhabLassoSharCloseQlike}{0.1265}  % QLIKE lasso shar c16
\newcommand{\rhabLassoSharCloseMse}{1.84}  % MSE x 1e10 lasso shar c16
\newcommand{\rhabLassoSharClosePct}{$-$0.3}  % QLIKE change % vs har
\newcommand{\rhabLassoSharCloseDm}{$-$0.44}  % DM vs the reference arm
\newcommand{\rhabLassoSharCloseP}{0.657}  % p-value
\newcommand{\rhabLassoSharCloseLateQlike}{0.1142}  % QLIKE lasso shar c16_2020
\newcommand{\rhabLassoSharCloseLateMse}{3.95}  % MSE x 1e10 lasso shar c16_2020
\newcommand{\rhabLassoSharCloseLatePct}{0.4}  % QLIKE change % vs har
\newcommand{\rhabLassoSharCloseLateDm}{0.42}  % DM vs the reference arm
\newcommand{\rhabLassoSharCloseLateP}{0.673}  % p-value
\newcommand{\rhabLassoSharqlAllQlike}{0.1840}  % QLIKE lasso sharql all
\newcommand{\rhabLassoSharqlAllMse}{1.69}  % MSE x 1e10 lasso sharql all
\newcommand{\rhabLassoSharqlAllPct}{$-$2.8}  % QLIKE change % vs har
\newcommand{\rhabLassoSharqlAllDm}{$-$6.63}  % DM vs the reference arm
\newcommand{\rhabLassoSharqlAllP}{\ensuremath{<}0.001}  % p-value
\newcommand{\rhabLassoSharqlCloseQlike}{0.1218}  % QLIKE lasso sharql c16
\newcommand{\rhabLassoSharqlCloseMse}{2.68}  % MSE x 1e10 lasso sharql c16
\newcommand{\rhabLassoSharqlClosePct}{$-$4.0}  % QLIKE change % vs har
\newcommand{\rhabLassoSharqlCloseDm}{$-$3.14}  % DM vs the reference arm
\newcommand{\rhabLassoSharqlCloseP}{0.002}  % p-value
\newcommand{\rhabLassoSharqlCloseLateQlike}{0.1125}  % QLIKE lasso sharql c16_2020
\newcommand{\rhabLassoSharqlCloseLateMse}{5.98}  % MSE x 1e10 lasso sharql c16_2020
\newcommand{\rhabLassoSharqlCloseLatePct}{$-$1.0}  % QLIKE change % vs har
\newcommand{\rhabLassoSharqlCloseLateDm}{$-$0.44}  % DM vs the reference arm
\newcommand{\rhabLassoSharqlCloseLateP}{0.662}  % p-value
\newcommand{\rhabLgbmFullAllQlike}{0.1613}  % QLIKE lgbm full all
\newcommand{\rhabLgbmFullAllMse}{1.41}  % MSE x 1e10 lgbm full all
\newcommand{\rhabLgbmFullAllPct}{$-$14.0}  % QLIKE change % vs har
\newcommand{\rhabLgbmFullAllDm}{$-$10.32}  % DM vs the reference arm
\newcommand{\rhabLgbmFullAllP}{\ensuremath{<}0.001}  % p-value
\newcommand{\rhabLgbmFullCloseQlike}{0.1201}  % QLIKE lgbm full c16
\newcommand{\rhabLgbmFullCloseMse}{2.34}  % MSE x 1e10 lgbm full c16
\newcommand{\rhabLgbmFullClosePct}{$-$2.7}  % QLIKE change % vs har
\newcommand{\rhabLgbmFullCloseDm}{$-$1.81}  % DM vs the reference arm
\newcommand{\rhabLgbmFullCloseP}{0.071}  % p-value
\newcommand{\rhabLgbmFullCloseLateQlike}{0.1089}  % QLIKE lgbm full c16_2020
\newcommand{\rhabLgbmFullCloseLateMse}{5.09}  % MSE x 1e10 lgbm full c16_2020
\newcommand{\rhabLgbmFullCloseLatePct}{$-$1.1}  % QLIKE change % vs har
\newcommand{\rhabLgbmFullCloseLateDm}{$-$0.39}  % DM vs the reference arm
\newcommand{\rhabLgbmFullCloseLateP}{0.697}  % p-value
\newcommand{\rhabLgbmHarAllQlike}{0.1875}  % QLIKE lgbm har all
\newcommand{\rhabLgbmHarAllMse}{1.18}  % MSE x 1e10 lgbm har all
\newcommand{\rhabLgbmHarCloseQlike}{0.1235}  % QLIKE lgbm har c16
\newcommand{\rhabLgbmHarCloseMse}{1.77}  % MSE x 1e10 lgbm har c16
\newcommand{\rhabLgbmHarCloseLateQlike}{0.1101}  % QLIKE lgbm har c16_2020
\newcommand{\rhabLgbmHarCloseLateMse}{3.68}  % MSE x 1e10 lgbm har c16_2020
\newcommand{\rhabOlsAddRmAllQlike}{0.1858}  % QLIKE ols add_rm all
\newcommand{\rhabOlsAddRmAllMse}{1.79}  % MSE x 1e10 ols add_rm all
\newcommand{\rhabOlsAddRmAllPct}{$-$1.9}  % QLIKE change % vs har
\newcommand{\rhabOlsAddRmAllDm}{$-$2.00}  % DM vs the reference arm
\newcommand{\rhabOlsAddRmAllP}{0.045}  % p-value
\newcommand{\rhabOlsAddRmCloseQlike}{0.1252}  % QLIKE ols add_rm c16
\newcommand{\rhabOlsAddRmCloseMse}{3.02}  % MSE x 1e10 ols add_rm c16
\newcommand{\rhabOlsAddRmClosePct}{$-$1.3}  % QLIKE change % vs har
\newcommand{\rhabOlsAddRmCloseDm}{$-$0.56}  % DM vs the reference arm
\newcommand{\rhabOlsAddRmCloseP}{0.576}  % p-value
\newcommand{\rhabOlsAddRmCloseLateQlike}{0.1221}  % QLIKE ols add_rm c16_2020
\newcommand{\rhabOlsAddRmCloseLateMse}{6.80}  % MSE x 1e10 ols add_rm c16_2020
\newcommand{\rhabOlsAddRmCloseLatePct}{7.5}  % QLIKE change % vs har
\newcommand{\rhabOlsAddRmCloseLateDm}{1.43}  % DM vs the reference arm
\newcommand{\rhabOlsAddRmCloseLateP}{0.152}  % p-value
\newcommand{\rhabOlsHarAllQlike}{0.1893}  % QLIKE ols har all
\newcommand{\rhabOlsHarAllMse}{1.11}  % MSE x 1e10 ols har all
\newcommand{\rhabOlsHarCloseQlike}{0.1269}  % QLIKE ols har c16
\newcommand{\rhabOlsHarCloseMse}{1.97}  % MSE x 1e10 ols har c16
\newcommand{\rhabOlsHarCloseLateQlike}{0.1137}  % QLIKE ols har c16_2020
\newcommand{\rhabOlsHarCloseLateMse}{4.24}  % MSE x 1e10 ols har c16_2020
\newcommand{\rhabOlsHarqAllQlike}{0.1893}  % QLIKE ols harq all
\newcommand{\rhabOlsHarqAllMse}{1.17}  % MSE x 1e10 ols harq all
\newcommand{\rhabOlsHarqAllPct}{0.0}  % QLIKE change % vs har
\newcommand{\rhabOlsHarqAllDm}{0.61}  % DM vs the reference arm
\newcommand{\rhabOlsHarqAllP}{0.543}  % p-value
\newcommand{\rhabOlsHarqCloseQlike}{0.1267}  % QLIKE ols harq c16
\newcommand{\rhabOlsHarqCloseMse}{1.78}  % MSE x 1e10 ols harq c16
\newcommand{\rhabOlsHarqClosePct}{$-$0.2}  % QLIKE change % vs har
\newcommand{\rhabOlsHarqCloseDm}{$-$1.38}  % DM vs the reference arm
\newcommand{\rhabOlsHarqCloseP}{0.166}  % p-value
\newcommand{\rhabOlsHarqCloseLateQlike}{0.1133}  % QLIKE ols harq c16_2020
\newcommand{\rhabOlsHarqCloseLateMse}{3.80}  % MSE x 1e10 ols harq c16_2020
\newcommand{\rhabOlsHarqCloseLatePct}{$-$0.3}  % QLIKE change % vs har
\newcommand{\rhabOlsHarqCloseLateDm}{$-$1.19}  % DM vs the reference arm
\newcommand{\rhabOlsHarqCloseLateP}{0.235}  % p-value
\newcommand{\rhabOlsLharAllQlike}{0.1842}  % QLIKE ols lhar all
\newcommand{\rhabOlsLharAllMse}{1.71}  % MSE x 1e10 ols lhar all
\newcommand{\rhabOlsLharAllPct}{$-$2.7}  % QLIKE change % vs har
\newcommand{\rhabOlsLharAllDm}{$-$6.19}  % DM vs the reference arm
\newcommand{\rhabOlsLharAllP}{\ensuremath{<}0.001}  % p-value
\newcommand{\rhabOlsLharCloseQlike}{0.1220}  % QLIKE ols lhar c16
\newcommand{\rhabOlsLharCloseMse}{2.81}  % MSE x 1e10 ols lhar c16
\newcommand{\rhabOlsLharClosePct}{$-$3.9}  % QLIKE change % vs har
\newcommand{\rhabOlsLharCloseDm}{$-$2.66}  % DM vs the reference arm
\newcommand{\rhabOlsLharCloseP}{0.008}  % p-value
\newcommand{\rhabOlsLharCloseLateQlike}{0.1132}  % QLIKE ols lhar c16_2020
\newcommand{\rhabOlsLharCloseLateMse}{6.30}  % MSE x 1e10 ols lhar c16_2020
\newcommand{\rhabOlsLharCloseLatePct}{$-$0.4}  % QLIKE change % vs har
\newcommand{\rhabOlsLharCloseLateDm}{$-$0.15}  % DM vs the reference arm
\newcommand{\rhabOlsLharCloseLateP}{0.883}  % p-value
\newcommand{\rhabOlsSharAllQlike}{0.1877}  % QLIKE ols shar all
\newcommand{\rhabOlsSharAllMse}{1.40}  % MSE x 1e10 ols shar all
\newcommand{\rhabOlsSharAllPct}{$-$0.8}  % QLIKE change % vs har
\newcommand{\rhabOlsSharAllDm}{$-$3.48}  % DM vs the reference arm
\newcommand{\rhabOlsSharAllP}{\ensuremath{<}0.001}  % p-value
\newcommand{\rhabOlsSharCloseQlike}{0.1265}  % QLIKE ols shar c16
\newcommand{\rhabOlsSharCloseMse}{1.84}  % MSE x 1e10 ols shar c16
\newcommand{\rhabOlsSharClosePct}{$-$0.3}  % QLIKE change % vs har
\newcommand{\rhabOlsSharCloseDm}{$-$0.42}  % DM vs the reference arm
\newcommand{\rhabOlsSharCloseP}{0.674}  % p-value
\newcommand{\rhabOlsSharCloseLateQlike}{0.1142}  % QLIKE ols shar c16_2020
\newcommand{\rhabOlsSharCloseLateMse}{3.95}  % MSE x 1e10 ols shar c16_2020
\newcommand{\rhabOlsSharCloseLatePct}{0.5}  % QLIKE change % vs har
\newcommand{\rhabOlsSharCloseLateDm}{0.44}  % DM vs the reference arm
\newcommand{\rhabOlsSharCloseLateP}{0.659}  % p-value
\newcommand{\rhabOlsSharqlAllQlike}{0.1841}  % QLIKE ols sharql all
\newcommand{\rhabOlsSharqlAllMse}{1.69}  % MSE x 1e10 ols sharql all
\newcommand{\rhabOlsSharqlAllPct}{$-$2.7}  % QLIKE change % vs har
\newcommand{\rhabOlsSharqlAllDm}{$-$6.53}  % DM vs the reference arm
\newcommand{\rhabOlsSharqlAllP}{\ensuremath{<}0.001}  % p-value
\newcommand{\rhabOlsSharqlCloseQlike}{0.1218}  % QLIKE ols sharql c16
\newcommand{\rhabOlsSharqlCloseMse}{2.68}  % MSE x 1e10 ols sharql c16
\newcommand{\rhabOlsSharqlClosePct}{$-$4.0}  % QLIKE change % vs har
\newcommand{\rhabOlsSharqlCloseDm}{$-$3.05}  % DM vs the reference arm
\newcommand{\rhabOlsSharqlCloseP}{0.002}  % p-value
\newcommand{\rhabOlsSharqlCloseLateQlike}{0.1126}  % QLIKE ols sharql c16_2020
\newcommand{\rhabOlsSharqlCloseLateMse}{5.99}  % MSE x 1e10 ols sharql c16_2020
\newcommand{\rhabOlsSharqlCloseLatePct}{$-$0.9}  % QLIKE change % vs har
\newcommand{\rhabOlsSharqlCloseLateDm}{$-$0.38}  % DM vs the reference arm
\newcommand{\rhabOlsSharqlCloseLateP}{0.703}  % p-value
\newcommand{\rhabRfFullAllQlike}{0.1685}  % QLIKE rf full all
\newcommand{\rhabRfFullAllMse}{1.32}  % MSE x 1e10 rf full all
\newcommand{\rhabRfFullAllPct}{$-$11.1}  % QLIKE change % vs har
\newcommand{\rhabRfFullAllDm}{$-$9.85}  % DM vs the reference arm
\newcommand{\rhabRfFullAllP}{\ensuremath{<}0.001}  % p-value
\newcommand{\rhabRfFullCloseQlike}{0.1248}  % QLIKE rf full c16
\newcommand{\rhabRfFullCloseMse}{1.99}  % MSE x 1e10 rf full c16
\newcommand{\rhabRfFullClosePct}{$-$1.4}  % QLIKE change % vs har
\newcommand{\rhabRfFullCloseDm}{$-$0.92}  % DM vs the reference arm
\newcommand{\rhabRfFullCloseP}{0.357}  % p-value
\newcommand{\rhabRfFullCloseLateQlike}{0.1123}  % QLIKE rf full c16_2020
\newcommand{\rhabRfFullCloseLateMse}{4.26}  % MSE x 1e10 rf full c16_2020
\newcommand{\rhabRfFullCloseLatePct}{$-$0.7}  % QLIKE change % vs har
\newcommand{\rhabRfFullCloseLateDm}{$-$0.26}  % DM vs the reference arm
\newcommand{\rhabRfFullCloseLateP}{0.794}  % p-value
\newcommand{\rhabRfHarAllQlike}{0.1895}  % QLIKE rf har all
\newcommand{\rhabRfHarAllMse}{1.19}  % MSE x 1e10 rf har all
\newcommand{\rhabRfHarCloseQlike}{0.1266}  % QLIKE rf har c16
\newcommand{\rhabRfHarCloseMse}{1.63}  % MSE x 1e10 rf har c16
\newcommand{\rhabRfHarCloseLateQlike}{0.1131}  % QLIKE rf har c16_2020
\newcommand{\rhabRfHarCloseLateMse}{3.39}  % MSE x 1e10 rf har c16_2020
\newcommand{\rhabRidgeAddEwAllQlike}{0.1876}  % QLIKE ridge add_ew all
\newcommand{\rhabRidgeAddEwAllMse}{1.42}  % MSE x 1e10 ridge add_ew all
\newcommand{\rhabRidgeAddEwAllPct}{$-$0.9}  % QLIKE change % vs har
\newcommand{\rhabRidgeAddEwAllDm}{$-$2.26}  % DM vs the reference arm
\newcommand{\rhabRidgeAddEwAllP}{0.024}  % p-value
\newcommand{\rhabRidgeAddEwCloseQlike}{0.1256}  % QLIKE ridge add_ew c16
\newcommand{\rhabRidgeAddEwCloseMse}{2.08}  % MSE x 1e10 ridge add_ew c16
\newcommand{\rhabRidgeAddEwClosePct}{$-$1.0}  % QLIKE change % vs har
\newcommand{\rhabRidgeAddEwCloseDm}{$-$0.89}  % DM vs the reference arm
\newcommand{\rhabRidgeAddEwCloseP}{0.374}  % p-value
\newcommand{\rhabRidgeAddEwCloseLateQlike}{0.1125}  % QLIKE ridge add_ew c16_2020
\newcommand{\rhabRidgeAddEwCloseLateMse}{4.51}  % MSE x 1e10 ridge add_ew c16_2020
\newcommand{\rhabRidgeAddEwCloseLatePct}{$-$1.0}  % QLIKE change % vs har
\newcommand{\rhabRidgeAddEwCloseLateDm}{$-$0.62}  % DM vs the reference arm
\newcommand{\rhabRidgeAddEwCloseLateP}{0.537}  % p-value
\newcommand{\rhabRidgeAddFomcAllQlike}{0.1847}  % QLIKE ridge add_fomc all
\newcommand{\rhabRidgeAddFomcAllMse}{1.13}  % MSE x 1e10 ridge add_fomc all
\newcommand{\rhabRidgeAddFomcAllPct}{$-$2.4}  % QLIKE change % vs har
\newcommand{\rhabRidgeAddFomcAllDm}{$-$5.79}  % DM vs the reference arm
\newcommand{\rhabRidgeAddFomcAllP}{\ensuremath{<}0.001}  % p-value
\newcommand{\rhabRidgeAddFomcCloseQlike}{0.1276}  % QLIKE ridge add_fomc c16
\newcommand{\rhabRidgeAddFomcCloseMse}{1.84}  % MSE x 1e10 ridge add_fomc c16
\newcommand{\rhabRidgeAddFomcClosePct}{0.6}  % QLIKE change % vs har
\newcommand{\rhabRidgeAddFomcCloseDm}{2.17}  % DM vs the reference arm
\newcommand{\rhabRidgeAddFomcCloseP}{0.030}  % p-value
\newcommand{\rhabRidgeAddFomcCloseLateQlike}{0.1135}  % QLIKE ridge add_fomc c16_2020
\newcommand{\rhabRidgeAddFomcCloseLateMse}{3.94}  % MSE x 1e10 ridge add_fomc c16_2020
\newcommand{\rhabRidgeAddFomcCloseLatePct}{$-$0.2}  % QLIKE change % vs har
\newcommand{\rhabRidgeAddFomcCloseLateDm}{$-$0.35}  % DM vs the reference arm
\newcommand{\rhabRidgeAddFomcCloseLateP}{0.729}  % p-value
\newcommand{\rhabRidgeAddIvAllQlike}{0.1897}  % QLIKE ridge add_iv all
\newcommand{\rhabRidgeAddIvAllMse}{1.28}  % MSE x 1e10 ridge add_iv all
\newcommand{\rhabRidgeAddIvAllPct}{0.2}  % QLIKE change % vs har
\newcommand{\rhabRidgeAddIvAllDm}{0.67}  % DM vs the reference arm
\newcommand{\rhabRidgeAddIvAllP}{0.506}  % p-value
\newcommand{\rhabRidgeAddIvCloseQlike}{0.1266}  % QLIKE ridge add_iv c16
\newcommand{\rhabRidgeAddIvCloseMse}{1.79}  % MSE x 1e10 ridge add_iv c16
\newcommand{\rhabRidgeAddIvClosePct}{$-$0.2}  % QLIKE change % vs har
\newcommand{\rhabRidgeAddIvCloseDm}{$-$0.20}  % DM vs the reference arm
\newcommand{\rhabRidgeAddIvCloseP}{0.845}  % p-value
\newcommand{\rhabRidgeAddIvCloseLateQlike}{0.1153}  % QLIKE ridge add_iv c16_2020
\newcommand{\rhabRidgeAddIvCloseLateMse}{3.82}  % MSE x 1e10 ridge add_iv c16_2020
\newcommand{\rhabRidgeAddIvCloseLatePct}{1.4}  % QLIKE change % vs har
\newcommand{\rhabRidgeAddIvCloseLateDm}{0.70}  % DM vs the reference arm
\newcommand{\rhabRidgeAddIvCloseLateP}{0.482}  % p-value
\newcommand{\rhabRidgeAddLiqAllQlike}{0.1885}  % QLIKE ridge add_liq all
\newcommand{\rhabRidgeAddLiqAllMse}{1.18}  % MSE x 1e10 ridge add_liq all
\newcommand{\rhabRidgeAddLiqAllPct}{$-$0.4}  % QLIKE change % vs har
\newcommand{\rhabRidgeAddLiqAllDm}{$-$0.98}  % DM vs the reference arm
\newcommand{\rhabRidgeAddLiqAllP}{0.326}  % p-value
\newcommand{\rhabRidgeAddLiqCloseQlike}{0.1304}  % QLIKE ridge add_liq c16
\newcommand{\rhabRidgeAddLiqCloseMse}{2.01}  % MSE x 1e10 ridge add_liq c16
\newcommand{\rhabRidgeAddLiqClosePct}{2.8}  % QLIKE change % vs har
\newcommand{\rhabRidgeAddLiqCloseDm}{1.91}  % DM vs the reference arm
\newcommand{\rhabRidgeAddLiqCloseP}{0.056}  % p-value
\newcommand{\rhabRidgeAddLiqCloseLateQlike}{0.1204}  % QLIKE ridge add_liq c16_2020
\newcommand{\rhabRidgeAddLiqCloseLateMse}{4.35}  % MSE x 1e10 ridge add_liq c16_2020
\newcommand{\rhabRidgeAddLiqCloseLatePct}{6.0}  % QLIKE change % vs har
\newcommand{\rhabRidgeAddLiqCloseLateDm}{2.09}  % DM vs the reference arm
\newcommand{\rhabRidgeAddLiqCloseLateP}{0.036}  % p-value
\newcommand{\rhabRidgeAddRmAllQlike}{0.1843}  % QLIKE ridge add_rm all
\newcommand{\rhabRidgeAddRmAllMse}{1.69}  % MSE x 1e10 ridge add_rm all
\newcommand{\rhabRidgeAddRmAllPct}{$-$2.6}  % QLIKE change % vs har
\newcommand{\rhabRidgeAddRmAllDm}{$-$4.14}  % DM vs the reference arm
\newcommand{\rhabRidgeAddRmAllP}{\ensuremath{<}0.001}  % p-value
\newcommand{\rhabRidgeAddRmCloseQlike}{0.1230}  % QLIKE ridge add_rm c16
\newcommand{\rhabRidgeAddRmCloseMse}{2.81}  % MSE x 1e10 ridge add_rm c16
\newcommand{\rhabRidgeAddRmClosePct}{$-$3.1}  % QLIKE change % vs har
\newcommand{\rhabRidgeAddRmCloseDm}{$-$1.80}  % DM vs the reference arm
\newcommand{\rhabRidgeAddRmCloseP}{0.071}  % p-value
\newcommand{\rhabRidgeAddRmCloseLateQlike}{0.1156}  % QLIKE ridge add_rm c16_2020
\newcommand{\rhabRidgeAddRmCloseLateMse}{6.29}  % MSE x 1e10 ridge add_rm c16_2020
\newcommand{\rhabRidgeAddRmCloseLatePct}{1.7}  % QLIKE change % vs har
\newcommand{\rhabRidgeAddRmCloseLateDm}{0.56}  % DM vs the reference arm
\newcommand{\rhabRidgeAddRmCloseLateP}{0.577}  % p-value
\newcommand{\rhabRidgeAddSentAllQlike}{0.1897}  % QLIKE ridge add_sent all
\newcommand{\rhabRidgeAddSentAllMse}{1.11}  % MSE x 1e10 ridge add_sent all
\newcommand{\rhabRidgeAddSentAllPct}{0.2}  % QLIKE change % vs har
\newcommand{\rhabRidgeAddSentAllDm}{0.82}  % DM vs the reference arm
\newcommand{\rhabRidgeAddSentAllP}{0.411}  % p-value
\newcommand{\rhabRidgeAddSentCloseQlike}{0.1310}  % QLIKE ridge add_sent c16
\newcommand{\rhabRidgeAddSentCloseMse}{1.98}  % MSE x 1e10 ridge add_sent c16
\newcommand{\rhabRidgeAddSentClosePct}{3.2}  % QLIKE change % vs har
\newcommand{\rhabRidgeAddSentCloseDm}{4.41}  % DM vs the reference arm
\newcommand{\rhabRidgeAddSentCloseP}{\ensuremath{<}0.001}  % p-value
\newcommand{\rhabRidgeAddSentCloseLateQlike}{0.1214}  % QLIKE ridge add_sent c16_2020
\newcommand{\rhabRidgeAddSentCloseLateMse}{4.27}  % MSE x 1e10 ridge add_sent c16_2020
\newcommand{\rhabRidgeAddSentCloseLatePct}{6.8}  % QLIKE change % vs har
\newcommand{\rhabRidgeAddSentCloseLateDm}{3.65}  % DM vs the reference arm
\newcommand{\rhabRidgeAddSentCloseLateP}{\ensuremath{<}0.001}  % p-value
\newcommand{\rhabRidgeAddVdAllQlike}{0.1903}  % QLIKE ridge add_vd all
\newcommand{\rhabRidgeAddVdAllMse}{1.11}  % MSE x 1e10 ridge add_vd all
\newcommand{\rhabRidgeAddVdAllPct}{0.5}  % QLIKE change % vs har
\newcommand{\rhabRidgeAddVdAllDm}{3.10}  % DM vs the reference arm
\newcommand{\rhabRidgeAddVdAllP}{0.002}  % p-value
\newcommand{\rhabRidgeAddVdCloseQlike}{0.1276}  % QLIKE ridge add_vd c16
\newcommand{\rhabRidgeAddVdCloseMse}{1.95}  % MSE x 1e10 ridge add_vd c16
\newcommand{\rhabRidgeAddVdClosePct}{0.6}  % QLIKE change % vs har
\newcommand{\rhabRidgeAddVdCloseDm}{1.55}  % DM vs the reference arm
\newcommand{\rhabRidgeAddVdCloseP}{0.120}  % p-value
\newcommand{\rhabRidgeAddVdCloseLateQlike}{0.1169}  % QLIKE ridge add_vd c16_2020
\newcommand{\rhabRidgeAddVdCloseLateMse}{4.21}  % MSE x 1e10 ridge add_vd c16_2020
\newcommand{\rhabRidgeAddVdCloseLatePct}{2.8}  % QLIKE change % vs har
\newcommand{\rhabRidgeAddVdCloseLateDm}{3.10}  % DM vs the reference arm
\newcommand{\rhabRidgeAddVdCloseLateP}{0.002}  % p-value
\newcommand{\rhabRidgeAddVwAllQlike}{0.1876}  % QLIKE ridge add_vw all
\newcommand{\rhabRidgeAddVwAllMse}{1.87}  % MSE x 1e10 ridge add_vw all
\newcommand{\rhabRidgeAddVwAllPct}{$-$0.9}  % QLIKE change % vs har
\newcommand{\rhabRidgeAddVwAllDm}{$-$1.58}  % DM vs the reference arm
\newcommand{\rhabRidgeAddVwAllP}{0.114}  % p-value
\newcommand{\rhabRidgeAddVwCloseQlike}{0.1264}  % QLIKE ridge add_vw c16
\newcommand{\rhabRidgeAddVwCloseMse}{2.27}  % MSE x 1e10 ridge add_vw c16
\newcommand{\rhabRidgeAddVwClosePct}{$-$0.4}  % QLIKE change % vs har
\newcommand{\rhabRidgeAddVwCloseDm}{$-$0.28}  % DM vs the reference arm
\newcommand{\rhabRidgeAddVwCloseP}{0.781}  % p-value
\newcommand{\rhabRidgeAddVwCloseLateQlike}{0.1139}  % QLIKE ridge add_vw c16_2020
\newcommand{\rhabRidgeAddVwCloseLateMse}{5.00}  % MSE x 1e10 ridge add_vw c16_2020
\newcommand{\rhabRidgeAddVwCloseLatePct}{0.2}  % QLIKE change % vs har
\newcommand{\rhabRidgeAddVwCloseLateDm}{0.12}  % DM vs the reference arm
\newcommand{\rhabRidgeAddVwCloseLateP}{0.907}  % p-value
\newcommand{\rhabRidgeFullAllQlike}{0.1791}  % QLIKE ridge full all
\newcommand{\rhabRidgeFullAllMse}{2.24}  % MSE x 1e10 ridge full all
\newcommand{\rhabRidgeFullAllPct}{$-$5.4}  % QLIKE change % vs har
\newcommand{\rhabRidgeFullAllDm}{$-$5.63}  % DM vs the reference arm
\newcommand{\rhabRidgeFullAllP}{\ensuremath{<}0.001}  % p-value
\newcommand{\rhabRidgeFullCloseQlike}{0.1279}  % QLIKE ridge full c16
\newcommand{\rhabRidgeFullCloseMse}{3.52}  % MSE x 1e10 ridge full c16
\newcommand{\rhabRidgeFullClosePct}{0.8}  % QLIKE change % vs har
\newcommand{\rhabRidgeFullCloseDm}{0.31}  % DM vs the reference arm
\newcommand{\rhabRidgeFullCloseP}{0.759}  % p-value
\newcommand{\rhabRidgeFullCloseLateQlike}{0.1260}  % QLIKE ridge full c16_2020
\newcommand{\rhabRidgeFullCloseLateMse}{7.99}  % MSE x 1e10 ridge full c16_2020
\newcommand{\rhabRidgeFullCloseLatePct}{10.9}  % QLIKE change % vs har
\newcommand{\rhabRidgeFullCloseLateDm}{2.31}  % DM vs the reference arm
\newcommand{\rhabRidgeFullCloseLateP}{0.021}  % p-value
\newcommand{\rhabRidgeHarAllQlike}{0.1893}  % QLIKE ridge har all
\newcommand{\rhabRidgeHarAllMse}{1.11}  % MSE x 1e10 ridge har all
\newcommand{\rhabRidgeHarCloseQlike}{0.1269}  % QLIKE ridge har c16
\newcommand{\rhabRidgeHarCloseMse}{1.97}  % MSE x 1e10 ridge har c16
\newcommand{\rhabRidgeHarCloseLateQlike}{0.1137}  % QLIKE ridge har c16_2020
\newcommand{\rhabRidgeHarCloseLateMse}{4.24}  % MSE x 1e10 ridge har c16_2020
\newcommand{\rhabRidgeHarqAllQlike}{0.1893}  % QLIKE ridge harq all
\newcommand{\rhabRidgeHarqAllMse}{1.16}  % MSE x 1e10 ridge harq all
\newcommand{\rhabRidgeHarqAllPct}{0.0}  % QLIKE change % vs har
\newcommand{\rhabRidgeHarqAllDm}{0.76}  % DM vs the reference arm
\newcommand{\rhabRidgeHarqAllP}{0.446}  % p-value
\newcommand{\rhabRidgeHarqCloseQlike}{0.1267}  % QLIKE ridge harq c16
\newcommand{\rhabRidgeHarqCloseMse}{1.79}  % MSE x 1e10 ridge harq c16
\newcommand{\rhabRidgeHarqClosePct}{$-$0.2}  % QLIKE change % vs har
\newcommand{\rhabRidgeHarqCloseDm}{$-$1.42}  % DM vs the reference arm
\newcommand{\rhabRidgeHarqCloseP}{0.157}  % p-value
\newcommand{\rhabRidgeHarqCloseLateQlike}{0.1133}  % QLIKE ridge harq c16_2020
\newcommand{\rhabRidgeHarqCloseLateMse}{3.81}  % MSE x 1e10 ridge harq c16_2020
\newcommand{\rhabRidgeHarqCloseLatePct}{$-$0.3}  % QLIKE change % vs har
\newcommand{\rhabRidgeHarqCloseLateDm}{$-$1.20}  % DM vs the reference arm
\newcommand{\rhabRidgeHarqCloseLateP}{0.228}  % p-value
\newcommand{\rhabRidgeLharAllQlike}{0.1841}  % QLIKE ridge lhar all
\newcommand{\rhabRidgeLharAllMse}{1.72}  % MSE x 1e10 ridge lhar all
\newcommand{\rhabRidgeLharAllPct}{$-$2.7}  % QLIKE change % vs har
\newcommand{\rhabRidgeLharAllDm}{$-$6.56}  % DM vs the reference arm
\newcommand{\rhabRidgeLharAllP}{\ensuremath{<}0.001}  % p-value
\newcommand{\rhabRidgeLharCloseQlike}{0.1216}  % QLIKE ridge lhar c16
\newcommand{\rhabRidgeLharCloseMse}{2.78}  % MSE x 1e10 ridge lhar c16
\newcommand{\rhabRidgeLharClosePct}{$-$4.2}  % QLIKE change % vs har
\newcommand{\rhabRidgeLharCloseDm}{$-$3.28}  % DM vs the reference arm
\newcommand{\rhabRidgeLharCloseP}{0.001}  % p-value
\newcommand{\rhabRidgeLharCloseLateQlike}{0.1121}  % QLIKE ridge lhar c16_2020
\newcommand{\rhabRidgeLharCloseLateMse}{6.23}  % MSE x 1e10 ridge lhar c16_2020
\newcommand{\rhabRidgeLharCloseLatePct}{$-$1.4}  % QLIKE change % vs har
\newcommand{\rhabRidgeLharCloseLateDm}{$-$0.59}  % DM vs the reference arm
\newcommand{\rhabRidgeLharCloseLateP}{0.554}  % p-value
\newcommand{\rhabRidgeLoEwAllQlike}{0.1789}  % QLIKE ridge lo_ew all
\newcommand{\rhabRidgeLoEwAllMse}{2.38}  % MSE x 1e10 ridge lo_ew all
\newcommand{\rhabRidgeLoEwAllPct}{$-$0.1}  % QLIKE change % vs full
\newcommand{\rhabRidgeLoEwAllDm}{$-$1.20}  % DM vs the reference arm
\newcommand{\rhabRidgeLoEwAllP}{0.230}  % p-value
\newcommand{\rhabRidgeLoEwCloseQlike}{0.1285}  % QLIKE ridge lo_ew c16
\newcommand{\rhabRidgeLoEwCloseMse}{3.74}  % MSE x 1e10 ridge lo_ew c16
\newcommand{\rhabRidgeLoEwClosePct}{0.5}  % QLIKE change % vs full
\newcommand{\rhabRidgeLoEwCloseDm}{1.23}  % DM vs the reference arm
\newcommand{\rhabRidgeLoEwCloseP}{0.218}  % p-value
\newcommand{\rhabRidgeLoEwCloseLateQlike}{0.1262}  % QLIKE ridge lo_ew c16_2020
\newcommand{\rhabRidgeLoEwCloseLateMse}{8.53}  % MSE x 1e10 ridge lo_ew c16_2020
\newcommand{\rhabRidgeLoEwCloseLatePct}{0.2}  % QLIKE change % vs full
\newcommand{\rhabRidgeLoEwCloseLateDm}{0.22}  % DM vs the reference arm
\newcommand{\rhabRidgeLoEwCloseLateP}{0.828}  % p-value
\newcommand{\rhabRidgeLoFomcAllQlike}{0.1833}  % QLIKE ridge lo_fomc all
\newcommand{\rhabRidgeLoFomcAllMse}{2.25}  % MSE x 1e10 ridge lo_fomc all
\newcommand{\rhabRidgeLoFomcAllPct}{2.3}  % QLIKE change % vs full
\newcommand{\rhabRidgeLoFomcAllDm}{5.55}  % DM vs the reference arm
\newcommand{\rhabRidgeLoFomcAllP}{\ensuremath{<}0.001}  % p-value
\newcommand{\rhabRidgeLoFomcCloseQlike}{0.1275}  % QLIKE ridge lo_fomc c16
\newcommand{\rhabRidgeLoFomcCloseMse}{3.58}  % MSE x 1e10 ridge lo_fomc c16
\newcommand{\rhabRidgeLoFomcClosePct}{$-$0.3}  % QLIKE change % vs full
\newcommand{\rhabRidgeLoFomcCloseDm}{$-$1.23}  % DM vs the reference arm
\newcommand{\rhabRidgeLoFomcCloseP}{0.219}  % p-value
\newcommand{\rhabRidgeLoFomcCloseLateQlike}{0.1257}  % QLIKE ridge lo_fomc c16_2020
\newcommand{\rhabRidgeLoFomcCloseLateMse}{8.15}  % MSE x 1e10 ridge lo_fomc c16_2020
\newcommand{\rhabRidgeLoFomcCloseLatePct}{$-$0.2}  % QLIKE change % vs full
\newcommand{\rhabRidgeLoFomcCloseLateDm}{$-$0.49}  % DM vs the reference arm
\newcommand{\rhabRidgeLoFomcCloseLateP}{0.622}  % p-value
\newcommand{\rhabRidgeLoIvAllQlike}{0.1791}  % QLIKE ridge lo_iv all
\newcommand{\rhabRidgeLoIvAllMse}{2.30}  % MSE x 1e10 ridge lo_iv all
\newcommand{\rhabRidgeLoIvAllPct}{0.0}  % QLIKE change % vs full
\newcommand{\rhabRidgeLoIvAllDm}{0.03}  % DM vs the reference arm
\newcommand{\rhabRidgeLoIvAllP}{0.980}  % p-value
\newcommand{\rhabRidgeLoIvCloseQlike}{0.1267}  % QLIKE ridge lo_iv c16
\newcommand{\rhabRidgeLoIvCloseMse}{3.61}  % MSE x 1e10 ridge lo_iv c16
\newcommand{\rhabRidgeLoIvClosePct}{$-$1.0}  % QLIKE change % vs full
\newcommand{\rhabRidgeLoIvCloseDm}{$-$1.02}  % DM vs the reference arm
\newcommand{\rhabRidgeLoIvCloseP}{0.309}  % p-value
\newcommand{\rhabRidgeLoIvCloseLateQlike}{0.1228}  % QLIKE ridge lo_iv c16_2020
\newcommand{\rhabRidgeLoIvCloseLateMse}{8.20}  % MSE x 1e10 ridge lo_iv c16_2020
\newcommand{\rhabRidgeLoIvCloseLatePct}{$-$2.5}  % QLIKE change % vs full
\newcommand{\rhabRidgeLoIvCloseLateDm}{$-$1.41}  % DM vs the reference arm
\newcommand{\rhabRidgeLoIvCloseLateP}{0.159}  % p-value
\newcommand{\rhabRidgeLoLiqAllQlike}{0.1794}  % QLIKE ridge lo_liq all
\newcommand{\rhabRidgeLoLiqAllMse}{2.29}  % MSE x 1e10 ridge lo_liq all
\newcommand{\rhabRidgeLoLiqAllPct}{0.2}  % QLIKE change % vs full
\newcommand{\rhabRidgeLoLiqAllDm}{0.75}  % DM vs the reference arm
\newcommand{\rhabRidgeLoLiqAllP}{0.454}  % p-value
\newcommand{\rhabRidgeLoLiqCloseQlike}{0.1269}  % QLIKE ridge lo_liq c16
\newcommand{\rhabRidgeLoLiqCloseMse}{3.93}  % MSE x 1e10 ridge lo_liq c16
\newcommand{\rhabRidgeLoLiqClosePct}{$-$0.8}  % QLIKE change % vs full
\newcommand{\rhabRidgeLoLiqCloseDm}{$-$1.14}  % DM vs the reference arm
\newcommand{\rhabRidgeLoLiqCloseP}{0.256}  % p-value
\newcommand{\rhabRidgeLoLiqCloseLateQlike}{0.1251}  % QLIKE ridge lo_liq c16_2020
\newcommand{\rhabRidgeLoLiqCloseLateMse}{8.99}  % MSE x 1e10 ridge lo_liq c16_2020
\newcommand{\rhabRidgeLoLiqCloseLatePct}{$-$0.7}  % QLIKE change % vs full
\newcommand{\rhabRidgeLoLiqCloseLateDm}{$-$0.77}  % DM vs the reference arm
\newcommand{\rhabRidgeLoLiqCloseLateP}{0.442}  % p-value
\newcommand{\rhabRidgeLoRmAllQlike}{0.1822}  % QLIKE ridge lo_rm all
\newcommand{\rhabRidgeLoRmAllMse}{1.82}  % MSE x 1e10 ridge lo_rm all
\newcommand{\rhabRidgeLoRmAllPct}{1.7}  % QLIKE change % vs full
\newcommand{\rhabRidgeLoRmAllDm}{3.75}  % DM vs the reference arm
\newcommand{\rhabRidgeLoRmAllP}{\ensuremath{<}0.001}  % p-value
\newcommand{\rhabRidgeLoRmCloseQlike}{0.1292}  % QLIKE ridge lo_rm c16
\newcommand{\rhabRidgeLoRmCloseMse}{2.26}  % MSE x 1e10 ridge lo_rm c16
\newcommand{\rhabRidgeLoRmClosePct}{1.0}  % QLIKE change % vs full
\newcommand{\rhabRidgeLoRmCloseDm}{0.73}  % DM vs the reference arm
\newcommand{\rhabRidgeLoRmCloseP}{0.466}  % p-value
\newcommand{\rhabRidgeLoRmCloseLateQlike}{0.1245}  % QLIKE ridge lo_rm c16_2020
\newcommand{\rhabRidgeLoRmCloseLateMse}{4.97}  % MSE x 1e10 ridge lo_rm c16_2020
\newcommand{\rhabRidgeLoRmCloseLatePct}{$-$1.2}  % QLIKE change % vs full
\newcommand{\rhabRidgeLoRmCloseLateDm}{$-$0.55}  % DM vs the reference arm
\newcommand{\rhabRidgeLoRmCloseLateP}{0.579}  % p-value
\newcommand{\rhabRidgeLoSentAllQlike}{0.1793}  % QLIKE ridge lo_sent all
\newcommand{\rhabRidgeLoSentAllMse}{2.24}  % MSE x 1e10 ridge lo_sent all
\newcommand{\rhabRidgeLoSentAllPct}{0.1}  % QLIKE change % vs full
\newcommand{\rhabRidgeLoSentAllDm}{0.55}  % DM vs the reference arm
\newcommand{\rhabRidgeLoSentAllP}{0.584}  % p-value
\newcommand{\rhabRidgeLoSentCloseQlike}{0.1261}  % QLIKE ridge lo_sent c16
\newcommand{\rhabRidgeLoSentCloseMse}{3.51}  % MSE x 1e10 ridge lo_sent c16
\newcommand{\rhabRidgeLoSentClosePct}{$-$1.4}  % QLIKE change % vs full
\newcommand{\rhabRidgeLoSentCloseDm}{$-$3.16}  % DM vs the reference arm
\newcommand{\rhabRidgeLoSentCloseP}{0.002}  % p-value
\newcommand{\rhabRidgeLoSentCloseLateQlike}{0.1218}  % QLIKE ridge lo_sent c16_2020
\newcommand{\rhabRidgeLoSentCloseLateMse}{7.97}  % MSE x 1e10 ridge lo_sent c16_2020
\newcommand{\rhabRidgeLoSentCloseLatePct}{$-$3.3}  % QLIKE change % vs full
\newcommand{\rhabRidgeLoSentCloseLateDm}{$-$3.21}  % DM vs the reference arm
\newcommand{\rhabRidgeLoSentCloseLateP}{0.001}  % p-value
\newcommand{\rhabRidgeLoVdAllQlike}{0.1789}  % QLIKE ridge lo_vd all
\newcommand{\rhabRidgeLoVdAllMse}{2.24}  % MSE x 1e10 ridge lo_vd all
\newcommand{\rhabRidgeLoVdAllPct}{$-$0.1}  % QLIKE change % vs full
\newcommand{\rhabRidgeLoVdAllDm}{$-$0.98}  % DM vs the reference arm
\newcommand{\rhabRidgeLoVdAllP}{0.326}  % p-value
\newcommand{\rhabRidgeLoVdCloseQlike}{0.1266}  % QLIKE ridge lo_vd c16
\newcommand{\rhabRidgeLoVdCloseMse}{3.51}  % MSE x 1e10 ridge lo_vd c16
\newcommand{\rhabRidgeLoVdClosePct}{$-$1.0}  % QLIKE change % vs full
\newcommand{\rhabRidgeLoVdCloseDm}{$-$2.17}  % DM vs the reference arm
\newcommand{\rhabRidgeLoVdCloseP}{0.030}  % p-value
\newcommand{\rhabRidgeLoVdCloseLateQlike}{0.1232}  % QLIKE ridge lo_vd c16_2020
\newcommand{\rhabRidgeLoVdCloseLateMse}{7.98}  % MSE x 1e10 ridge lo_vd c16_2020
\newcommand{\rhabRidgeLoVdCloseLatePct}{$-$2.2}  % QLIKE change % vs full
\newcommand{\rhabRidgeLoVdCloseLateDm}{$-$2.26}  % DM vs the reference arm
\newcommand{\rhabRidgeLoVdCloseLateP}{0.024}  % p-value
\newcommand{\rhabRidgeLoVwAllQlike}{0.1789}  % QLIKE ridge lo_vw all
\newcommand{\rhabRidgeLoVwAllMse}{1.97}  % MSE x 1e10 ridge lo_vw all
\newcommand{\rhabRidgeLoVwAllPct}{$-$0.1}  % QLIKE change % vs full
\newcommand{\rhabRidgeLoVwAllDm}{$-$0.70}  % DM vs the reference arm
\newcommand{\rhabRidgeLoVwAllP}{0.484}  % p-value
\newcommand{\rhabRidgeLoVwCloseQlike}{0.1274}  % QLIKE ridge lo_vw c16
\newcommand{\rhabRidgeLoVwCloseMse}{3.50}  % MSE x 1e10 ridge lo_vw c16
\newcommand{\rhabRidgeLoVwClosePct}{$-$0.4}  % QLIKE change % vs full
\newcommand{\rhabRidgeLoVwCloseDm}{$-$0.90}  % DM vs the reference arm
\newcommand{\rhabRidgeLoVwCloseP}{0.366}  % p-value
\newcommand{\rhabRidgeLoVwCloseLateQlike}{0.1246}  % QLIKE ridge lo_vw c16_2020
\newcommand{\rhabRidgeLoVwCloseLateMse}{7.95}  % MSE x 1e10 ridge lo_vw c16_2020
\newcommand{\rhabRidgeLoVwCloseLatePct}{$-$1.1}  % QLIKE change % vs full
\newcommand{\rhabRidgeLoVwCloseLateDm}{$-$1.96}  % DM vs the reference arm
\newcommand{\rhabRidgeLoVwCloseLateP}{0.050}  % p-value
\newcommand{\rhabRidgeSharAllQlike}{0.1876}  % QLIKE ridge shar all
\newcommand{\rhabRidgeSharAllMse}{1.35}  % MSE x 1e10 ridge shar all
\newcommand{\rhabRidgeSharAllPct}{$-$0.9}  % QLIKE change % vs har
\newcommand{\rhabRidgeSharAllDm}{$-$3.95}  % DM vs the reference arm
\newcommand{\rhabRidgeSharAllP}{\ensuremath{<}0.001}  % p-value
\newcommand{\rhabRidgeSharCloseQlike}{0.1264}  % QLIKE ridge shar c16
\newcommand{\rhabRidgeSharCloseMse}{1.81}  % MSE x 1e10 ridge shar c16
\newcommand{\rhabRidgeSharClosePct}{$-$0.4}  % QLIKE change % vs har
\newcommand{\rhabRidgeSharCloseDm}{$-$0.69}  % DM vs the reference arm
\newcommand{\rhabRidgeSharCloseP}{0.488}  % p-value
\newcommand{\rhabRidgeSharCloseLateQlike}{0.1139}  % QLIKE ridge shar c16_2020
\newcommand{\rhabRidgeSharCloseLateMse}{3.87}  % MSE x 1e10 ridge shar c16_2020
\newcommand{\rhabRidgeSharCloseLatePct}{0.2}  % QLIKE change % vs har
\newcommand{\rhabRidgeSharCloseLateDm}{0.23}  % DM vs the reference arm
\newcommand{\rhabRidgeSharCloseLateP}{0.821}  % p-value
\newcommand{\rhabRidgeSharqlAllQlike}{0.1840}  % QLIKE ridge sharql all
\newcommand{\rhabRidgeSharqlAllMse}{1.68}  % MSE x 1e10 ridge sharql all
\newcommand{\rhabRidgeSharqlAllPct}{$-$2.8}  % QLIKE change % vs har
\newcommand{\rhabRidgeSharqlAllDm}{$-$6.82}  % DM vs the reference arm
\newcommand{\rhabRidgeSharqlAllP}{\ensuremath{<}0.001}  % p-value
\newcommand{\rhabRidgeSharqlCloseQlike}{0.1216}  % QLIKE ridge sharql c16
\newcommand{\rhabRidgeSharqlCloseMse}{2.65}  % MSE x 1e10 ridge sharql c16
\newcommand{\rhabRidgeSharqlClosePct}{$-$4.2}  % QLIKE change % vs har
\newcommand{\rhabRidgeSharqlCloseDm}{$-$3.49}  % DM vs the reference arm
\newcommand{\rhabRidgeSharqlCloseP}{\ensuremath{<}0.001}  % p-value
\newcommand{\rhabRidgeSharqlCloseLateQlike}{0.1119}  % QLIKE ridge sharql c16_2020
\newcommand{\rhabRidgeSharqlCloseLateMse}{5.92}  % MSE x 1e10 ridge sharql c16_2020
\newcommand{\rhabRidgeSharqlCloseLatePct}{$-$1.6}  % QLIKE change % vs har
\newcommand{\rhabRidgeSharqlCloseLateDm}{$-$0.73}  % DM vs the reference arm
\newcommand{\rhabRidgeSharqlCloseLateP}{0.468}  % p-value
\newcommand{\rhabXgbFullAllQlike}{0.1627}  % QLIKE xgb full all
\newcommand{\rhabXgbFullAllMse}{1.34}  % MSE x 1e10 xgb full all
\newcommand{\rhabXgbFullAllPct}{$-$13.8}  % QLIKE change % vs har
\newcommand{\rhabXgbFullAllDm}{$-$9.94}  % DM vs the reference arm
\newcommand{\rhabXgbFullAllP}{\ensuremath{<}0.001}  % p-value
\newcommand{\rhabXgbFullCloseQlike}{0.1192}  % QLIKE xgb full c16
\newcommand{\rhabXgbFullCloseMse}{2.20}  % MSE x 1e10 xgb full c16
\newcommand{\rhabXgbFullClosePct}{$-$3.4}  % QLIKE change % vs har
\newcommand{\rhabXgbFullCloseDm}{$-$2.18}  % DM vs the reference arm
\newcommand{\rhabXgbFullCloseP}{0.030}  % p-value
\newcommand{\rhabXgbFullCloseLateQlike}{0.1079}  % QLIKE xgb full c16_2020
\newcommand{\rhabXgbFullCloseLateMse}{4.79}  % MSE x 1e10 xgb full c16_2020
\newcommand{\rhabXgbFullCloseLatePct}{$-$2.1}  % QLIKE change % vs har
\newcommand{\rhabXgbFullCloseLateDm}{$-$0.76}  % DM vs the reference arm
\newcommand{\rhabXgbFullCloseLateP}{0.444}  % p-value
\newcommand{\rhabXgbHarAllQlike}{0.1888}  % QLIKE xgb har all
\newcommand{\rhabXgbHarAllMse}{1.16}  % MSE x 1e10 xgb har all
\newcommand{\rhabXgbHarCloseQlike}{0.1234}  % QLIKE xgb har c16
\newcommand{\rhabXgbHarCloseMse}{1.73}  % MSE x 1e10 xgb har c16
\newcommand{\rhabXgbHarCloseLateQlike}{0.1102}  % QLIKE xgb har c16_2020
\newcommand{\rhabXgbHarCloseLateMse}{3.62}  % MSE x 1e10 xgb har c16_2020
\newcommand{\rhabRidgeFullLiveSentQlike}{0.1833}  % QLIKE ridge full live_sent (subsample)
\newcommand{\rhabRidgeFullLiveSentPct}{$-$1.9}  % QLIKE change % vs the reference (subsample)
\newcommand{\rhabRidgeFullLiveSentDm}{$-$0.98}  % DM vs the reference arm (subsample)
\newcommand{\rhabRidgeFullLiveSentP}{0.327}  % p-value
\newcommand{\rhabRidgeLoSentLiveSentQlike}{0.1842}  % QLIKE ridge lo_sent live_sent (subsample)
\newcommand{\rhabRidgeLoSentLiveSentPct}{0.5}  % QLIKE change % vs the reference (subsample)
\newcommand{\rhabRidgeLoSentLiveSentDm}{0.71}  % DM vs the reference arm (subsample)
\newcommand{\rhabRidgeLoSentLiveSentP}{0.477}  % p-value
\newcommand{\rhabRidgeAddSentLiveSentQlike}{0.1846}  % QLIKE ridge add_sent live_sent (subsample)
\newcommand{\rhabRidgeAddSentLiveSentPct}{$-$1.2}  % QLIKE change % vs the reference (subsample)
\newcommand{\rhabRidgeAddSentLiveSentDm}{$-$1.18}  % DM vs the reference arm (subsample)
\newcommand{\rhabRidgeAddSentLiveSentP}{0.238}  % p-value
\newcommand{\rhabRidgeHarLiveSentQlike}{0.1868}  % QLIKE ridge har live_sent (subsample)
\newcommand{\rhabLassoFullLiveSentQlike}{0.1780}  % QLIKE lasso full live_sent (subsample)
\newcommand{\rhabLassoFullLiveSentPct}{$-$4.7}  % QLIKE change % vs the reference (subsample)
\newcommand{\rhabLassoFullLiveSentDm}{$-$1.83}  % DM vs the reference arm (subsample)
\newcommand{\rhabLassoFullLiveSentP}{0.068}  % p-value
\newcommand{\rhabLassoLoSentLiveSentQlike}{0.1786}  % QLIKE lasso lo_sent live_sent (subsample)
\newcommand{\rhabLassoLoSentLiveSentPct}{0.3}  % QLIKE change % vs the reference (subsample)
\newcommand{\rhabLassoLoSentLiveSentDm}{0.41}  % DM vs the reference arm (subsample)
\newcommand{\rhabLassoLoSentLiveSentP}{0.685}  % p-value
\newcommand{\rhabLassoAddSentLiveSentQlike}{0.1850}  % QLIKE lasso add_sent live_sent (subsample)
\newcommand{\rhabLassoAddSentLiveSentPct}{$-$1.0}  % QLIKE change % vs the reference (subsample)
\newcommand{\rhabLassoAddSentLiveSentDm}{$-$0.96}  % DM vs the reference arm (subsample)
\newcommand{\rhabLassoAddSentLiveSentP}{0.339}  % p-value
\newcommand{\rhabLassoHarLiveSentQlike}{0.1868}  % QLIKE lasso har live_sent (subsample)
\newcommand{\rhabEnetFullLiveSentQlike}{0.1818}  % QLIKE enet full live_sent (subsample)
\newcommand{\rhabEnetFullLiveSentPct}{$-$2.7}  % QLIKE change % vs the reference (subsample)
\newcommand{\rhabEnetFullLiveSentDm}{$-$1.55}  % DM vs the reference arm (subsample)
\newcommand{\rhabEnetFullLiveSentP}{0.121}  % p-value
\newcommand{\rhabEnetLoSentLiveSentQlike}{0.1833}  % QLIKE enet lo_sent live_sent (subsample)
\newcommand{\rhabEnetLoSentLiveSentPct}{0.8}  % QLIKE change % vs the reference (subsample)
\newcommand{\rhabEnetLoSentLiveSentDm}{1.43}  % DM vs the reference arm (subsample)
\newcommand{\rhabEnetLoSentLiveSentP}{0.154}  % p-value
\newcommand{\rhabEnetAddSentLiveSentQlike}{0.1845}  % QLIKE enet add_sent live_sent (subsample)
\newcommand{\rhabEnetAddSentLiveSentPct}{$-$1.3}  % QLIKE change % vs the reference (subsample)
\newcommand{\rhabEnetAddSentLiveSentDm}{$-$1.48}  % DM vs the reference arm (subsample)
\newcommand{\rhabEnetAddSentLiveSentP}{0.139}  % p-value
\newcommand{\rhabEnetHarLiveSentQlike}{0.1868}  % QLIKE enet har live_sent (subsample)
\newcommand{\rhabLgbmFullLiveSentQlike}{0.1564}  % QLIKE lgbm full live_sent (subsample)
\newcommand{\rhabLgbmFullLiveSentPct}{$-$13.8}  % QLIKE change % vs the reference (subsample)
\newcommand{\rhabLgbmFullLiveSentDm}{$-$3.76}  % DM vs the reference arm (subsample)
\newcommand{\rhabLgbmFullLiveSentP}{\ensuremath{<}0.001}  % p-value
\newcommand{\rhabLgbmHarLiveSentQlike}{0.1814}  % QLIKE lgbm har live_sent (subsample)
\newcommand{\rhabXgbFullLiveSentQlike}{0.1559}  % QLIKE xgb full live_sent (subsample)
\newcommand{\rhabXgbFullLiveSentPct}{$-$15.6}  % QLIKE change % vs the reference (subsample)
\newcommand{\rhabXgbFullLiveSentDm}{$-$3.69}  % DM vs the reference arm (subsample)
\newcommand{\rhabXgbFullLiveSentP}{\ensuremath{<}0.001}  % p-value
\newcommand{\rhabXgbHarLiveSentQlike}{0.1846}  % QLIKE xgb har live_sent (subsample)
\newcommand{\rhabRfFullLiveSentQlike}{0.1676}  % QLIKE rf full live_sent (subsample)
\newcommand{\rhabRfFullLiveSentPct}{$-$10.8}  % QLIKE change % vs the reference (subsample)
\newcommand{\rhabRfFullLiveSentDm}{$-$4.09}  % DM vs the reference arm (subsample)
\newcommand{\rhabRfFullLiveSentP}{\ensuremath{<}0.001}  % p-value
\newcommand{\rhabRfHarLiveSentQlike}{0.1878}  % QLIKE rf har live_sent (subsample)
\newcommand{\rhabRidgeFullLiveSentCloseQlike}{0.1222}  % QLIKE ridge full live16_sent (subsample)
\newcommand{\rhabRidgeFullLiveSentClosePct}{12.8}  % QLIKE change % vs the reference (subsample)
\newcommand{\rhabRidgeFullLiveSentCloseDm}{2.03}  % DM vs the reference arm (subsample)
\newcommand{\rhabRidgeFullLiveSentCloseP}{0.042}  % p-value
\newcommand{\rhabRidgeLoSentLiveSentCloseQlike}{0.1141}  % QLIKE ridge lo_sent live16_sent (subsample)
\newcommand{\rhabRidgeLoSentLiveSentClosePct}{$-$6.6}  % QLIKE change % vs the reference (subsample)
\newcommand{\rhabRidgeLoSentLiveSentCloseDm}{$-$3.48}  % DM vs the reference arm (subsample)
\newcommand{\rhabRidgeLoSentLiveSentCloseP}{\ensuremath{<}0.001}  % p-value
\newcommand{\rhabRidgeAddSentLiveSentCloseQlike}{0.1188}  % QLIKE ridge add_sent live16_sent (subsample)
\newcommand{\rhabRidgeAddSentLiveSentClosePct}{9.7}  % QLIKE change % vs the reference (subsample)
\newcommand{\rhabRidgeAddSentLiveSentCloseDm}{2.63}  % DM vs the reference arm (subsample)
\newcommand{\rhabRidgeAddSentLiveSentCloseP}{0.009}  % p-value
\newcommand{\rhabRidgeHarLiveSentCloseQlike}{0.1083}  % QLIKE ridge har live16_sent (subsample)
\newcommand{\rhabLassoFullLiveSentCloseQlike}{0.1234}  % QLIKE lasso full live16_sent (subsample)
\newcommand{\rhabLassoFullLiveSentClosePct}{13.9}  % QLIKE change % vs the reference (subsample)
\newcommand{\rhabLassoFullLiveSentCloseDm}{2.17}  % DM vs the reference arm (subsample)
\newcommand{\rhabLassoFullLiveSentCloseP}{0.030}  % p-value
\newcommand{\rhabLassoLoSentLiveSentCloseQlike}{0.1155}  % QLIKE lasso lo_sent live16_sent (subsample)
\newcommand{\rhabLassoLoSentLiveSentClosePct}{$-$6.3}  % QLIKE change % vs the reference (subsample)
\newcommand{\rhabLassoLoSentLiveSentCloseDm}{$-$2.92}  % DM vs the reference arm (subsample)
\newcommand{\rhabLassoLoSentLiveSentCloseP}{0.003}  % p-value
\newcommand{\rhabLassoAddSentLiveSentCloseQlike}{0.1195}  % QLIKE lasso add_sent live16_sent (subsample)
\newcommand{\rhabLassoAddSentLiveSentClosePct}{10.3}  % QLIKE change % vs the reference (subsample)
\newcommand{\rhabLassoAddSentLiveSentCloseDm}{2.76}  % DM vs the reference arm (subsample)
\newcommand{\rhabLassoAddSentLiveSentCloseP}{0.006}  % p-value
\newcommand{\rhabLassoHarLiveSentCloseQlike}{0.1083}  % QLIKE lasso har live16_sent (subsample)
\newcommand{\rhabEnetFullLiveSentCloseQlike}{0.1171}  % QLIKE enet full live16_sent (subsample)
\newcommand{\rhabEnetFullLiveSentClosePct}{8.1}  % QLIKE change % vs the reference (subsample)
\newcommand{\rhabEnetFullLiveSentCloseDm}{1.35}  % DM vs the reference arm (subsample)
\newcommand{\rhabEnetFullLiveSentCloseP}{0.176}  % p-value
\newcommand{\rhabEnetLoSentLiveSentCloseQlike}{0.1134}  % QLIKE enet lo_sent live16_sent (subsample)
\newcommand{\rhabEnetLoSentLiveSentClosePct}{$-$3.2}  % QLIKE change % vs the reference (subsample)
\newcommand{\rhabEnetLoSentLiveSentCloseDm}{$-$2.23}  % DM vs the reference arm (subsample)
\newcommand{\rhabEnetLoSentLiveSentCloseP}{0.026}  % p-value
\newcommand{\rhabEnetAddSentLiveSentCloseQlike}{0.1156}  % QLIKE enet add_sent live16_sent (subsample)
\newcommand{\rhabEnetAddSentLiveSentClosePct}{6.7}  % QLIKE change % vs the reference (subsample)
\newcommand{\rhabEnetAddSentLiveSentCloseDm}{1.91}  % DM vs the reference arm (subsample)
\newcommand{\rhabEnetAddSentLiveSentCloseP}{0.056}  % p-value
\newcommand{\rhabEnetHarLiveSentCloseQlike}{0.1083}  % QLIKE enet har live16_sent (subsample)
\newcommand{\rhabLgbmFullLiveSentCloseQlike}{0.0995}  % QLIKE lgbm full live16_sent (subsample)
\newcommand{\rhabLgbmFullLiveSentClosePct}{$-$2.5}  % QLIKE change % vs the reference (subsample)
\newcommand{\rhabLgbmFullLiveSentCloseDm}{$-$0.57}  % DM vs the reference arm (subsample)
\newcommand{\rhabLgbmFullLiveSentCloseP}{0.565}  % p-value
\newcommand{\rhabLgbmHarLiveSentCloseQlike}{0.1021}  % QLIKE lgbm har live16_sent (subsample)
\newcommand{\rhabXgbFullLiveSentCloseQlike}{0.0997}  % QLIKE xgb full live16_sent (subsample)
\newcommand{\rhabXgbFullLiveSentClosePct}{$-$2.8}  % QLIKE change % vs the reference (subsample)
\newcommand{\rhabXgbFullLiveSentCloseDm}{$-$0.61}  % DM vs the reference arm (subsample)
\newcommand{\rhabXgbFullLiveSentCloseP}{0.541}  % p-value
\newcommand{\rhabXgbHarLiveSentCloseQlike}{0.1025}  % QLIKE xgb har live16_sent (subsample)
\newcommand{\rhabRfFullLiveSentCloseQlike}{0.1026}  % QLIKE rf full live16_sent (subsample)
\newcommand{\rhabRfFullLiveSentClosePct}{$-$2.6}  % QLIKE change % vs the reference (subsample)
\newcommand{\rhabRfFullLiveSentCloseDm}{$-$0.65}  % DM vs the reference arm (subsample)
\newcommand{\rhabRfFullLiveSentCloseP}{0.513}  % p-value
\newcommand{\rhabRfHarLiveSentCloseQlike}{0.1053}  % QLIKE rf har live16_sent (subsample)
\newcommand{\rhabRidgeFullLiveIvQlike}{0.1940}  % QLIKE ridge full live_iv (subsample)
\newcommand{\rhabRidgeFullLiveIvPct}{$-$4.7}  % QLIKE change % vs the reference (subsample)
\newcommand{\rhabRidgeFullLiveIvDm}{$-$3.62}  % DM vs the reference arm (subsample)
\newcommand{\rhabRidgeFullLiveIvP}{\ensuremath{<}0.001}  % p-value
\newcommand{\rhabRidgeLoIvLiveIvQlike}{0.1942}  % QLIKE ridge lo_iv live_iv (subsample)
\newcommand{\rhabRidgeLoIvLiveIvPct}{0.1}  % QLIKE change % vs the reference (subsample)
\newcommand{\rhabRidgeLoIvLiveIvDm}{0.25}  % DM vs the reference arm (subsample)
\newcommand{\rhabRidgeLoIvLiveIvP}{0.806}  % p-value
\newcommand{\rhabRidgeAddIvLiveIvQlike}{0.2037}  % QLIKE ridge add_iv live_iv (subsample)
\newcommand{\rhabRidgeAddIvLiveIvPct}{0.0}  % QLIKE change % vs the reference (subsample)
\newcommand{\rhabRidgeAddIvLiveIvDm}{0.02}  % DM vs the reference arm (subsample)
\newcommand{\rhabRidgeAddIvLiveIvP}{0.986}  % p-value
\newcommand{\rhabRidgeHarLiveIvQlike}{0.2037}  % QLIKE ridge har live_iv (subsample)
\newcommand{\rhabLassoFullLiveIvQlike}{0.1924}  % QLIKE lasso full live_iv (subsample)
\newcommand{\rhabLassoFullLiveIvPct}{$-$5.5}  % QLIKE change % vs the reference (subsample)
\newcommand{\rhabLassoFullLiveIvDm}{$-$2.98}  % DM vs the reference arm (subsample)
\newcommand{\rhabLassoFullLiveIvP}{0.003}  % p-value
\newcommand{\rhabLassoLoIvLiveIvQlike}{0.1931}  % QLIKE lasso lo_iv live_iv (subsample)
\newcommand{\rhabLassoLoIvLiveIvPct}{0.4}  % QLIKE change % vs the reference (subsample)
\newcommand{\rhabLassoLoIvLiveIvDm}{0.47}  % DM vs the reference arm (subsample)
\newcommand{\rhabLassoLoIvLiveIvP}{0.641}  % p-value
\newcommand{\rhabLassoAddIvLiveIvQlike}{0.2019}  % QLIKE lasso add_iv live_iv (subsample)
\newcommand{\rhabLassoAddIvLiveIvPct}{$-$0.9}  % QLIKE change % vs the reference (subsample)
\newcommand{\rhabLassoAddIvLiveIvDm}{$-$1.39}  % DM vs the reference arm (subsample)
\newcommand{\rhabLassoAddIvLiveIvP}{0.165}  % p-value
\newcommand{\rhabLassoHarLiveIvQlike}{0.2037}  % QLIKE lasso har live_iv (subsample)
\newcommand{\rhabEnetFullLiveIvQlike}{0.1921}  % QLIKE enet full live_iv (subsample)
\newcommand{\rhabEnetFullLiveIvPct}{$-$5.7}  % QLIKE change % vs the reference (subsample)
\newcommand{\rhabEnetFullLiveIvDm}{$-$4.87}  % DM vs the reference arm (subsample)
\newcommand{\rhabEnetFullLiveIvP}{\ensuremath{<}0.001}  % p-value
\newcommand{\rhabEnetLoIvLiveIvQlike}{0.1924}  % QLIKE enet lo_iv live_iv (subsample)
\newcommand{\rhabEnetLoIvLiveIvPct}{0.1}  % QLIKE change % vs the reference (subsample)
\newcommand{\rhabEnetLoIvLiveIvDm}{0.23}  % DM vs the reference arm (subsample)
\newcommand{\rhabEnetLoIvLiveIvP}{0.816}  % p-value
\newcommand{\rhabEnetAddIvLiveIvQlike}{0.2033}  % QLIKE enet add_iv live_iv (subsample)
\newcommand{\rhabEnetAddIvLiveIvPct}{$-$0.2}  % QLIKE change % vs the reference (subsample)
\newcommand{\rhabEnetAddIvLiveIvDm}{$-$0.36}  % DM vs the reference arm (subsample)
\newcommand{\rhabEnetAddIvLiveIvP}{0.722}  % p-value
\newcommand{\rhabEnetHarLiveIvQlike}{0.2037}  % QLIKE enet har live_iv (subsample)
\newcommand{\rhabLgbmFullLiveIvQlike}{0.1739}  % QLIKE lgbm full live_iv (subsample)
\newcommand{\rhabLgbmFullLiveIvPct}{$-$13.3}  % QLIKE change % vs the reference (subsample)
\newcommand{\rhabLgbmFullLiveIvDm}{$-$7.31}  % DM vs the reference arm (subsample)
\newcommand{\rhabLgbmFullLiveIvP}{\ensuremath{<}0.001}  % p-value
\newcommand{\rhabLgbmHarLiveIvQlike}{0.2006}  % QLIKE lgbm har live_iv (subsample)
\newcommand{\rhabXgbFullLiveIvQlike}{0.1739}  % QLIKE xgb full live_iv (subsample)
\newcommand{\rhabXgbFullLiveIvPct}{$-$14.1}  % QLIKE change % vs the reference (subsample)
\newcommand{\rhabXgbFullLiveIvDm}{$-$7.16}  % DM vs the reference arm (subsample)
\newcommand{\rhabXgbFullLiveIvP}{\ensuremath{<}0.001}  % p-value
\newcommand{\rhabXgbHarLiveIvQlike}{0.2024}  % QLIKE xgb har live_iv (subsample)
\newcommand{\rhabRfFullLiveIvQlike}{0.1818}  % QLIKE rf full live_iv (subsample)
\newcommand{\rhabRfFullLiveIvPct}{$-$10.4}  % QLIKE change % vs the reference (subsample)
\newcommand{\rhabRfFullLiveIvDm}{$-$7.47}  % DM vs the reference arm (subsample)
\newcommand{\rhabRfFullLiveIvP}{\ensuremath{<}0.001}  % p-value
\newcommand{\rhabRfHarLiveIvQlike}{0.2029}  % QLIKE rf har live_iv (subsample)
\newcommand{\rhabRidgeFullLiveIvCloseQlike}{0.1380}  % QLIKE ridge full live16_iv (subsample)
\newcommand{\rhabRidgeFullLiveIvClosePct}{6.0}  % QLIKE change % vs the reference (subsample)
\newcommand{\rhabRidgeFullLiveIvCloseDm}{1.50}  % DM vs the reference arm (subsample)
\newcommand{\rhabRidgeFullLiveIvCloseP}{0.134}  % p-value
\newcommand{\rhabRidgeLoIvLiveIvCloseQlike}{0.1355}  % QLIKE ridge lo_iv live16_iv (subsample)
\newcommand{\rhabRidgeLoIvLiveIvClosePct}{$-$1.8}  % QLIKE change % vs the reference (subsample)
\newcommand{\rhabRidgeLoIvLiveIvCloseDm}{$-$1.22}  % DM vs the reference arm (subsample)
\newcommand{\rhabRidgeLoIvLiveIvCloseP}{0.222}  % p-value
\newcommand{\rhabRidgeAddIvLiveIvCloseQlike}{0.1306}  % QLIKE ridge add_iv live16_iv (subsample)
\newcommand{\rhabRidgeAddIvLiveIvClosePct}{0.3}  % QLIKE change % vs the reference (subsample)
\newcommand{\rhabRidgeAddIvLiveIvCloseDm}{0.19}  % DM vs the reference arm (subsample)
\newcommand{\rhabRidgeAddIvLiveIvCloseP}{0.853}  % p-value
\newcommand{\rhabRidgeHarLiveIvCloseQlike}{0.1302}  % QLIKE ridge har live16_iv (subsample)
\newcommand{\rhabLassoFullLiveIvCloseQlike}{0.1425}  % QLIKE lasso full live16_iv (subsample)
\newcommand{\rhabLassoFullLiveIvClosePct}{9.4}  % QLIKE change % vs the reference (subsample)
\newcommand{\rhabLassoFullLiveIvCloseDm}{1.93}  % DM vs the reference arm (subsample)
\newcommand{\rhabLassoFullLiveIvCloseP}{0.054}  % p-value
\newcommand{\rhabLassoLoIvLiveIvCloseQlike}{0.1378}  % QLIKE lasso lo_iv live16_iv (subsample)
\newcommand{\rhabLassoLoIvLiveIvClosePct}{$-$3.3}  % QLIKE change % vs the reference (subsample)
\newcommand{\rhabLassoLoIvLiveIvCloseDm}{$-$1.78}  % DM vs the reference arm (subsample)
\newcommand{\rhabLassoLoIvLiveIvCloseP}{0.075}  % p-value
\newcommand{\rhabLassoAddIvLiveIvCloseQlike}{0.1300}  % QLIKE lasso add_iv live16_iv (subsample)
\newcommand{\rhabLassoAddIvLiveIvClosePct}{$-$0.1}  % QLIKE change % vs the reference (subsample)
\newcommand{\rhabLassoAddIvLiveIvCloseDm}{$-$0.08}  % DM vs the reference arm (subsample)
\newcommand{\rhabLassoAddIvLiveIvCloseP}{0.940}  % p-value
\newcommand{\rhabLassoHarLiveIvCloseQlike}{0.1302}  % QLIKE lasso har live16_iv (subsample)
\newcommand{\rhabEnetFullLiveIvCloseQlike}{0.1355}  % QLIKE enet full live16_iv (subsample)
\newcommand{\rhabEnetFullLiveIvClosePct}{4.1}  % QLIKE change % vs the reference (subsample)
\newcommand{\rhabEnetFullLiveIvCloseDm}{1.12}  % DM vs the reference arm (subsample)
\newcommand{\rhabEnetFullLiveIvCloseP}{0.262}  % p-value
\newcommand{\rhabEnetLoIvLiveIvCloseQlike}{0.1321}  % QLIKE enet lo_iv live16_iv (subsample)
\newcommand{\rhabEnetLoIvLiveIvClosePct}{$-$2.5}  % QLIKE change % vs the reference (subsample)
\newcommand{\rhabEnetLoIvLiveIvCloseDm}{$-$1.58}  % DM vs the reference arm (subsample)
\newcommand{\rhabEnetLoIvLiveIvCloseP}{0.115}  % p-value
\newcommand{\rhabEnetAddIvLiveIvCloseQlike}{0.1308}  % QLIKE enet add_iv live16_iv (subsample)
\newcommand{\rhabEnetAddIvLiveIvClosePct}{0.5}  % QLIKE change % vs the reference (subsample)
\newcommand{\rhabEnetAddIvLiveIvCloseDm}{0.31}  % DM vs the reference arm (subsample)
\newcommand{\rhabEnetAddIvLiveIvCloseP}{0.756}  % p-value
\newcommand{\rhabEnetHarLiveIvCloseQlike}{0.1302}  % QLIKE enet har live16_iv (subsample)
\newcommand{\rhabLgbmFullLiveIvCloseQlike}{0.1242}  % QLIKE lgbm full live16_iv (subsample)
\newcommand{\rhabLgbmFullLiveIvClosePct}{$-$1.6}  % QLIKE change % vs the reference (subsample)
\newcommand{\rhabLgbmFullLiveIvCloseDm}{$-$0.74}  % DM vs the reference arm (subsample)
\newcommand{\rhabLgbmFullLiveIvCloseP}{0.461}  % p-value
\newcommand{\rhabLgbmHarLiveIvCloseQlike}{0.1263}  % QLIKE lgbm har live16_iv (subsample)
\newcommand{\rhabXgbFullLiveIvCloseQlike}{0.1234}  % QLIKE xgb full live16_iv (subsample)
\newcommand{\rhabXgbFullLiveIvClosePct}{$-$2.3}  % QLIKE change % vs the reference (subsample)
\newcommand{\rhabXgbFullLiveIvCloseDm}{$-$1.01}  % DM vs the reference arm (subsample)
\newcommand{\rhabXgbFullLiveIvCloseP}{0.312}  % p-value
\newcommand{\rhabXgbHarLiveIvCloseQlike}{0.1263}  % QLIKE xgb har live16_iv (subsample)
\newcommand{\rhabRfFullLiveIvCloseQlike}{0.1304}  % QLIKE rf full live16_iv (subsample)
\newcommand{\rhabRfFullLiveIvClosePct}{$-$0.4}  % QLIKE change % vs the reference (subsample)
\newcommand{\rhabRfFullLiveIvCloseDm}{$-$0.16}  % DM vs the reference arm (subsample)
\newcommand{\rhabRfFullLiveIvCloseP}{0.870}  % p-value
\newcommand{\rhabRfHarLiveIvCloseQlike}{0.1309}  % QLIKE rf har live16_iv (subsample)
\newcommand{\rhabRidgeFullLiveIvbQlike}{0.1757}  % QLIKE ridge full live_ivb (subsample)
\newcommand{\rhabRidgeFullLiveIvbPct}{$-$3.6}  % QLIKE change % vs the reference (subsample)
\newcommand{\rhabRidgeFullLiveIvbDm}{$-$2.12}  % DM vs the reference arm (subsample)
\newcommand{\rhabRidgeFullLiveIvbP}{0.034}  % p-value
\newcommand{\rhabRidgeLoIvLiveIvbQlike}{0.1746}  % QLIKE ridge lo_iv live_ivb (subsample)
\newcommand{\rhabRidgeLoIvLiveIvbPct}{$-$0.6}  % QLIKE change % vs the reference (subsample)
\newcommand{\rhabRidgeLoIvLiveIvbDm}{$-$0.79}  % DM vs the reference arm (subsample)
\newcommand{\rhabRidgeLoIvLiveIvbP}{0.431}  % p-value
\newcommand{\rhabRidgeAddIvLiveIvbQlike}{0.1850}  % QLIKE ridge add_iv live_ivb (subsample)
\newcommand{\rhabRidgeAddIvLiveIvbPct}{1.4}  % QLIKE change % vs the reference (subsample)
\newcommand{\rhabRidgeAddIvLiveIvbDm}{1.50}  % DM vs the reference arm (subsample)
\newcommand{\rhabRidgeAddIvLiveIvbP}{0.134}  % p-value
\newcommand{\rhabRidgeHarLiveIvbQlike}{0.1824}  % QLIKE ridge har live_ivb (subsample)
\newcommand{\rhabLassoFullLiveIvbQlike}{0.1718}  % QLIKE lasso full live_ivb (subsample)
\newcommand{\rhabLassoFullLiveIvbPct}{$-$5.8}  % QLIKE change % vs the reference (subsample)
\newcommand{\rhabLassoFullLiveIvbDm}{$-$2.58}  % DM vs the reference arm (subsample)
\newcommand{\rhabLassoFullLiveIvbP}{0.010}  % p-value
\newcommand{\rhabLassoLoIvLiveIvbQlike}{0.1700}  % QLIKE lasso lo_iv live_ivb (subsample)
\newcommand{\rhabLassoLoIvLiveIvbPct}{$-$1.0}  % QLIKE change % vs the reference (subsample)
\newcommand{\rhabLassoLoIvLiveIvbDm}{$-$0.96}  % DM vs the reference arm (subsample)
\newcommand{\rhabLassoLoIvLiveIvbP}{0.339}  % p-value
\newcommand{\rhabLassoAddIvLiveIvbQlike}{0.1844}  % QLIKE lasso add_iv live_ivb (subsample)
\newcommand{\rhabLassoAddIvLiveIvbPct}{1.1}  % QLIKE change % vs the reference (subsample)
\newcommand{\rhabLassoAddIvLiveIvbDm}{1.04}  % DM vs the reference arm (subsample)
\newcommand{\rhabLassoAddIvLiveIvbP}{0.300}  % p-value
\newcommand{\rhabLassoHarLiveIvbQlike}{0.1824}  % QLIKE lasso har live_ivb (subsample)
\newcommand{\rhabEnetFullLiveIvbQlike}{0.1745}  % QLIKE enet full live_ivb (subsample)
\newcommand{\rhabEnetFullLiveIvbPct}{$-$4.3}  % QLIKE change % vs the reference (subsample)
\newcommand{\rhabEnetFullLiveIvbDm}{$-$2.67}  % DM vs the reference arm (subsample)
\newcommand{\rhabEnetFullLiveIvbP}{0.008}  % p-value
\newcommand{\rhabEnetLoIvLiveIvbQlike}{0.1732}  % QLIKE enet lo_iv live_ivb (subsample)
\newcommand{\rhabEnetLoIvLiveIvbPct}{$-$0.7}  % QLIKE change % vs the reference (subsample)
\newcommand{\rhabEnetLoIvLiveIvbDm}{$-$0.88}  % DM vs the reference arm (subsample)
\newcommand{\rhabEnetLoIvLiveIvbP}{0.377}  % p-value
\newcommand{\rhabEnetAddIvLiveIvbQlike}{0.1844}  % QLIKE enet add_iv live_ivb (subsample)
\newcommand{\rhabEnetAddIvLiveIvbPct}{1.1}  % QLIKE change % vs the reference (subsample)
\newcommand{\rhabEnetAddIvLiveIvbDm}{1.30}  % DM vs the reference arm (subsample)
\newcommand{\rhabEnetAddIvLiveIvbP}{0.195}  % p-value
\newcommand{\rhabEnetHarLiveIvbQlike}{0.1824}  % QLIKE enet har live_ivb (subsample)
\newcommand{\rhabLgbmFullLiveIvbQlike}{0.1518}  % QLIKE lgbm full live_ivb (subsample)
\newcommand{\rhabLgbmFullLiveIvbPct}{$-$14.6}  % QLIKE change % vs the reference (subsample)
\newcommand{\rhabLgbmFullLiveIvbDm}{$-$5.14}  % DM vs the reference arm (subsample)
\newcommand{\rhabLgbmFullLiveIvbP}{\ensuremath{<}0.001}  % p-value
\newcommand{\rhabLgbmHarLiveIvbQlike}{0.1778}  % QLIKE lgbm har live_ivb (subsample)
\newcommand{\rhabXgbFullLiveIvbQlike}{0.1511}  % QLIKE xgb full live_ivb (subsample)
\newcommand{\rhabXgbFullLiveIvbPct}{$-$16.3}  % QLIKE change % vs the reference (subsample)
\newcommand{\rhabXgbFullLiveIvbDm}{$-$4.98}  % DM vs the reference arm (subsample)
\newcommand{\rhabXgbFullLiveIvbP}{\ensuremath{<}0.001}  % p-value
\newcommand{\rhabXgbHarLiveIvbQlike}{0.1805}  % QLIKE xgb har live_ivb (subsample)
\newcommand{\rhabRfFullLiveIvbQlike}{0.1621}  % QLIKE rf full live_ivb (subsample)
\newcommand{\rhabRfFullLiveIvbPct}{$-$11.3}  % QLIKE change % vs the reference (subsample)
\newcommand{\rhabRfFullLiveIvbDm}{$-$5.43}  % DM vs the reference arm (subsample)
\newcommand{\rhabRfFullLiveIvbP}{\ensuremath{<}0.001}  % p-value
\newcommand{\rhabRfHarLiveIvbQlike}{0.1829}  % QLIKE rf har live_ivb (subsample)
\newcommand{\rhabRidgeFullLiveIvbCloseQlike}{0.1180}  % QLIKE ridge full live16_ivb (subsample)
\newcommand{\rhabRidgeFullLiveIvbClosePct}{15.1}  % QLIKE change % vs the reference (subsample)
\newcommand{\rhabRidgeFullLiveIvbCloseDm}{2.54}  % DM vs the reference arm (subsample)
\newcommand{\rhabRidgeFullLiveIvbCloseP}{0.011}  % p-value
\newcommand{\rhabRidgeLoIvLiveIvbCloseQlike}{0.1123}  % QLIKE ridge lo_iv live16_ivb (subsample)
\newcommand{\rhabRidgeLoIvLiveIvbClosePct}{$-$4.9}  % QLIKE change % vs the reference (subsample)
\newcommand{\rhabRidgeLoIvLiveIvbCloseDm}{$-$1.73}  % DM vs the reference arm (subsample)
\newcommand{\rhabRidgeLoIvLiveIvbCloseP}{0.083}  % p-value
\newcommand{\rhabRidgeAddIvLiveIvbCloseQlike}{0.1073}  % QLIKE ridge add_iv live16_ivb (subsample)
\newcommand{\rhabRidgeAddIvLiveIvbClosePct}{4.6}  % QLIKE change % vs the reference (subsample)
\newcommand{\rhabRidgeAddIvLiveIvbCloseDm}{1.47}  % DM vs the reference arm (subsample)
\newcommand{\rhabRidgeAddIvLiveIvbCloseP}{0.142}  % p-value
\newcommand{\rhabRidgeHarLiveIvbCloseQlike}{0.1026}  % QLIKE ridge har live16_ivb (subsample)
\newcommand{\rhabLassoFullLiveIvbCloseQlike}{0.1202}  % QLIKE lasso full live16_ivb (subsample)
\newcommand{\rhabLassoFullLiveIvbClosePct}{17.2}  % QLIKE change % vs the reference (subsample)
\newcommand{\rhabLassoFullLiveIvbCloseDm}{2.74}  % DM vs the reference arm (subsample)
\newcommand{\rhabLassoFullLiveIvbCloseP}{0.006}  % p-value
\newcommand{\rhabLassoLoIvLiveIvbCloseQlike}{0.1120}  % QLIKE lasso lo_iv live16_ivb (subsample)
\newcommand{\rhabLassoLoIvLiveIvbClosePct}{$-$6.8}  % QLIKE change % vs the reference (subsample)
\newcommand{\rhabLassoLoIvLiveIvbCloseDm}{$-$2.01}  % DM vs the reference arm (subsample)
\newcommand{\rhabLassoLoIvLiveIvbCloseP}{0.044}  % p-value
\newcommand{\rhabLassoAddIvLiveIvbCloseQlike}{0.1066}  % QLIKE lasso add_iv live16_ivb (subsample)
\newcommand{\rhabLassoAddIvLiveIvbClosePct}{3.9}  % QLIKE change % vs the reference (subsample)
\newcommand{\rhabLassoAddIvLiveIvbCloseDm}{1.22}  % DM vs the reference arm (subsample)
\newcommand{\rhabLassoAddIvLiveIvbCloseP}{0.221}  % p-value
\newcommand{\rhabLassoHarLiveIvbCloseQlike}{0.1026}  % QLIKE lasso har live16_ivb (subsample)
\newcommand{\rhabEnetFullLiveIvbCloseQlike}{0.1147}  % QLIKE enet full live16_ivb (subsample)
\newcommand{\rhabEnetFullLiveIvbClosePct}{11.8}  % QLIKE change % vs the reference (subsample)
\newcommand{\rhabEnetFullLiveIvbCloseDm}{2.03}  % DM vs the reference arm (subsample)
\newcommand{\rhabEnetFullLiveIvbCloseP}{0.043}  % p-value
\newcommand{\rhabEnetLoIvLiveIvbCloseQlike}{0.1073}  % QLIKE enet lo_iv live16_ivb (subsample)
\newcommand{\rhabEnetLoIvLiveIvbClosePct}{$-$6.5}  % QLIKE change % vs the reference (subsample)
\newcommand{\rhabEnetLoIvLiveIvbCloseDm}{$-$2.14}  % DM vs the reference arm (subsample)
\newcommand{\rhabEnetLoIvLiveIvbCloseP}{0.033}  % p-value
\newcommand{\rhabEnetAddIvLiveIvbCloseQlike}{0.1076}  % QLIKE enet add_iv live16_ivb (subsample)
\newcommand{\rhabEnetAddIvLiveIvbClosePct}{4.9}  % QLIKE change % vs the reference (subsample)
\newcommand{\rhabEnetAddIvLiveIvbCloseDm}{1.60}  % DM vs the reference arm (subsample)
\newcommand{\rhabEnetAddIvLiveIvbCloseP}{0.110}  % p-value
\newcommand{\rhabEnetHarLiveIvbCloseQlike}{0.1026}  % QLIKE enet har live16_ivb (subsample)
\newcommand{\rhabLgbmFullLiveIvbCloseQlike}{0.0964}  % QLIKE lgbm full live16_ivb (subsample)
\newcommand{\rhabLgbmFullLiveIvbClosePct}{$-$1.0}  % QLIKE change % vs the reference (subsample)
\newcommand{\rhabLgbmFullLiveIvbCloseDm}{$-$0.29}  % DM vs the reference arm (subsample)
\newcommand{\rhabLgbmFullLiveIvbCloseP}{0.775}  % p-value
\newcommand{\rhabLgbmHarLiveIvbCloseQlike}{0.0974}  % QLIKE lgbm har live16_ivb (subsample)
\newcommand{\rhabXgbFullLiveIvbCloseQlike}{0.0962}  % QLIKE xgb full live16_ivb (subsample)
\newcommand{\rhabXgbFullLiveIvbClosePct}{$-$1.5}  % QLIKE change % vs the reference (subsample)
\newcommand{\rhabXgbFullLiveIvbCloseDm}{$-$0.40}  % DM vs the reference arm (subsample)
\newcommand{\rhabXgbFullLiveIvbCloseP}{0.689}  % p-value
\newcommand{\rhabXgbHarLiveIvbCloseQlike}{0.0976}  % QLIKE xgb har live16_ivb (subsample)
\newcommand{\rhabRfFullLiveIvbCloseQlike}{0.1005}  % QLIKE rf full live16_ivb (subsample)
\newcommand{\rhabRfFullLiveIvbClosePct}{0.2}  % QLIKE change % vs the reference (subsample)
\newcommand{\rhabRfFullLiveIvbCloseDm}{0.05}  % DM vs the reference arm (subsample)
\newcommand{\rhabRfFullLiveIvbCloseP}{0.958}  % p-value
\newcommand{\rhabRfHarLiveIvbCloseQlike}{0.1003}  % QLIKE rf har live16_ivb (subsample)
\newcommand{\rhabRidgeFullLiveVdQlike}{0.1942}  % QLIKE ridge full live_vd (subsample)
\newcommand{\rhabRidgeFullLiveVdPct}{4.2}  % QLIKE change % vs the reference (subsample)
\newcommand{\rhabRidgeFullLiveVdDm}{1.33}  % DM vs the reference arm (subsample)
\newcommand{\rhabRidgeFullLiveVdP}{0.183}  % p-value
\newcommand{\rhabRidgeLoVdLiveVdQlike}{0.1887}  % QLIKE ridge lo_vd live_vd (subsample)
\newcommand{\rhabRidgeLoVdLiveVdPct}{$-$2.8}  % QLIKE change % vs the reference (subsample)
\newcommand{\rhabRidgeLoVdLiveVdDm}{$-$3.17}  % DM vs the reference arm (subsample)
\newcommand{\rhabRidgeLoVdLiveVdP}{0.002}  % p-value
\newcommand{\rhabRidgeAddVdLiveVdQlike}{0.1944}  % QLIKE ridge add_vd live_vd (subsample)
\newcommand{\rhabRidgeAddVdLiveVdPct}{4.3}  % QLIKE change % vs the reference (subsample)
\newcommand{\rhabRidgeAddVdLiveVdDm}{2.52}  % DM vs the reference arm (subsample)
\newcommand{\rhabRidgeAddVdLiveVdP}{0.012}  % p-value
\newcommand{\rhabRidgeHarLiveVdQlike}{0.1864}  % QLIKE ridge har live_vd (subsample)
\newcommand{\rhabLassoFullLiveVdQlike}{0.1898}  % QLIKE lasso full live_vd (subsample)
\newcommand{\rhabLassoFullLiveVdPct}{1.9}  % QLIKE change % vs the reference (subsample)
\newcommand{\rhabLassoFullLiveVdDm}{0.45}  % DM vs the reference arm (subsample)
\newcommand{\rhabLassoFullLiveVdP}{0.651}  % p-value
\newcommand{\rhabLassoLoVdLiveVdQlike}{0.1846}  % QLIKE lasso lo_vd live_vd (subsample)
\newcommand{\rhabLassoLoVdLiveVdPct}{$-$2.7}  % QLIKE change % vs the reference (subsample)
\newcommand{\rhabLassoLoVdLiveVdDm}{$-$2.51}  % DM vs the reference arm (subsample)
\newcommand{\rhabLassoLoVdLiveVdP}{0.012}  % p-value
\newcommand{\rhabLassoAddVdLiveVdQlike}{0.1941}  % QLIKE lasso add_vd live_vd (subsample)
\newcommand{\rhabLassoAddVdLiveVdPct}{4.2}  % QLIKE change % vs the reference (subsample)
\newcommand{\rhabLassoAddVdLiveVdDm}{2.49}  % DM vs the reference arm (subsample)
\newcommand{\rhabLassoAddVdLiveVdP}{0.013}  % p-value
\newcommand{\rhabLassoHarLiveVdQlike}{0.1864}  % QLIKE lasso har live_vd (subsample)
\newcommand{\rhabEnetFullLiveVdQlike}{0.1908}  % QLIKE enet full live_vd (subsample)
\newcommand{\rhabEnetFullLiveVdPct}{2.4}  % QLIKE change % vs the reference (subsample)
\newcommand{\rhabEnetFullLiveVdDm}{0.85}  % DM vs the reference arm (subsample)
\newcommand{\rhabEnetFullLiveVdP}{0.397}  % p-value
\newcommand{\rhabEnetLoVdLiveVdQlike}{0.1869}  % QLIKE enet lo_vd live_vd (subsample)
\newcommand{\rhabEnetLoVdLiveVdPct}{$-$2.1}  % QLIKE change % vs the reference (subsample)
\newcommand{\rhabEnetLoVdLiveVdDm}{$-$3.47}  % DM vs the reference arm (subsample)
\newcommand{\rhabEnetLoVdLiveVdP}{\ensuremath{<}0.001}  % p-value
\newcommand{\rhabEnetAddVdLiveVdQlike}{0.1916}  % QLIKE enet add_vd live_vd (subsample)
\newcommand{\rhabEnetAddVdLiveVdPct}{2.8}  % QLIKE change % vs the reference (subsample)
\newcommand{\rhabEnetAddVdLiveVdDm}{2.20}  % DM vs the reference arm (subsample)
\newcommand{\rhabEnetAddVdLiveVdP}{0.028}  % p-value
\newcommand{\rhabEnetHarLiveVdQlike}{0.1864}  % QLIKE enet har live_vd (subsample)
\newcommand{\rhabLgbmFullLiveVdQlike}{0.1612}  % QLIKE lgbm full live_vd (subsample)
\newcommand{\rhabLgbmFullLiveVdPct}{$-$10.9}  % QLIKE change % vs the reference (subsample)
\newcommand{\rhabLgbmFullLiveVdDm}{$-$2.01}  % DM vs the reference arm (subsample)
\newcommand{\rhabLgbmFullLiveVdP}{0.044}  % p-value
\newcommand{\rhabLgbmHarLiveVdQlike}{0.1809}  % QLIKE lgbm har live_vd (subsample)
\newcommand{\rhabXgbFullLiveVdQlike}{0.1616}  % QLIKE xgb full live_vd (subsample)
\newcommand{\rhabXgbFullLiveVdPct}{$-$12.1}  % QLIKE change % vs the reference (subsample)
\newcommand{\rhabXgbFullLiveVdDm}{$-$2.10}  % DM vs the reference arm (subsample)
\newcommand{\rhabXgbFullLiveVdP}{0.036}  % p-value
\newcommand{\rhabXgbHarLiveVdQlike}{0.1839}  % QLIKE xgb har live_vd (subsample)
\newcommand{\rhabRfFullLiveVdQlike}{0.1704}  % QLIKE rf full live_vd (subsample)
\newcommand{\rhabRfFullLiveVdPct}{$-$9.1}  % QLIKE change % vs the reference (subsample)
\newcommand{\rhabRfFullLiveVdDm}{$-$2.40}  % DM vs the reference arm (subsample)
\newcommand{\rhabRfFullLiveVdP}{0.016}  % p-value
\newcommand{\rhabRfHarLiveVdQlike}{0.1874}  % QLIKE rf har live_vd (subsample)
\newcommand{\rhabRidgeFullLiveVdCloseQlike}{0.1269}  % QLIKE ridge full live16_vd (subsample)
\newcommand{\rhabRidgeFullLiveVdClosePct}{45.8}  % QLIKE change % vs the reference (subsample)
\newcommand{\rhabRidgeFullLiveVdCloseDm}{3.73}  % DM vs the reference arm (subsample)
\newcommand{\rhabRidgeFullLiveVdCloseP}{\ensuremath{<}0.001}  % p-value
\newcommand{\rhabRidgeLoVdLiveVdCloseQlike}{0.1048}  % QLIKE ridge lo_vd live16_vd (subsample)
\newcommand{\rhabRidgeLoVdLiveVdClosePct}{$-$17.4}  % QLIKE change % vs the reference (subsample)
\newcommand{\rhabRidgeLoVdLiveVdCloseDm}{$-$4.97}  % DM vs the reference arm (subsample)
\newcommand{\rhabRidgeLoVdLiveVdCloseP}{\ensuremath{<}0.001}  % p-value
\newcommand{\rhabRidgeAddVdLiveVdCloseQlike}{0.1018}  % QLIKE ridge add_vd live16_vd (subsample)
\newcommand{\rhabRidgeAddVdLiveVdClosePct}{17.0}  % QLIKE change % vs the reference (subsample)
\newcommand{\rhabRidgeAddVdLiveVdCloseDm}{3.66}  % DM vs the reference arm (subsample)
\newcommand{\rhabRidgeAddVdLiveVdCloseP}{\ensuremath{<}0.001}  % p-value
\newcommand{\rhabRidgeHarLiveVdCloseQlike}{0.0870}  % QLIKE ridge har live16_vd (subsample)
\newcommand{\rhabLassoFullLiveVdCloseQlike}{0.1271}  % QLIKE lasso full live16_vd (subsample)
\newcommand{\rhabLassoFullLiveVdClosePct}{46.0}  % QLIKE change % vs the reference (subsample)
\newcommand{\rhabLassoFullLiveVdCloseDm}{3.87}  % DM vs the reference arm (subsample)
\newcommand{\rhabLassoFullLiveVdCloseP}{\ensuremath{<}0.001}  % p-value
\newcommand{\rhabLassoLoVdLiveVdCloseQlike}{0.1097}  % QLIKE lasso lo_vd live16_vd (subsample)
\newcommand{\rhabLassoLoVdLiveVdClosePct}{$-$13.7}  % QLIKE change % vs the reference (subsample)
\newcommand{\rhabLassoLoVdLiveVdCloseDm}{$-$4.18}  % DM vs the reference arm (subsample)
\newcommand{\rhabLassoLoVdLiveVdCloseP}{\ensuremath{<}0.001}  % p-value
\newcommand{\rhabLassoAddVdLiveVdCloseQlike}{0.1018}  % QLIKE lasso add_vd live16_vd (subsample)
\newcommand{\rhabLassoAddVdLiveVdClosePct}{17.0}  % QLIKE change % vs the reference (subsample)
\newcommand{\rhabLassoAddVdLiveVdCloseDm}{3.75}  % DM vs the reference arm (subsample)
\newcommand{\rhabLassoAddVdLiveVdCloseP}{\ensuremath{<}0.001}  % p-value
\newcommand{\rhabLassoHarLiveVdCloseQlike}{0.0870}  % QLIKE lasso har live16_vd (subsample)
\newcommand{\rhabEnetFullLiveVdCloseQlike}{0.1227}  % QLIKE enet full live16_vd (subsample)
\newcommand{\rhabEnetFullLiveVdClosePct}{41.0}  % QLIKE change % vs the reference (subsample)
\newcommand{\rhabEnetFullLiveVdCloseDm}{3.59}  % DM vs the reference arm (subsample)
\newcommand{\rhabEnetFullLiveVdCloseP}{\ensuremath{<}0.001}  % p-value
\newcommand{\rhabEnetLoVdLiveVdCloseQlike}{0.1029}  % QLIKE enet lo_vd live16_vd (subsample)
\newcommand{\rhabEnetLoVdLiveVdClosePct}{$-$16.2}  % QLIKE change % vs the reference (subsample)
\newcommand{\rhabEnetLoVdLiveVdCloseDm}{$-$4.00}  % DM vs the reference arm (subsample)
\newcommand{\rhabEnetLoVdLiveVdCloseP}{\ensuremath{<}0.001}  % p-value
\newcommand{\rhabEnetAddVdLiveVdCloseQlike}{0.1005}  % QLIKE enet add_vd live16_vd (subsample)
\newcommand{\rhabEnetAddVdLiveVdClosePct}{15.5}  % QLIKE change % vs the reference (subsample)
\newcommand{\rhabEnetAddVdLiveVdCloseDm}{3.55}  % DM vs the reference arm (subsample)
\newcommand{\rhabEnetAddVdLiveVdCloseP}{\ensuremath{<}0.001}  % p-value
\newcommand{\rhabEnetHarLiveVdCloseQlike}{0.0870}  % QLIKE enet har live16_vd (subsample)
\newcommand{\rhabLgbmFullLiveVdCloseQlike}{0.0802}  % QLIKE lgbm full live16_vd (subsample)
\newcommand{\rhabLgbmFullLiveVdClosePct}{$-$1.5}  % QLIKE change % vs the reference (subsample)
\newcommand{\rhabLgbmFullLiveVdCloseDm}{$-$0.29}  % DM vs the reference arm (subsample)
\newcommand{\rhabLgbmFullLiveVdCloseP}{0.769}  % p-value
\newcommand{\rhabLgbmHarLiveVdCloseQlike}{0.0814}  % QLIKE lgbm har live16_vd (subsample)
\newcommand{\rhabXgbFullLiveVdCloseQlike}{0.0785}  % QLIKE xgb full live16_vd (subsample)
\newcommand{\rhabXgbFullLiveVdClosePct}{$-$4.3}  % QLIKE change % vs the reference (subsample)
\newcommand{\rhabXgbFullLiveVdCloseDm}{$-$0.89}  % DM vs the reference arm (subsample)
\newcommand{\rhabXgbFullLiveVdCloseP}{0.375}  % p-value
\newcommand{\rhabXgbHarLiveVdCloseQlike}{0.0820}  % QLIKE xgb har live16_vd (subsample)
\newcommand{\rhabRfFullLiveVdCloseQlike}{0.0822}  % QLIKE rf full live16_vd (subsample)
\newcommand{\rhabRfFullLiveVdClosePct}{$-$1.0}  % QLIKE change % vs the reference (subsample)
\newcommand{\rhabRfFullLiveVdCloseDm}{$-$0.21}  % DM vs the reference arm (subsample)
\newcommand{\rhabRfFullLiveVdCloseP}{0.833}  % p-value
\newcommand{\rhabRfHarLiveVdCloseQlike}{0.0831}  % QLIKE rf har live16_vd (subsample)
\newcommand{\rhabLiveSentStart}{19~April 2022}  % first bar of the live window sent
\newcommand{\rhabLiveSentEnd}{30~April 2024}  % last bar of the live window sent
\newcommand{\rhabLiveSentDays}{508}  % trading days of the live window sent
\newcommand{\rhabLiveIvStart}{19~April 2018}  % first bar of the live window iv
\newcommand{\rhabLiveIvEnd}{30~April 2024}  % last bar of the live window iv
\newcommand{\rhabLiveIvDays}{1{,}506}  % trading days of the live window iv
\newcommand{\rhabLiveIvbStart}{6~July 2021}  % first bar of the live window ivb
\newcommand{\rhabLiveIvbEnd}{30~April 2024}  % last bar of the live window ivb
\newcommand{\rhabLiveIvbDays}{706}  % trading days of the live window ivb
\newcommand{\rhabLiveVdStart}{16~November 2022}  % first bar of the live window vd
\newcommand{\rhabLiveVdEnd}{31~August 2023}  % last bar of the live window vd
\newcommand{\rhabLiveVdDays}{196}  % trading days of the live window vd
\newcommand{\rhabRidgeFullVsOlsAddRmAllDm}{$-$4.70}  % DM ridge full vs OLS add_rm
\newcommand{\rhabRidgeFullVsOlsAddRmAllP}{\ensuremath{<}0.001}  % p-value
\newcommand{\rhabRidgeFullVsOlsAddRmCloseDm}{1.16}  % DM ridge full vs OLS add_rm
\newcommand{\rhabRidgeFullVsOlsAddRmCloseP}{0.244}  % p-value
\newcommand{\rhabRidgeFullVsOlsAddRmCloseLateDm}{0.95}  % DM ridge full vs OLS add_rm
\newcommand{\rhabRidgeFullVsOlsAddRmCloseLateP}{0.341}  % p-value
\newcommand{\rhabRidgeFullVsOlsSharAllDm}{$-$4.77}  % DM ridge full vs OLS shar
\newcommand{\rhabRidgeFullVsOlsSharAllP}{\ensuremath{<}0.001}  % p-value
\newcommand{\rhabRidgeFullVsOlsSharCloseDm}{0.41}  % DM ridge full vs OLS shar
\newcommand{\rhabRidgeFullVsOlsSharCloseP}{0.682}  % p-value
\newcommand{\rhabRidgeFullVsOlsSharCloseLateDm}{2.34}  % DM ridge full vs OLS shar
\newcommand{\rhabRidgeFullVsOlsSharCloseLateP}{0.019}  % p-value
\newcommand{\rhabRidgeFullVsOlsHarqAllDm}{$-$5.66}  % DM ridge full vs OLS harq
\newcommand{\rhabRidgeFullVsOlsHarqAllP}{\ensuremath{<}0.001}  % p-value
\newcommand{\rhabRidgeFullVsOlsHarqCloseDm}{0.37}  % DM ridge full vs OLS harq
\newcommand{\rhabRidgeFullVsOlsHarqCloseP}{0.715}  % p-value
\newcommand{\rhabRidgeFullVsOlsHarqCloseLateDm}{2.35}  % DM ridge full vs OLS harq
\newcommand{\rhabRidgeFullVsOlsHarqCloseLateP}{0.019}  % p-value
\newcommand{\rhabRidgeFullVsOlsLharAllDm}{$-$3.23}  % DM ridge full vs OLS lhar
\newcommand{\rhabRidgeFullVsOlsLharAllP}{0.001}  % p-value
\newcommand{\rhabRidgeFullVsOlsLharCloseDm}{2.40}  % DM ridge full vs OLS lhar
\newcommand{\rhabRidgeFullVsOlsLharCloseP}{0.016}  % p-value
\newcommand{\rhabRidgeFullVsOlsLharCloseLateDm}{2.96}  % DM ridge full vs OLS lhar
\newcommand{\rhabRidgeFullVsOlsLharCloseLateP}{0.003}  % p-value
\newcommand{\rhabRidgeFullVsOlsSharqlAllDm}{$-$3.04}  % DM ridge full vs OLS sharql
\newcommand{\rhabRidgeFullVsOlsSharqlAllP}{0.002}  % p-value
\newcommand{\rhabRidgeFullVsOlsSharqlCloseDm}{2.44}  % DM ridge full vs OLS sharql
\newcommand{\rhabRidgeFullVsOlsSharqlCloseP}{0.015}  % p-value
\newcommand{\rhabRidgeFullVsOlsSharqlCloseLateDm}{3.13}  % DM ridge full vs OLS sharql
\newcommand{\rhabRidgeFullVsOlsSharqlCloseLateP}{0.002}  % p-value
\newcommand{\rhabRidgeFullVsOlsHarAllDm}{$-$5.63}  % DM ridge full vs OLS har
\newcommand{\rhabRidgeFullVsOlsHarAllP}{\ensuremath{<}0.001}  % p-value
\newcommand{\rhabRidgeFullVsOlsHarCloseDm}{0.31}  % DM ridge full vs OLS har
\newcommand{\rhabRidgeFullVsOlsHarCloseP}{0.759}  % p-value
\newcommand{\rhabRidgeFullVsOlsHarCloseLateDm}{2.31}  % DM ridge full vs OLS har
\newcommand{\rhabRidgeFullVsOlsHarCloseLateP}{0.021}  % p-value
\newcommand{\rhabLassoFullVsOlsAddRmAllDm}{$-$4.53}  % DM lasso full vs OLS add_rm
\newcommand{\rhabLassoFullVsOlsAddRmAllP}{\ensuremath{<}0.001}  % p-value
\newcommand{\rhabLassoFullVsOlsAddRmCloseDm}{2.40}  % DM lasso full vs OLS add_rm
\newcommand{\rhabLassoFullVsOlsAddRmCloseP}{0.016}  % p-value
\newcommand{\rhabLassoFullVsOlsAddRmCloseLateDm}{2.26}  % DM lasso full vs OLS add_rm
\newcommand{\rhabLassoFullVsOlsAddRmCloseLateP}{0.024}  % p-value
\newcommand{\rhabLassoFullVsOlsSharAllDm}{$-$4.07}  % DM lasso full vs OLS shar
\newcommand{\rhabLassoFullVsOlsSharAllP}{\ensuremath{<}0.001}  % p-value
\newcommand{\rhabLassoFullVsOlsSharCloseDm}{1.12}  % DM lasso full vs OLS shar
\newcommand{\rhabLassoFullVsOlsSharCloseP}{0.261}  % p-value
\newcommand{\rhabLassoFullVsOlsSharCloseLateDm}{2.54}  % DM lasso full vs OLS shar
\newcommand{\rhabLassoFullVsOlsSharCloseLateP}{0.011}  % p-value
\newcommand{\rhabLassoFullVsOlsHarqAllDm}{$-$4.75}  % DM lasso full vs OLS harq
\newcommand{\rhabLassoFullVsOlsHarqAllP}{\ensuremath{<}0.001}  % p-value
\newcommand{\rhabLassoFullVsOlsHarqCloseDm}{1.08}  % DM lasso full vs OLS harq
\newcommand{\rhabLassoFullVsOlsHarqCloseP}{0.279}  % p-value
\newcommand{\rhabLassoFullVsOlsHarqCloseLateDm}{2.55}  % DM lasso full vs OLS harq
\newcommand{\rhabLassoFullVsOlsHarqCloseLateP}{0.011}  % p-value
\newcommand{\rhabLassoFullVsOlsLharAllDm}{$-$2.95}  % DM lasso full vs OLS lhar
\newcommand{\rhabLassoFullVsOlsLharAllP}{0.003}  % p-value
\newcommand{\rhabLassoFullVsOlsLharCloseDm}{3.00}  % DM lasso full vs OLS lhar
\newcommand{\rhabLassoFullVsOlsLharCloseP}{0.003}  % p-value
\newcommand{\rhabLassoFullVsOlsLharCloseLateDm}{3.28}  % DM lasso full vs OLS lhar
\newcommand{\rhabLassoFullVsOlsLharCloseLateP}{0.001}  % p-value
\newcommand{\rhabLassoFullVsOlsSharqlAllDm}{$-$2.81}  % DM lasso full vs OLS sharql
\newcommand{\rhabLassoFullVsOlsSharqlAllP}{0.005}  % p-value
\newcommand{\rhabLassoFullVsOlsSharqlCloseDm}{2.96}  % DM lasso full vs OLS sharql
\newcommand{\rhabLassoFullVsOlsSharqlCloseP}{0.003}  % p-value
\newcommand{\rhabLassoFullVsOlsSharqlCloseLateDm}{3.33}  % DM lasso full vs OLS sharql
\newcommand{\rhabLassoFullVsOlsSharqlCloseLateP}{\ensuremath{<}0.001}  % p-value
\newcommand{\rhabLassoFullVsOlsHarAllDm}{$-$4.72}  % DM lasso full vs OLS har
\newcommand{\rhabLassoFullVsOlsHarAllP}{\ensuremath{<}0.001}  % p-value
\newcommand{\rhabLassoFullVsOlsHarCloseDm}{1.04}  % DM lasso full vs OLS har
\newcommand{\rhabLassoFullVsOlsHarCloseP}{0.300}  % p-value
\newcommand{\rhabLassoFullVsOlsHarCloseLateDm}{2.51}  % DM lasso full vs OLS har
\newcommand{\rhabLassoFullVsOlsHarCloseLateP}{0.012}  % p-value
\newcommand{\rhabEnetFullVsOlsAddRmAllDm}{$-$5.61}  % DM enet full vs OLS add_rm
\newcommand{\rhabEnetFullVsOlsAddRmAllP}{\ensuremath{<}0.001}  % p-value
\newcommand{\rhabEnetFullVsOlsAddRmCloseDm}{0.46}  % DM enet full vs OLS add_rm
\newcommand{\rhabEnetFullVsOlsAddRmCloseP}{0.649}  % p-value
\newcommand{\rhabEnetFullVsOlsAddRmCloseLateDm}{0.29}  % DM enet full vs OLS add_rm
\newcommand{\rhabEnetFullVsOlsAddRmCloseLateP}{0.770}  % p-value
\newcommand{\rhabEnetFullVsOlsSharAllDm}{$-$5.97}  % DM enet full vs OLS shar
\newcommand{\rhabEnetFullVsOlsSharAllP}{\ensuremath{<}0.001}  % p-value
\newcommand{\rhabEnetFullVsOlsSharCloseDm}{$-$0.11}  % DM enet full vs OLS shar
\newcommand{\rhabEnetFullVsOlsSharCloseP}{0.914}  % p-value
\newcommand{\rhabEnetFullVsOlsSharCloseLateDm}{1.92}  % DM enet full vs OLS shar
\newcommand{\rhabEnetFullVsOlsSharCloseLateP}{0.055}  % p-value
\newcommand{\rhabEnetFullVsOlsHarqAllDm}{$-$7.03}  % DM enet full vs OLS harq
\newcommand{\rhabEnetFullVsOlsHarqAllP}{\ensuremath{<}0.001}  % p-value
\newcommand{\rhabEnetFullVsOlsHarqCloseDm}{$-$0.15}  % DM enet full vs OLS harq
\newcommand{\rhabEnetFullVsOlsHarqCloseP}{0.883}  % p-value
\newcommand{\rhabEnetFullVsOlsHarqCloseLateDm}{1.96}  % DM enet full vs OLS harq
\newcommand{\rhabEnetFullVsOlsHarqCloseLateP}{0.050}  % p-value
\newcommand{\rhabEnetFullVsOlsLharAllDm}{$-$4.58}  % DM enet full vs OLS lhar
\newcommand{\rhabEnetFullVsOlsLharAllP}{\ensuremath{<}0.001}  % p-value
\newcommand{\rhabEnetFullVsOlsLharCloseDm}{1.86}  % DM enet full vs OLS lhar
\newcommand{\rhabEnetFullVsOlsLharCloseP}{0.063}  % p-value
\newcommand{\rhabEnetFullVsOlsLharCloseLateDm}{2.46}  % DM enet full vs OLS lhar
\newcommand{\rhabEnetFullVsOlsLharCloseLateP}{0.014}  % p-value
\newcommand{\rhabEnetFullVsOlsSharqlAllDm}{$-$4.34}  % DM enet full vs OLS sharql
\newcommand{\rhabEnetFullVsOlsSharqlAllP}{\ensuremath{<}0.001}  % p-value
\newcommand{\rhabEnetFullVsOlsSharqlCloseDm}{1.92}  % DM enet full vs OLS sharql
\newcommand{\rhabEnetFullVsOlsSharqlCloseP}{0.055}  % p-value
\newcommand{\rhabEnetFullVsOlsSharqlCloseLateDm}{2.65}  % DM enet full vs OLS sharql
\newcommand{\rhabEnetFullVsOlsSharqlCloseLateP}{0.008}  % p-value
\newcommand{\rhabEnetFullVsOlsHarAllDm}{$-$6.98}  % DM enet full vs OLS har
\newcommand{\rhabEnetFullVsOlsHarAllP}{\ensuremath{<}0.001}  % p-value
\newcommand{\rhabEnetFullVsOlsHarCloseDm}{$-$0.21}  % DM enet full vs OLS har
\newcommand{\rhabEnetFullVsOlsHarCloseP}{0.832}  % p-value
\newcommand{\rhabEnetFullVsOlsHarCloseLateDm}{1.91}  % DM enet full vs OLS har
\newcommand{\rhabEnetFullVsOlsHarCloseLateP}{0.056}  % p-value
\newcommand{\rhabLgbmFullVsOlsAddRmAllDm}{$-$7.94}  % DM lgbm full vs OLS add_rm
\newcommand{\rhabLgbmFullVsOlsAddRmAllP}{\ensuremath{<}0.001}  % p-value
\newcommand{\rhabLgbmFullVsOlsAddRmCloseDm}{$-$1.78}  % DM lgbm full vs OLS add_rm
\newcommand{\rhabLgbmFullVsOlsAddRmCloseP}{0.076}  % p-value
\newcommand{\rhabLgbmFullVsOlsAddRmCloseLateDm}{$-$2.44}  % DM lgbm full vs OLS add_rm
\newcommand{\rhabLgbmFullVsOlsAddRmCloseLateP}{0.015}  % p-value
\newcommand{\rhabLgbmFullVsOlsSharAllDm}{$-$9.26}  % DM lgbm full vs OLS shar
\newcommand{\rhabLgbmFullVsOlsSharAllP}{\ensuremath{<}0.001}  % p-value
\newcommand{\rhabLgbmFullVsOlsSharCloseDm}{$-$3.86}  % DM lgbm full vs OLS shar
\newcommand{\rhabLgbmFullVsOlsSharCloseP}{\ensuremath{<}0.001}  % p-value
\newcommand{\rhabLgbmFullVsOlsSharCloseLateDm}{$-$2.34}  % DM lgbm full vs OLS shar
\newcommand{\rhabLgbmFullVsOlsSharCloseLateP}{0.019}  % p-value
\newcommand{\rhabLgbmFullVsOlsHarqAllDm}{$-$10.44}  % DM lgbm full vs OLS harq
\newcommand{\rhabLgbmFullVsOlsHarqAllP}{\ensuremath{<}0.001}  % p-value
\newcommand{\rhabLgbmFullVsOlsHarqCloseDm}{$-$3.77}  % DM lgbm full vs OLS harq
\newcommand{\rhabLgbmFullVsOlsHarqCloseP}{\ensuremath{<}0.001}  % p-value
\newcommand{\rhabLgbmFullVsOlsHarqCloseLateDm}{$-$1.70}  % DM lgbm full vs OLS harq
\newcommand{\rhabLgbmFullVsOlsHarqCloseLateP}{0.088}  % p-value
\newcommand{\rhabLgbmFullVsOlsLharAllDm}{$-$8.32}  % DM lgbm full vs OLS lhar
\newcommand{\rhabLgbmFullVsOlsLharAllP}{\ensuremath{<}0.001}  % p-value
\newcommand{\rhabLgbmFullVsOlsLharCloseDm}{$-$1.02}  % DM lgbm full vs OLS lhar
\newcommand{\rhabLgbmFullVsOlsLharCloseP}{0.310}  % p-value
\newcommand{\rhabLgbmFullVsOlsLharCloseLateDm}{$-$1.42}  % DM lgbm full vs OLS lhar
\newcommand{\rhabLgbmFullVsOlsLharCloseLateP}{0.155}  % p-value
\newcommand{\rhabLgbmFullVsOlsSharqlAllDm}{$-$8.22}  % DM lgbm full vs OLS sharql
\newcommand{\rhabLgbmFullVsOlsSharqlAllP}{\ensuremath{<}0.001}  % p-value
\newcommand{\rhabLgbmFullVsOlsSharqlCloseDm}{$-$1.03}  % DM lgbm full vs OLS sharql
\newcommand{\rhabLgbmFullVsOlsSharqlCloseP}{0.302}  % p-value
\newcommand{\rhabLgbmFullVsOlsSharqlCloseLateDm}{$-$1.45}  % DM lgbm full vs OLS sharql
\newcommand{\rhabLgbmFullVsOlsSharqlCloseLateP}{0.148}  % p-value
\newcommand{\rhabLgbmFullVsOlsHarAllDm}{$-$10.41}  % DM lgbm full vs OLS har
\newcommand{\rhabLgbmFullVsOlsHarAllP}{\ensuremath{<}0.001}  % p-value
\newcommand{\rhabLgbmFullVsOlsHarCloseDm}{$-$3.89}  % DM lgbm full vs OLS har
\newcommand{\rhabLgbmFullVsOlsHarCloseP}{\ensuremath{<}0.001}  % p-value
\newcommand{\rhabLgbmFullVsOlsHarCloseLateDm}{$-$1.86}  % DM lgbm full vs OLS har
\newcommand{\rhabLgbmFullVsOlsHarCloseLateP}{0.063}  % p-value
\newcommand{\rhabXgbFullVsOlsAddRmAllDm}{$-$7.47}  % DM xgb full vs OLS add_rm
\newcommand{\rhabXgbFullVsOlsAddRmAllP}{\ensuremath{<}0.001}  % p-value
\newcommand{\rhabXgbFullVsOlsAddRmCloseDm}{$-$2.16}  % DM xgb full vs OLS add_rm
\newcommand{\rhabXgbFullVsOlsAddRmCloseP}{0.031}  % p-value
\newcommand{\rhabXgbFullVsOlsAddRmCloseLateDm}{$-$2.61}  % DM xgb full vs OLS add_rm
\newcommand{\rhabXgbFullVsOlsAddRmCloseLateP}{0.009}  % p-value
\newcommand{\rhabXgbFullVsOlsSharAllDm}{$-$8.81}  % DM xgb full vs OLS shar
\newcommand{\rhabXgbFullVsOlsSharAllP}{\ensuremath{<}0.001}  % p-value
\newcommand{\rhabXgbFullVsOlsSharCloseDm}{$-$4.87}  % DM xgb full vs OLS shar
\newcommand{\rhabXgbFullVsOlsSharCloseP}{\ensuremath{<}0.001}  % p-value
\newcommand{\rhabXgbFullVsOlsSharCloseLateDm}{$-$2.99}  % DM xgb full vs OLS shar
\newcommand{\rhabXgbFullVsOlsSharCloseLateP}{0.003}  % p-value
\newcommand{\rhabXgbFullVsOlsHarqAllDm}{$-$9.96}  % DM xgb full vs OLS harq
\newcommand{\rhabXgbFullVsOlsHarqAllP}{\ensuremath{<}0.001}  % p-value
\newcommand{\rhabXgbFullVsOlsHarqCloseDm}{$-$4.77}  % DM xgb full vs OLS harq
\newcommand{\rhabXgbFullVsOlsHarqCloseP}{\ensuremath{<}0.001}  % p-value
\newcommand{\rhabXgbFullVsOlsHarqCloseLateDm}{$-$2.27}  % DM xgb full vs OLS harq
\newcommand{\rhabXgbFullVsOlsHarqCloseLateP}{0.023}  % p-value
\newcommand{\rhabXgbFullVsOlsLharAllDm}{$-$7.85}  % DM xgb full vs OLS lhar
\newcommand{\rhabXgbFullVsOlsLharAllP}{\ensuremath{<}0.001}  % p-value
\newcommand{\rhabXgbFullVsOlsLharCloseDm}{$-$1.70}  % DM xgb full vs OLS lhar
\newcommand{\rhabXgbFullVsOlsLharCloseP}{0.090}  % p-value
\newcommand{\rhabXgbFullVsOlsLharCloseLateDm}{$-$1.85}  % DM xgb full vs OLS lhar
\newcommand{\rhabXgbFullVsOlsLharCloseLateP}{0.065}  % p-value
\newcommand{\rhabXgbFullVsOlsSharqlAllDm}{$-$7.76}  % DM xgb full vs OLS sharql
\newcommand{\rhabXgbFullVsOlsSharqlAllP}{\ensuremath{<}0.001}  % p-value
\newcommand{\rhabXgbFullVsOlsSharqlCloseDm}{$-$1.82}  % DM xgb full vs OLS sharql
\newcommand{\rhabXgbFullVsOlsSharqlCloseP}{0.069}  % p-value
\newcommand{\rhabXgbFullVsOlsSharqlCloseLateDm}{$-$1.95}  % DM xgb full vs OLS sharql
\newcommand{\rhabXgbFullVsOlsSharqlCloseLateP}{0.051}  % p-value
\newcommand{\rhabXgbFullVsOlsHarAllDm}{$-$9.93}  % DM xgb full vs OLS har
\newcommand{\rhabXgbFullVsOlsHarAllP}{\ensuremath{<}0.001}  % p-value
\newcommand{\rhabXgbFullVsOlsHarCloseDm}{$-$4.88}  % DM xgb full vs OLS har
\newcommand{\rhabXgbFullVsOlsHarCloseP}{\ensuremath{<}0.001}  % p-value
\newcommand{\rhabXgbFullVsOlsHarCloseLateDm}{$-$2.44}  % DM xgb full vs OLS har
\newcommand{\rhabXgbFullVsOlsHarCloseLateP}{0.015}  % p-value
\newcommand{\rhabRfFullVsOlsAddRmAllDm}{$-$6.68}  % DM rf full vs OLS add_rm
\newcommand{\rhabRfFullVsOlsAddRmAllP}{\ensuremath{<}0.001}  % p-value
\newcommand{\rhabRfFullVsOlsAddRmCloseDm}{$-$0.14}  % DM rf full vs OLS add_rm
\newcommand{\rhabRfFullVsOlsAddRmCloseP}{0.892}  % p-value
\newcommand{\rhabRfFullVsOlsAddRmCloseLateDm}{$-$1.75}  % DM rf full vs OLS add_rm
\newcommand{\rhabRfFullVsOlsAddRmCloseLateP}{0.080}  % p-value
\newcommand{\rhabRfFullVsOlsSharAllDm}{$-$8.45}  % DM rf full vs OLS shar
\newcommand{\rhabRfFullVsOlsSharAllP}{\ensuremath{<}0.001}  % p-value
\newcommand{\rhabRfFullVsOlsSharCloseDm}{$-$0.98}  % DM rf full vs OLS shar
\newcommand{\rhabRfFullVsOlsSharCloseP}{0.328}  % p-value
\newcommand{\rhabRfFullVsOlsSharCloseLateDm}{$-$0.76}  % DM rf full vs OLS shar
\newcommand{\rhabRfFullVsOlsSharCloseLateP}{0.444}  % p-value
\newcommand{\rhabRfFullVsOlsHarqAllDm}{$-$9.73}  % DM rf full vs OLS harq
\newcommand{\rhabRfFullVsOlsHarqAllP}{\ensuremath{<}0.001}  % p-value
\newcommand{\rhabRfFullVsOlsHarqCloseDm}{$-$1.05}  % DM rf full vs OLS harq
\newcommand{\rhabRfFullVsOlsHarqCloseP}{0.293}  % p-value
\newcommand{\rhabRfFullVsOlsHarqCloseLateDm}{$-$0.36}  % DM rf full vs OLS harq
\newcommand{\rhabRfFullVsOlsHarqCloseLateP}{0.715}  % p-value
\newcommand{\rhabRfFullVsOlsLharAllDm}{$-$7.08}  % DM rf full vs OLS lhar
\newcommand{\rhabRfFullVsOlsLharAllP}{\ensuremath{<}0.001}  % p-value
\newcommand{\rhabRfFullVsOlsLharCloseDm}{1.54}  % DM rf full vs OLS lhar
\newcommand{\rhabRfFullVsOlsLharCloseP}{0.124}  % p-value
\newcommand{\rhabRfFullVsOlsLharCloseLateDm}{$-$0.26}  % DM rf full vs OLS lhar
\newcommand{\rhabRfFullVsOlsLharCloseLateP}{0.794}  % p-value
\newcommand{\rhabRfFullVsOlsSharqlAllDm}{$-$6.99}  % DM rf full vs OLS sharql
\newcommand{\rhabRfFullVsOlsSharqlAllP}{\ensuremath{<}0.001}  % p-value
\newcommand{\rhabRfFullVsOlsSharqlCloseDm}{1.76}  % DM rf full vs OLS sharql
\newcommand{\rhabRfFullVsOlsSharqlCloseP}{0.078}  % p-value
\newcommand{\rhabRfFullVsOlsSharqlCloseLateDm}{$-$0.09}  % DM rf full vs OLS sharql
\newcommand{\rhabRfFullVsOlsSharqlCloseLateP}{0.926}  % p-value
\newcommand{\rhabRfFullVsOlsHarAllDm}{$-$9.69}  % DM rf full vs OLS har
\newcommand{\rhabRfFullVsOlsHarAllP}{\ensuremath{<}0.001}  % p-value
\newcommand{\rhabRfFullVsOlsHarCloseDm}{$-$1.16}  % DM rf full vs OLS har
\newcommand{\rhabRfFullVsOlsHarCloseP}{0.246}  % p-value
\newcommand{\rhabRfFullVsOlsHarCloseLateDm}{$-$0.49}  % DM rf full vs OLS har
\newcommand{\rhabRfFullVsOlsHarCloseLateP}{0.621}  % p-value
\newcommand{\rhabNClassesAddRm}{three}  % classes with an add_rm arm
\newcommand{\rhabNClassesBeatAddRm}{none}  % full set beats its own add_rm, DM p < 0.05
\newcommand{\rhabBestOlsBench}{SHAR + HARQ + LHAR}  % lowest QLIKE OLS augmented HAR
\newcommand{\rhabLiveAttnFirst}{23~March 2009}  % first print stocktwits_attention
\newcommand{\rhabLiveAttnLast}{30~April 2024}  % last print stocktwits_attention
\newcommand{\rhabLiveAttnFromKone}{11~February 2021}  % first non-zeroed value, k = 1
\newcommand{\rhabLiveAttnFromKlast}{13~May 2021}  % the same, k = 3125
\newcommand{\rhabLiveAttnSharePct}{31}  % % of evaluation bars live, k = 1
\newcommand{\rhabLiveNSentFirst}{1~June 2010}  % first print stocktwits_sentcount
\newcommand{\rhabLiveNSentLast}{30~April 2024}  % last print stocktwits_sentcount
\newcommand{\rhabLiveNSentFromKone}{19~April 2022}  % first non-zeroed value, k = 1
\newcommand{\rhabLiveNSentFromKlast}{19~July 2022}  % the same, k = 3125
\newcommand{\rhabLiveNSentSharePct}{20}  % % of evaluation bars live, k = 1
\newcommand{\rhabLiveSentFirst}{1~June 2010}  % first print stocktwits_sentiment
\newcommand{\rhabLiveSentLast}{30~April 2024}  % last print stocktwits_sentiment
\newcommand{\rhabLiveSentFromKone}{19~April 2022}  % first non-zeroed value, k = 1
\newcommand{\rhabLiveSentFromKlast}{19~July 2022}  % the same, k = 3125
\newcommand{\rhabLiveSentSharePct}{20}  % % of evaluation bars live, k = 1
\newcommand{\rhabLiveVixFirst}{22~May 2006}  % first print vix
\newcommand{\rhabLiveVixLast}{30~April 2024}  % last print vix
\newcommand{\rhabLiveVixFromKone}{19~April 2018}  % first non-zeroed value, k = 1
\newcommand{\rhabLiveVixFromKlast}{18~July 2018}  % the same, k = 3125
\newcommand{\rhabLiveVixSharePct}{58}  % % of evaluation bars live, k = 1
\newcommand{\rhabLiveVixthreemFirst}{12~August 2009}  % first print vix3m
\newcommand{\rhabLiveVixthreemLast}{30~April 2024}  % last print vix3m
\newcommand{\rhabLiveVixthreemFromKone}{6~July 2021}  % first non-zeroed value, k = 1
\newcommand{\rhabLiveVixthreemFromKlast}{1~October 2021}  % the same, k = 3125
\newcommand{\rhabLiveVixthreemSharePct}{27}  % % of evaluation bars live, k = 1
\newcommand{\rhabLiveVvixFirst}{26~March 2012}  % first print vvix
\newcommand{\rhabLiveVvixLast}{30~April 2024}  % last print vvix
\newcommand{\rhabLiveVvixFromKone}{12~February 2024}  % first non-zeroed value, k = 1
\newcommand{\rhabLiveVvixFromKlast}{never}  % the same, k = 3125
\newcommand{\rhabLiveVvixSharePct}{2}  % % of evaluation bars live, k = 1
\newcommand{\rhabLiveVdAllOcFirst}{3~January 2011}  % first print voldemand_all_open_and_close
\newcommand{\rhabLiveVdAllOcLast}{31~August 2023}  % last print voldemand_all_open_and_close
\newcommand{\rhabLiveVdAllOcFromKone}{16~November 2022}  % first non-zeroed value, k = 1
\newcommand{\rhabLiveVdAllOcFromKlast}{20~February 2023}  % the same, k = 3125
\newcommand{\rhabLiveVdAllOcSharePct}{8}  % % of evaluation bars live, k = 1
\newcommand{\rhabLiveVdAllOFirst}{3~January 2011}  % first print voldemand_all_open_only
\newcommand{\rhabLiveVdAllOLast}{31~August 2023}  % last print voldemand_all_open_only
\newcommand{\rhabLiveVdAllOFromKone}{16~November 2022}  % first non-zeroed value, k = 1
\newcommand{\rhabLiveVdAllOFromKlast}{20~February 2023}  % the same, k = 3125
\newcommand{\rhabLiveVdAllOSharePct}{8}  % % of evaluation bars live, k = 1
\newcommand{\rhabLiveVdSpxOcFirst}{3~January 2011}  % first print voldemand_spx_open_and_close
\newcommand{\rhabLiveVdSpxOcLast}{31~August 2023}  % last print voldemand_spx_open_and_close
\newcommand{\rhabLiveVdSpxOcFromKone}{16~November 2022}  % first non-zeroed value, k = 1
\newcommand{\rhabLiveVdSpxOcFromKlast}{20~February 2023}  % the same, k = 3125
\newcommand{\rhabLiveVdSpxOcSharePct}{8}  % % of evaluation bars live, k = 1
\newcommand{\rhabLiveVdSpxOFirst}{3~January 2011}  % first print voldemand_spx_open_only
\newcommand{\rhabLiveVdSpxOLast}{31~August 2023}  % last print voldemand_spx_open_only
\newcommand{\rhabLiveVdSpxOFromKone}{16~November 2022}  % first non-zeroed value, k = 1
\newcommand{\rhabLiveVdSpxOFromKlast}{20~February 2023}  % the same, k = 3125
\newcommand{\rhabLiveVdSpxOSharePct}{8}  % % of evaluation bars live, k = 1
\fi
\ifdefined\rhcxWMain\else% AUTO-GENERATED by writeup/make_rehaul_complexity.py -- do not edit; re-run the script.
% \rhcx... macros of the draft complexity / scaling subsection (writeup/drafts/2026-10-26_complexity.tex): the
% window x feature-count grid of results/ladder_2026-10-26/grid (RULE.txt ADDENDUM 1).  Names: Gain<Model><Span>
% <W><k> (QLIKE gain over the same window's k = 0 rung, x 10^3), Kstar / Kninety (best k, k at 90 % of the full
% set's gain), ...Read (the reading across windows), Model Lgbm Xgb Ridge Lasso Enet, Span Bar (every bar) / Close
% (16:00), W One Two Main Long, k One Three Five Ten Twenty Full.
% PRE-PATCH design: the FOMC calendar features before the pending FOMC fix (RULE.txt ADDENDUM 2).
\providecommand{\rhcxWOne}{}\renewcommand{\rhcxWOne}{1{,}000}  % rolling window, trading days (13,000 bars)
\providecommand{\rhcxTOne}{}\renewcommand{\rhcxTOne}{13{,}000}  % window length in bars, T = 13 W
\providecommand{\rhcxWTwo}{}\renewcommand{\rhcxWTwo}{2{,}000}  % rolling window, trading days (26,000 bars)
\providecommand{\rhcxTTwo}{}\renewcommand{\rhcxTTwo}{26{,}000}  % window length in bars, T = 13 W
\providecommand{\rhcxWMain}{}\renewcommand{\rhcxWMain}{3{,}000}  % rolling window, trading days (39,000 bars)
\providecommand{\rhcxTMain}{}\renewcommand{\rhcxTMain}{39{,}000}  % window length in bars, T = 13 W
\providecommand{\rhcxWLong}{}\renewcommand{\rhcxWLong}{3{,}500}  % rolling window, trading days (45,500 bars)
\providecommand{\rhcxTLong}{}\renewcommand{\rhcxTLong}{45{,}500}  % window length in bars, T = 13 W
\providecommand{\rhcxNBars}{}\renewcommand{\rhcxNBars}{33{,}527}  % scored bars (Table 1's), every window
\providecommand{\rhcxNDays}{}\renewcommand{\rhcxNDays}{2{,}579}  % evaluation days
\providecommand{\rhcxNOrigins}{}\renewcommand{\rhcxNOrigins}{126}  % monthly origins (refit, held 21 sessions)
\providecommand{\rhcxNSeries}{}\renewcommand{\rhcxNSeries}{46}  % exogenous series
\providecommand{\rhcxPFull}{}\renewcommand{\rhcxPFull}{564}  % predictors of the full set (12 + 12 x 46)
\providecommand{\rhcxPHar}{}\renewcommand{\rhcxPHar}{12}  % predictors of the k = 0 rung (HAR benchmark set)
\providecommand{\rhcxCMax}{}\renewcommand{\rhcxCMax}{0.043}  % largest complexity c = P / T (full set, W = 1,000)
\providecommand{\rhcxCMin}{}\renewcommand{\rhcxCMin}{\ensuremath{2.6\times 10^{-4}}}  % smallest c (k = 0, W = 3,500)
\providecommand{\rhcxCMaxMain}{}\renewcommand{\rhcxCMaxMain}{0.014}  % c of the full set at W = 3,000
\providecommand{\rhcxNGates}{}\renewcommand{\rhcxNGates}{20}  % gates (window x model)
\providecommand{\rhcxNGatesPass}{}\renewcommand{\rhcxNGatesPass}{20}  % gates passed
\providecommand{\rhcxHacLag}{}\renewcommand{\rhcxHacLag}{8}  % Newey-West lag of the DM tests
\providecommand{\rhcxPenRidgeOne}{}\renewcommand{\rhcxPenRidgeOne}{\ensuremath{10^{4}}}  % ridge frozen penalty at W 1000 (all_features)
\providecommand{\rhcxPenLassoOne}{}\renewcommand{\rhcxPenLassoOne}{\ensuremath{10^{-3}}}  % lasso frozen penalty at W 1000 (all_features)
\providecommand{\rhcxPenEnetOne}{}\renewcommand{\rhcxPenEnetOne}{\ensuremath{10^{-2}}}  % enet frozen penalty at W 1000 (all_features)
\providecommand{\rhcxPenRidgeTwo}{}\renewcommand{\rhcxPenRidgeTwo}{\ensuremath{10^{3}}}  % ridge frozen penalty at W 2000 (all_features)
\providecommand{\rhcxPenLassoTwo}{}\renewcommand{\rhcxPenLassoTwo}{\ensuremath{10^{-4}}}  % lasso frozen penalty at W 2000 (all_features)
\providecommand{\rhcxPenEnetTwo}{}\renewcommand{\rhcxPenEnetTwo}{\ensuremath{10^{-3}}}  % enet frozen penalty at W 2000 (all_features)
\providecommand{\rhcxPenRidgeMain}{}\renewcommand{\rhcxPenRidgeMain}{\ensuremath{10^{3}}}  % ridge frozen penalty at W 3000 (all_features)
\providecommand{\rhcxPenLassoMain}{}\renewcommand{\rhcxPenLassoMain}{\ensuremath{10^{-4}}}  % lasso frozen penalty at W 3000 (all_features)
\providecommand{\rhcxPenEnetMain}{}\renewcommand{\rhcxPenEnetMain}{\ensuremath{10^{-3}}}  % enet frozen penalty at W 3000 (all_features)
\providecommand{\rhcxPenRidgeLong}{}\renewcommand{\rhcxPenRidgeLong}{\ensuremath{10^{3}}}  % ridge frozen penalty at W 3500 (all_features)
\providecommand{\rhcxPenLassoLong}{}\renewcommand{\rhcxPenLassoLong}{\ensuremath{10^{-4}}}  % lasso frozen penalty at W 3500 (all_features)
\providecommand{\rhcxPenEnetLong}{}\renewcommand{\rhcxPenEnetLong}{\ensuremath{10^{-3}}}  % enet frozen penalty at W 3500 (all_features)
\providecommand{\rhcxKstarLgbmBarOne}{}\renewcommand{\rhcxKstarLgbmBarOne}{46}  % lgbm all W 1000: rung with the lowest QLIKE
\providecommand{\rhcxKninetyLgbmBarOne}{}\renewcommand{\rhcxKninetyLgbmBarOne}{20}  % lgbm all W 1000: smallest k with 90% of the full set's gain
\providecommand{\rhcxGainLgbmBarOneOne}{}\renewcommand{\rhcxGainLgbmBarOneOne}{\ensuremath{-}0.92}  % lgbm all W 1000 k 1: QLIKE gain over k = 0, x 10^3
\providecommand{\rhcxGainLgbmBarOneThree}{}\renewcommand{\rhcxGainLgbmBarOneThree}{3.45}  % lgbm all W 1000 k 3: QLIKE gain over k = 0, x 10^3
\providecommand{\rhcxGainLgbmBarOneFive}{}\renewcommand{\rhcxGainLgbmBarOneFive}{6.06}  % lgbm all W 1000 k 5: QLIKE gain over k = 0, x 10^3
\providecommand{\rhcxGainLgbmBarOneTen}{}\renewcommand{\rhcxGainLgbmBarOneTen}{6.00}  % lgbm all W 1000 k 10: QLIKE gain over k = 0, x 10^3
\providecommand{\rhcxGainLgbmBarOneTwenty}{}\renewcommand{\rhcxGainLgbmBarOneTwenty}{15.27}  % lgbm all W 1000 k 20: QLIKE gain over k = 0, x 10^3
\providecommand{\rhcxGainLgbmBarOneFull}{}\renewcommand{\rhcxGainLgbmBarOneFull}{15.85}  % lgbm all W 1000 k 46: QLIKE gain over k = 0, x 10^3
\providecommand{\rhcxGainFullPLgbmBarOne}{}\renewcommand{\rhcxGainFullPLgbmBarOne}{p < 0.001}  % lgbm all W 1000: DM p, full set vs k = 0
\providecommand{\rhcxGainBestPLgbmBarOne}{}\renewcommand{\rhcxGainBestPLgbmBarOne}{p < 0.001}  % lgbm all W 1000: DM p, best rung vs k = 0
\providecommand{\rhcxKstarLgbmBarTwo}{}\renewcommand{\rhcxKstarLgbmBarTwo}{46}  % lgbm all W 2000: rung with the lowest QLIKE
\providecommand{\rhcxKninetyLgbmBarTwo}{}\renewcommand{\rhcxKninetyLgbmBarTwo}{10}  % lgbm all W 2000: smallest k with 90% of the full set's gain
\providecommand{\rhcxGainLgbmBarTwoOne}{}\renewcommand{\rhcxGainLgbmBarTwoOne}{0.25}  % lgbm all W 2000 k 1: QLIKE gain over k = 0, x 10^3
\providecommand{\rhcxGainLgbmBarTwoThree}{}\renewcommand{\rhcxGainLgbmBarTwoThree}{5.22}  % lgbm all W 2000 k 3: QLIKE gain over k = 0, x 10^3
\providecommand{\rhcxGainLgbmBarTwoFive}{}\renewcommand{\rhcxGainLgbmBarTwoFive}{7.21}  % lgbm all W 2000 k 5: QLIKE gain over k = 0, x 10^3
\providecommand{\rhcxGainLgbmBarTwoTen}{}\renewcommand{\rhcxGainLgbmBarTwoTen}{22.21}  % lgbm all W 2000 k 10: QLIKE gain over k = 0, x 10^3
\providecommand{\rhcxGainLgbmBarTwoTwenty}{}\renewcommand{\rhcxGainLgbmBarTwoTwenty}{23.05}  % lgbm all W 2000 k 20: QLIKE gain over k = 0, x 10^3
\providecommand{\rhcxGainLgbmBarTwoFull}{}\renewcommand{\rhcxGainLgbmBarTwoFull}{23.08}  % lgbm all W 2000 k 46: QLIKE gain over k = 0, x 10^3
\providecommand{\rhcxGainFullPLgbmBarTwo}{}\renewcommand{\rhcxGainFullPLgbmBarTwo}{p < 0.001}  % lgbm all W 2000: DM p, full set vs k = 0
\providecommand{\rhcxGainBestPLgbmBarTwo}{}\renewcommand{\rhcxGainBestPLgbmBarTwo}{p < 0.001}  % lgbm all W 2000: DM p, best rung vs k = 0
\providecommand{\rhcxKstarLgbmBarMain}{}\renewcommand{\rhcxKstarLgbmBarMain}{46}  % lgbm all W 3000: rung with the lowest QLIKE
\providecommand{\rhcxKninetyLgbmBarMain}{}\renewcommand{\rhcxKninetyLgbmBarMain}{10}  % lgbm all W 3000: smallest k with 90% of the full set's gain
\providecommand{\rhcxGainLgbmBarMainOne}{}\renewcommand{\rhcxGainLgbmBarMainOne}{0.89}  % lgbm all W 3000 k 1: QLIKE gain over k = 0, x 10^3
\providecommand{\rhcxGainLgbmBarMainThree}{}\renewcommand{\rhcxGainLgbmBarMainThree}{8.58}  % lgbm all W 3000 k 3: QLIKE gain over k = 0, x 10^3
\providecommand{\rhcxGainLgbmBarMainFive}{}\renewcommand{\rhcxGainLgbmBarMainFive}{8.73}  % lgbm all W 3000 k 5: QLIKE gain over k = 0, x 10^3
\providecommand{\rhcxGainLgbmBarMainTen}{}\renewcommand{\rhcxGainLgbmBarMainTen}{24.08}  % lgbm all W 3000 k 10: QLIKE gain over k = 0, x 10^3
\providecommand{\rhcxGainLgbmBarMainTwenty}{}\renewcommand{\rhcxGainLgbmBarMainTwenty}{24.46}  % lgbm all W 3000 k 20: QLIKE gain over k = 0, x 10^3
\providecommand{\rhcxGainLgbmBarMainFull}{}\renewcommand{\rhcxGainLgbmBarMainFull}{25.47}  % lgbm all W 3000 k 46: QLIKE gain over k = 0, x 10^3
\providecommand{\rhcxGainFullPLgbmBarMain}{}\renewcommand{\rhcxGainFullPLgbmBarMain}{p < 0.001}  % lgbm all W 3000: DM p, full set vs k = 0
\providecommand{\rhcxGainBestPLgbmBarMain}{}\renewcommand{\rhcxGainBestPLgbmBarMain}{p < 0.001}  % lgbm all W 3000: DM p, best rung vs k = 0
\providecommand{\rhcxKstarLgbmBarLong}{}\renewcommand{\rhcxKstarLgbmBarLong}{46}  % lgbm all W 3500: rung with the lowest QLIKE
\providecommand{\rhcxKninetyLgbmBarLong}{}\renewcommand{\rhcxKninetyLgbmBarLong}{10}  % lgbm all W 3500: smallest k with 90% of the full set's gain
\providecommand{\rhcxGainLgbmBarLongOne}{}\renewcommand{\rhcxGainLgbmBarLongOne}{0.42}  % lgbm all W 3500 k 1: QLIKE gain over k = 0, x 10^3
\providecommand{\rhcxGainLgbmBarLongThree}{}\renewcommand{\rhcxGainLgbmBarLongThree}{8.22}  % lgbm all W 3500 k 3: QLIKE gain over k = 0, x 10^3
\providecommand{\rhcxGainLgbmBarLongFive}{}\renewcommand{\rhcxGainLgbmBarLongFive}{7.78}  % lgbm all W 3500 k 5: QLIKE gain over k = 0, x 10^3
\providecommand{\rhcxGainLgbmBarLongTen}{}\renewcommand{\rhcxGainLgbmBarLongTen}{24.67}  % lgbm all W 3500 k 10: QLIKE gain over k = 0, x 10^3
\providecommand{\rhcxGainLgbmBarLongTwenty}{}\renewcommand{\rhcxGainLgbmBarLongTwenty}{24.81}  % lgbm all W 3500 k 20: QLIKE gain over k = 0, x 10^3
\providecommand{\rhcxGainLgbmBarLongFull}{}\renewcommand{\rhcxGainLgbmBarLongFull}{25.21}  % lgbm all W 3500 k 46: QLIKE gain over k = 0, x 10^3
\providecommand{\rhcxGainFullPLgbmBarLong}{}\renewcommand{\rhcxGainFullPLgbmBarLong}{p < 0.001}  % lgbm all W 3500: DM p, full set vs k = 0
\providecommand{\rhcxGainBestPLgbmBarLong}{}\renewcommand{\rhcxGainBestPLgbmBarLong}{p < 0.001}  % lgbm all W 3500: DM p, best rung vs k = 0
\providecommand{\rhcxKstarReadLgbmBar}{}\renewcommand{\rhcxKstarReadLgbmBar}{is flat}  % lgbm all: best k across W (RULE add. 1 (7))
\providecommand{\rhcxKninetyReadLgbmBar}{}\renewcommand{\rhcxKninetyReadLgbmBar}{falls}  % lgbm all: k90 across W
\providecommand{\rhcxFullBestLgbmBar}{}\renewcommand{\rhcxFullBestLgbmBar}{4}  % lgbm all: windows (of 4) where k = 46 is best
\providecommand{\rhcxSlopePosLgbmBar}{}\renewcommand{\rhcxSlopePosLgbmBar}{4}  % lgbm all: windows with a positive gain-log(c) slope
\providecommand{\rhcxKstarLgbmCloseOne}{}\renewcommand{\rhcxKstarLgbmCloseOne}{5}  % lgbm c16 W 1000: rung with the lowest QLIKE
\providecommand{\rhcxKninetyLgbmCloseOne}{}\renewcommand{\rhcxKninetyLgbmCloseOne}{5}  % lgbm c16 W 1000: smallest k with 90% of the full set's gain
\providecommand{\rhcxGainLgbmCloseOneOne}{}\renewcommand{\rhcxGainLgbmCloseOneOne}{\ensuremath{-}1.95}  % lgbm c16 W 1000 k 1: QLIKE gain over k = 0, x 10^3
\providecommand{\rhcxGainLgbmCloseOneThree}{}\renewcommand{\rhcxGainLgbmCloseOneThree}{1.45}  % lgbm c16 W 1000 k 3: QLIKE gain over k = 0, x 10^3
\providecommand{\rhcxGainLgbmCloseOneFive}{}\renewcommand{\rhcxGainLgbmCloseOneFive}{3.77}  % lgbm c16 W 1000 k 5: QLIKE gain over k = 0, x 10^3
\providecommand{\rhcxGainLgbmCloseOneTen}{}\renewcommand{\rhcxGainLgbmCloseOneTen}{2.77}  % lgbm c16 W 1000 k 10: QLIKE gain over k = 0, x 10^3
\providecommand{\rhcxGainLgbmCloseOneTwenty}{}\renewcommand{\rhcxGainLgbmCloseOneTwenty}{1.22}  % lgbm c16 W 1000 k 20: QLIKE gain over k = 0, x 10^3
\providecommand{\rhcxGainLgbmCloseOneFull}{}\renewcommand{\rhcxGainLgbmCloseOneFull}{1.62}  % lgbm c16 W 1000 k 46: QLIKE gain over k = 0, x 10^3
\providecommand{\rhcxGainFullPLgbmCloseOne}{}\renewcommand{\rhcxGainFullPLgbmCloseOne}{p = 0.490}  % lgbm c16 W 1000: DM p, full set vs k = 0
\providecommand{\rhcxGainBestPLgbmCloseOne}{}\renewcommand{\rhcxGainBestPLgbmCloseOne}{p = 0.084}  % lgbm c16 W 1000: DM p, best rung vs k = 0
\providecommand{\rhcxKstarLgbmCloseTwo}{}\renewcommand{\rhcxKstarLgbmCloseTwo}{5}  % lgbm c16 W 2000: rung with the lowest QLIKE
\providecommand{\rhcxKninetyLgbmCloseTwo}{}\renewcommand{\rhcxKninetyLgbmCloseTwo}{3}  % lgbm c16 W 2000: smallest k with 90% of the full set's gain
\providecommand{\rhcxGainLgbmCloseTwoOne}{}\renewcommand{\rhcxGainLgbmCloseTwoOne}{\ensuremath{-}0.59}  % lgbm c16 W 2000 k 1: QLIKE gain over k = 0, x 10^3
\providecommand{\rhcxGainLgbmCloseTwoThree}{}\renewcommand{\rhcxGainLgbmCloseTwoThree}{1.95}  % lgbm c16 W 2000 k 3: QLIKE gain over k = 0, x 10^3
\providecommand{\rhcxGainLgbmCloseTwoFive}{}\renewcommand{\rhcxGainLgbmCloseTwoFive}{2.45}  % lgbm c16 W 2000 k 5: QLIKE gain over k = 0, x 10^3
\providecommand{\rhcxGainLgbmCloseTwoTen}{}\renewcommand{\rhcxGainLgbmCloseTwoTen}{0.62}  % lgbm c16 W 2000 k 10: QLIKE gain over k = 0, x 10^3
\providecommand{\rhcxGainLgbmCloseTwoTwenty}{}\renewcommand{\rhcxGainLgbmCloseTwoTwenty}{1.20}  % lgbm c16 W 2000 k 20: QLIKE gain over k = 0, x 10^3
\providecommand{\rhcxGainLgbmCloseTwoFull}{}\renewcommand{\rhcxGainLgbmCloseTwoFull}{0.45}  % lgbm c16 W 2000 k 46: QLIKE gain over k = 0, x 10^3
\providecommand{\rhcxGainFullPLgbmCloseTwo}{}\renewcommand{\rhcxGainFullPLgbmCloseTwo}{p = 0.828}  % lgbm c16 W 2000: DM p, full set vs k = 0
\providecommand{\rhcxGainBestPLgbmCloseTwo}{}\renewcommand{\rhcxGainBestPLgbmCloseTwo}{p = 0.224}  % lgbm c16 W 2000: DM p, best rung vs k = 0
\providecommand{\rhcxKstarLgbmCloseMain}{}\renewcommand{\rhcxKstarLgbmCloseMain}{3}  % lgbm c16 W 3000: rung with the lowest QLIKE
\providecommand{\rhcxKninetyLgbmCloseMain}{}\renewcommand{\rhcxKninetyLgbmCloseMain}{3}  % lgbm c16 W 3000: smallest k with 90% of the full set's gain
\providecommand{\rhcxGainLgbmCloseMainOne}{}\renewcommand{\rhcxGainLgbmCloseMainOne}{0.49}  % lgbm c16 W 3000 k 1: QLIKE gain over k = 0, x 10^3
\providecommand{\rhcxGainLgbmCloseMainThree}{}\renewcommand{\rhcxGainLgbmCloseMainThree}{4.35}  % lgbm c16 W 3000 k 3: QLIKE gain over k = 0, x 10^3
\providecommand{\rhcxGainLgbmCloseMainFive}{}\renewcommand{\rhcxGainLgbmCloseMainFive}{3.90}  % lgbm c16 W 3000 k 5: QLIKE gain over k = 0, x 10^3
\providecommand{\rhcxGainLgbmCloseMainTen}{}\renewcommand{\rhcxGainLgbmCloseMainTen}{2.57}  % lgbm c16 W 3000 k 10: QLIKE gain over k = 0, x 10^3
\providecommand{\rhcxGainLgbmCloseMainTwenty}{}\renewcommand{\rhcxGainLgbmCloseMainTwenty}{2.31}  % lgbm c16 W 3000 k 20: QLIKE gain over k = 0, x 10^3
\providecommand{\rhcxGainLgbmCloseMainFull}{}\renewcommand{\rhcxGainLgbmCloseMainFull}{3.01}  % lgbm c16 W 3000 k 46: QLIKE gain over k = 0, x 10^3
\providecommand{\rhcxGainFullPLgbmCloseMain}{}\renewcommand{\rhcxGainFullPLgbmCloseMain}{p = 0.105}  % lgbm c16 W 3000: DM p, full set vs k = 0
\providecommand{\rhcxGainBestPLgbmCloseMain}{}\renewcommand{\rhcxGainBestPLgbmCloseMain}{p = 0.005}  % lgbm c16 W 3000: DM p, best rung vs k = 0
\providecommand{\rhcxKstarLgbmCloseLong}{}\renewcommand{\rhcxKstarLgbmCloseLong}{3}  % lgbm c16 W 3500: rung with the lowest QLIKE
\providecommand{\rhcxKninetyLgbmCloseLong}{}\renewcommand{\rhcxKninetyLgbmCloseLong}{3}  % lgbm c16 W 3500: smallest k with 90% of the full set's gain
\providecommand{\rhcxGainLgbmCloseLongOne}{}\renewcommand{\rhcxGainLgbmCloseLongOne}{0.17}  % lgbm c16 W 3500 k 1: QLIKE gain over k = 0, x 10^3
\providecommand{\rhcxGainLgbmCloseLongThree}{}\renewcommand{\rhcxGainLgbmCloseLongThree}{3.63}  % lgbm c16 W 3500 k 3: QLIKE gain over k = 0, x 10^3
\providecommand{\rhcxGainLgbmCloseLongFive}{}\renewcommand{\rhcxGainLgbmCloseLongFive}{2.38}  % lgbm c16 W 3500 k 5: QLIKE gain over k = 0, x 10^3
\providecommand{\rhcxGainLgbmCloseLongTen}{}\renewcommand{\rhcxGainLgbmCloseLongTen}{2.32}  % lgbm c16 W 3500 k 10: QLIKE gain over k = 0, x 10^3
\providecommand{\rhcxGainLgbmCloseLongTwenty}{}\renewcommand{\rhcxGainLgbmCloseLongTwenty}{1.76}  % lgbm c16 W 3500 k 20: QLIKE gain over k = 0, x 10^3
\providecommand{\rhcxGainLgbmCloseLongFull}{}\renewcommand{\rhcxGainLgbmCloseLongFull}{2.14}  % lgbm c16 W 3500 k 46: QLIKE gain over k = 0, x 10^3
\providecommand{\rhcxGainFullPLgbmCloseLong}{}\renewcommand{\rhcxGainFullPLgbmCloseLong}{p = 0.241}  % lgbm c16 W 3500: DM p, full set vs k = 0
\providecommand{\rhcxGainBestPLgbmCloseLong}{}\renewcommand{\rhcxGainBestPLgbmCloseLong}{p = 0.019}  % lgbm c16 W 3500: DM p, best rung vs k = 0
\providecommand{\rhcxKstarReadLgbmClose}{}\renewcommand{\rhcxKstarReadLgbmClose}{falls}  % lgbm c16: best k across W (RULE add. 1 (7))
\providecommand{\rhcxKninetyReadLgbmClose}{}\renewcommand{\rhcxKninetyReadLgbmClose}{falls}  % lgbm c16: k90 across W
\providecommand{\rhcxFullBestLgbmClose}{}\renewcommand{\rhcxFullBestLgbmClose}{0}  % lgbm c16: windows (of 4) where k = 46 is best
\providecommand{\rhcxSlopePosLgbmClose}{}\renewcommand{\rhcxSlopePosLgbmClose}{4}  % lgbm c16: windows with a positive gain-log(c) slope
\providecommand{\rhcxKstarXgbBarOne}{}\renewcommand{\rhcxKstarXgbBarOne}{20}  % xgb all W 1000: rung with the lowest QLIKE
\providecommand{\rhcxKninetyXgbBarOne}{}\renewcommand{\rhcxKninetyXgbBarOne}{10}  % xgb all W 1000: smallest k with 90% of the full set's gain
\providecommand{\rhcxGainXgbBarOneOne}{}\renewcommand{\rhcxGainXgbBarOneOne}{\ensuremath{-}0.95}  % xgb all W 1000 k 1: QLIKE gain over k = 0, x 10^3
\providecommand{\rhcxGainXgbBarOneThree}{}\renewcommand{\rhcxGainXgbBarOneThree}{1.85}  % xgb all W 1000 k 3: QLIKE gain over k = 0, x 10^3
\providecommand{\rhcxGainXgbBarOneFive}{}\renewcommand{\rhcxGainXgbBarOneFive}{6.46}  % xgb all W 1000 k 5: QLIKE gain over k = 0, x 10^3
\providecommand{\rhcxGainXgbBarOneTen}{}\renewcommand{\rhcxGainXgbBarOneTen}{20.33}  % xgb all W 1000 k 10: QLIKE gain over k = 0, x 10^3
\providecommand{\rhcxGainXgbBarOneTwenty}{}\renewcommand{\rhcxGainXgbBarOneTwenty}{20.59}  % xgb all W 1000 k 20: QLIKE gain over k = 0, x 10^3
\providecommand{\rhcxGainXgbBarOneFull}{}\renewcommand{\rhcxGainXgbBarOneFull}{20.39}  % xgb all W 1000 k 46: QLIKE gain over k = 0, x 10^3
\providecommand{\rhcxGainFullPXgbBarOne}{}\renewcommand{\rhcxGainFullPXgbBarOne}{p < 0.001}  % xgb all W 1000: DM p, full set vs k = 0
\providecommand{\rhcxGainBestPXgbBarOne}{}\renewcommand{\rhcxGainBestPXgbBarOne}{p < 0.001}  % xgb all W 1000: DM p, best rung vs k = 0
\providecommand{\rhcxKstarXgbBarTwo}{}\renewcommand{\rhcxKstarXgbBarTwo}{20}  % xgb all W 2000: rung with the lowest QLIKE
\providecommand{\rhcxKninetyXgbBarTwo}{}\renewcommand{\rhcxKninetyXgbBarTwo}{10}  % xgb all W 2000: smallest k with 90% of the full set's gain
\providecommand{\rhcxGainXgbBarTwoOne}{}\renewcommand{\rhcxGainXgbBarTwoOne}{0.36}  % xgb all W 2000 k 1: QLIKE gain over k = 0, x 10^3
\providecommand{\rhcxGainXgbBarTwoThree}{}\renewcommand{\rhcxGainXgbBarTwoThree}{5.21}  % xgb all W 2000 k 3: QLIKE gain over k = 0, x 10^3
\providecommand{\rhcxGainXgbBarTwoFive}{}\renewcommand{\rhcxGainXgbBarTwoFive}{8.44}  % xgb all W 2000 k 5: QLIKE gain over k = 0, x 10^3
\providecommand{\rhcxGainXgbBarTwoTen}{}\renewcommand{\rhcxGainXgbBarTwoTen}{22.28}  % xgb all W 2000 k 10: QLIKE gain over k = 0, x 10^3
\providecommand{\rhcxGainXgbBarTwoTwenty}{}\renewcommand{\rhcxGainXgbBarTwoTwenty}{22.89}  % xgb all W 2000 k 20: QLIKE gain over k = 0, x 10^3
\providecommand{\rhcxGainXgbBarTwoFull}{}\renewcommand{\rhcxGainXgbBarTwoFull}{22.87}  % xgb all W 2000 k 46: QLIKE gain over k = 0, x 10^3
\providecommand{\rhcxGainFullPXgbBarTwo}{}\renewcommand{\rhcxGainFullPXgbBarTwo}{p < 0.001}  % xgb all W 2000: DM p, full set vs k = 0
\providecommand{\rhcxGainBestPXgbBarTwo}{}\renewcommand{\rhcxGainBestPXgbBarTwo}{p < 0.001}  % xgb all W 2000: DM p, best rung vs k = 0
\providecommand{\rhcxKstarXgbBarMain}{}\renewcommand{\rhcxKstarXgbBarMain}{46}  % xgb all W 3000: rung with the lowest QLIKE
\providecommand{\rhcxKninetyXgbBarMain}{}\renewcommand{\rhcxKninetyXgbBarMain}{10}  % xgb all W 3000: smallest k with 90% of the full set's gain
\providecommand{\rhcxGainXgbBarMainOne}{}\renewcommand{\rhcxGainXgbBarMainOne}{1.05}  % xgb all W 3000 k 1: QLIKE gain over k = 0, x 10^3
\providecommand{\rhcxGainXgbBarMainThree}{}\renewcommand{\rhcxGainXgbBarMainThree}{6.97}  % xgb all W 3000 k 3: QLIKE gain over k = 0, x 10^3
\providecommand{\rhcxGainXgbBarMainFive}{}\renewcommand{\rhcxGainXgbBarMainFive}{9.36}  % xgb all W 3000 k 5: QLIKE gain over k = 0, x 10^3
\providecommand{\rhcxGainXgbBarMainTen}{}\renewcommand{\rhcxGainXgbBarMainTen}{23.85}  % xgb all W 3000 k 10: QLIKE gain over k = 0, x 10^3
\providecommand{\rhcxGainXgbBarMainTwenty}{}\renewcommand{\rhcxGainXgbBarMainTwenty}{24.37}  % xgb all W 3000 k 20: QLIKE gain over k = 0, x 10^3
\providecommand{\rhcxGainXgbBarMainFull}{}\renewcommand{\rhcxGainXgbBarMainFull}{24.60}  % xgb all W 3000 k 46: QLIKE gain over k = 0, x 10^3
\providecommand{\rhcxGainFullPXgbBarMain}{}\renewcommand{\rhcxGainFullPXgbBarMain}{p < 0.001}  % xgb all W 3000: DM p, full set vs k = 0
\providecommand{\rhcxGainBestPXgbBarMain}{}\renewcommand{\rhcxGainBestPXgbBarMain}{p < 0.001}  % xgb all W 3000: DM p, best rung vs k = 0
\providecommand{\rhcxKstarXgbBarLong}{}\renewcommand{\rhcxKstarXgbBarLong}{46}  % xgb all W 3500: rung with the lowest QLIKE
\providecommand{\rhcxKninetyXgbBarLong}{}\renewcommand{\rhcxKninetyXgbBarLong}{10}  % xgb all W 3500: smallest k with 90% of the full set's gain
\providecommand{\rhcxGainXgbBarLongOne}{}\renewcommand{\rhcxGainXgbBarLongOne}{1.15}  % xgb all W 3500 k 1: QLIKE gain over k = 0, x 10^3
\providecommand{\rhcxGainXgbBarLongThree}{}\renewcommand{\rhcxGainXgbBarLongThree}{8.52}  % xgb all W 3500 k 3: QLIKE gain over k = 0, x 10^3
\providecommand{\rhcxGainXgbBarLongFive}{}\renewcommand{\rhcxGainXgbBarLongFive}{8.23}  % xgb all W 3500 k 5: QLIKE gain over k = 0, x 10^3
\providecommand{\rhcxGainXgbBarLongTen}{}\renewcommand{\rhcxGainXgbBarLongTen}{24.04}  % xgb all W 3500 k 10: QLIKE gain over k = 0, x 10^3
\providecommand{\rhcxGainXgbBarLongTwenty}{}\renewcommand{\rhcxGainXgbBarLongTwenty}{24.01}  % xgb all W 3500 k 20: QLIKE gain over k = 0, x 10^3
\providecommand{\rhcxGainXgbBarLongFull}{}\renewcommand{\rhcxGainXgbBarLongFull}{24.37}  % xgb all W 3500 k 46: QLIKE gain over k = 0, x 10^3
\providecommand{\rhcxGainFullPXgbBarLong}{}\renewcommand{\rhcxGainFullPXgbBarLong}{p < 0.001}  % xgb all W 3500: DM p, full set vs k = 0
\providecommand{\rhcxGainBestPXgbBarLong}{}\renewcommand{\rhcxGainBestPXgbBarLong}{p < 0.001}  % xgb all W 3500: DM p, best rung vs k = 0
\providecommand{\rhcxKstarReadXgbBar}{}\renewcommand{\rhcxKstarReadXgbBar}{rises}  % xgb all: best k across W (RULE add. 1 (7))
\providecommand{\rhcxKninetyReadXgbBar}{}\renewcommand{\rhcxKninetyReadXgbBar}{is flat}  % xgb all: k90 across W
\providecommand{\rhcxFullBestXgbBar}{}\renewcommand{\rhcxFullBestXgbBar}{2}  % xgb all: windows (of 4) where k = 46 is best
\providecommand{\rhcxSlopePosXgbBar}{}\renewcommand{\rhcxSlopePosXgbBar}{4}  % xgb all: windows with a positive gain-log(c) slope
\providecommand{\rhcxKstarXgbCloseOne}{}\renewcommand{\rhcxKstarXgbCloseOne}{5}  % xgb c16 W 1000: rung with the lowest QLIKE
\providecommand{\rhcxKninetyXgbCloseOne}{}\renewcommand{\rhcxKninetyXgbCloseOne}{5}  % xgb c16 W 1000: smallest k with 90% of the full set's gain
\providecommand{\rhcxGainXgbCloseOneOne}{}\renewcommand{\rhcxGainXgbCloseOneOne}{\ensuremath{-}1.91}  % xgb c16 W 1000 k 1: QLIKE gain over k = 0, x 10^3
\providecommand{\rhcxGainXgbCloseOneThree}{}\renewcommand{\rhcxGainXgbCloseOneThree}{\ensuremath{-}1.79}  % xgb c16 W 1000 k 3: QLIKE gain over k = 0, x 10^3
\providecommand{\rhcxGainXgbCloseOneFive}{}\renewcommand{\rhcxGainXgbCloseOneFive}{2.01}  % xgb c16 W 1000 k 5: QLIKE gain over k = 0, x 10^3
\providecommand{\rhcxGainXgbCloseOneTen}{}\renewcommand{\rhcxGainXgbCloseOneTen}{\ensuremath{-}0.19}  % xgb c16 W 1000 k 10: QLIKE gain over k = 0, x 10^3
\providecommand{\rhcxGainXgbCloseOneTwenty}{}\renewcommand{\rhcxGainXgbCloseOneTwenty}{0.19}  % xgb c16 W 1000 k 20: QLIKE gain over k = 0, x 10^3
\providecommand{\rhcxGainXgbCloseOneFull}{}\renewcommand{\rhcxGainXgbCloseOneFull}{0.25}  % xgb c16 W 1000 k 46: QLIKE gain over k = 0, x 10^3
\providecommand{\rhcxGainFullPXgbCloseOne}{}\renewcommand{\rhcxGainFullPXgbCloseOne}{p = 0.909}  % xgb c16 W 1000: DM p, full set vs k = 0
\providecommand{\rhcxGainBestPXgbCloseOne}{}\renewcommand{\rhcxGainBestPXgbCloseOne}{p = 0.325}  % xgb c16 W 1000: DM p, best rung vs k = 0
\providecommand{\rhcxKstarXgbCloseTwo}{}\renewcommand{\rhcxKstarXgbCloseTwo}{5}  % xgb c16 W 2000: rung with the lowest QLIKE
\providecommand{\rhcxKninetyXgbCloseTwo}{}\renewcommand{\rhcxKninetyXgbCloseTwo}{5}  % xgb c16 W 2000: smallest k with 90% of the full set's gain
\providecommand{\rhcxGainXgbCloseTwoOne}{}\renewcommand{\rhcxGainXgbCloseTwoOne}{\ensuremath{-}0.66}  % xgb c16 W 2000 k 1: QLIKE gain over k = 0, x 10^3
\providecommand{\rhcxGainXgbCloseTwoThree}{}\renewcommand{\rhcxGainXgbCloseTwoThree}{1.42}  % xgb c16 W 2000 k 3: QLIKE gain over k = 0, x 10^3
\providecommand{\rhcxGainXgbCloseTwoFive}{}\renewcommand{\rhcxGainXgbCloseTwoFive}{2.08}  % xgb c16 W 2000 k 5: QLIKE gain over k = 0, x 10^3
\providecommand{\rhcxGainXgbCloseTwoTen}{}\renewcommand{\rhcxGainXgbCloseTwoTen}{1.30}  % xgb c16 W 2000 k 10: QLIKE gain over k = 0, x 10^3
\providecommand{\rhcxGainXgbCloseTwoTwenty}{}\renewcommand{\rhcxGainXgbCloseTwoTwenty}{1.90}  % xgb c16 W 2000 k 20: QLIKE gain over k = 0, x 10^3
\providecommand{\rhcxGainXgbCloseTwoFull}{}\renewcommand{\rhcxGainXgbCloseTwoFull}{1.59}  % xgb c16 W 2000 k 46: QLIKE gain over k = 0, x 10^3
\providecommand{\rhcxGainFullPXgbCloseTwo}{}\renewcommand{\rhcxGainFullPXgbCloseTwo}{p = 0.423}  % xgb c16 W 2000: DM p, full set vs k = 0
\providecommand{\rhcxGainBestPXgbCloseTwo}{}\renewcommand{\rhcxGainBestPXgbCloseTwo}{p = 0.291}  % xgb c16 W 2000: DM p, best rung vs k = 0
\providecommand{\rhcxKstarXgbCloseMain}{}\renewcommand{\rhcxKstarXgbCloseMain}{5}  % xgb c16 W 3000: rung with the lowest QLIKE
\providecommand{\rhcxKninetyXgbCloseMain}{}\renewcommand{\rhcxKninetyXgbCloseMain}{5}  % xgb c16 W 3000: smallest k with 90% of the full set's gain
\providecommand{\rhcxGainXgbCloseMainOne}{}\renewcommand{\rhcxGainXgbCloseMainOne}{0.47}  % xgb c16 W 3000 k 1: QLIKE gain over k = 0, x 10^3
\providecommand{\rhcxGainXgbCloseMainThree}{}\renewcommand{\rhcxGainXgbCloseMainThree}{1.29}  % xgb c16 W 3000 k 3: QLIKE gain over k = 0, x 10^3
\providecommand{\rhcxGainXgbCloseMainFive}{}\renewcommand{\rhcxGainXgbCloseMainFive}{3.81}  % xgb c16 W 3000 k 5: QLIKE gain over k = 0, x 10^3
\providecommand{\rhcxGainXgbCloseMainTen}{}\renewcommand{\rhcxGainXgbCloseMainTen}{2.24}  % xgb c16 W 3000 k 10: QLIKE gain over k = 0, x 10^3
\providecommand{\rhcxGainXgbCloseMainTwenty}{}\renewcommand{\rhcxGainXgbCloseMainTwenty}{3.46}  % xgb c16 W 3000 k 20: QLIKE gain over k = 0, x 10^3
\providecommand{\rhcxGainXgbCloseMainFull}{}\renewcommand{\rhcxGainXgbCloseMainFull}{3.06}  % xgb c16 W 3000 k 46: QLIKE gain over k = 0, x 10^3
\providecommand{\rhcxGainFullPXgbCloseMain}{}\renewcommand{\rhcxGainFullPXgbCloseMain}{p = 0.120}  % xgb c16 W 3000: DM p, full set vs k = 0
\providecommand{\rhcxGainBestPXgbCloseMain}{}\renewcommand{\rhcxGainBestPXgbCloseMain}{p = 0.050}  % xgb c16 W 3000: DM p, best rung vs k = 0
\providecommand{\rhcxKstarXgbCloseLong}{}\renewcommand{\rhcxKstarXgbCloseLong}{3}  % xgb c16 W 3500: rung with the lowest QLIKE
\providecommand{\rhcxKninetyXgbCloseLong}{}\renewcommand{\rhcxKninetyXgbCloseLong}{3}  % xgb c16 W 3500: smallest k with 90% of the full set's gain
\providecommand{\rhcxGainXgbCloseLongOne}{}\renewcommand{\rhcxGainXgbCloseLongOne}{1.13}  % xgb c16 W 3500 k 1: QLIKE gain over k = 0, x 10^3
\providecommand{\rhcxGainXgbCloseLongThree}{}\renewcommand{\rhcxGainXgbCloseLongThree}{4.34}  % xgb c16 W 3500 k 3: QLIKE gain over k = 0, x 10^3
\providecommand{\rhcxGainXgbCloseLongFive}{}\renewcommand{\rhcxGainXgbCloseLongFive}{4.22}  % xgb c16 W 3500 k 5: QLIKE gain over k = 0, x 10^3
\providecommand{\rhcxGainXgbCloseLongTen}{}\renewcommand{\rhcxGainXgbCloseLongTen}{3.09}  % xgb c16 W 3500 k 10: QLIKE gain over k = 0, x 10^3
\providecommand{\rhcxGainXgbCloseLongTwenty}{}\renewcommand{\rhcxGainXgbCloseLongTwenty}{3.06}  % xgb c16 W 3500 k 20: QLIKE gain over k = 0, x 10^3
\providecommand{\rhcxGainXgbCloseLongFull}{}\renewcommand{\rhcxGainXgbCloseLongFull}{3.04}  % xgb c16 W 3500 k 46: QLIKE gain over k = 0, x 10^3
\providecommand{\rhcxGainFullPXgbCloseLong}{}\renewcommand{\rhcxGainFullPXgbCloseLong}{p = 0.100}  % xgb c16 W 3500: DM p, full set vs k = 0
\providecommand{\rhcxGainBestPXgbCloseLong}{}\renewcommand{\rhcxGainBestPXgbCloseLong}{p = 0.008}  % xgb c16 W 3500: DM p, best rung vs k = 0
\providecommand{\rhcxKstarReadXgbClose}{}\renewcommand{\rhcxKstarReadXgbClose}{falls}  % xgb c16: best k across W (RULE add. 1 (7))
\providecommand{\rhcxKninetyReadXgbClose}{}\renewcommand{\rhcxKninetyReadXgbClose}{falls}  % xgb c16: k90 across W
\providecommand{\rhcxFullBestXgbClose}{}\renewcommand{\rhcxFullBestXgbClose}{0}  % xgb c16: windows (of 4) where k = 46 is best
\providecommand{\rhcxSlopePosXgbClose}{}\renewcommand{\rhcxSlopePosXgbClose}{4}  % xgb c16: windows with a positive gain-log(c) slope
\providecommand{\rhcxKstarRidgeBarOne}{}\renewcommand{\rhcxKstarRidgeBarOne}{46}  % ridge all W 1000: rung with the lowest QLIKE
\providecommand{\rhcxKninetyRidgeBarOne}{}\renewcommand{\rhcxKninetyRidgeBarOne}{20}  % ridge all W 1000: smallest k with 90% of the full set's gain
\providecommand{\rhcxGainRidgeBarOneOne}{}\renewcommand{\rhcxGainRidgeBarOneOne}{4.23}  % ridge all W 1000 k 1: QLIKE gain over k = 0, x 10^3
\providecommand{\rhcxGainRidgeBarOneThree}{}\renewcommand{\rhcxGainRidgeBarOneThree}{6.12}  % ridge all W 1000 k 3: QLIKE gain over k = 0, x 10^3
\providecommand{\rhcxGainRidgeBarOneFive}{}\renewcommand{\rhcxGainRidgeBarOneFive}{5.78}  % ridge all W 1000 k 5: QLIKE gain over k = 0, x 10^3
\providecommand{\rhcxGainRidgeBarOneTen}{}\renewcommand{\rhcxGainRidgeBarOneTen}{6.26}  % ridge all W 1000 k 10: QLIKE gain over k = 0, x 10^3
\providecommand{\rhcxGainRidgeBarOneTwenty}{}\renewcommand{\rhcxGainRidgeBarOneTwenty}{9.24}  % ridge all W 1000 k 20: QLIKE gain over k = 0, x 10^3
\providecommand{\rhcxGainRidgeBarOneFull}{}\renewcommand{\rhcxGainRidgeBarOneFull}{9.82}  % ridge all W 1000 k 46: QLIKE gain over k = 0, x 10^3
\providecommand{\rhcxGainFullPRidgeBarOne}{}\renewcommand{\rhcxGainFullPRidgeBarOne}{p < 0.001}  % ridge all W 1000: DM p, full set vs k = 0
\providecommand{\rhcxGainBestPRidgeBarOne}{}\renewcommand{\rhcxGainBestPRidgeBarOne}{p < 0.001}  % ridge all W 1000: DM p, best rung vs k = 0
\providecommand{\rhcxKstarRidgeBarTwo}{}\renewcommand{\rhcxKstarRidgeBarTwo}{5}  % ridge all W 2000: rung with the lowest QLIKE
\providecommand{\rhcxKninetyRidgeBarTwo}{}\renewcommand{\rhcxKninetyRidgeBarTwo}{3}  % ridge all W 2000: smallest k with 90% of the full set's gain
\providecommand{\rhcxGainRidgeBarTwoOne}{}\renewcommand{\rhcxGainRidgeBarTwoOne}{2.40}  % ridge all W 2000 k 1: QLIKE gain over k = 0, x 10^3
\providecommand{\rhcxGainRidgeBarTwoThree}{}\renewcommand{\rhcxGainRidgeBarTwoThree}{9.09}  % ridge all W 2000 k 3: QLIKE gain over k = 0, x 10^3
\providecommand{\rhcxGainRidgeBarTwoFive}{}\renewcommand{\rhcxGainRidgeBarTwoFive}{10.30}  % ridge all W 2000 k 5: QLIKE gain over k = 0, x 10^3
\providecommand{\rhcxGainRidgeBarTwoTen}{}\renewcommand{\rhcxGainRidgeBarTwoTen}{8.58}  % ridge all W 2000 k 10: QLIKE gain over k = 0, x 10^3
\providecommand{\rhcxGainRidgeBarTwoTwenty}{}\renewcommand{\rhcxGainRidgeBarTwoTwenty}{9.79}  % ridge all W 2000 k 20: QLIKE gain over k = 0, x 10^3
\providecommand{\rhcxGainRidgeBarTwoFull}{}\renewcommand{\rhcxGainRidgeBarTwoFull}{6.84}  % ridge all W 2000 k 46: QLIKE gain over k = 0, x 10^3
\providecommand{\rhcxGainFullPRidgeBarTwo}{}\renewcommand{\rhcxGainFullPRidgeBarTwo}{p = 0.004}  % ridge all W 2000: DM p, full set vs k = 0
\providecommand{\rhcxGainBestPRidgeBarTwo}{}\renewcommand{\rhcxGainBestPRidgeBarTwo}{p < 0.001}  % ridge all W 2000: DM p, best rung vs k = 0
\providecommand{\rhcxKstarRidgeBarMain}{}\renewcommand{\rhcxKstarRidgeBarMain}{10}  % ridge all W 3000: rung with the lowest QLIKE
\providecommand{\rhcxKninetyRidgeBarMain}{}\renewcommand{\rhcxKninetyRidgeBarMain}{3}  % ridge all W 3000: smallest k with 90% of the full set's gain
\providecommand{\rhcxGainRidgeBarMainOne}{}\renewcommand{\rhcxGainRidgeBarMainOne}{3.80}  % ridge all W 3000 k 1: QLIKE gain over k = 0, x 10^3
\providecommand{\rhcxGainRidgeBarMainThree}{}\renewcommand{\rhcxGainRidgeBarMainThree}{5.07}  % ridge all W 3000 k 3: QLIKE gain over k = 0, x 10^3
\providecommand{\rhcxGainRidgeBarMainFive}{}\renewcommand{\rhcxGainRidgeBarMainFive}{11.42}  % ridge all W 3000 k 5: QLIKE gain over k = 0, x 10^3
\providecommand{\rhcxGainRidgeBarMainTen}{}\renewcommand{\rhcxGainRidgeBarMainTen}{11.71}  % ridge all W 3000 k 10: QLIKE gain over k = 0, x 10^3
\providecommand{\rhcxGainRidgeBarMainTwenty}{}\renewcommand{\rhcxGainRidgeBarMainTwenty}{10.81}  % ridge all W 3000 k 20: QLIKE gain over k = 0, x 10^3
\providecommand{\rhcxGainRidgeBarMainFull}{}\renewcommand{\rhcxGainRidgeBarMainFull}{5.41}  % ridge all W 3000 k 46: QLIKE gain over k = 0, x 10^3
\providecommand{\rhcxGainFullPRidgeBarMain}{}\renewcommand{\rhcxGainFullPRidgeBarMain}{p = 0.015}  % ridge all W 3000: DM p, full set vs k = 0
\providecommand{\rhcxGainBestPRidgeBarMain}{}\renewcommand{\rhcxGainBestPRidgeBarMain}{p < 0.001}  % ridge all W 3000: DM p, best rung vs k = 0
\providecommand{\rhcxKstarRidgeBarLong}{}\renewcommand{\rhcxKstarRidgeBarLong}{10}  % ridge all W 3500: rung with the lowest QLIKE
\providecommand{\rhcxKninetyRidgeBarLong}{}\renewcommand{\rhcxKninetyRidgeBarLong}{10}  % ridge all W 3500: smallest k with 90% of the full set's gain
\providecommand{\rhcxGainRidgeBarLongOne}{}\renewcommand{\rhcxGainRidgeBarLongOne}{2.51}  % ridge all W 3500 k 1: QLIKE gain over k = 0, x 10^3
\providecommand{\rhcxGainRidgeBarLongThree}{}\renewcommand{\rhcxGainRidgeBarLongThree}{5.39}  % ridge all W 3500 k 3: QLIKE gain over k = 0, x 10^3
\providecommand{\rhcxGainRidgeBarLongFive}{}\renewcommand{\rhcxGainRidgeBarLongFive}{8.87}  % ridge all W 3500 k 5: QLIKE gain over k = 0, x 10^3
\providecommand{\rhcxGainRidgeBarLongTen}{}\renewcommand{\rhcxGainRidgeBarLongTen}{12.53}  % ridge all W 3500 k 10: QLIKE gain over k = 0, x 10^3
\providecommand{\rhcxGainRidgeBarLongTwenty}{}\renewcommand{\rhcxGainRidgeBarLongTwenty}{11.90}  % ridge all W 3500 k 20: QLIKE gain over k = 0, x 10^3
\providecommand{\rhcxGainRidgeBarLongFull}{}\renewcommand{\rhcxGainRidgeBarLongFull}{11.15}  % ridge all W 3500 k 46: QLIKE gain over k = 0, x 10^3
\providecommand{\rhcxGainFullPRidgeBarLong}{}\renewcommand{\rhcxGainFullPRidgeBarLong}{p < 0.001}  % ridge all W 3500: DM p, full set vs k = 0
\providecommand{\rhcxGainBestPRidgeBarLong}{}\renewcommand{\rhcxGainBestPRidgeBarLong}{p < 0.001}  % ridge all W 3500: DM p, best rung vs k = 0
\providecommand{\rhcxKstarReadRidgeBar}{}\renewcommand{\rhcxKstarReadRidgeBar}{is not monotone}  % ridge all: best k across W (RULE add. 1 (7))
\providecommand{\rhcxKninetyReadRidgeBar}{}\renewcommand{\rhcxKninetyReadRidgeBar}{is not monotone}  % ridge all: k90 across W
\providecommand{\rhcxFullBestRidgeBar}{}\renewcommand{\rhcxFullBestRidgeBar}{1}  % ridge all: windows (of 4) where k = 46 is best
\providecommand{\rhcxSlopePosRidgeBar}{}\renewcommand{\rhcxSlopePosRidgeBar}{4}  % ridge all: windows with a positive gain-log(c) slope
\providecommand{\rhcxKstarRidgeCloseOne}{}\renewcommand{\rhcxKstarRidgeCloseOne}{10}  % ridge c16 W 1000: rung with the lowest QLIKE
\providecommand{\rhcxKninetyRidgeCloseOne}{}\renewcommand{\rhcxKninetyRidgeCloseOne}{3}  % ridge c16 W 1000: smallest k with 90% of the full set's gain
\providecommand{\rhcxGainRidgeCloseOneOne}{}\renewcommand{\rhcxGainRidgeCloseOneOne}{3.08}  % ridge c16 W 1000 k 1: QLIKE gain over k = 0, x 10^3
\providecommand{\rhcxGainRidgeCloseOneThree}{}\renewcommand{\rhcxGainRidgeCloseOneThree}{5.03}  % ridge c16 W 1000 k 3: QLIKE gain over k = 0, x 10^3
\providecommand{\rhcxGainRidgeCloseOneFive}{}\renewcommand{\rhcxGainRidgeCloseOneFive}{4.92}  % ridge c16 W 1000 k 5: QLIKE gain over k = 0, x 10^3
\providecommand{\rhcxGainRidgeCloseOneTen}{}\renewcommand{\rhcxGainRidgeCloseOneTen}{5.36}  % ridge c16 W 1000 k 10: QLIKE gain over k = 0, x 10^3
\providecommand{\rhcxGainRidgeCloseOneTwenty}{}\renewcommand{\rhcxGainRidgeCloseOneTwenty}{3.56}  % ridge c16 W 1000 k 20: QLIKE gain over k = 0, x 10^3
\providecommand{\rhcxGainRidgeCloseOneFull}{}\renewcommand{\rhcxGainRidgeCloseOneFull}{4.14}  % ridge c16 W 1000 k 46: QLIKE gain over k = 0, x 10^3
\providecommand{\rhcxGainFullPRidgeCloseOne}{}\renewcommand{\rhcxGainFullPRidgeCloseOne}{p = 0.074}  % ridge c16 W 1000: DM p, full set vs k = 0
\providecommand{\rhcxGainBestPRidgeCloseOne}{}\renewcommand{\rhcxGainBestPRidgeCloseOne}{p = 0.003}  % ridge c16 W 1000: DM p, best rung vs k = 0
\providecommand{\rhcxKstarRidgeCloseTwo}{}\renewcommand{\rhcxKstarRidgeCloseTwo}{5}  % ridge c16 W 2000: rung with the lowest QLIKE
\providecommand{\rhcxKninetyRidgeCloseTwo}{}\renewcommand{\rhcxKninetyRidgeCloseTwo}{46}  % ridge c16 W 2000: smallest k with 90% of the full set's gain
\providecommand{\rhcxGainRidgeCloseTwoOne}{}\renewcommand{\rhcxGainRidgeCloseTwoOne}{\ensuremath{-}1.34}  % ridge c16 W 2000 k 1: QLIKE gain over k = 0, x 10^3
\providecommand{\rhcxGainRidgeCloseTwoThree}{}\renewcommand{\rhcxGainRidgeCloseTwoThree}{1.50}  % ridge c16 W 2000 k 3: QLIKE gain over k = 0, x 10^3
\providecommand{\rhcxGainRidgeCloseTwoFive}{}\renewcommand{\rhcxGainRidgeCloseTwoFive}{4.62}  % ridge c16 W 2000 k 5: QLIKE gain over k = 0, x 10^3
\providecommand{\rhcxGainRidgeCloseTwoTen}{}\renewcommand{\rhcxGainRidgeCloseTwoTen}{3.83}  % ridge c16 W 2000 k 10: QLIKE gain over k = 0, x 10^3
\providecommand{\rhcxGainRidgeCloseTwoTwenty}{}\renewcommand{\rhcxGainRidgeCloseTwoTwenty}{2.91}  % ridge c16 W 2000 k 20: QLIKE gain over k = 0, x 10^3
\providecommand{\rhcxGainRidgeCloseTwoFull}{}\renewcommand{\rhcxGainRidgeCloseTwoFull}{\ensuremath{-}8.48}  % ridge c16 W 2000 k 46: QLIKE gain over k = 0, x 10^3
\providecommand{\rhcxGainFullPRidgeCloseTwo}{}\renewcommand{\rhcxGainFullPRidgeCloseTwo}{p = 0.166}  % ridge c16 W 2000: DM p, full set vs k = 0
\providecommand{\rhcxGainBestPRidgeCloseTwo}{}\renewcommand{\rhcxGainBestPRidgeCloseTwo}{p = 0.037}  % ridge c16 W 2000: DM p, best rung vs k = 0
\providecommand{\rhcxKstarRidgeCloseMain}{}\renewcommand{\rhcxKstarRidgeCloseMain}{5}  % ridge c16 W 3000: rung with the lowest QLIKE
\providecommand{\rhcxKninetyRidgeCloseMain}{}\renewcommand{\rhcxKninetyRidgeCloseMain}{46}  % ridge c16 W 3000: smallest k with 90% of the full set's gain
\providecommand{\rhcxGainRidgeCloseMainOne}{}\renewcommand{\rhcxGainRidgeCloseMainOne}{\ensuremath{-}0.81}  % ridge c16 W 3000 k 1: QLIKE gain over k = 0, x 10^3
\providecommand{\rhcxGainRidgeCloseMainThree}{}\renewcommand{\rhcxGainRidgeCloseMainThree}{\ensuremath{-}0.06}  % ridge c16 W 3000 k 3: QLIKE gain over k = 0, x 10^3
\providecommand{\rhcxGainRidgeCloseMainFive}{}\renewcommand{\rhcxGainRidgeCloseMainFive}{5.09}  % ridge c16 W 3000 k 5: QLIKE gain over k = 0, x 10^3
\providecommand{\rhcxGainRidgeCloseMainTen}{}\renewcommand{\rhcxGainRidgeCloseMainTen}{4.70}  % ridge c16 W 3000 k 10: QLIKE gain over k = 0, x 10^3
\providecommand{\rhcxGainRidgeCloseMainTwenty}{}\renewcommand{\rhcxGainRidgeCloseMainTwenty}{3.66}  % ridge c16 W 3000 k 20: QLIKE gain over k = 0, x 10^3
\providecommand{\rhcxGainRidgeCloseMainFull}{}\renewcommand{\rhcxGainRidgeCloseMainFull}{\ensuremath{-}4.75}  % ridge c16 W 3000 k 46: QLIKE gain over k = 0, x 10^3
\providecommand{\rhcxGainFullPRidgeCloseMain}{}\renewcommand{\rhcxGainFullPRidgeCloseMain}{p = 0.236}  % ridge c16 W 3000: DM p, full set vs k = 0
\providecommand{\rhcxGainBestPRidgeCloseMain}{}\renewcommand{\rhcxGainBestPRidgeCloseMain}{p = 0.010}  % ridge c16 W 3000: DM p, best rung vs k = 0
\providecommand{\rhcxKstarRidgeCloseLong}{}\renewcommand{\rhcxKstarRidgeCloseLong}{10}  % ridge c16 W 3500: rung with the lowest QLIKE
\providecommand{\rhcxKninetyRidgeCloseLong}{}\renewcommand{\rhcxKninetyRidgeCloseLong}{10}  % ridge c16 W 3500: smallest k with 90% of the full set's gain
\providecommand{\rhcxGainRidgeCloseLongOne}{}\renewcommand{\rhcxGainRidgeCloseLongOne}{\ensuremath{-}1.26}  % ridge c16 W 3500 k 1: QLIKE gain over k = 0, x 10^3
\providecommand{\rhcxGainRidgeCloseLongThree}{}\renewcommand{\rhcxGainRidgeCloseLongThree}{0.08}  % ridge c16 W 3500 k 3: QLIKE gain over k = 0, x 10^3
\providecommand{\rhcxGainRidgeCloseLongFive}{}\renewcommand{\rhcxGainRidgeCloseLongFive}{0.88}  % ridge c16 W 3500 k 5: QLIKE gain over k = 0, x 10^3
\providecommand{\rhcxGainRidgeCloseLongTen}{}\renewcommand{\rhcxGainRidgeCloseLongTen}{5.10}  % ridge c16 W 3500 k 10: QLIKE gain over k = 0, x 10^3
\providecommand{\rhcxGainRidgeCloseLongTwenty}{}\renewcommand{\rhcxGainRidgeCloseLongTwenty}{4.32}  % ridge c16 W 3500 k 20: QLIKE gain over k = 0, x 10^3
\providecommand{\rhcxGainRidgeCloseLongFull}{}\renewcommand{\rhcxGainRidgeCloseLongFull}{4.25}  % ridge c16 W 3500 k 46: QLIKE gain over k = 0, x 10^3
\providecommand{\rhcxGainFullPRidgeCloseLong}{}\renewcommand{\rhcxGainFullPRidgeCloseLong}{p = 0.079}  % ridge c16 W 3500: DM p, full set vs k = 0
\providecommand{\rhcxGainBestPRidgeCloseLong}{}\renewcommand{\rhcxGainBestPRidgeCloseLong}{p = 0.014}  % ridge c16 W 3500: DM p, best rung vs k = 0
\providecommand{\rhcxKstarReadRidgeClose}{}\renewcommand{\rhcxKstarReadRidgeClose}{is not monotone}  % ridge c16: best k across W (RULE add. 1 (7))
\providecommand{\rhcxKninetyReadRidgeClose}{}\renewcommand{\rhcxKninetyReadRidgeClose}{is not monotone}  % ridge c16: k90 across W
\providecommand{\rhcxFullBestRidgeClose}{}\renewcommand{\rhcxFullBestRidgeClose}{0}  % ridge c16: windows (of 4) where k = 46 is best
\providecommand{\rhcxSlopePosRidgeClose}{}\renewcommand{\rhcxSlopePosRidgeClose}{2}  % ridge c16: windows with a positive gain-log(c) slope
\providecommand{\rhcxKstarLassoBarOne}{}\renewcommand{\rhcxKstarLassoBarOne}{20}  % lasso all W 1000: rung with the lowest QLIKE
\providecommand{\rhcxKninetyLassoBarOne}{}\renewcommand{\rhcxKninetyLassoBarOne}{20}  % lasso all W 1000: smallest k with 90% of the full set's gain
\providecommand{\rhcxGainLassoBarOneOne}{}\renewcommand{\rhcxGainLassoBarOneOne}{3.36}  % lasso all W 1000 k 1: QLIKE gain over k = 0, x 10^3
\providecommand{\rhcxGainLassoBarOneThree}{}\renewcommand{\rhcxGainLassoBarOneThree}{6.43}  % lasso all W 1000 k 3: QLIKE gain over k = 0, x 10^3
\providecommand{\rhcxGainLassoBarOneFive}{}\renewcommand{\rhcxGainLassoBarOneFive}{5.55}  % lasso all W 1000 k 5: QLIKE gain over k = 0, x 10^3
\providecommand{\rhcxGainLassoBarOneTen}{}\renewcommand{\rhcxGainLassoBarOneTen}{4.21}  % lasso all W 1000 k 10: QLIKE gain over k = 0, x 10^3
\providecommand{\rhcxGainLassoBarOneTwenty}{}\renewcommand{\rhcxGainLassoBarOneTwenty}{9.97}  % lasso all W 1000 k 20: QLIKE gain over k = 0, x 10^3
\providecommand{\rhcxGainLassoBarOneFull}{}\renewcommand{\rhcxGainLassoBarOneFull}{8.26}  % lasso all W 1000 k 46: QLIKE gain over k = 0, x 10^3
\providecommand{\rhcxGainFullPLassoBarOne}{}\renewcommand{\rhcxGainFullPLassoBarOne}{p < 0.001}  % lasso all W 1000: DM p, full set vs k = 0
\providecommand{\rhcxGainBestPLassoBarOne}{}\renewcommand{\rhcxGainBestPLassoBarOne}{p < 0.001}  % lasso all W 1000: DM p, best rung vs k = 0
\providecommand{\rhcxKstarLassoBarTwo}{}\renewcommand{\rhcxKstarLassoBarTwo}{10}  % lasso all W 2000: rung with the lowest QLIKE
\providecommand{\rhcxKninetyLassoBarTwo}{}\renewcommand{\rhcxKninetyLassoBarTwo}{1}  % lasso all W 2000: smallest k with 90% of the full set's gain
\providecommand{\rhcxGainLassoBarTwoOne}{}\renewcommand{\rhcxGainLassoBarTwoOne}{6.47}  % lasso all W 2000 k 1: QLIKE gain over k = 0, x 10^3
\providecommand{\rhcxGainLassoBarTwoThree}{}\renewcommand{\rhcxGainLassoBarTwoThree}{7.49}  % lasso all W 2000 k 3: QLIKE gain over k = 0, x 10^3
\providecommand{\rhcxGainLassoBarTwoFive}{}\renewcommand{\rhcxGainLassoBarTwoFive}{8.78}  % lasso all W 2000 k 5: QLIKE gain over k = 0, x 10^3
\providecommand{\rhcxGainLassoBarTwoTen}{}\renewcommand{\rhcxGainLassoBarTwoTen}{11.50}  % lasso all W 2000 k 10: QLIKE gain over k = 0, x 10^3
\providecommand{\rhcxGainLassoBarTwoTwenty}{}\renewcommand{\rhcxGainLassoBarTwoTwenty}{9.80}  % lasso all W 2000 k 20: QLIKE gain over k = 0, x 10^3
\providecommand{\rhcxGainLassoBarTwoFull}{}\renewcommand{\rhcxGainLassoBarTwoFull}{4.37}  % lasso all W 2000 k 46: QLIKE gain over k = 0, x 10^3
\providecommand{\rhcxGainFullPLassoBarTwo}{}\renewcommand{\rhcxGainFullPLassoBarTwo}{p = 0.161}  % lasso all W 2000: DM p, full set vs k = 0
\providecommand{\rhcxGainBestPLassoBarTwo}{}\renewcommand{\rhcxGainBestPLassoBarTwo}{p < 0.001}  % lasso all W 2000: DM p, best rung vs k = 0
\providecommand{\rhcxKstarLassoBarMain}{}\renewcommand{\rhcxKstarLassoBarMain}{10}  % lasso all W 3000: rung with the lowest QLIKE
\providecommand{\rhcxKninetyLassoBarMain}{}\renewcommand{\rhcxKninetyLassoBarMain}{1}  % lasso all W 3000: smallest k with 90% of the full set's gain
\providecommand{\rhcxGainLassoBarMainOne}{}\renewcommand{\rhcxGainLassoBarMainOne}{6.08}  % lasso all W 3000 k 1: QLIKE gain over k = 0, x 10^3
\providecommand{\rhcxGainLassoBarMainThree}{}\renewcommand{\rhcxGainLassoBarMainThree}{7.07}  % lasso all W 3000 k 3: QLIKE gain over k = 0, x 10^3
\providecommand{\rhcxGainLassoBarMainFive}{}\renewcommand{\rhcxGainLassoBarMainFive}{11.05}  % lasso all W 3000 k 5: QLIKE gain over k = 0, x 10^3
\providecommand{\rhcxGainLassoBarMainTen}{}\renewcommand{\rhcxGainLassoBarMainTen}{13.42}  % lasso all W 3000 k 10: QLIKE gain over k = 0, x 10^3
\providecommand{\rhcxGainLassoBarMainTwenty}{}\renewcommand{\rhcxGainLassoBarMainTwenty}{12.74}  % lasso all W 3000 k 20: QLIKE gain over k = 0, x 10^3
\providecommand{\rhcxGainLassoBarMainFull}{}\renewcommand{\rhcxGainLassoBarMainFull}{5.81}  % lasso all W 3000 k 46: QLIKE gain over k = 0, x 10^3
\providecommand{\rhcxGainFullPLassoBarMain}{}\renewcommand{\rhcxGainFullPLassoBarMain}{p = 0.059}  % lasso all W 3000: DM p, full set vs k = 0
\providecommand{\rhcxGainBestPLassoBarMain}{}\renewcommand{\rhcxGainBestPLassoBarMain}{p < 0.001}  % lasso all W 3000: DM p, best rung vs k = 0
\providecommand{\rhcxKstarLassoBarLong}{}\renewcommand{\rhcxKstarLassoBarLong}{10}  % lasso all W 3500: rung with the lowest QLIKE
\providecommand{\rhcxKninetyLassoBarLong}{}\renewcommand{\rhcxKninetyLassoBarLong}{10}  % lasso all W 3500: smallest k with 90% of the full set's gain
\providecommand{\rhcxGainLassoBarLongOne}{}\renewcommand{\rhcxGainLassoBarLongOne}{6.18}  % lasso all W 3500 k 1: QLIKE gain over k = 0, x 10^3
\providecommand{\rhcxGainLassoBarLongThree}{}\renewcommand{\rhcxGainLassoBarLongThree}{6.94}  % lasso all W 3500 k 3: QLIKE gain over k = 0, x 10^3
\providecommand{\rhcxGainLassoBarLongFive}{}\renewcommand{\rhcxGainLassoBarLongFive}{6.66}  % lasso all W 3500 k 5: QLIKE gain over k = 0, x 10^3
\providecommand{\rhcxGainLassoBarLongTen}{}\renewcommand{\rhcxGainLassoBarLongTen}{14.76}  % lasso all W 3500 k 10: QLIKE gain over k = 0, x 10^3
\providecommand{\rhcxGainLassoBarLongTwenty}{}\renewcommand{\rhcxGainLassoBarLongTwenty}{14.27}  % lasso all W 3500 k 20: QLIKE gain over k = 0, x 10^3
\providecommand{\rhcxGainLassoBarLongFull}{}\renewcommand{\rhcxGainLassoBarLongFull}{13.34}  % lasso all W 3500 k 46: QLIKE gain over k = 0, x 10^3
\providecommand{\rhcxGainFullPLassoBarLong}{}\renewcommand{\rhcxGainFullPLassoBarLong}{p < 0.001}  % lasso all W 3500: DM p, full set vs k = 0
\providecommand{\rhcxGainBestPLassoBarLong}{}\renewcommand{\rhcxGainBestPLassoBarLong}{p < 0.001}  % lasso all W 3500: DM p, best rung vs k = 0
\providecommand{\rhcxKstarReadLassoBar}{}\renewcommand{\rhcxKstarReadLassoBar}{falls}  % lasso all: best k across W (RULE add. 1 (7))
\providecommand{\rhcxKninetyReadLassoBar}{}\renewcommand{\rhcxKninetyReadLassoBar}{is not monotone}  % lasso all: k90 across W
\providecommand{\rhcxFullBestLassoBar}{}\renewcommand{\rhcxFullBestLassoBar}{0}  % lasso all: windows (of 4) where k = 46 is best
\providecommand{\rhcxSlopePosLassoBar}{}\renewcommand{\rhcxSlopePosLassoBar}{4}  % lasso all: windows with a positive gain-log(c) slope
\providecommand{\rhcxKstarLassoCloseOne}{}\renewcommand{\rhcxKstarLassoCloseOne}{46}  % lasso c16 W 1000: rung with the lowest QLIKE
\providecommand{\rhcxKninetyLassoCloseOne}{}\renewcommand{\rhcxKninetyLassoCloseOne}{46}  % lasso c16 W 1000: smallest k with 90% of the full set's gain
\providecommand{\rhcxGainLassoCloseOneOne}{}\renewcommand{\rhcxGainLassoCloseOneOne}{1.53}  % lasso c16 W 1000 k 1: QLIKE gain over k = 0, x 10^3
\providecommand{\rhcxGainLassoCloseOneThree}{}\renewcommand{\rhcxGainLassoCloseOneThree}{3.17}  % lasso c16 W 1000 k 3: QLIKE gain over k = 0, x 10^3
\providecommand{\rhcxGainLassoCloseOneFive}{}\renewcommand{\rhcxGainLassoCloseOneFive}{\ensuremath{-}1.77}  % lasso c16 W 1000 k 5: QLIKE gain over k = 0, x 10^3
\providecommand{\rhcxGainLassoCloseOneTen}{}\renewcommand{\rhcxGainLassoCloseOneTen}{2.68}  % lasso c16 W 1000 k 10: QLIKE gain over k = 0, x 10^3
\providecommand{\rhcxGainLassoCloseOneTwenty}{}\renewcommand{\rhcxGainLassoCloseOneTwenty}{3.82}  % lasso c16 W 1000 k 20: QLIKE gain over k = 0, x 10^3
\providecommand{\rhcxGainLassoCloseOneFull}{}\renewcommand{\rhcxGainLassoCloseOneFull}{4.33}  % lasso c16 W 1000 k 46: QLIKE gain over k = 0, x 10^3
\providecommand{\rhcxGainFullPLassoCloseOne}{}\renewcommand{\rhcxGainFullPLassoCloseOne}{p = 0.160}  % lasso c16 W 1000: DM p, full set vs k = 0
\providecommand{\rhcxGainBestPLassoCloseOne}{}\renewcommand{\rhcxGainBestPLassoCloseOne}{p = 0.160}  % lasso c16 W 1000: DM p, best rung vs k = 0
\providecommand{\rhcxKstarLassoCloseTwo}{}\renewcommand{\rhcxKstarLassoCloseTwo}{10}  % lasso c16 W 2000: rung with the lowest QLIKE
\providecommand{\rhcxKninetyLassoCloseTwo}{}\renewcommand{\rhcxKninetyLassoCloseTwo}{46}  % lasso c16 W 2000: smallest k with 90% of the full set's gain
\providecommand{\rhcxGainLassoCloseTwoOne}{}\renewcommand{\rhcxGainLassoCloseTwoOne}{\ensuremath{-}1.20}  % lasso c16 W 2000 k 1: QLIKE gain over k = 0, x 10^3
\providecommand{\rhcxGainLassoCloseTwoThree}{}\renewcommand{\rhcxGainLassoCloseTwoThree}{\ensuremath{-}2.00}  % lasso c16 W 2000 k 3: QLIKE gain over k = 0, x 10^3
\providecommand{\rhcxGainLassoCloseTwoFive}{}\renewcommand{\rhcxGainLassoCloseTwoFive}{\ensuremath{-}1.11}  % lasso c16 W 2000 k 5: QLIKE gain over k = 0, x 10^3
\providecommand{\rhcxGainLassoCloseTwoTen}{}\renewcommand{\rhcxGainLassoCloseTwoTen}{1.69}  % lasso c16 W 2000 k 10: QLIKE gain over k = 0, x 10^3
\providecommand{\rhcxGainLassoCloseTwoTwenty}{}\renewcommand{\rhcxGainLassoCloseTwoTwenty}{1.08}  % lasso c16 W 2000 k 20: QLIKE gain over k = 0, x 10^3
\providecommand{\rhcxGainLassoCloseTwoFull}{}\renewcommand{\rhcxGainLassoCloseTwoFull}{\ensuremath{-}17.71}  % lasso c16 W 2000 k 46: QLIKE gain over k = 0, x 10^3
\providecommand{\rhcxGainFullPLassoCloseTwo}{}\renewcommand{\rhcxGainFullPLassoCloseTwo}{p = 0.040}  % lasso c16 W 2000: DM p, full set vs k = 0
\providecommand{\rhcxGainBestPLassoCloseTwo}{}\renewcommand{\rhcxGainBestPLassoCloseTwo}{p = 0.491}  % lasso c16 W 2000: DM p, best rung vs k = 0
\providecommand{\rhcxKstarLassoCloseMain}{}\renewcommand{\rhcxKstarLassoCloseMain}{20}  % lasso c16 W 3000: rung with the lowest QLIKE
\providecommand{\rhcxKninetyLassoCloseMain}{}\renewcommand{\rhcxKninetyLassoCloseMain}{46}  % lasso c16 W 3000: smallest k with 90% of the full set's gain
\providecommand{\rhcxGainLassoCloseMainOne}{}\renewcommand{\rhcxGainLassoCloseMainOne}{\ensuremath{-}1.77}  % lasso c16 W 3000 k 1: QLIKE gain over k = 0, x 10^3
\providecommand{\rhcxGainLassoCloseMainThree}{}\renewcommand{\rhcxGainLassoCloseMainThree}{\ensuremath{-}3.15}  % lasso c16 W 3000 k 3: QLIKE gain over k = 0, x 10^3
\providecommand{\rhcxGainLassoCloseMainFive}{}\renewcommand{\rhcxGainLassoCloseMainFive}{\ensuremath{-}0.81}  % lasso c16 W 3000 k 5: QLIKE gain over k = 0, x 10^3
\providecommand{\rhcxGainLassoCloseMainTen}{}\renewcommand{\rhcxGainLassoCloseMainTen}{2.73}  % lasso c16 W 3000 k 10: QLIKE gain over k = 0, x 10^3
\providecommand{\rhcxGainLassoCloseMainTwenty}{}\renewcommand{\rhcxGainLassoCloseMainTwenty}{2.86}  % lasso c16 W 3000 k 20: QLIKE gain over k = 0, x 10^3
\providecommand{\rhcxGainLassoCloseMainFull}{}\renewcommand{\rhcxGainLassoCloseMainFull}{\ensuremath{-}8.63}  % lasso c16 W 3000 k 46: QLIKE gain over k = 0, x 10^3
\providecommand{\rhcxGainFullPLassoCloseMain}{}\renewcommand{\rhcxGainFullPLassoCloseMain}{p = 0.086}  % lasso c16 W 3000: DM p, full set vs k = 0
\providecommand{\rhcxGainBestPLassoCloseMain}{}\renewcommand{\rhcxGainBestPLassoCloseMain}{p = 0.233}  % lasso c16 W 3000: DM p, best rung vs k = 0
\providecommand{\rhcxKstarLassoCloseLong}{}\renewcommand{\rhcxKstarLassoCloseLong}{10}  % lasso c16 W 3500: rung with the lowest QLIKE
\providecommand{\rhcxKninetyLassoCloseLong}{}\renewcommand{\rhcxKninetyLassoCloseLong}{10}  % lasso c16 W 3500: smallest k with 90% of the full set's gain
\providecommand{\rhcxGainLassoCloseLongOne}{}\renewcommand{\rhcxGainLassoCloseLongOne}{\ensuremath{-}1.33}  % lasso c16 W 3500 k 1: QLIKE gain over k = 0, x 10^3
\providecommand{\rhcxGainLassoCloseLongThree}{}\renewcommand{\rhcxGainLassoCloseLongThree}{\ensuremath{-}3.66}  % lasso c16 W 3500 k 3: QLIKE gain over k = 0, x 10^3
\providecommand{\rhcxGainLassoCloseLongFive}{}\renewcommand{\rhcxGainLassoCloseLongFive}{\ensuremath{-}3.07}  % lasso c16 W 3500 k 5: QLIKE gain over k = 0, x 10^3
\providecommand{\rhcxGainLassoCloseLongTen}{}\renewcommand{\rhcxGainLassoCloseLongTen}{4.13}  % lasso c16 W 3500 k 10: QLIKE gain over k = 0, x 10^3
\providecommand{\rhcxGainLassoCloseLongTwenty}{}\renewcommand{\rhcxGainLassoCloseLongTwenty}{3.73}  % lasso c16 W 3500 k 20: QLIKE gain over k = 0, x 10^3
\providecommand{\rhcxGainLassoCloseLongFull}{}\renewcommand{\rhcxGainLassoCloseLongFull}{2.78}  % lasso c16 W 3500 k 46: QLIKE gain over k = 0, x 10^3
\providecommand{\rhcxGainFullPLassoCloseLong}{}\renewcommand{\rhcxGainFullPLassoCloseLong}{p = 0.294}  % lasso c16 W 3500: DM p, full set vs k = 0
\providecommand{\rhcxGainBestPLassoCloseLong}{}\renewcommand{\rhcxGainBestPLassoCloseLong}{p = 0.054}  % lasso c16 W 3500: DM p, best rung vs k = 0
\providecommand{\rhcxKstarReadLassoClose}{}\renewcommand{\rhcxKstarReadLassoClose}{is not monotone}  % lasso c16: best k across W (RULE add. 1 (7))
\providecommand{\rhcxKninetyReadLassoClose}{}\renewcommand{\rhcxKninetyReadLassoClose}{falls}  % lasso c16: k90 across W
\providecommand{\rhcxFullBestLassoClose}{}\renewcommand{\rhcxFullBestLassoClose}{1}  % lasso c16: windows (of 4) where k = 46 is best
\providecommand{\rhcxSlopePosLassoClose}{}\renewcommand{\rhcxSlopePosLassoClose}{2}  % lasso c16: windows with a positive gain-log(c) slope
\providecommand{\rhcxKstarEnetBarOne}{}\renewcommand{\rhcxKstarEnetBarOne}{46}  % enet all W 1000: rung with the lowest QLIKE
\providecommand{\rhcxKninetyEnetBarOne}{}\renewcommand{\rhcxKninetyEnetBarOne}{10}  % enet all W 1000: smallest k with 90% of the full set's gain
\providecommand{\rhcxGainEnetBarOneOne}{}\renewcommand{\rhcxGainEnetBarOneOne}{3.92}  % enet all W 1000 k 1: QLIKE gain over k = 0, x 10^3
\providecommand{\rhcxGainEnetBarOneThree}{}\renewcommand{\rhcxGainEnetBarOneThree}{6.72}  % enet all W 1000 k 3: QLIKE gain over k = 0, x 10^3
\providecommand{\rhcxGainEnetBarOneFive}{}\renewcommand{\rhcxGainEnetBarOneFive}{6.98}  % enet all W 1000 k 5: QLIKE gain over k = 0, x 10^3
\providecommand{\rhcxGainEnetBarOneTen}{}\renewcommand{\rhcxGainEnetBarOneTen}{10.72}  % enet all W 1000 k 10: QLIKE gain over k = 0, x 10^3
\providecommand{\rhcxGainEnetBarOneTwenty}{}\renewcommand{\rhcxGainEnetBarOneTwenty}{11.01}  % enet all W 1000 k 20: QLIKE gain over k = 0, x 10^3
\providecommand{\rhcxGainEnetBarOneFull}{}\renewcommand{\rhcxGainEnetBarOneFull}{11.22}  % enet all W 1000 k 46: QLIKE gain over k = 0, x 10^3
\providecommand{\rhcxGainFullPEnetBarOne}{}\renewcommand{\rhcxGainFullPEnetBarOne}{p < 0.001}  % enet all W 1000: DM p, full set vs k = 0
\providecommand{\rhcxGainBestPEnetBarOne}{}\renewcommand{\rhcxGainBestPEnetBarOne}{p < 0.001}  % enet all W 1000: DM p, best rung vs k = 0
\providecommand{\rhcxKstarEnetBarTwo}{}\renewcommand{\rhcxKstarEnetBarTwo}{10}  % enet all W 2000: rung with the lowest QLIKE
\providecommand{\rhcxKninetyEnetBarTwo}{}\renewcommand{\rhcxKninetyEnetBarTwo}{5}  % enet all W 2000: smallest k with 90% of the full set's gain
\providecommand{\rhcxGainEnetBarTwoOne}{}\renewcommand{\rhcxGainEnetBarTwoOne}{4.04}  % enet all W 2000 k 1: QLIKE gain over k = 0, x 10^3
\providecommand{\rhcxGainEnetBarTwoThree}{}\renewcommand{\rhcxGainEnetBarTwoThree}{5.55}  % enet all W 2000 k 3: QLIKE gain over k = 0, x 10^3
\providecommand{\rhcxGainEnetBarTwoFive}{}\renewcommand{\rhcxGainEnetBarTwoFive}{10.57}  % enet all W 2000 k 5: QLIKE gain over k = 0, x 10^3
\providecommand{\rhcxGainEnetBarTwoTen}{}\renewcommand{\rhcxGainEnetBarTwoTen}{11.11}  % enet all W 2000 k 10: QLIKE gain over k = 0, x 10^3
\providecommand{\rhcxGainEnetBarTwoTwenty}{}\renewcommand{\rhcxGainEnetBarTwoTwenty}{10.51}  % enet all W 2000 k 20: QLIKE gain over k = 0, x 10^3
\providecommand{\rhcxGainEnetBarTwoFull}{}\renewcommand{\rhcxGainEnetBarTwoFull}{8.14}  % enet all W 2000 k 46: QLIKE gain over k = 0, x 10^3
\providecommand{\rhcxGainFullPEnetBarTwo}{}\renewcommand{\rhcxGainFullPEnetBarTwo}{p < 0.001}  % enet all W 2000: DM p, full set vs k = 0
\providecommand{\rhcxGainBestPEnetBarTwo}{}\renewcommand{\rhcxGainBestPEnetBarTwo}{p < 0.001}  % enet all W 2000: DM p, best rung vs k = 0
\providecommand{\rhcxKstarEnetBarMain}{}\renewcommand{\rhcxKstarEnetBarMain}{10}  % enet all W 3000: rung with the lowest QLIKE
\providecommand{\rhcxKninetyEnetBarMain}{}\renewcommand{\rhcxKninetyEnetBarMain}{3}  % enet all W 3000: smallest k with 90% of the full set's gain
\providecommand{\rhcxGainEnetBarMainOne}{}\renewcommand{\rhcxGainEnetBarMainOne}{\ensuremath{-}0.16}  % enet all W 3000 k 1: QLIKE gain over k = 0, x 10^3
\providecommand{\rhcxGainEnetBarMainThree}{}\renewcommand{\rhcxGainEnetBarMainThree}{7.98}  % enet all W 3000 k 3: QLIKE gain over k = 0, x 10^3
\providecommand{\rhcxGainEnetBarMainFive}{}\renewcommand{\rhcxGainEnetBarMainFive}{11.58}  % enet all W 3000 k 5: QLIKE gain over k = 0, x 10^3
\providecommand{\rhcxGainEnetBarMainTen}{}\renewcommand{\rhcxGainEnetBarMainTen}{11.78}  % enet all W 3000 k 10: QLIKE gain over k = 0, x 10^3
\providecommand{\rhcxGainEnetBarMainTwenty}{}\renewcommand{\rhcxGainEnetBarMainTwenty}{10.97}  % enet all W 3000 k 20: QLIKE gain over k = 0, x 10^3
\providecommand{\rhcxGainEnetBarMainFull}{}\renewcommand{\rhcxGainEnetBarMainFull}{8.80}  % enet all W 3000 k 46: QLIKE gain over k = 0, x 10^3
\providecommand{\rhcxGainFullPEnetBarMain}{}\renewcommand{\rhcxGainFullPEnetBarMain}{p < 0.001}  % enet all W 3000: DM p, full set vs k = 0
\providecommand{\rhcxGainBestPEnetBarMain}{}\renewcommand{\rhcxGainBestPEnetBarMain}{p < 0.001}  % enet all W 3000: DM p, best rung vs k = 0
\providecommand{\rhcxKstarEnetBarLong}{}\renewcommand{\rhcxKstarEnetBarLong}{10}  % enet all W 3500: rung with the lowest QLIKE
\providecommand{\rhcxKninetyEnetBarLong}{}\renewcommand{\rhcxKninetyEnetBarLong}{5}  % enet all W 3500: smallest k with 90% of the full set's gain
\providecommand{\rhcxGainEnetBarLongOne}{}\renewcommand{\rhcxGainEnetBarLongOne}{3.88}  % enet all W 3500 k 1: QLIKE gain over k = 0, x 10^3
\providecommand{\rhcxGainEnetBarLongThree}{}\renewcommand{\rhcxGainEnetBarLongThree}{8.42}  % enet all W 3500 k 3: QLIKE gain over k = 0, x 10^3
\providecommand{\rhcxGainEnetBarLongFive}{}\renewcommand{\rhcxGainEnetBarLongFive}{11.87}  % enet all W 3500 k 5: QLIKE gain over k = 0, x 10^3
\providecommand{\rhcxGainEnetBarLongTen}{}\renewcommand{\rhcxGainEnetBarLongTen}{12.25}  % enet all W 3500 k 10: QLIKE gain over k = 0, x 10^3
\providecommand{\rhcxGainEnetBarLongTwenty}{}\renewcommand{\rhcxGainEnetBarLongTwenty}{12.04}  % enet all W 3500 k 20: QLIKE gain over k = 0, x 10^3
\providecommand{\rhcxGainEnetBarLongFull}{}\renewcommand{\rhcxGainEnetBarLongFull}{12.08}  % enet all W 3500 k 46: QLIKE gain over k = 0, x 10^3
\providecommand{\rhcxGainFullPEnetBarLong}{}\renewcommand{\rhcxGainFullPEnetBarLong}{p < 0.001}  % enet all W 3500: DM p, full set vs k = 0
\providecommand{\rhcxGainBestPEnetBarLong}{}\renewcommand{\rhcxGainBestPEnetBarLong}{p < 0.001}  % enet all W 3500: DM p, best rung vs k = 0
\providecommand{\rhcxKstarReadEnetBar}{}\renewcommand{\rhcxKstarReadEnetBar}{falls}  % enet all: best k across W (RULE add. 1 (7))
\providecommand{\rhcxKninetyReadEnetBar}{}\renewcommand{\rhcxKninetyReadEnetBar}{is not monotone}  % enet all: k90 across W
\providecommand{\rhcxFullBestEnetBar}{}\renewcommand{\rhcxFullBestEnetBar}{1}  % enet all: windows (of 4) where k = 46 is best
\providecommand{\rhcxSlopePosEnetBar}{}\renewcommand{\rhcxSlopePosEnetBar}{4}  % enet all: windows with a positive gain-log(c) slope
\providecommand{\rhcxKstarEnetCloseOne}{}\renewcommand{\rhcxKstarEnetCloseOne}{46}  % enet c16 W 1000: rung with the lowest QLIKE
\providecommand{\rhcxKninetyEnetCloseOne}{}\renewcommand{\rhcxKninetyEnetCloseOne}{46}  % enet c16 W 1000: smallest k with 90% of the full set's gain
\providecommand{\rhcxGainEnetCloseOneOne}{}\renewcommand{\rhcxGainEnetCloseOneOne}{2.47}  % enet c16 W 1000 k 1: QLIKE gain over k = 0, x 10^3
\providecommand{\rhcxGainEnetCloseOneThree}{}\renewcommand{\rhcxGainEnetCloseOneThree}{4.19}  % enet c16 W 1000 k 3: QLIKE gain over k = 0, x 10^3
\providecommand{\rhcxGainEnetCloseOneFive}{}\renewcommand{\rhcxGainEnetCloseOneFive}{4.46}  % enet c16 W 1000 k 5: QLIKE gain over k = 0, x 10^3
\providecommand{\rhcxGainEnetCloseOneTen}{}\renewcommand{\rhcxGainEnetCloseOneTen}{4.97}  % enet c16 W 1000 k 10: QLIKE gain over k = 0, x 10^3
\providecommand{\rhcxGainEnetCloseOneTwenty}{}\renewcommand{\rhcxGainEnetCloseOneTwenty}{5.20}  % enet c16 W 1000 k 20: QLIKE gain over k = 0, x 10^3
\providecommand{\rhcxGainEnetCloseOneFull}{}\renewcommand{\rhcxGainEnetCloseOneFull}{5.87}  % enet c16 W 1000 k 46: QLIKE gain over k = 0, x 10^3
\providecommand{\rhcxGainFullPEnetCloseOne}{}\renewcommand{\rhcxGainFullPEnetCloseOne}{p = 0.004}  % enet c16 W 1000: DM p, full set vs k = 0
\providecommand{\rhcxGainBestPEnetCloseOne}{}\renewcommand{\rhcxGainBestPEnetCloseOne}{p = 0.004}  % enet c16 W 1000: DM p, best rung vs k = 0
\providecommand{\rhcxKstarEnetCloseTwo}{}\renewcommand{\rhcxKstarEnetCloseTwo}{10}  % enet c16 W 2000: rung with the lowest QLIKE
\providecommand{\rhcxKninetyEnetCloseTwo}{}\renewcommand{\rhcxKninetyEnetCloseTwo}{46}  % enet c16 W 2000: smallest k with 90% of the full set's gain
\providecommand{\rhcxGainEnetCloseTwoOne}{}\renewcommand{\rhcxGainEnetCloseTwoOne}{\ensuremath{-}0.77}  % enet c16 W 2000 k 1: QLIKE gain over k = 0, x 10^3
\providecommand{\rhcxGainEnetCloseTwoThree}{}\renewcommand{\rhcxGainEnetCloseTwoThree}{0.35}  % enet c16 W 2000 k 3: QLIKE gain over k = 0, x 10^3
\providecommand{\rhcxGainEnetCloseTwoFive}{}\renewcommand{\rhcxGainEnetCloseTwoFive}{4.82}  % enet c16 W 2000 k 5: QLIKE gain over k = 0, x 10^3
\providecommand{\rhcxGainEnetCloseTwoTen}{}\renewcommand{\rhcxGainEnetCloseTwoTen}{5.86}  % enet c16 W 2000 k 10: QLIKE gain over k = 0, x 10^3
\providecommand{\rhcxGainEnetCloseTwoTwenty}{}\renewcommand{\rhcxGainEnetCloseTwoTwenty}{2.34}  % enet c16 W 2000 k 20: QLIKE gain over k = 0, x 10^3
\providecommand{\rhcxGainEnetCloseTwoFull}{}\renewcommand{\rhcxGainEnetCloseTwoFull}{\ensuremath{-}4.89}  % enet c16 W 2000 k 46: QLIKE gain over k = 0, x 10^3
\providecommand{\rhcxGainFullPEnetCloseTwo}{}\renewcommand{\rhcxGainFullPEnetCloseTwo}{p = 0.325}  % enet c16 W 2000: DM p, full set vs k = 0
\providecommand{\rhcxGainBestPEnetCloseTwo}{}\renewcommand{\rhcxGainBestPEnetCloseTwo}{p = 0.009}  % enet c16 W 2000: DM p, best rung vs k = 0
\providecommand{\rhcxKstarEnetCloseMain}{}\renewcommand{\rhcxKstarEnetCloseMain}{10}  % enet c16 W 3000: rung with the lowest QLIKE
\providecommand{\rhcxKninetyEnetCloseMain}{}\renewcommand{\rhcxKninetyEnetCloseMain}{46}  % enet c16 W 3000: smallest k with 90% of the full set's gain
\providecommand{\rhcxGainEnetCloseMainOne}{}\renewcommand{\rhcxGainEnetCloseMainOne}{0.12}  % enet c16 W 3000 k 1: QLIKE gain over k = 0, x 10^3
\providecommand{\rhcxGainEnetCloseMainThree}{}\renewcommand{\rhcxGainEnetCloseMainThree}{1.80}  % enet c16 W 3000 k 3: QLIKE gain over k = 0, x 10^3
\providecommand{\rhcxGainEnetCloseMainFive}{}\renewcommand{\rhcxGainEnetCloseMainFive}{5.12}  % enet c16 W 3000 k 5: QLIKE gain over k = 0, x 10^3
\providecommand{\rhcxGainEnetCloseMainTen}{}\renewcommand{\rhcxGainEnetCloseMainTen}{5.55}  % enet c16 W 3000 k 10: QLIKE gain over k = 0, x 10^3
\providecommand{\rhcxGainEnetCloseMainTwenty}{}\renewcommand{\rhcxGainEnetCloseMainTwenty}{2.41}  % enet c16 W 3000 k 20: QLIKE gain over k = 0, x 10^3
\providecommand{\rhcxGainEnetCloseMainFull}{}\renewcommand{\rhcxGainEnetCloseMainFull}{\ensuremath{-}1.68}  % enet c16 W 3000 k 46: QLIKE gain over k = 0, x 10^3
\providecommand{\rhcxGainFullPEnetCloseMain}{}\renewcommand{\rhcxGainFullPEnetCloseMain}{p = 0.631}  % enet c16 W 3000: DM p, full set vs k = 0
\providecommand{\rhcxGainBestPEnetCloseMain}{}\renewcommand{\rhcxGainBestPEnetCloseMain}{p = 0.009}  % enet c16 W 3000: DM p, best rung vs k = 0
\providecommand{\rhcxKstarEnetCloseLong}{}\renewcommand{\rhcxKstarEnetCloseLong}{5}  % enet c16 W 3500: rung with the lowest QLIKE
\providecommand{\rhcxKninetyEnetCloseLong}{}\renewcommand{\rhcxKninetyEnetCloseLong}{5}  % enet c16 W 3500: smallest k with 90% of the full set's gain
\providecommand{\rhcxGainEnetCloseLongOne}{}\renewcommand{\rhcxGainEnetCloseLongOne}{\ensuremath{-}0.90}  % enet c16 W 3500 k 1: QLIKE gain over k = 0, x 10^3
\providecommand{\rhcxGainEnetCloseLongThree}{}\renewcommand{\rhcxGainEnetCloseLongThree}{2.24}  % enet c16 W 3500 k 3: QLIKE gain over k = 0, x 10^3
\providecommand{\rhcxGainEnetCloseLongFive}{}\renewcommand{\rhcxGainEnetCloseLongFive}{5.41}  % enet c16 W 3500 k 5: QLIKE gain over k = 0, x 10^3
\providecommand{\rhcxGainEnetCloseLongTen}{}\renewcommand{\rhcxGainEnetCloseLongTen}{5.04}  % enet c16 W 3500 k 10: QLIKE gain over k = 0, x 10^3
\providecommand{\rhcxGainEnetCloseLongTwenty}{}\renewcommand{\rhcxGainEnetCloseLongTwenty}{4.89}  % enet c16 W 3500 k 20: QLIKE gain over k = 0, x 10^3
\providecommand{\rhcxGainEnetCloseLongFull}{}\renewcommand{\rhcxGainEnetCloseLongFull}{4.83}  % enet c16 W 3500 k 46: QLIKE gain over k = 0, x 10^3
\providecommand{\rhcxGainFullPEnetCloseLong}{}\renewcommand{\rhcxGainFullPEnetCloseLong}{p = 0.037}  % enet c16 W 3500: DM p, full set vs k = 0
\providecommand{\rhcxGainBestPEnetCloseLong}{}\renewcommand{\rhcxGainBestPEnetCloseLong}{p = 0.004}  % enet c16 W 3500: DM p, best rung vs k = 0
\providecommand{\rhcxKstarReadEnetClose}{}\renewcommand{\rhcxKstarReadEnetClose}{falls}  % enet c16: best k across W (RULE add. 1 (7))
\providecommand{\rhcxKninetyReadEnetClose}{}\renewcommand{\rhcxKninetyReadEnetClose}{falls}  % enet c16: k90 across W
\providecommand{\rhcxFullBestEnetClose}{}\renewcommand{\rhcxFullBestEnetClose}{1}  % enet c16: windows (of 4) where k = 46 is best
\providecommand{\rhcxSlopePosEnetClose}{}\renewcommand{\rhcxSlopePosEnetClose}{3}  % enet c16: windows with a positive gain-log(c) slope
\providecommand{\rhcxNKstarRisesBar}{}\renewcommand{\rhcxNKstarRisesBar}{1}  % all: models whose best k RISES with W
\providecommand{\rhcxWhoKstarRisesBar}{}\renewcommand{\rhcxWhoKstarRisesBar}{XGBoost}  % all: those models
\providecommand{\rhcxNKstarFallsBar}{}\renewcommand{\rhcxNKstarFallsBar}{2}  % all: models whose best k FALLS with W
\providecommand{\rhcxWhoKstarFallsBar}{}\renewcommand{\rhcxWhoKstarFallsBar}{lasso and elastic net}  % all: those models
\providecommand{\rhcxNKstarFlatBar}{}\renewcommand{\rhcxNKstarFlatBar}{1}  % all: models whose best k FLAT with W
\providecommand{\rhcxWhoKstarFlatBar}{}\renewcommand{\rhcxWhoKstarFlatBar}{LightGBM}  % all: those models
\providecommand{\rhcxNKstarMixedBar}{}\renewcommand{\rhcxNKstarMixedBar}{1}  % all: models whose best k MIXED with W
\providecommand{\rhcxWhoKstarMixedBar}{}\renewcommand{\rhcxWhoKstarMixedBar}{ridge}  % all: those models
\providecommand{\rhcxNFullBestBar}{}\renewcommand{\rhcxNFullBestBar}{8}  % all: (model, window) cells where the full set is best, of 20
\providecommand{\rhcxNFullBestTreesBar}{}\renewcommand{\rhcxNFullBestTreesBar}{6}  % all: tree cells (of 8) where the full set is best
\providecommand{\rhcxNFullBestLinearBar}{}\renewcommand{\rhcxNFullBestLinearBar}{2}  % all: linear cells (of 12) where the full set is best
\providecommand{\rhcxNKstarRisesClose}{}\renewcommand{\rhcxNKstarRisesClose}{0}  % c16: models whose best k RISES with W
\providecommand{\rhcxWhoKstarRisesClose}{}\renewcommand{\rhcxWhoKstarRisesClose}{none}  % c16: those models
\providecommand{\rhcxNKstarFallsClose}{}\renewcommand{\rhcxNKstarFallsClose}{3}  % c16: models whose best k FALLS with W
\providecommand{\rhcxWhoKstarFallsClose}{}\renewcommand{\rhcxWhoKstarFallsClose}{LightGBM, XGBoost and elastic net}  % c16: those models
\providecommand{\rhcxNKstarFlatClose}{}\renewcommand{\rhcxNKstarFlatClose}{0}  % c16: models whose best k FLAT with W
\providecommand{\rhcxWhoKstarFlatClose}{}\renewcommand{\rhcxWhoKstarFlatClose}{none}  % c16: those models
\providecommand{\rhcxNKstarMixedClose}{}\renewcommand{\rhcxNKstarMixedClose}{2}  % c16: models whose best k MIXED with W
\providecommand{\rhcxWhoKstarMixedClose}{}\renewcommand{\rhcxWhoKstarMixedClose}{ridge and lasso}  % c16: those models
\providecommand{\rhcxNFullBestClose}{}\renewcommand{\rhcxNFullBestClose}{2}  % c16: (model, window) cells where the full set is best, of 20
\providecommand{\rhcxNFullBestTreesClose}{}\renewcommand{\rhcxNFullBestTreesClose}{0}  % c16: tree cells (of 8) where the full set is best
\providecommand{\rhcxNFullBestLinearClose}{}\renewcommand{\rhcxNFullBestLinearClose}{2}  % c16: linear cells (of 12) where the full set is best
\fi
\begin{itemize}
\item \textbf{Questions and measures.} Which predictors do the pooled models use, is the exogenous signal dense and
  weak or concentrated, and does a series that moves the forecasts also make them more accurate? We explain the
  full-set models of Table~\ref{tab:rh_oos} at \rhfipNOrigins{} monthly origins, \rhfipFirstOrigin{} to
  \rhfipLastOrigin: each origin's refit is held for its block of \rhfipOriginEvery{} trading days, and the blocks
  cover the \rhfipNRows{} bars of the span, including the shortened futures sessions on NYSE holidays, which
  Table~\ref{tab:rh_oos} does not score (it scores \rhorNIntB). A group's share is $s_A$, its share of a tree
  ensemble's mean absolute TreeSHAP value \citep{lundberg2020local}, or $S_A$, its share of a linear model's absolute
  coefficients on the window-standardized predictors, $|\hat\beta_{k,j}|\,\mathrm{sd}_{k,j}$. Standard groups: the
  HAR components, the calendar dummies, the futures' \rhfipNMom{} realized moments, the VIX and the futures' volume
  and trade count; non-standard: the other \rhfipNNonStd{} series. Models are compared by their own shares, across
  model types on the exogenous total; a share measures how much a series moves the forecast, not how much it
  improves its accuracy. No explained bar lies in a tuning sample: the penalties were chosen on the last pre-sample
  fold, ending \rhfipTuneLast, the tree configurations on an earlier development sample
  (Appendix~\ref{app:rh_hyper}). Construction, intervals and checks: Appendix~\ref{app:rh_fi_compute}.
\item \textbf{What the models use} (Table~\ref{tab:rh_fip_shares}). In the trees the HAR components carry
  \rhfipHarMin--\rhfipHarMax\% of the attribution and the futures' realized moments \rhfipMomMin--\rhfipMomMax\%; the
  leading exogenous series are the futures' own \rhfipTreeTopOne{} and \rhfipTreeTopTwo. The non-standard series hold
  \rhfipLgbmNs, \rhfipXgbNs{} and \rhfipRfNs\% (LightGBM, XGBoost, random forest), most of it in five series:
  \rhfipNsOne{} (\rhfipNsOneMin--\rhfipNsOneMax\%), \rhfipNsTwo{} (\rhfipNsTwoMin--\rhfipNsTwoMax\%), \rhfipNsThree{}
  (\rhfipNsThreeMin--\rhfipNsThreeMax\%), \rhfipNsFour{} (\rhfipNsFourMin--\rhfipNsFourMax\%) and \rhfipNsFive{}
  (\rhfipNsFiveMin--\rhfipNsFiveMax\%). The trees lean on them less than the linear models do: of the exogenous total
  they hold \rhfipLgbmNsEx, \rhfipXgbNsEx{} and \rhfipRfNsEx\% in the trees against \rhfipRidgeNsEx, \rhfipLassoNsEx{}
  and \rhfipEnetNsEx\% in the ridge, lasso and elastic net, the bootstrap interval of the difference lying below zero
  in \rhfipNDLess{} of \rhfipNDAll{} comparisons, and the linear models rank the FOMC calendar near the top (the ridge
  puts \rhfipFomcUntil{} at rank \rhfipRidgeRankFomcUntil{} of \rhfipNExog, the trees at
  \rhfipTreeRankFomcUntilMin--\rhfipTreeRankFomcUntilMax). The trees' exogenous attribution is concentrated: its
  effective number of series, $N_{\mathrm{eff}} = 1/\sum_A s_A^2$, is \rhfipLgbmNeff, \rhfipXgbNeff{} and
  \rhfipRfNeff, against \rhfipRidgeNeff{} for the ridge, and \rhfipLgbmNHalf, \rhfipXgbNHalf{} and \rhfipRfNHalf{}
  series hold half of it. Rankings over time, conditional shares, response shapes and interactions:
  Appendices~\ref{app:rh_fi_rank}--\ref{app:rh_fi_shapes}.
\item \textbf{How many series carry the gain} (Figure~\ref{fig:rh_ladder}). Following
  \citet{MaischZhangUlrichZimmer2026}, we refit each model on its top $k$ exogenous series only, $k = 0$ (the HAR
  benchmark set), 1, 3, 5, 10, 20 and \rhfipNExog{} (the full set), at the same origins and hyperparameters, and score
  every rung on the \rhorNIntB{} bars of Table~\ref{tab:rh_oos}. Each model ranks the series once, by its own measure
  on a fit explained on \rhorPenYears, before the span, so the ranking has no look-ahead. Rung $k$ retains
  $R_k = (L_0 - L_k)/(L_0 - L_{\rhfipNExog})$ of the full set's QLIKE gain; the reading rule was fixed before scoring.
  \begin{itemize}
  \item Trees: the top 10 series keep \rhldTreeRTenMin{}--\rhldTreeRTenMax{}\% of the gain, the top
    5 only \rhldTreeRFiveMin{}--\rhldTreeRFiveMax{}\%; \rhfipFomcUntil{}, ranked
    \rhldTreeRankFomcUntilMin{}--\rhldTreeRankFomcUntilMax{} by the trees, enters between the two
    rungs. The other series add little (top 10 against the full set: DM $p = \rhldLgbmPTenFull{}$,
    \rhldXgbPTenFull{} and \rhldRfPTenFull{}).
  \item Linear models: with 5 or 10 series they keep \rhldLinRFiveMin{}--\rhldLinRFiveMax{}\% and
    \rhldLinRTenMin{}--\rhldLinRTenMax{}\% of the gain, more than the full set (top 10 against the
    full set: DM $p \le \rhldLinPTenFullMax{}$), and the lasso's single top series,
    \rhldLassoTopOne{}, matches the full set. Part of the shortfall is the cost of holding a fit on
    \rhfipNCols{} predictors for a month: re-estimated at every bar (Table~\ref{tab:rh_oos}), the full sets are as
    accurate as the monthly top-10 rungs (post hoc; DM $p$ from \rhldLinPTenVsWalkMin{} to
    \rhldLinPTenVsWalkMax{}).
  \item Ridge against lasso: a dense and weak signal favors the ridge, which keeps and shrinks every series, over the
    lasso, which selects \citep[see][]{GiannoneLenzaPrimiceri2021}. In Table~\ref{tab:rh_oos} the full set changes
    QLIKE relative to the HAR benchmark by \rhorVsOlsPctLassoFull\% in the lasso and \rhorVsOlsPctRidgeFull\% in the
    ridge (lasso against ridge: DM $= \rhldDmLassoRidge{}$, $p = \rhldPLassoRidge{}$).
  \end{itemize}
  The exogenous signal is concentrated, not dense and weak, in both model classes. Across windows of \rhcxWOne{} to
  \rhcxWLong{} trading days the best number of series rises with the window, as a dense and weak signal predicts,
  only for \rhcxWhoKstarRisesBar{} (Appendix~\ref{app:rh_fi_scaling}); the window test of the earlier draft is in
  Appendix~\ref{app:rh_fi_window}.
\item \textbf{The late-starting series while live} (Figure~\ref{fig:rh_fi_live}). Under the late-start rule
  (Section~\ref{sec:rh_predictors}) the trees see mostly the availability averages of StockTwits, VIX3M and
  volatility demand, which record when a series prints, not what it prints. In the full years in which their values
  are live, the yearly shares rise: StockTwits from \rhfilStwPreMin--\rhfilStwPreMax\% (\rhfilStwPreYears) to
  \rhfilStwLiveMin--\rhfilStwLiveMax\% (\rhfilStwLiveYears), VIX3M from \rhfilVixThreePreMin--\rhfilVixThreePreMax\%
  to \rhfilVixThreeLiveMin--\rhfilVixThreeLiveMax\%, volatility demand from \rhfilVdPreMin--\rhfilVdPreMax\% to
  \rhfilVdLiveMin--\rhfilVdLiveMax\% (\rhfilVdLiveYears). The pooled shares thus understate their use, but the use
  buys little accuracy: on its live window, adding StockTwits to the HAR benchmark changes the trees' QLIKE by
  \rhabLgbmAddSentLiveSentPct, \rhabXgbAddSentLiveSentPct{} and \rhabRfAddSentLiveSentPct\%, none significantly
  (Table~\ref{tab:rh_ablation}, Panel~E).
\item \textbf{Use is not accuracy.} The FOMC calendar holds \rhfipFomcBucketMin--\rhfipFomcBucketMax\% of the trees'
  attribution (Table~\ref{tab:rh_fip_shares}) yet carries most of their accuracy gain (Table~\ref{tab:rh_ablation}):
  added alone to the HAR benchmark it changes QLIKE by \rhabLgbmAddFomcAllPct, \rhabXgbAddFomcAllPct{} and
  \rhabRfAddFomcAllPct\%, against \rhabLgbmFullAllPct, \rhabXgbFullAllPct{} and \rhabRfFullAllPct\% for the full set
  (\rhabTreeFomcGainShareMin{}--\rhabTreeFomcGainShareMax{}\% of the gain), and removing it raises
  LightGBM's and XGBoost's QLIKE by \rhabLgbmLoFomcAllPct{} and \rhabXgbLoFomcAllPct\%. It moves the forecast strongly
  on few bars: on FOMC statement days \rhfipFomcUntil{} holds \rhfipLgbmUntilFomc\% of LightGBM's attribution, on
  other days \rhfipLgbmUntilOther\% (Appendix~\ref{app:rh_fi_cond}). The futures' realized moments are the converse:
  second in attribution, yet removing them raises the trees' QLIKE by only \rhabLgbmLoRmAllPct, \rhabXgbLoRmAllPct{}
  and \rhabRfLoRmAllPct\%, as other predictors carry much of the same information. What a series adds to accuracy is
  measured by the ablations (Section~\ref{sec:rh_results_ablation}), not by its share.
\item \textbf{In short.} The trees rely on the HAR components and the futures' realized moments, and their exogenous
  attribution is concentrated. The signal is not dense and weak: about 10 series carry the trees' gain and 5 to 10
  the linear models', and the lasso is at least as accurate as the ridge. The alternative data are used more while
  live, without a measurable gain, and the FOMC calendar carries most of the trees' gain on a small share of their
  attribution.
\end{itemize}
\iffalse
\begin{itemize}
\item \textbf{Models:} the closing-interval models of record on the full predictor set ($p = \rhfiP$), trained on
  the last 30 minutes of each trading day (Section~\ref{sec:rh_econ_last30}), not on the intraday scheme.
  \begin{itemize}
  \item Lasso and ridge with the penalty strength $\lambda$ selected on a pre-sample: $\lambda = \rhfiLassoAlpha$ and
    $\rhfiRidgeAlpha$ (scikit-learn parametrization). Penalty weight $\omega_j$, the factor on predictor $j$'s penalty
    term: $\underline{\omega} = \rhfiBackbonePf$ on the HAR benchmark block, $1$ on the exogenous predictors.
    Re-estimated every trading day; the scheduled re-selection of $\lambda$ never replaced the pre-sample choice.
  \item Not the earlier study's single-penalty specification (one $\lambda$ for every predictor, re-selected every
    \rhfiOldBlock{} trading days): the paper does not use it.
  \item Tree ensembles at fixed hyperparameters, the best entries of the development-sample search (LightGBM and
    XGBoost: the top Optuna trial; random forest: trial \rhdmRfTrial); constant and duplicate predictors dropped in
    each estimation window.
  \end{itemize}
\item \textbf{Notation.}
  \begin{itemize}
  \item Re-estimation $k = 1, \dots, K$, $K = \rhfiNRefits$, one every \rhfiRefitEvery{} trading days. Re-estimation
    $k$ uses the window $\mathcal W_{t_k}$, the \rhfiWindow{} trading days before its first forecast day $t_k$, and
    forecasts its out-of-sample block $\mathcal O_k$, the trading days from $t_k$ to the next re-estimation
    ($\le \rhfiRefitEvery$).
    The blocks make up this section's forecast days $\mathcal O = \bigcup_k \mathcal O_k$:
    $|\mathcal O| = \rhfiNOos$ trading days, \rhfiFirst{} to \rhfiLast.
  \item Predictors $\bm x_t \in \mathbb{R}^p$; forecast $\hat y_t = \hat f_k(\bm x_t)$, $t \in \mathcal O_k$, on the
    scale of the target $y_t$, the diurnally adjusted realized volatility of the closing half hour.
  \item Predictor group $A \subseteq \{1, \dots, p\}$: all predictors built from one series (its averages at every
    horizon, its missing-data and activity indicators). \rhfiNUnits{} groups: \rhfiNExog{} exogenous series in
    \rhfiNBuckets{} categories (Section~\ref{sec:rh_predictors}), the HAR components, the calendar dummies.
  \item Linear models: $\hat{\bm\beta}_k \in \mathbb{R}^p$, the coefficient estimates behind the forecast for $t_k$;
    $\hat{\bm\beta}_{k,A}$, its subvector on group $A$.
  \item Trees: SHAP values $\phi_{t,j}$ \citep{lundberg2017unified} and SHAP interaction values $\phi_{t,jl}$
    \citep{lundberg2020local}; they add up to the forecast,
    \[
      \hat y_t = \phi_0 + \sum_{j=1}^{p} \phi_{t,j}, \qquad
      \phi_{t,j} = \sum_{l=1}^{p} \phi_{t,jl}, \qquad
      \phi_{t,A} = \sum_{j \in A} \phi_{t,j}.
    \]
    For $j \ne l$, $\phi_{t,jl} = \tfrac12\bigl(\phi_{t,l} \mid j \text{ present} - \phi_{t,l} \mid j \text{ absent}\bigr)$,
    and $\phi_{t,jj} = \phi_{t,j} - \sum_{l \ne j} \phi_{t,jl}$.
  \item Path-dependent TreeSHAP \citep{lundberg2020local}: the training observations per node (node covers) stand in
    for the background sample.
  \item Number triples below: LightGBM, XGBoost, random forest.
  \end{itemize}
\item \textbf{Summary table:} Table~\ref{tab:rh_fi_summary} ranks the predictor groups model by model, as in the
  variable-importance rankings of \citet{gu2020empirical}, each model by its own aggregate measure below: the lasso's
  selection count $K_A$, the ridge's coefficient shares $S_A$ and $S^{\mathrm{sd}}_A$, and the trees' mean SHAP rank
  $\bar r_A$. Under each column: Kendall's $W$ of the rankings across re-estimations, and the number of exogenous
  series with a predictor that beats the random probe.
% AUTO-GENERATED by writeup/make_paper_beats.py -- do not edit; re-run the script.
% Review-layer stub of generated/rehaul_fi_summary.tex: the caption and label of each float, the body omitted.
\providecommand{\ser}[2]{#2}\providecommand{\serx}[3]{#3}
\ifdefined\rhfiNRefits\else% AUTO-GENERATED by writeup/make_rehaul_fi.py -- do not edit; re-run the script.
% \rhfi... macros for writeup/rehaul/08_feature_importance.tex (results/rehaul_fi/).
\providecommand{\rhfiNRefits}{}\renewcommand{\rhfiNRefits}{147}  % refits, one every 10 sessions
\providecommand{\rhfiNOos}{}\renewcommand{\rhfiNOos}{1{,}469}  % scored held-out sessions
\providecommand{\rhfiP}{}\renewcommand{\rhfiP}{628}  % columns, all features
\providecommand{\rhfiNUnits}{}\renewcommand{\rhfiNUnits}{49}  % units: 47 exogenous series + HAR ladder + calendar
\providecommand{\rhfiNExog}{}\renewcommand{\rhfiNExog}{47}  % exogenous series
\providecommand{\rhfiFirst}{}\renewcommand{\rhfiFirst}{25~June~2018}  % first scored session
\providecommand{\rhfiLast}{}\renewcommand{\rhfiLast}{30~April~2024}  % last scored session
\providecommand{\rhfiWindow}{}\renewcommand{\rhfiWindow}{2{,}000}  % sessions in a training window
\providecommand{\rhfiRefitEvery}{}\renewcommand{\rhfiRefitEvery}{10}  % sessions between refits
\providecommand{\rhfiBackbonePf}{}\renewcommand{\rhfiBackbonePf}{0.01}  % backbone penalty factor
\providecommand{\rhfiNBuckets}{}\renewcommand{\rhfiNBuckets}{8}  % exogenous buckets
\providecommand{\rhfiOldBlock}{}\renewcommand{\rhfiOldBlock}{250}  % sessions between the old arm's penalty re-choices
\providecommand{\rhfiNPairs}{}\renewcommand{\rhfiNPairs}{1{,}176}  % series pairs
\providecommand{\rhfiNProbeCols}{}\renewcommand{\rhfiNProbeCols}{631}  % columns of a probe fit
\providecommand{\rhfiTopWPairs}{}\renewcommand{\rhfiTopWPairs}{50}  % pairs in the interaction W
\providecommand{\rhfiProbeLevels}{}\renewcommand{\rhfiProbeLevels}{20}  % values of the few-valued noise column
\providecommand{\rhfiRidgeAlpha}{}\renewcommand{\rhfiRidgeAlpha}{\ensuremath{10^{3}}}
\providecommand{\rhfiLassoAlpha}{}\renewcommand{\rhfiLassoAlpha}{\ensuremath{10^{-2}}}
\providecommand{\rhfiLassoNzMed}{}\renewcommand{\rhfiLassoNzMed}{30}  % median nonzero coefficients per refit
\providecommand{\rhfiLassoNzMin}{}\renewcommand{\rhfiLassoNzMin}{27}
\providecommand{\rhfiLassoNzMax}{}\renewcommand{\rhfiLassoNzMax}{35}
\providecommand{\rhfiLassoExogNzMed}{}\renewcommand{\rhfiLassoExogNzMed}{14}  % ... exogenous columns
\providecommand{\rhfiLassoColsEver}{}\renewcommand{\rhfiLassoColsEver}{60}  % columns nonzero in >= 1 refit
\providecommand{\rhfiLassoColsMajority}{}\renewcommand{\rhfiLassoColsMajority}{29}  % columns nonzero in >= half the refits
\providecommand{\rhfiLassoExogEver}{}\renewcommand{\rhfiLassoExogEver}{27}  % exogenous series with a nonzero column in >= 1 refit
\providecommand{\rhfiLassoExogMajority}{}\renewcommand{\rhfiLassoExogMajority}{11}  % ... in >= half the refits
\providecommand{\rhfiLassoExogAlways}{}\renewcommand{\rhfiLassoExogAlways}{3}  % ... in every refit
\providecommand{\rhfiLassoAlwaysList}{}\renewcommand{\rhfiLassoAlwaysList}{\texttt{sumret}, \texttt{sumabsret\_vwstock}, \texttt{sumautocov}}  % exogenous series selected in every refit
\providecommand{\rhfiLassoProbeSel}{}\renewcommand{\rhfiLassoProbeSel}{1}  % continuous probe selected, of the probe fits
\providecommand{\rhfiLassoProbeSelLevels}{}\renewcommand{\rhfiLassoProbeSelLevels}{0}
\providecommand{\rhfiLassoProbeSelBinary}{}\renewcommand{\rhfiLassoProbeSelBinary}{0}
\providecommand{\rhfiNProbe}{}\renewcommand{\rhfiNProbe}{49}  % probe refits (every third)
\providecommand{\rhfiLassoW}{}\renewcommand{\rhfiLassoW}{0.74}  % Kendall W of the per-refit series ranks by sum |beta|
\providecommand{\rhfiLassoBackboneShare}{}\renewcommand{\rhfiLassoBackboneShare}{88}  % median share of sum |beta| on the backbone, %
\providecommand{\rhfiOldLassoSpearman}{}\renewcommand{\rhfiOldLassoSpearman}{0.70}  % Spearman, selection share: pre-tuned vs re-tuned one-penalty lasso
\providecommand{\rhfiOldLassoNzMed}{}\renewcommand{\rhfiOldLassoNzMed}{23}  % re-tuned lasso: median nonzero per refit
\providecommand{\rhfiOldLassoExogMajority}{}\renewcommand{\rhfiOldLassoExogMajority}{18}
\providecommand{\rhfiRidgeTopOne}{}\renewcommand{\rhfiRidgeTopOne}{Calendar}
\providecommand{\rhfiRidgeTopOneShare}{}\renewcommand{\rhfiRidgeTopOneShare}{31.0}  % % of sum |beta|
\providecommand{\rhfiRidgeTopTwo}{}\renewcommand{\rhfiRidgeTopTwo}{HAR ladder}
\providecommand{\rhfiRidgeTopTwoShare}{}\renewcommand{\rhfiRidgeTopTwoShare}{22.8}  % % of sum |beta|
\providecommand{\rhfiRidgeTopThree}{}\renewcommand{\rhfiRidgeTopThree}{\texttt{sumret}}
\providecommand{\rhfiRidgeTopThreeShare}{}\renewcommand{\rhfiRidgeTopThreeShare}{3.6}  % % of sum |beta|
\providecommand{\rhfiRidgeTopFour}{}\renewcommand{\rhfiRidgeTopFour}{\texttt{sumabsret\_vwstock}}
\providecommand{\rhfiRidgeTopFourShare}{}\renewcommand{\rhfiRidgeTopFourShare}{2.2}  % % of sum |beta|
\providecommand{\rhfiRidgeSdTopOne}{}\renewcommand{\rhfiRidgeSdTopOne}{HAR ladder}
\providecommand{\rhfiRidgeSdTopOneShare}{}\renewcommand{\rhfiRidgeSdTopOneShare}{30.6}  % % of sum |beta| sd
\providecommand{\rhfiRidgeSdTopTwo}{}\renewcommand{\rhfiRidgeSdTopTwo}{Calendar}
\providecommand{\rhfiRidgeSdTopTwoShare}{}\renewcommand{\rhfiRidgeSdTopTwoShare}{9.3}  % % of sum |beta| sd
\providecommand{\rhfiRidgeSdTopThree}{}\renewcommand{\rhfiRidgeSdTopThree}{\texttt{sumret}}
\providecommand{\rhfiRidgeSdTopThreeShare}{}\renewcommand{\rhfiRidgeSdTopThreeShare}{4.5}  % % of sum |beta| sd
\providecommand{\rhfiRidgeTopExog}{}\renewcommand{\rhfiRidgeTopExog}{\texttt{sumret}}
\providecommand{\rhfiRidgeTopExogShare}{}\renewcommand{\rhfiRidgeTopExogShare}{3.6}
\providecommand{\rhfiRidgeBackboneShare}{}\renewcommand{\rhfiRidgeBackboneShare}{54}  % median share of sum |beta| on the backbone, %
\providecommand{\rhfiRidgeW}{}\renewcommand{\rhfiRidgeW}{0.85}
\providecommand{\rhfiRidgeAllSelected}{}\renewcommand{\rhfiRidgeAllSelected}{49}  % units with N_A = K under the ridge
\providecommand{\rhfiLassoProbeRate}{}\renewcommand{\rhfiLassoProbeRate}{2.0}  % continuous probe's selection rate, %
\providecommand{\rhfiLassoColsAboveProbe}{}\renewcommand{\rhfiLassoColsAboveProbe}{54}  % columns selected more often than the continuous probe
\providecommand{\rhfiLassoSeriesAboveProbe}{}\renewcommand{\rhfiLassoSeriesAboveProbe}{23}  % exogenous series with such a column
\providecommand{\rhfiLassoTopCol}{}\renewcommand{\rhfiLassoTopCol}{\texttt{is\_month\_end}}
\providecommand{\rhfiRidgeSpearmanSd}{}\renewcommand{\rhfiRidgeSpearmanSd}{0.85}  % Spearman, series sum|beta| vs sum|beta| sd
\providecommand{\rhfiRidgeProbeBeat}{}\renewcommand{\rhfiRidgeProbeBeat}{107}  % columns with |beta| above the continuous probe's in most probe fits
\providecommand{\rhfiRidgeProbeBeatSeries}{}\renewcommand{\rhfiRidgeProbeBeatSeries}{38}
\providecommand{\rhfiScaleValueN}{}\renewcommand{\rhfiScaleValueN}{282}
\providecommand{\rhfiScaleValueSd}{}\renewcommand{\rhfiScaleValueSd}{0.84}  % median window sd, value ladders
\providecommand{\rhfiScaleValueSdLo}{}\renewcommand{\rhfiScaleValueSdLo}{0.03}
\providecommand{\rhfiScaleValueSdHi}{}\renewcommand{\rhfiScaleValueSdHi}{1.76}
\providecommand{\rhfiScaleValueIqr}{}\renewcommand{\rhfiScaleValueIqr}{0.94}  % median window IQR, value ladders
\providecommand{\rhfiScaleFlagN}{}\renewcommand{\rhfiScaleFlagN}{324}
\providecommand{\rhfiScaleFlagSd}{}\renewcommand{\rhfiScaleFlagSd}{0.022}  % median window sd, availability / activity ladders
\providecommand{\rhfiScaleFlagSdHi}{}\renewcommand{\rhfiScaleFlagSdHi}{0.17}
\providecommand{\rhfiScaleCalSd}{}\renewcommand{\rhfiScaleCalSd}{0.39}  % median window sd, calendar
\providecommand{\rhfiScaleHarSd}{}\renewcommand{\rhfiScaleHarSd}{1.07}  % median window sd, HAR ladder
\providecommand{\rhfiLgbmW}{}\renewcommand{\rhfiLgbmW}{0.57}  % lgbm: Kendall W, per-refit series ranks
\providecommand{\rhfiLgbmWCols}{}\renewcommand{\rhfiLgbmWCols}{0.91}  % lgbm: Kendall W, column ranks
\providecommand{\rhfiLgbmTopOne}{}\renewcommand{\rhfiLgbmTopOne}{HAR ladder}
\providecommand{\rhfiLgbmTopOneRank}{}\renewcommand{\rhfiLgbmTopOneRank}{1.0}  % mean rank
\providecommand{\rhfiLgbmTopTwo}{}\renewcommand{\rhfiLgbmTopTwo}{\texttt{sumret}}
\providecommand{\rhfiLgbmTopTwoRank}{}\renewcommand{\rhfiLgbmTopTwoRank}{2.9}  % mean rank
\providecommand{\rhfiLgbmTopThree}{}\renewcommand{\rhfiLgbmTopThree}{\texttt{sumabsret}}
\providecommand{\rhfiLgbmTopThreeRank}{}\renewcommand{\rhfiLgbmTopThreeRank}{4.8}  % mean rank
\providecommand{\rhfiLgbmSpBorda}{}\renewcommand{\rhfiLgbmSpBorda}{0.90}  % Spearman, mean rank vs top-10 Borda
\providecommand{\rhfiLgbmSpGlobal}{}\renewcommand{\rhfiLgbmSpGlobal}{0.96}  % Spearman, mean rank vs global mean |phi|
\providecommand{\rhfiLgbmSpLocal}{}\renewcommand{\rhfiLgbmSpLocal}{1.000}  % Spearman, mean ranks local refit vs stored
\providecommand{\rhfiLgbmProbeRank}{}\renewcommand{\rhfiLgbmProbeRank}{83}  % median rank of the continuous probe, 631 columns
\providecommand{\rhfiLgbmProbeCols}{}\renewcommand{\rhfiLgbmProbeCols}{60}  % columns above the continuous probe in most probe fits
\providecommand{\rhfiLgbmProbeSeries}{}\renewcommand{\rhfiLgbmProbeSeries}{27}  % exogenous series with such a column
\providecommand{\rhfiLgbmTopCol}{}\renewcommand{\rhfiLgbmTopCol}{\texttt{har\_ma\_1}}  % column with the lowest mean rank
\providecommand{\rhfiLgbmTopColRank}{}\renewcommand{\rhfiLgbmTopColRank}{1.3}
\providecommand{\rhfiLgbmTopColTwo}{}\renewcommand{\rhfiLgbmTopColTwo}{\texttt{har\_ma\_5}}
\providecommand{\rhfiLgbmTopColTwoRank}{}\renewcommand{\rhfiLgbmTopColTwoRank}{1.8}
\providecommand{\rhfiLgbmPredGap}{}\renewcommand{\rhfiLgbmPredGap}{0}  % max |local refit - stored refit| forecast
\providecommand{\rhfiLgbmInterW}{}\renewcommand{\rhfiLgbmInterW}{0.45}  % Kendall W of per-refit pair ranks, top-50 pairs
\providecommand{\rhfiLgbmInterRatio}{}\renewcommand{\rhfiLgbmInterRatio}{90}  % sum of series-pair interactions / sum main, %
\providecommand{\rhfiLgbmInterSpGlobal}{}\renewcommand{\rhfiLgbmInterSpGlobal}{0.99}
\providecommand{\rhfiLgbmInterFloor}{}\renewcommand{\rhfiLgbmInterFloor}{0.49}  % probe floor x 1e3
\providecommand{\rhfiLgbmInterAboveFloor}{}\renewcommand{\rhfiLgbmInterAboveFloor}{116}  % series pairs above the probe floor
\providecommand{\rhfiLgbmInterTopA}{}\renewcommand{\rhfiLgbmInterTopA}{HAR ladder}
\providecommand{\rhfiLgbmInterTopB}{}\renewcommand{\rhfiLgbmInterTopB}{\texttt{sumabsret}}
\providecommand{\rhfiLgbmInterTopVal}{}\renewcommand{\rhfiLgbmInterTopVal}{12.6}  % x 1e3
\providecommand{\rhfiLgbmHTop}{}\renewcommand{\rhfiLgbmHTop}{0.33}  % H^2 of the top series pair (pooled), %
\providecommand{\rhfiLgbmHMin}{}\renewcommand{\rhfiLgbmHMin}{0.02}  % min H^2 over the top-10 series pairs, %
\providecommand{\rhfiLgbmHMax}{}\renewcommand{\rhfiLgbmHMax}{0.33}  % max H^2 over the top-10 series pairs, %
\providecommand{\rhfiLgbmHMinTop}{}\renewcommand{\rhfiLgbmHMinTop}{0.19}  % top pair: interaction PD var / smaller main PD var
\providecommand{\rhfiLgbmHMinLo}{}\renewcommand{\rhfiLgbmHMinLo}{0.05}  % ... min over the top-10 series pairs
\providecommand{\rhfiLgbmHMinHi}{}\renewcommand{\rhfiLgbmHMinHi}{0.54}  % ... max
\providecommand{\rhfiLgbmHColTop}{}\renewcommand{\rhfiLgbmHColTop}{0.36}  % H^2 of the top column pair, %
\providecommand{\rhfiLgbmHMinColTop}{}\renewcommand{\rhfiLgbmHMinColTop}{0.02}  % ... min-normalised
\providecommand{\rhfiLgbmMainTopVal}{}\renewcommand{\rhfiLgbmMainTopVal}{195}  % largest global mean |phi_A| x 1e3
\providecommand{\rhfiLgbmWithinTop}{}\renewcommand{\rhfiLgbmWithinTop}{HAR ladder}
\providecommand{\rhfiLgbmWithinTopVal}{}\renewcommand{\rhfiLgbmWithinTopVal}{18.7}
\providecommand{\rhfiXgbW}{}\renewcommand{\rhfiXgbW}{0.60}  % xgb: Kendall W, per-refit series ranks
\providecommand{\rhfiXgbWCols}{}\renewcommand{\rhfiXgbWCols}{0.90}  % xgb: Kendall W, column ranks
\providecommand{\rhfiXgbTopOne}{}\renewcommand{\rhfiXgbTopOne}{HAR ladder}
\providecommand{\rhfiXgbTopOneRank}{}\renewcommand{\rhfiXgbTopOneRank}{1.0}  % mean rank
\providecommand{\rhfiXgbTopTwo}{}\renewcommand{\rhfiXgbTopTwo}{\texttt{sumret}}
\providecommand{\rhfiXgbTopTwoRank}{}\renewcommand{\rhfiXgbTopTwoRank}{3.0}  % mean rank
\providecommand{\rhfiXgbTopThree}{}\renewcommand{\rhfiXgbTopThree}{\texttt{sumabsret}}
\providecommand{\rhfiXgbTopThreeRank}{}\renewcommand{\rhfiXgbTopThreeRank}{3.9}  % mean rank
\providecommand{\rhfiXgbSpBorda}{}\renewcommand{\rhfiXgbSpBorda}{0.87}  % Spearman, mean rank vs top-10 Borda
\providecommand{\rhfiXgbSpGlobal}{}\renewcommand{\rhfiXgbSpGlobal}{0.97}  % Spearman, mean rank vs global mean |phi|
\providecommand{\rhfiXgbSpLocal}{}\renewcommand{\rhfiXgbSpLocal}{0.993}  % Spearman, mean ranks local refit vs stored
\providecommand{\rhfiXgbProbeRank}{}\renewcommand{\rhfiXgbProbeRank}{110}  % median rank of the continuous probe, 631 columns
\providecommand{\rhfiXgbProbeCols}{}\renewcommand{\rhfiXgbProbeCols}{78}  % columns above the continuous probe in most probe fits
\providecommand{\rhfiXgbProbeSeries}{}\renewcommand{\rhfiXgbProbeSeries}{35}  % exogenous series with such a column
\providecommand{\rhfiXgbTopCol}{}\renewcommand{\rhfiXgbTopCol}{\texttt{har\_ma\_1}}  % column with the lowest mean rank
\providecommand{\rhfiXgbTopColRank}{}\renewcommand{\rhfiXgbTopColRank}{1.3}
\providecommand{\rhfiXgbTopColTwo}{}\renewcommand{\rhfiXgbTopColTwo}{\texttt{har\_ma\_5}}
\providecommand{\rhfiXgbTopColTwoRank}{}\renewcommand{\rhfiXgbTopColTwoRank}{1.8}
\providecommand{\rhfiXgbPredGap}{}\renewcommand{\rhfiXgbPredGap}{0.17}  % max |local refit - stored refit| forecast
\providecommand{\rhfiXgbInterW}{}\renewcommand{\rhfiXgbInterW}{0.47}  % Kendall W of per-refit pair ranks, top-50 pairs
\providecommand{\rhfiXgbInterRatio}{}\renewcommand{\rhfiXgbInterRatio}{74}  % sum of series-pair interactions / sum main, %
\providecommand{\rhfiXgbInterSpGlobal}{}\renewcommand{\rhfiXgbInterSpGlobal}{0.99}
\providecommand{\rhfiXgbInterFloor}{}\renewcommand{\rhfiXgbInterFloor}{0.41}  % probe floor x 1e3
\providecommand{\rhfiXgbInterAboveFloor}{}\renewcommand{\rhfiXgbInterAboveFloor}{115}  % series pairs above the probe floor
\providecommand{\rhfiXgbInterTopA}{}\renewcommand{\rhfiXgbInterTopA}{HAR ladder}
\providecommand{\rhfiXgbInterTopB}{}\renewcommand{\rhfiXgbInterTopB}{\texttt{sumabsret}}
\providecommand{\rhfiXgbInterTopVal}{}\renewcommand{\rhfiXgbInterTopVal}{15.2}  % x 1e3
\providecommand{\rhfiXgbHTop}{}\renewcommand{\rhfiXgbHTop}{0.37}  % H^2 of the top series pair (pooled), %
\providecommand{\rhfiXgbHMin}{}\renewcommand{\rhfiXgbHMin}{0.02}  % min H^2 over the top-10 series pairs, %
\providecommand{\rhfiXgbHMax}{}\renewcommand{\rhfiXgbHMax}{0.68}  % max H^2 over the top-10 series pairs, %
\providecommand{\rhfiXgbHMinTop}{}\renewcommand{\rhfiXgbHMinTop}{0.39}  % top pair: interaction PD var / smaller main PD var
\providecommand{\rhfiXgbHMinLo}{}\renewcommand{\rhfiXgbHMinLo}{0.08}  % ... min over the top-10 series pairs
\providecommand{\rhfiXgbHMinHi}{}\renewcommand{\rhfiXgbHMinHi}{0.80}  % ... max
\providecommand{\rhfiXgbHColTop}{}\renewcommand{\rhfiXgbHColTop}{0.39}  % H^2 of the top column pair, %
\providecommand{\rhfiXgbHMinColTop}{}\renewcommand{\rhfiXgbHMinColTop}{0.02}  % ... min-normalised
\providecommand{\rhfiXgbMainTopVal}{}\renewcommand{\rhfiXgbMainTopVal}{192}  % largest global mean |phi_A| x 1e3
\providecommand{\rhfiXgbWithinTop}{}\renewcommand{\rhfiXgbWithinTop}{HAR ladder}
\providecommand{\rhfiXgbWithinTopVal}{}\renewcommand{\rhfiXgbWithinTopVal}{19.1}
\providecommand{\rhfiRfW}{}\renewcommand{\rhfiRfW}{0.62}  % rf: Kendall W, per-refit series ranks
\providecommand{\rhfiRfWCols}{}\renewcommand{\rhfiRfWCols}{0.93}  % rf: Kendall W, column ranks
\providecommand{\rhfiRfTopOne}{}\renewcommand{\rhfiRfTopOne}{HAR ladder}
\providecommand{\rhfiRfTopOneRank}{}\renewcommand{\rhfiRfTopOneRank}{1.0}  % mean rank
\providecommand{\rhfiRfTopTwo}{}\renewcommand{\rhfiRfTopTwo}{\texttt{sumabsret}}
\providecommand{\rhfiRfTopTwoRank}{}\renewcommand{\rhfiRfTopTwoRank}{2.6}  % mean rank
\providecommand{\rhfiRfTopThree}{}\renewcommand{\rhfiRfTopThree}{\texttt{sumret}}
\providecommand{\rhfiRfTopThreeRank}{}\renewcommand{\rhfiRfTopThreeRank}{4.0}  % mean rank
\providecommand{\rhfiRfSpBorda}{}\renewcommand{\rhfiRfSpBorda}{0.92}  % Spearman, mean rank vs top-10 Borda
\providecommand{\rhfiRfSpGlobal}{}\renewcommand{\rhfiRfSpGlobal}{0.98}  % Spearman, mean rank vs global mean |phi|
\providecommand{\rhfiRfSpLocal}{}\renewcommand{\rhfiRfSpLocal}{0.996}  % Spearman, mean ranks local refit vs stored
\providecommand{\rhfiRfProbeRank}{}\renewcommand{\rhfiRfProbeRank}{97}  % median rank of the continuous probe, 631 columns
\providecommand{\rhfiRfProbeCols}{}\renewcommand{\rhfiRfProbeCols}{71}  % columns above the continuous probe in most probe fits
\providecommand{\rhfiRfProbeSeries}{}\renewcommand{\rhfiRfProbeSeries}{29}  % exogenous series with such a column
\providecommand{\rhfiRfTopCol}{}\renewcommand{\rhfiRfTopCol}{\texttt{har\_ma\_1}}  % column with the lowest mean rank
\providecommand{\rhfiRfTopColRank}{}\renewcommand{\rhfiRfTopColRank}{1.3}
\providecommand{\rhfiRfTopColTwo}{}\renewcommand{\rhfiRfTopColTwo}{\texttt{har\_ma\_5}}
\providecommand{\rhfiRfTopColTwoRank}{}\renewcommand{\rhfiRfTopColTwoRank}{1.7}
\providecommand{\rhfiRfPredGap}{}\renewcommand{\rhfiRfPredGap}{0.11}  % max |local refit - stored refit| forecast
\providecommand{\rhfiRfInterW}{}\renewcommand{\rhfiRfInterW}{0.57}  % Kendall W of per-refit pair ranks, top-50 pairs
\providecommand{\rhfiRfInterRatio}{}\renewcommand{\rhfiRfInterRatio}{72}  % sum of series-pair interactions / sum main, %
\providecommand{\rhfiRfInterSpGlobal}{}\renewcommand{\rhfiRfInterSpGlobal}{1.00}
\providecommand{\rhfiRfInterFloor}{}\renewcommand{\rhfiRfInterFloor}{0.50}  % probe floor x 1e3
\providecommand{\rhfiRfInterAboveFloor}{}\renewcommand{\rhfiRfInterAboveFloor}{52}  % series pairs above the probe floor
\providecommand{\rhfiRfInterTopA}{}\renewcommand{\rhfiRfInterTopA}{HAR ladder}
\providecommand{\rhfiRfInterTopB}{}\renewcommand{\rhfiRfInterTopB}{\texttt{sumabsret}}
\providecommand{\rhfiRfInterTopVal}{}\renewcommand{\rhfiRfInterTopVal}{23.8}  % x 1e3
\providecommand{\rhfiRfHTop}{}\renewcommand{\rhfiRfHTop}{0.85}  % H^2 of the top series pair (pooled), %
\providecommand{\rhfiRfHMin}{}\renewcommand{\rhfiRfHMin}{0.02}  % min H^2 over the top-10 series pairs, %
\providecommand{\rhfiRfHMax}{}\renewcommand{\rhfiRfHMax}{0.85}  % max H^2 over the top-10 series pairs, %
\providecommand{\rhfiRfHMinTop}{}\renewcommand{\rhfiRfHMinTop}{1.50}  % top pair: interaction PD var / smaller main PD var
\providecommand{\rhfiRfHMinLo}{}\renewcommand{\rhfiRfHMinLo}{0.17}  % ... min over the top-10 series pairs
\providecommand{\rhfiRfHMinHi}{}\renewcommand{\rhfiRfHMinHi}{5.46}  % ... max
\providecommand{\rhfiRfHColTop}{}\renewcommand{\rhfiRfHColTop}{1.16}  % H^2 of the top column pair, %
\providecommand{\rhfiRfHMinColTop}{}\renewcommand{\rhfiRfHMinColTop}{0.06}  % ... min-normalised
\providecommand{\rhfiRfMainTopVal}{}\renewcommand{\rhfiRfMainTopVal}{224}  % largest global mean |phi_A| x 1e3
\providecommand{\rhfiRfWithinTop}{}\renewcommand{\rhfiRfWithinTop}{HAR ladder}
\providecommand{\rhfiRfWithinTopVal}{}\renewcommand{\rhfiRfWithinTopVal}{34.3}
\providecommand{\rhfiSpTreesMin}{}\renewcommand{\rhfiSpTreesMin}{0.83}  % Spearman of mean ranks, trees
\providecommand{\rhfiSpTreesMax}{}\renewcommand{\rhfiSpTreesMax}{0.96}
\providecommand{\rhfiSpTreeLinMin}{}\renewcommand{\rhfiSpTreeLinMin}{0.17}  % Spearman of mean ranks, tree vs linear
\providecommand{\rhfiSpTreeLinMax}{}\renewcommand{\rhfiSpTreeLinMax}{0.55}
\providecommand{\rhfiSpLassoRidge}{}\renewcommand{\rhfiSpLassoRidge}{0.47}
\providecommand{\rhfiHRefits}{}\renewcommand{\rhfiHRefits}{15}  % refits for H
\providecommand{\rhfiHRows}{}\renewcommand{\rhfiHRows}{100}  % evaluation points per refit
\providecommand{\rhfiHTopSeries}{}\renewcommand{\rhfiHTopSeries}{10}
\providecommand{\rhfiFitSecRf}{}\renewcommand{\rhfiFitSecRf}{212}  % mean seconds per forest refit (11 processes sharing the CPU)
\providecommand{\rhfiShapSecRf}{}\renewcommand{\rhfiShapSecRf}{182}  % mean seconds per forest refit, SHAP + interactions
\providecommand{\rhfiShapSecLgbm}{}\renewcommand{\rhfiShapSecLgbm}{0.9}  % ... LightGBM
\providecommand{\rhfiShapSecXgb}{}\renewcommand{\rhfiShapSecXgb}{0.4}  % ... XGBoost
\providecommand{\rhfiWallTreesMin}{}\renewcommand{\rhfiWallTreesMin}{133}  % wall minutes, tree stage on 11 processes
\providecommand{\rhfiWallHMin}{}\renewcommand{\rhfiWallHMin}{7}  % wall minutes, H stage
\providecommand{\rhfiWallLinMin}{}\renewcommand{\rhfiWallLinMin}{7}  % wall minutes, linear stage
\providecommand{\rhfiWorkers}{}\renewcommand{\rhfiWorkers}{11}
\fi
\begin{table}[tbp]
\centering
\caption{Variable importance by model}
\label{tab:rh_fi_summary}
\textit{[body omitted in this review layer]}
\end{table}

\item \textbf{Lasso, aggregate:} the selection count of a group and of a predictor, the number of the $K$
  re-estimations with a nonzero coefficient,
  \[
    K_A = \sum_{k=1}^{K} \mathbf{1}\bigl\{\exists j \in A:\ \hat\beta_{k,j} \ne 0\bigr\}, \qquad
    K_j = \sum_{k=1}^{K} \mathbf{1}\bigl\{\hat\beta_{k,j} \ne 0\bigr\};
  \]
  $K_A / K$ is the selection frequency, the share of estimation windows with a nonzero coefficient.
  \begin{itemize}
  \item Per re-estimation: \rhfiLassoNzMed{} nonzero coefficients (median; range \rhfiLassoNzMin--\rhfiLassoNzMax), of
    which \rhfiLassoExogNzMed{} on exogenous predictors.
  \item Predictors: \rhfiLassoColsEver{} of \rhfiP{} selected at least once, \rhfiLassoColsMajority{} in at least half
    of the re-estimations.
  \item Exogenous series: \rhfiLassoExogEver{} of \rhfiNExog{} selected at least once; \rhfiLassoExogMajority{} in at
    least half of the re-estimations; \rhfiLassoExogAlways{} in every re-estimation (\rhfiLassoAlwaysList).
  \item HAR benchmark block: selected in every re-estimation ($\omega_j = \rhfiBackbonePf$);
    \rhfiLassoBackboneShare\,\% of $\sum_j |\hat\beta_{k,j}|$ (median re-estimation).
  \item Random-probe benchmark: the continuous probe, an i.i.d.\ uniform noise predictor (defined under Known issues),
    is selected in \rhfiLassoProbeRate\,\% of its fits; \rhfiLassoColsAboveProbe{} predictors, in
    \rhfiLassoSeriesAboveProbe{} exogenous series, are selected more often.
  \end{itemize}
\item \textbf{Lasso, time-varying:} the magnitude and the selection indicator of each group at each re-estimation,
  \[
    \lVert\hat{\bm\beta}_{k,A}\rVert_1 = \sum_{j \in A} \bigl|\hat\beta_{k,j}\bigr|, \qquad
    \mathbf{1}\bigl\{\lVert\hat{\bm\beta}_{k,A}\rVert_1 > 0\bigr\}
  \]
  Rank stability across re-estimations: Kendall's coefficient of
  concordance $W$ of the rankings by $\lVert\hat{\bm\beta}_{k,A}\rVert_1$ at each re-estimation ($W = 1$: the same
  ranking at every re-estimation; formula below), $W = \rhfiLassoW$.
\item \textbf{Ridge, aggregate:} the group's share of the absolute coefficient mass,
  \[
    S_A = \frac{\sum_{k} \sum_{j \in A} |\hat\beta_{k,j}|}{\sum_{k} \sum_{j=1}^{p} |\hat\beta_{k,j}|} .
  \]
  \begin{itemize}
  \item Scale check: only the averages of series values are robust-scaled.
    \begin{itemize}
    \item Averages of series values (\rhfiScaleValueN{} predictors): standard deviation over the estimation window
      \rhfiScaleValueSd{} (median; 5th--95th percentiles: \rhfiScaleValueSdLo--\rhfiScaleValueSdHi), interquartile
      range \rhfiScaleValueIqr.
    \item Averaged missing-data and activity indicators (\rhfiScaleFlagN): standard deviation \rhfiScaleFlagSd;
      calendar dummies \rhfiScaleCalSd; HAR components \rhfiScaleHarSd.
    \item So the raw $|\hat\beta_{k,j}|$ are in different units. Check with the group's share of the absolute
      standardized coefficients,
      \[
        S^{\mathrm{sd}}_A = \frac{\sum_{k} \sum_{j \in A} |\hat\beta_{k,j}|\, \mathrm{sd}_{k,j}}
          {\sum_{k} \sum_{j=1}^{p} |\hat\beta_{k,j}|\, \mathrm{sd}_{k,j}},
      \]
      with $\mathrm{sd}_{k,j}$ the standard deviation of predictor $j$ over $\mathcal W_{t_k}$;
      Spearman rank correlation of $S^{\mathrm{sd}}_A$
      with $S_A$ across groups \rhfiRidgeSpearmanSd.
    \item Leading groups on that scale: \rhfiRidgeSdTopOne{} \rhfiRidgeSdTopOneShare\,\%, \rhfiRidgeSdTopTwo{}
      \rhfiRidgeSdTopTwoShare\,\%, \rhfiRidgeSdTopThree{} \rhfiRidgeSdTopThreeShare\,\%.
    \end{itemize}
  \item Leading groups: \rhfiRidgeTopOne{} \rhfiRidgeTopOneShare\,\%, \rhfiRidgeTopTwo{} \rhfiRidgeTopTwoShare\,\%,
    \rhfiRidgeTopThree{} \rhfiRidgeTopThreeShare\,\%; HAR benchmark block \rhfiRidgeBackboneShare\,\% (median
    re-estimation).
  \item Ridge sets no coefficient to zero: all \rhfiRidgeAllSelected{} groups have $K_A = K$, so the selection frequency
    does not discriminate.
  \item Random-probe benchmark: \rhfiRidgeProbeBeat{} predictors, in \rhfiRidgeProbeBeatSeries{} exogenous series, have
    $|\hat\beta_{k,j}|$ above the continuous probe's (robust-scaled) in most probe fits.
  \end{itemize}
\item \textbf{Ridge, time-varying:} the share at re-estimation $k$,
  $S_{k,A} = \lVert\hat{\bm\beta}_{k,A}\rVert_1 / \sum_{A'} \lVert\hat{\bm\beta}_{k,A'}\rVert_1$, $A'$ over all groups
  (mean rank by $\lVert\hat{\bm\beta}_{k,A}\rVert_1$); Kendall's
  $W = \rhfiRidgeW$.
\item \textbf{Trees, ranking at each re-estimation:} the group's mean absolute SHAP value on the out-of-sample block
  of re-estimation $k$, then its rank,
  \[
    I_{k,A} = \frac{1}{|\mathcal O_k|} \sum_{t \in \mathcal O_k} |\phi_{t,A}|, \qquad
    r_{k,A} = 1 + \#\{A' : I_{k,A'} > I_{k,A}\} \quad \text{(ties averaged)}.
  \]
\item \textbf{Trees, global ranking:} aggregate the $K$ rankings over time,
  \[
    \bar r_A = \frac{1}{K} \sum_{k=1}^{K} r_{k,A}, \qquad
    \mathrm{BC}_A = \sum_{k=1}^{K} \max\bigl(0,\ 11 - r_{k,A}\bigr).
  \]
  \begin{itemize}
  \item $\mathrm{BC}_A$: truncated Borda count (a full Borda count is affine in $\bar r_A$).
  \item Kendall's coefficient of concordance \citep{KendallBabingtonSmith1939} over the $K$ rankings of the $n$
    groups, with rank sums
    $R_A = \sum_k r_{k,A}$ and their mean $\bar R$:
    \[
      W = \frac{12 \sum_{A} \bigl(R_A - \bar R\bigr)^2}
        {K^2 (n^3 - n) - K \sum_{k} \sum_{g \in \mathcal E_k} \bigl(|g|^3 - |g|\bigr)},
    \]
    $\mathcal E_k$ the sets of tied (equally ranked) groups at re-estimation $k$; $W = 1$: the same ranking at every
    re-estimation.
  \item LightGBM: \rhfiLgbmTopOne{} ($\bar r_A = \rhfiLgbmTopOneRank$), \rhfiLgbmTopTwo{} (\rhfiLgbmTopTwoRank),
    \rhfiLgbmTopThree{} (\rhfiLgbmTopThreeRank); Kendall's $W = \rhfiLgbmW$.
  \item XGBoost: \rhfiXgbTopOne{} (\rhfiXgbTopOneRank), \rhfiXgbTopTwo{} (\rhfiXgbTopTwoRank), \rhfiXgbTopThree{}
    (\rhfiXgbTopThreeRank); Kendall's $W = \rhfiXgbW$.
  \item Random forest: \rhfiRfTopOne{} (\rhfiRfTopOneRank), \rhfiRfTopTwo{} (\rhfiRfTopTwoRank), \rhfiRfTopThree{}
    (\rhfiRfTopThreeRank); Kendall's $W = \rhfiRfW$.
  \item Mean rank vs truncated Borda count, Spearman rank correlation over groups: \rhfiLgbmSpBorda, \rhfiXgbSpBorda,
    \rhfiRfSpBorda.
  \item Aggregate of the ranks vs rank of the aggregate, the mean absolute SHAP value over $\mathcal O$,
    $\bar I_A = |\mathcal O|^{-1} \sum_{t \in \mathcal O} |\phi_{t,A}|$: \rhfiLgbmSpGlobal, \rhfiXgbSpGlobal, \rhfiRfSpGlobal.
  \item Predictor level ($n = \rhfiP$): Kendall's $W = \rhfiLgbmWCols$, \rhfiXgbWCols, \rhfiRfWCols.
  \item Agreement of $\bar r_A$ across models (Spearman rank correlation): trees \rhfiSpTreesMin--\rhfiSpTreesMax;
    tree vs linear \rhfiSpTreeLinMin--\rhfiSpTreeLinMax; lasso vs ridge \rhfiSpLassoRidge.
  \end{itemize}
\item \textbf{Interaction effects, global:} the mean absolute SHAP interaction value of two groups over
  $\mathcal O$,
  \[
    \phi_{t,AB} = \sum_{j \in A} \sum_{l \in B} \bigl(\phi_{t,jl} + \phi_{t,lj}\bigr), \qquad
    \bar I_{AB} = \frac{1}{|\mathcal O|} \sum_{t \in \mathcal O} |\phi_{t,AB}|, \quad A \ne B;
  \]
  for two predictors, $|\phi_{t,jl} + \phi_{t,lj}|$, $j \ne l$.
  \begin{itemize}
  \item Top pair: LightGBM \rhfiLgbmInterTopA{} $\times$ \rhfiLgbmInterTopB{} (\rhfiLgbmInterTopVal);
    XGBoost \rhfiXgbInterTopA{} $\times$ \rhfiXgbInterTopB{} (\rhfiXgbInterTopVal); random forest
    \rhfiRfInterTopA{} $\times$ \rhfiRfInterTopB{} (\rhfiRfInterTopVal); $\times 10^{-3}$, in units of $y_t$.
  \item Scale: largest $\bar I_A$ \rhfiLgbmMainTopVal, \rhfiXgbMainTopVal, \rhfiRfMainTopVal{} ($\times 10^{-3}$); all
    series-pair $\bar I_{AB}$ summed, as \% of $\sum_A \bar I_A$: \rhfiLgbmInterRatio, \rhfiXgbInterRatio,
    \rhfiRfInterRatio.
  \item Within a group ($\sum_{j \ne l \in A} \phi_{t,jl}$): largest \rhfiLgbmWithinTop{} (\rhfiLgbmWithinTopVal),
    \rhfiXgbWithinTop{} (\rhfiXgbWithinTopVal), \rhfiRfWithinTop{} (\rhfiRfWithinTopVal); $\times 10^{-3}$.
  \item Noise floor: the continuous probe's largest $\bar I$ with any group, median over probe fits:
    \rhfiLgbmInterFloor, \rhfiXgbInterFloor, \rhfiRfInterFloor{} ($\times 10^{-3}$); series pairs above it:
    \rhfiLgbmInterAboveFloor, \rhfiXgbInterAboveFloor, \rhfiRfInterAboveFloor{} of \rhfiNPairs.
  \end{itemize}
\item \textbf{Interaction effects, cross-check:} the $H$-statistic of \citet{FriedmanPopescu2008} on each model's top
  \rhfiHTopSeries{} series pairs, from the partial dependence functions $\mathrm{PD}_S$, $S \in \{A, B, AB\}$,
  $AB \equiv A \cup B$:
  \[
    H^2_{AB} = \frac{\sum_{t} \bigl[\mathrm{PD}_{AB}(\bm x_t) - \mathrm{PD}_A(\bm x_t) - \mathrm{PD}_B(\bm x_t)\bigr]^2}
                    {\sum_{t} \mathrm{PD}_{AB}(\bm x_t)^2}, \qquad
    \mathrm{PD}_S(\bm x) = \frac{1}{|\mathcal H_k|} \sum_{t' \in \mathcal H_k} \hat f_k\bigl(\bm x_S, \bm x_{t',-S}\bigr),
  \]
  $t \in \mathcal H_k$; each $\mathrm{PD}_S$ centered over $\mathcal H_k$.
  \begin{itemize}
  \item Points: $\mathcal H_k$, the \rhfiHRows{} trading days from $t_k$ on, both evaluation points and background
    sample, at \rhfiHRefits{} re-estimations (every tenth); sums pooled over re-estimations.
  \item The denominator contains the HAR components' own partial dependence whenever they are in the pair. A second
    normalization, by the smaller single-group partial dependence:
    \[
      H^2_{\min,AB} = \frac{\sum_{t} \bigl[\mathrm{PD}_{AB}(\bm x_t) - \mathrm{PD}_A(\bm x_t) - \mathrm{PD}_B(\bm x_t)\bigr]^2}
                           {\min\bigl\{\sum_{t} \mathrm{PD}_A(\bm x_t)^2,\ \sum_{t} \mathrm{PD}_B(\bm x_t)^2\bigr\}} .
    \]
  \item Top pair, $H^2$ (\%): \rhfiLgbmHTop, \rhfiXgbHTop, \rhfiRfHTop; over the top \rhfiHTopSeries:
    \rhfiLgbmHMin--\rhfiLgbmHMax, \rhfiXgbHMin--\rhfiXgbHMax, \rhfiRfHMin--\rhfiRfHMax.
  \item Top pair, $H^2_{\min}$: \rhfiLgbmHMinTop, \rhfiXgbHMinTop, \rhfiRfHMinTop; over the top
    \rhfiHTopSeries: \rhfiLgbmHMinLo--\rhfiLgbmHMinHi, \rhfiXgbHMinLo--\rhfiXgbHMinHi, \rhfiRfHMinLo--\rhfiRfHMinHi.
  \item Top predictor pair, $H^2$ (\%) and $H^2_{\min}$: LightGBM \rhfiLgbmHColTop, \rhfiLgbmHMinColTop; XGBoost
    \rhfiXgbHColTop, \rhfiXgbHMinColTop; random forest \rhfiRfHColTop, \rhfiRfHMinColTop.
  \end{itemize}
\item \textbf{Interaction effects, time-varying:} $\phi_{t,AB}$ on the out-of-sample block of each re-estimation,
  ranked at each re-estimation and aggregated as the SHAP values are,
  \[
    I_{k,AB} = \frac{1}{|\mathcal O_k|} \sum_{t \in \mathcal O_k} |\phi_{t,AB}| \ \to\ r_{k,AB} \ \to\ \bar r_{AB},\ W .
  \]
  \begin{itemize}
  \item Kendall's $W$ over each model's top \rhfiTopWPairs{} pairs: \rhfiLgbmInterW, \rhfiXgbInterW, \rhfiRfInterW.
  \item Mean rank vs $\bar I_{AB}$ (Spearman rank correlation, all pairs): \rhfiLgbmInterSpGlobal,
    \rhfiXgbInterSpGlobal, \rhfiRfInterSpGlobal.
  \end{itemize}
\item \textbf{Known issues, reconciled:}
  \begin{itemize}
  \item Mean decrease in impurity (MDI) favors predictors with many distinct values
    \citep{StroblBoulesteixZeileisHothorn2007}: MDI is not used; SHAP values and coefficients only.
  \item Random probes, at every third re-estimation (\rhfiNProbe{} fits): three i.i.d.\ noise predictors, continuous
    ($\mathrm{U}(0,1)$), discrete uniform on \rhfiProbeLevels{} levels, and binary (Bernoulli, probability one half).
    \begin{itemize}
    \item Lasso: the probes are selected in \rhfiLassoProbeSel, \rhfiLassoProbeSelLevels, \rhfiLassoProbeSelBinary{}
      of \rhfiNProbe{} fits (robust-scaled, $\omega_j = 1$).
    \item Ridge: \rhfiRidgeProbeBeat{} predictors have $|\hat\beta_{k,j}|$ above the continuous probe's in most probe
      fits.
    \item Trees: the continuous probe's median rank among \rhfiNProbeCols{} predictors by mean $|\phi_{t,j}|$:
      \rhfiLgbmProbeRank,
      \rhfiXgbProbeRank, \rhfiRfProbeRank; exogenous series with a predictor above it in most probe fits:
      \rhfiLgbmProbeSeries, \rhfiXgbProbeSeries, \rhfiRfProbeSeries.
    \end{itemize}
  \item Dilution across correlated predictors: $\phi_{t,A}$ and $\lVert\hat{\bm\beta}_{k,A}\rVert_1$ sum over the
    group's predictors, so credit split among a series' averages at different horizons stays inside the group.
  \item The earlier study's single-penalty lasso ($\lambda$ re-selected): median \rhfiOldLassoNzMed{} nonzero
    coefficients per re-estimation; \rhfiOldLassoExogMajority{} exogenous series selected in at least half of the
    re-estimations;
    Spearman rank correlation of its $K_A$ with that of the lasso of record \rhfiOldLassoSpearman.
  \item Reproduction: re-running the estimation reproduces the archived LightGBM models exactly; XGBoost and the
    random forest differ across library versions (maximum absolute forecast difference \rhfiXgbPredGap{} and
    \rhfiRfPredGap, in units of $y_t$). SHAP values: from the archived models; interaction values: from the reproduced
    ones (Spearman rank correlation of $\bar r_A$, reproduced vs archived: \rhfiLgbmSpLocal, \rhfiXgbSpLocal,
    \rhfiRfSpLocal).
  \end{itemize}
\item \textbf{Compute} (one laptop, \rhfiWorkers{} processes):
  \begin{itemize}
  \item Reproducing the tree models and computing the SHAP interaction values: \rhfiWallTreesMin{} minutes. Per
    re-estimation: random forest
    estimation \rhfiFitSecRf\,s, its SHAP and interaction values \rhfiShapSecRf\,s; LightGBM \rhfiShapSecLgbm\,s,
    XGBoost \rhfiShapSecXgb\,s.
  \item Linear coefficient paths and probes: \rhfiWallLinMin{} minutes; $H^2$: \rhfiWallHMin{} minutes.
  \end{itemize}
\end{itemize}
\fi


